% ═══════════════════════════════════════════════════════════════
% Paper 2 Technical Supplement
% The Structure of Admissible Physics: End-to-End Derivations
% ═══════════════════════════════════════════════════════════════

\documentclass[11pt,a4paper]{article}
\usepackage[margin=1in]{geometry}
\usepackage{amsmath,amssymb,amsthm,mathtools}
\usepackage{enumitem}
\usepackage{booktabs}
\usepackage{array}
\usepackage[colorlinks=true,linkcolor=black,citecolor=black,urlcolor=black]{hyperref}
\usepackage[most]{tcolorbox}
\tcbuselibrary{skins,breakable}
\usepackage{fancyhdr}
\usepackage{caption}

% ── Theorem environments ──
\newtheorem{theorem}{Theorem}[section]
\newtheorem{lemma}[theorem]{Lemma}
\newtheorem{proposition}[theorem]{Proposition}
\newtheorem{corollary}[theorem]{Corollary}
\newtheorem{theorem_app}{Theorem}[section]
\newtheorem{lemma_app}[theorem_app]{Lemma}
\theoremstyle{definition}
\newtheorem{definition}[theorem]{Definition}
\newtheorem{remark}[theorem]{Remark}

% ── Box style (matches Paper 1 / Paper 2) ──
\tcbset{
    apfbox/.style={
        enhanced,
        colback=black!3,
        frame hidden,
        borderline={0.5pt}{0pt}{black},
        arc=0pt,
        left=8pt, right=8pt, top=8pt, bottom=8pt,
        fonttitle=\bfseries\color{black},
        coltitle=black,
        detach title,
        before upper={\tcbtitle\par\smallskip},
        before skip=14pt,
        after skip=14pt,
        grow to left by=-8pt,
        grow to right by=-8pt,
        sharp corners,
        breakable,
        pad after break=4pt,
        pad before break=4pt,
        lines before break=3,
    }
}

% ── Shortcuts ──
\newcommand{\eps}{\varepsilon}
\newcommand{\End}{\mathrm{End}}
\newcommand{\SU}{\mathrm{SU}}
\newcommand{\U}{\mathrm{U}}
\newcommand{\id}{\mathrm{id}}
\DeclareMathOperator{\spn}{span}
\DeclareMathOperator{\conv}{conv}
\DeclareMathOperator{\tr}{tr}
\DeclareMathOperator{\diag}{diag}
\newcommand{\coderef}[2]{\smallskip\noindent\textit{Code reference:} \texttt{#1} in \texttt{#2}.}
\newenvironment{varlist}{\begin{itemize}[nosep,leftmargin=2em]}{\end{itemize}}

\pagestyle{fancy}
\fancyhf{}
\fancyhead[L]{\small\itshape Paper 2 Technical Supplement}
\fancyhead[R]{\small\itshape E.\,S.\ Brooke}
\fancyfoot[C]{\thepage}
\renewcommand{\headrulewidth}{0.4pt}


% ═══════════════════════════════════════════════════════════════
\begin{document}
% ═══════════════════════════════════════════════════════════════

\begin{center}
{\LARGE\bfseries Technical Supplement:\\[4pt]
The Structure of Admissible Physics}

\bigskip
{\large E.\,S.\ Brooke}\\[4pt]
{\small Companion to Paper~2: ``The Structure of Admissible Physics''}

\bigskip
{\small March 2026}
\end{center}

\bigskip


% ══════════════════════════════════════════════════════════════
% ABSTRACT
% ══════════════════════════════════════════════════════════════

\begin{abstract}
\noindent
This supplement provides self-contained proofs for every theorem stated in Paper~2 whose derivation depends on material from companion papers or requires detail beyond what the main text presents.  The goal is that a reader with access to Paper~1, its Technical Supplement, and this document can verify the structural chain from the \emph{Principle of Least Enforcement Cost} (PLEC) --- the canonical variational formulation of the enforceability-of-distinction foundation, decomposed into four structurally necessary components (A1 finite capacity, MD positive floor, A2 argmin, BW non-degeneracy) --- to the Standard Model gauge group and fermion content without consulting any other paper in the series.  The cosmological density fractions are derived conditionally: the structural partition $42 + 3 + 16 = 61$ is proved here; its identification with the Friedmann-equation density fractions is conditional on the horizon identification established in Paper~6.  The one-loop $\beta$-function sign is verified here using an imported formula; its derivation from capacity arguments is in Paper~4A.

Part~I establishes the core derivation chain: cost function uniqueness ($L_{\mathrm{Cauchy}}$), the bridge from superadditivity to gauge structure ($L_\Delta \to T_{\mathrm{alg}} \to$ gauge identification), the three carrier prerequisites that close the Theorem~R dependency on formerly external lemmas ($T_{7B}$, $L_{\mathrm{irr,uniform}}$, $B1'$), and the derivation of Lorentzian spacetime from the cost kernel: causal ordering from irreversibility, $d = 4$ from Lovelock uniqueness and DOF minimality, and the Einstein field equations ($\Delta_{\mathrm{sig}}$, $T_8$, $T_{9,\mathrm{grav}}$).  Part~II derives the gauge template and matter content: the exhaustive Lie algebra classification ($L_{\mathrm{gauge,unique}}$), the 1{,}680-template fermion scan with seven formally specified filters ($T_{\mathrm{field}}$), anomaly cancellation, and hypercharge uniqueness.  Part~III covers capacity counting and cosmology: the MECE partition, horizon equipartition ($L_{\mathrm{equip}}$), the bridge to the Friedmann equations, saturation independence, and uncertainty analysis.  Appendix~A compares the framework with existing reconstruction and ``SM from principles'' programs.

All upstream results (the four PLEC components A1, MD, A2, BW; constitutive semantics SP, K3, FD1--FD6; and the full Paper~1 derivation chain) are imported from the Paper~1 Technical Supplement and are restated in the glossary below.  Appendix~C provides complete, self-contained proofs of all ten imported results, so that no external document is needed to verify any theorem in this supplement.  No result depends on Papers~3--7; forward references to those papers appear only as scope pointers and are flagged as such.
\end{abstract}


% ══════════════════════════════════════════════════════════════
% FRONT-MATTER: WHAT IS ASSUMED, WHAT IS PROVED, WHAT IS DEFERRED
% ══════════════════════════════════════════════════════════════

\medskip
\begin{tcolorbox}[apfbox, title={What Is Assumed, What Is Proved, What Is Deferred}]

\textbf{Imported from Paper~1 Supplement (proved there; re-proved in Appendix~C for self-containment).}  The first four items are the structurally necessary components of PLEC (none derivable from the others; essentiality proofs in Paper~1 Supplement, \S1).
\begin{enumerate}[nosep,label=\arabic*.,leftmargin=2em]
  \item \textbf{A1} (PLEC upper bound): enforcement capacity is finite and positive.
  \item \textbf{MD} (PLEC positive floor): uniform positive cost per unit of enforcement complexity.
  \item \textbf{A2} (PLEC argmin selection): the realized configuration is the $\arg\min$ of total enforcement cost over the admissible set.
  \item \textbf{BW} (PLEC non-degeneracy): the cost spectrum witnesses a budget-window triple.
  \item \textbf{SP, K3, FD1--FD6}: constitutive admissibility semantics.
  \item $T_{\mathrm{form}}$: finite-dimensional substrate.
  \item $T_{\mathrm{embed}}$: canonical vector-space embedding.
  \item $T_{\mathrm{sep}}$: sector decomposition $S_\Gamma = \bigoplus_d M_d \oplus \Pi$.
  \item $L_{\eps^*}$: uniform cost floor $\eps \ge \eps^* > 0$.
  \item $L_\Delta$: superadditivity of co-located enforcement costs.
  \item $L_{\mathrm{blk}}$, $L_\Pi$: direct-sum failure and codespace rotation.
  \item $T_{\mathrm{alg}}$: finite-dimensional noncommutative $*$-algebra.
  \item $T_{\mathrm{GNS}}$: GNS Hilbert-space representation.
  \item T2a--T2c: Wedderburn--Artin decomposition, field selection $\mathbb{F} = \mathbb{C}$.
  \item $T_{\mathrm{Born}}$: Born rule.  $T_M$: interface monogamy.
  \item $L_{\mathrm{loc}}$: budget additivity for independent interfaces.
  \item $L_{\mathrm{irr}}$: irreversibility from A1 + $L_{\mathrm{nc}}$ + $L_{\mathrm{loc}}$.
  \item T3: locality $\to$ gauge structure (Skolem--Noether $+$ Doplicher--Roberts).
  \item $L_\omega$: sector-orthogonal bilinear form.
  \item $T_{\mathrm{adj}}$: self-adjoint enforcement projections.
\end{enumerate}

\smallskip\noindent\textbf{Core results proved in this supplement (self-contained, no downstream dependence).}
\begin{enumerate}[nosep,label=\arabic*.,leftmargin=2em]
  \item $L_{\mathrm{Cauchy}}$: cost function uniqueness (\S\ref{sec:cauchy_supp}).
  \item Local support and commuting algebras from disjoint substrates (Definition~\ref{def:local_support}, Lemma~\ref{lem:loc_commut}).
  \item Gauge identification theorem: internal vs.\ kinematic automorphisms; $G = \mathrm{Inn}(\mathcal{A}) = \prod_k \mathrm{PU}(n_k)$ (\S\ref{sec:gauge_ident}).
  \item Exact sequence $1 \to Z \to \mathrm{U}(\mathcal{A}) \to \mathrm{Inn}(\mathcal{A}) \to 1$ and the source exclusivity of abelian factors from $T_{\mathrm{rel}}$ (Proposition~\ref{prop:exact_seq}, Lemma~\ref{lem:abelian_source}).
  \item Unique abelian factor: structural theorem (template-independent, Proposition~\ref{prop:unique_U1}) and SM specialization (Theorem~\ref{thm:unique_U1_consolidated}).
  \item Lift from $\mathrm{PU}(n)$ to $\SU(n)$ and global structure $[\SU(3) \times \SU(2) \times \U(1)_Y]/\mathbb{Z}_6$ (Propositions~\ref{prop:PU_to_SU}, \ref{prop:global_structure}).
  \item $T_{7B}$: positive-definite cost kernel at a single interface (\S\ref{sec:T7B}).
  \item $L_{\mathrm{irr,uniform}}$: gauge-sector irreversibility at saturated interfaces (\S\ref{sec:Lirr_uniform}).
  \item $B1'$: complex carrier from oriented composite robustness (\S\ref{sec:B1prime}).
  \item $L_{\mathrm{gauge,unique}}$: gauge template uniqueness (\S\ref{sec:gauge_class_supp}).
  \item $T_{\mathrm{field}}$: fermion content by exhaustive 1{,}680-template scan with seven filters (\S\ref{sec:fermion_scan}).
  \item $L_{\mathrm{hypercharge}}$: hypercharge uniqueness from anomaly system (\S\ref{sec:hypercharge}).
  \item $L_{\mathrm{count}}$, $P_{\mathrm{exhaust}}$: capacity counting and MECE partition (\S\ref{sec:capacity_count}, \S\ref{sec:mece}).
\end{enumerate}

\smallskip\noindent\textbf{Downstream results: proved here conditional on inter-paper bridges.}

These results use the core chain above together with imported mathematical theorems (Hawking--King--McCarthy, Malament, Kolmogorov, Lovelock) whose exact hypotheses are verified in the text.  They are included to show the derivation path, but their full rigor depends on structure established in companion papers; the epistemic status of each is flagged explicitly.
\begin{enumerate}[nosep,label=\arabic*.,leftmargin=2em]
  \item $\Delta_{\mathrm{sig}}$: Lorentzian signature from causal structure (\S\ref{sec:signature}).  \textit{Status:} hypothesis verification for HKM/Malament provided; the causal partial order $\prec$ is derived from $L_{\mathrm{irr}}$; the construction of the spacetime metric from the cost kernel is presented (\S\ref{sec:cost_to_lorentz}) but the full dynamical treatment requires Paper~6.
  \item $T_8$: spacetime dimension $d = 4$ (\S\ref{sec:d4}).  \textit{Status:} four-condition proof complete; the argument uses the linearized GR DOF formula (imported) and Lovelock's theorem (hypotheses verified in Lemma~\ref{lem:lovelock_hyp}).
  \item $T_{9,\mathrm{grav}}$: Einstein field equations (\S\ref{sec:einstein}).  \textit{Status:} Lovelock hypotheses verified; the five admissibility conditions are derived; the full coupling to matter stress-energy requires Paper~6.
  \item $L_{\mathrm{equip}}$: horizon equipartition (\S\ref{sec:equip}).  \textit{Status:} the maximum-entropy argument is self-contained; the identification of the enforcement horizon with the cosmological de~Sitter horizon is established in Paper~6.  Without that bridge, the cosmological prediction $\Omega_\Lambda = 42/61$ is conditional.
  \item $L_{\mathrm{sat,part}}$: saturation independence (\S\ref{sec:sat_indep}).  \textit{Status:} self-contained, but its cosmological application requires $L_{\mathrm{equip}}$'s horizon bridge.
\end{enumerate}

\smallskip\noindent\textbf{Not proved here (honest flags).}
\begin{enumerate}[nosep,label=\arabic*.,leftmargin=2em]
  \item $L_{\mathrm{irr,uniform}}$ is a \emph{sufficient-condition} theorem: it guarantees gauge-sector irreversibility at interfaces satisfying the saturation bound~(\ref{eq:saturation_bound}).  It does not prove irreversibility at \emph{all} interfaces; the corollary shows it holds generically at near-capacity interfaces.
  \item Quantitative $\beta$-function derivation: structural necessity of asymptotic freedom is proved here; the one-loop coefficient is derived in Paper~4A ($L_{\mathrm{beta,capacity}}$).  A schematic sign calculation is included (\S\ref{sec:beta_schematic}).
  \item Cosmological horizon identification: the structural partition $42+3+16=61$ is proved here; the identification of the enforcement horizon with the cosmological de~Sitter horizon is established in Paper~6.
  \item Mass hierarchy, CKM/PMNS mixing, CP violation: require the generation structure derived in Papers~4A/4B.
  \item Confinement scale $\Lambda_{\mathrm{QCD}}$, string tension, hadronization dynamics: require quantitative RG treatment (Paper~4A).
\end{enumerate}
\end{tcolorbox}

\medskip
\begin{tcolorbox}[apfbox, title={Framing convention: ``forces'' and ``selects'' denote uniqueness, not processes}]
\footnotesize
The supplement's theorems and proofs use the verbs ``forces,'' ``selects,'' and ``rejects'' as standard mathematical shorthand for uniqueness results, $\arg\min$ properties, and inadmissibility results respectively.  Under the APF ontology, these verbs do not denote processes: nothing is forced by an agent, nothing is selected by a mechanism, nothing is rejected by a filter.  Each theorem asserts a property of the set of admissible configurations---that it has a unique member, a minimum member, or no member satisfying some condition.  For example, ``capacity minimality selects $N_c = 3$'' is shorthand for ``$N_c = 3$ is the $\arg\min$ of dim(G) over R1--R3-admissible color groups,'' and the alternative values (SU(4), SU(5), E$_6$, \ldots) are \emph{admissible} in the sense that each has cost below the capacity bound, but have higher cost than SU(3) and are therefore not what is.  The companion note \emph{What Is, Is Minimum-Cost Admissible} spells this out in full.  Readers may substitute the descriptive reading throughout the supplement without any change to the proofs.
\end{tcolorbox}

\smallskip\noindent\textit{Self-contained Standard Model datum.}\quad
Section~\ref{sec:hypercharge} derives the unique anomaly-free hypercharge
assignments for one generation of the Standard Model fermion template
purely from the tools assembled in this supplement (gauge template from
\S\ref{sec:gauge_class_supp}, fermion content from \S\ref{sec:fermion_scan},
anomaly conditions from \S\ref{sec:hypercharge}).  The result
($Y_Q = 1/6$, $Y_u = 2/3$, $Y_d = -1/3$, $Y_L = -1/2$, $Y_e = -1$)
with charge quantization as a derived consequence constitutes a
nontrivial, experimentally verified prediction obtained without any
input beyond A1 and the derivation chain.

\medskip\noindent\textbf{Glossary of imported objects.}  Each entry gives the object's name, one-sentence definition, source in the Paper~1 Supplement (P1S), and where it is first used in this supplement.

\smallskip
\begin{center}\small
\renewcommand{\arraystretch}{1.15}
\begin{tabular}{@{}lp{5.8cm}ll@{}}
\toprule
\textbf{Symbol} & \textbf{Definition} & \textbf{Source} & \textbf{Used at} \\
\midrule
$\Gamma$ & Enforcement interface (physical boundary) & P1S FD1 & Throughout \\
$S_\Gamma$ & Substrate: set of configurations at $\Gamma$ & P1S FD1 & Throughout \\
$\mathcal{D}(\Gamma)$ & Enforceable distinctions (binary predicates, $\eps > 0$) & P1S FD1 & \S\ref{sec:cauchy_supp} \\
$\eps(d)$ & Enforcement cost: $\inf_{\mathcal{P}(d)} \kappa$ & P1S FD4 & \S\ref{sec:cauchy_supp} \\
$C(\Gamma)$ & Total enforcement budget at $\Gamma$ & P1S A1 & Throughout \\
$\eps^*$ & Uniform cost floor: $\eps(d) \ge \eps^*$ for all $d$ & P1S $L_{\eps^*}$ & \S\ref{sec:cauchy_supp} \\
$M_d$ & Anchor set: minimal perturbation-support subspace & P1S FD3 & \S\ref{sec:nc_chain} \\
$\Pi$ & Pool: $S_\Gamma \ominus \bigoplus_d M_d$ & P1S $T_{\mathrm{sep}}$ & \S\ref{sec:nc_chain} \\
$E_d$ & Enforcement projection: $B$-orth.\ proj.\ onto $M_d$ & P1S FD5/$T_{\mathrm{adj}}$ & \S\ref{sec:nc_chain} \\
$B(\cdot,\cdot)$ & Sector-orthogonal bilinear form on $V$ & P1S $L_\omega$ & \S\ref{sec:T7B} \\
$\mathcal{A}$ & Enforcement algebra: $\mathrm{Alg}_\mathbb{R}(\{E_d, E_{\{d_i,d_j\}}\})$ & P1S O3 & \S\ref{sec:gauge_ident} \\
$\omega$ & Faithful positive state on $\mathcal{A}$ & P1S O4 & \S\ref{sec:gauge_ident} \\
$\mathcal{H}_\omega$ & GNS Hilbert space from $(\mathcal{A}, \omega)$ & P1S $T_{\mathrm{GNS}}$ & \S\ref{sec:gauge_ident} \\
$\Delta(d_1,d_2)$ & Superadditive surplus: $\eps(\{d_1,d_2\}) - \eps(d_1) - \eps(d_2)$ & P1S $L_\Delta$ & \S\ref{sec:nc_chain} \\
\bottomrule
\end{tabular}
\end{center}

\smallskip\noindent\textbf{APF $\leftrightarrow$ standard terminology.}

\smallskip
\begin{center}\small
\renewcommand{\arraystretch}{1.15}
\begin{tabular}{@{}ll@{}}
\toprule
\textbf{APF term} & \textbf{Standard / GPT equivalent} \\
\midrule
Distinction $d$ & Effect / binary test / observable outcome \\
Enforcement projection $E_d$ & Sharp effect / compression / filter \\
Pool $\Pi$ & Shared latent support / coupling sector \\
Superadditivity surplus $\Delta$ & Nonclassicality witness / coupling measure \\
Sector decomposition ($T_{\mathrm{sep}}$) & Face decomposition of state space \\
Enforcement algebra $\mathcal{A}$ & Observable algebra / $C^*$-algebra of observables \\
GNS representation & Gelfand--Naimark--Segal construction \\
\bottomrule
\end{tabular}
\end{center}


% ══════════════════════════════════════════════════════════════
% DEPENDENCY TABLE
% ══════════════════════════════════════════════════════════════

\medskip\noindent\textbf{Dependency table.}  Each row gives a theorem proved in this supplement, its inputs (all previously proved or imported), and the downstream results that depend on it.

\smallskip
\begin{center}\footnotesize
\renewcommand{\arraystretch}{1.15}
\begin{tabular}{@{}lp{4.5cm}p{4.5cm}@{}}
\toprule
\textbf{Theorem} & \textbf{Inputs} & \textbf{Used by} \\
\midrule
$L_{\mathrm{Cauchy}}$ (\S\ref{sec:cauchy_supp}) & A1, $L_{\mathrm{loc}}$, $L_{\eps^*}$ & $L_{\mathrm{gauge,unique}}$ (cost ranking) \\
Gauge ident.\ (\S\ref{sec:gauge_ident}) & $T_{\mathrm{alg}}$, $T_{\mathrm{sep}}$, $L_{\mathrm{loc}}$, T3 & Theorem R (gauge $=$ internal Aut) \\
$T_{7B}$ (\S\ref{sec:T7B}) & A1, K2, K3, SP, $L_\omega$, $T_{\mathrm{adj}}$ & $L_{\mathrm{irr,uniform}}$ \\
$L_{\mathrm{irr,uniform}}$ (\S\ref{sec:Lirr_uniform}) & $L_{\mathrm{loc}}$, $L_{\mathrm{nc}}$, $L_{\mathrm{irr}}$, $T_{7B}$, saturation bound & Theorem R (R2 chain) \\
$B1'$ (\S\ref{sec:B1prime}) & A1 (minimality), $T_M$, rep.\ theory & Theorem R (R1 chain) \\
$L_{\mathrm{gauge,unique}}$ (\S\ref{sec:gauge_class_supp}) & Theorem R, $L_{\mathrm{Cauchy}}$, Lie classif. & $T_{\mathrm{field}}$, $L_{\mathrm{count}}$ \\
$T_{\mathrm{field}}$ (\S\ref{sec:fermion_scan}) & $L_{\mathrm{gauge,unique}}$, AF, anomalies & $L_{\mathrm{count}}$, $P_{\mathrm{exhaust}}$ \\
$L_{\mathrm{hypercharge}}$ (\S\ref{sec:hypercharge}) & $T_{\mathrm{field}}$, anomaly conditions & $L_{\mathrm{count}}$ (charge quantization) \\
$L_{\mathrm{count}}$ (\S\ref{sec:capacity_count}) & $T_{\mathrm{field}}$, $T_{\mathrm{Higgs}}$, $L_{\mathrm{gauge,unique}}$ & $P_{\mathrm{exhaust}}$, $L_{\mathrm{equip}}$ \\
$P_{\mathrm{exhaust}}$ (\S\ref{sec:mece}) & $L_{\mathrm{count}}$, $T_3$, $T_{\mathrm{conf}}$ & $L_{\mathrm{equip}}$ \\
$L_{\mathrm{equip}}$ (\S\ref{sec:equip}) & $P_{\mathrm{exhaust}}$, $L_{\mathrm{irr}}$, $L_{\eps^*}$ & $\Omega_\Lambda$, $\Omega_m$ predictions \\
$L_{\mathrm{sat,part}}$ (\S\ref{sec:sat_indep}) & $P_{\mathrm{exhaust}}$, $L_{\mathrm{equip}}$ & Robustness of predictions \\
\bottomrule
\end{tabular}
\end{center}


\tableofcontents
\newpage


% ╔═══════════════════════════════════════════════════════════════╗
% ║                    PART I: CORE DERIVATIONS                   ║
% ╚═══════════════════════════════════════════════════════════════╝

\part{Core Derivations}
\label{part:core}


% ══════════════════════════════════════════════════════════════
\section{Cost function uniqueness ($L_{\mathrm{Cauchy}}$)}
\label{sec:cauchy_supp}
% ══════════════════════════════════════════════════════════════

The enforcement cost function $F$ maps channel count to capacity committed.
Paper~2 \S4 states the uniqueness result; this section provides the full
formal treatment, specifying domain, regularity, and operational meaning
of each condition.

% ======================================================================
% Paper 2 Technical Supplement — §1
% Cost Function Uniqueness via Cauchy's Functional Equation
% ======================================================================
%
% PURPOSE.  Paper 2 §4 states that the enforcement cost function
% F(n) = n is unique under additivity, monotonicity, and unit
% normalization.  The reviewer asks (Q2): "What are the exact hypotheses
% under which Cauchy's functional equation yields F(d) = d in your
% setting?  What is the domain of d, and how is additivity defined
% operationally?"  This section provides the full formal treatment.
%
% UPSTREAM: A1, L_loc, L_cost (from Paper 1 Supplement).
% NO downstream results in Paper 2 require extension beyond N.
% The rational extension is needed only for Paper 4A (Weinberg angle).
% ======================================================================


\subsection{Domain and operational definitions}
\label{sec:cauchy_domain}

The enforcement cost function $F$ maps the \emph{number of enforcement
channels} consumed by a distinction to the capacity committed.  We now
specify its domain, its operational meaning, and the three conditions
that determine it uniquely.

\begin{definition}[Enforcement channel count]\label{def:channel_count}
At an interface $\Gamma$ with finite capacity $C(\Gamma)$ (A1), each
irreducible distinction $d \in \mathcal{D}_{\mathrm{irr}}(\Gamma)$ costs
$\varepsilon^*$ to enforce ($L_{\varepsilon^*}$, Paper~1 Supplement).
A composite distinction occupying $n$ independent irreducible channels
has cost $\varepsilon(\{d_1, \ldots, d_n\}) = n \cdot \varepsilon^*$
when the $n$ channels are at \emph{separate interfaces} ($L_{\mathrm{loc}}$),
or $n \cdot \varepsilon^* + \Delta$ when co-located ($L_\Delta$, with
$\Delta \ge 0$).  The \emph{channel count} is the integer
$n(d) := \varepsilon(d) / \varepsilon^*$ evaluated at separate interfaces,
where the surplus $\Delta = 0$.
\end{definition}

\noindent The domain of $F$ is therefore the positive integers:
$F: \mathbb{N} \to \mathbb{R}_{>0}$, where $F(n)$ gives the capacity
cost (in units of $\varepsilon^*$) of enforcing $n$ independent channels
at separate interfaces.

\begin{remark}[Why the domain is $\mathbb{N}$, not $\mathbb{R}^+$]
\label{rem:domain_N}
Enforcement channels are discrete.  This discreteness is not assumed ---
it is derived from A1 + MD:
\begin{enumerate}[label=(\roman*),nosep]
  \item A1 provides finite capacity $C(\Gamma) < \infty$.
  \item $L_{\varepsilon^*}$ (from MD) provides a positive cost floor
        $\varepsilon^* > 0$ for each irreducible channel.
  \item Therefore the maximum number of independently enforceable channels
        at any interface is $\lfloor C(\Gamma) / \varepsilon^* \rfloor
        < \infty$, an integer.
\end{enumerate}
The domain $\mathbb{N}$ is a \emph{consequence}, not a choice.  The
classical Cauchy problem on $\mathbb{R}^+$ --- with its pathological
non-measurable solutions --- does not arise because the domain is
already discrete.
\end{remark}


\subsection{The three conditions and their derivation}
\label{sec:cauchy_conditions}

\begin{theorem}[$L_{\mathrm{Cauchy}}$: Cost Function Uniqueness]\label{thm:LCauchy}
Let $F: \mathbb{N} \to \mathbb{R}_{>0}$ satisfy:
\begin{align}
  &\textbf{C1 (Additivity):}  &&F(a + b) = F(a) + F(b)
    \quad \text{for all } a, b \in \mathbb{N}, \label{eq:C1_supp} \\
  &\textbf{C2 (Monotonicity):} &&a > b \;\Rightarrow\; F(a) > F(b)
    \quad \text{for all } a, b \in \mathbb{N}, \label{eq:C2_supp} \\
  &\textbf{C3 (Unit):}          &&F(1) = 1. \label{eq:C3_supp}
\end{align}
Then $F(n) = n$ for all $n \in \mathbb{N}$.
\end{theorem}

\begin{proof}
By C1 (induction): $F(n) = F(1 + 1 + \cdots + 1) = n \cdot F(1)$.
By C3: $F(1) = 1$.  Therefore $F(n) = n$ for all $n \in \mathbb{N}$.
C2 is satisfied automatically: $n > m \Rightarrow F(n) = n > m = F(m)$.
\end{proof}

\noindent The proof on $\mathbb{N}$ is trivial once the three conditions
are established.  The substantive content is the \emph{derivation of the
conditions from A1}:

\begin{proposition}[Derivation of C1--C3 from admissibility]\label{prop:C1C3}
\leavevmode
\begin{enumerate}[label=\textbf{C\arabic*},nosep,leftmargin=3em]
  \item \textbf{(Additivity).}  Let $d_1, d_2$ be distinctions at
  \emph{separate} interfaces $\Gamma_1, \Gamma_2$.  By $L_{\mathrm{loc}}$
  (Paper~1 Supplement~\cite{BrookeSupplement}, \S4): enforcement at separate
  interfaces is independent, with independent budgets.  The total
  enforcement cost is $\varepsilon(d_1) + \varepsilon(d_2)$ --- costs
  add exactly when the enforcement regions are disjoint.  In channel-count
  language: $F(n(d_1) + n(d_2)) = F(n(d_1)) + F(n(d_2))$.

  The operational meaning of ``additivity'' is therefore:
  \emph{independently enforceable channels at separate interfaces
  contribute independently to the total capacity commitment.}  This is
  the content of $L_{\mathrm{loc}}$ applied to the cost functional.

  \textit{What C1 does not claim:} C1 does not assert additivity at a
  \emph{shared} interface.  At a shared interface, the superadditive
  surplus $\Delta > 0$ ($L_\Delta$) produces non-additive costs.
  C1 applies only to the separated-interface regime, which is precisely
  the regime relevant to the cost function $F$ (Definition~\ref{def:channel_count}).

  \item \textbf{(Monotonicity).}  More enforcement channels require
  more capacity.  This is immediate from A1: each channel costs
  $\varepsilon^* > 0$ ($L_{\varepsilon^*}$), so adding a channel
  increases the total cost.

  \item \textbf{(Unit).}  $F(1) = 1$ is a choice of capacity units:
  we measure cost in units of the single-channel cost $\varepsilon^*$.
  This is a normalization convention, analogous to $c = 1$ in natural
  units.  It introduces no physical content.
\end{enumerate}
\end{proposition}

\noindent\textit{Connection to resource theories and GPTs.}\quad
The conditions C1--C3 have close analogues in resource-theoretic
frameworks: C1 corresponds to the tensor-product structure of
independent resources; C2 to monotonicity under free operations; C3
to the choice of a gold standard (unit resource).  The operational
stance --- deriving algebraic structure from resource constraints ---
is shared with GPT-style reconstructions (Hardy~\cite{Hardy2001},
CDP~\cite{CDP2011}) and the ``accessible fragments'' approach.  The
difference is scope: those programs reconstruct the state space; this
program uses the same starting point to derive gauge content and
matter spectrum.


\subsection{Rational extension (for Paper~4A)}
\label{sec:cauchy_rational}

Paper~2's results ($L_{\mathrm{gauge,unique}}$, $T_{\mathrm{field}}$,
$L_{\mathrm{count}}$, $P_{\mathrm{exhaust}}$, $L_{\mathrm{equip}}$) use
$F$ only on $\mathbb{N}$.  The Weinberg angle derivation (Paper~4A)
evaluates $F$ at the rational argument $1/d$ (the reciprocal channel
count), producing $\gamma = F(d) + F(1/d) = d + 1/d$.  We record the
rational extension here for completeness.

\begin{proposition}[Rational extension of $F$, with operational premise]\label{prop:rational_ext}
\textbf{Premise} (\textbf{Sharing-admissibility, S}): A single
enforcement channel evaluated at fractional occupancy $1/q$ ($q \in \mathbb{N}$)
contributes $F(1/q)$ to each of $q$ separated-interface sectors that
jointly share the channel, and the $q$ shared contributions combine
additively under C1.  Equivalently, $F$ is extended to $\mathbb{Q}_{>0}$
as the unique additive extension consistent with C1 applied $q$ times
across sharing.

\smallskip\noindent\textbf{Conclusion.}  Under premise S, $F(p/q) = p/q$
for all $p, q \in \mathbb{N}$.
\end{proposition}

\begin{proof}
For any $q \in \mathbb{N}$: by S, $F$ is defined at $p/q$; by C1
applied $q$ times, $q \cdot F(p/q) = F(p/q + \cdots + p/q) = F(p) = p$.
Therefore $F(p/q) = p/q$.
\end{proof}

\smallskip\noindent\textit{Logical status of the sharing premise.}\quad
Premise S is not derivable from C1 alone: C1 as stated in the integer
domain ($a, b \in \mathbb{N}$) is vacuous at fractional arguments, so
``C1 extended to $\mathbb{Q}_{>0}$'' is definitionally a new domain
assignment.  S supplies the operational content: it interprets $F(1/q)$
as the per-sector share of a channel decomposed across $q$ separated
interfaces.  The interpretation is consistent with $L_{\mathrm{loc}}$
(locality of capacity) and $T_M$ (interface monogamy)---sharing at
separated interfaces does not incur superadditive surplus---but it is
a reading beyond C1's native domain.  The honest flag is that the
Weinberg-angle derivation in Paper~4A rests on C1~$+$~S, not on C1
alone.  S is therefore listed alongside MD, BW, and the constitutive
admissibility semantics (SP, K3, FD1--FD6) in the assumption inventory
when Paper~4A uses the rational extension.

\noindent\textit{Operational justification for the rational extension.}\quad
The rational argument $1/d$ arises when the cost functional is evaluated
on a \emph{fractional} channel occupancy: a single enforcement channel
shared among $d$ sectors contributes $1/d$ of a full channel to each
sector.  This is the capacity-accounting interpretation of the Weinberg
angle formula's ``inverse channel count'' term.  The extension from
$\mathbb{N}$ to $\mathbb{Q}_{>0}$ requires no additional regularity
assumption (no measurability, no continuity): it follows from C1 applied
$q$ times, as the proof above shows.

\smallskip\noindent\textit{Exact status of the rational extension in the axiomatics.}\quad
The rational extension is a \emph{theorem} of C1, not an additional
axiom.  The proof ($q \cdot F(p/q) = F(p) = p$, therefore
$F(p/q) = p/q$) uses only C1 (additivity) applied $q$ times.  No
``sharing axiom'' is introduced.  The physical interpretation ---
that $F(1/d)$ describes the cost of a channel shared among $d$
separated interfaces --- is a \emph{reading} of the mathematical
result, not a premise.  Sharing is admissible only in the
separated-interface regime (the domain of $F$): at a \emph{shared}
interface, the superadditive surplus $\Delta > 0$ prevents clean
fractional decomposition, which is precisely why $F$ is defined on
separated-interface costs and the surplus is tracked separately by
$L_\Delta$.  The rational extension evaluates $F(1/d)$ at
configurations where the sharing cost is accounted for at the pool
level, not at the channel level.

\smallskip\noindent\textit{No real extension needed.}\quad The classical
Cauchy problem on $\mathbb{R}^+$ admits pathological (non-measurable,
non-monotone, additive) solutions constructible via the Axiom of Choice.
These are excluded by C2 (Darboux, 1875): any Lebesgue-measurable (or
merely bounded on an interval) additive function $f: \mathbb{R}^+ \to
\mathbb{R}$ satisfying $f(x+y) = f(x) + f(y)$ and $f(1) = 1$ is
$f(x) = x$.  However, no result in this paper or in Paper~4A requires
the domain to extend beyond $\mathbb{Q}_{>0}$.  The Darboux theorem is
available as a safety net but is never load-bearing.

\coderef{check\_L\_Cauchy\_uniqueness}{L\_Cauchy\_uniqueness.py}


\subsection{What L$_{\mathrm{Cauchy}}$ eliminates}
\label{sec:cauchy_consequences}

Without $L_{\mathrm{Cauchy}}$, the cost function $F$ is a free parameter.
Any monotone function $F: \mathbb{N} \to \mathbb{R}_{>0}$ with $F(1) = 1$
would be compatible with A1.  With $L_{\mathrm{Cauchy}}$, the class
collapses to a single element:

\smallskip
\begin{center}
\small
\begin{tabular}{@{}lll@{}}
\toprule
\textbf{Candidate} & \textbf{Condition violated} & \textbf{Consequence if allowed} \\
\midrule
$F(n) = n^2$ & C1 ($9 = F(3) \ne F(1) + F(2) = 5$) & Weinberg angle shifts \\
$F(n) = \sqrt{n}$ & C1 ($\sqrt{3} \ne 1 + \sqrt{2}$) & Capacity non-additive \\
$F(n) = \log n$ & C1 ($\log 3 \ne 0 + \log 2$); C3 ($F(1) = 0$) & No unit cost \\
$F(n) = c$ (constant) & C2 ($F(2) = F(1)$); C1 ($c \ne 2c$) & Cost insensitive \\
$F(n) = n^{3/2}$ & C1 ($3\sqrt{3} \ne 1 + 2\sqrt{2}$) & Weinberg angle shifts \\
\bottomrule
\end{tabular}
\end{center}

\smallskip\noindent Each alternative is excluded by C1 alone (additivity).
The uniqueness is not a consequence of fine-tuning or accidental
cancellation: it is a theorem of algebra (the only additive function on
$\mathbb{N}$ with $F(1) = 1$ is the identity).  C2 is logically
redundant on $\mathbb{N}$ once C1 and C3 are established, but it is
stated explicitly because it is the physically motivated condition
(more channels cost more) and because it becomes non-redundant on
$\mathbb{R}^+$ where it excludes pathological solutions.


\newpage


% ══════════════════════════════════════════════════════════════
\section{From superadditivity to gauge structure}
\label{sec:nc_chain}
% ══════════════════════════════════════════════════════════════

This section restates the derivation chain $L_\Delta \to L_{\mathrm{blk}}
\to L_\Pi \to T_{\mathrm{alg}}$ (proved in the Paper~1 Supplement) in
condensed form, and then proves the gauge identification theorem that
separates internal (gauge) automorphisms from kinematic (diffeomorphism)
automorphisms.

The chain is summarized first; the proofs are in the Paper~1 Supplement
and are not reproduced.  The gauge identification theorem
(\S\ref{sec:gauge_ident}) is new to this supplement.


\subsection{The non-commutativity chain (summary)}
\label{sec:nc_summary}

The chain from superadditivity to non-commutativity has four links,
each proved in the Paper~1 Supplement (P1S).  We restate the key result
of each without re-proof.

\smallskip\noindent\textbf{Link 1: $L_\Delta$ (Superadditivity).}
For co-located distinctions $d_1, d_2$ at an interface $\Gamma$ with
non-trivial pool $\Pi \ne \{0\}$:
$\Delta(d_1, d_2) := \eps(\{d_1, d_2\}) - \eps(d_1) - \eps(d_2) > 0$.
\textit{Source:} P1S Theorem~5.1.  \textit{Inputs:} A1, K2, K3,
$T_{\mathrm{sep}}$, SP.

\smallskip\noindent\textbf{Link 2: $L_{\mathrm{blk}}$ (Direct-Sum Failure).}
If $\Delta > 0$, the minimum-cost joint defense cannot be a direct sum
of individual defenses: $E_{\{d_1,d_2\}} \ne E_{d_1} + E_{d_2}$.
\textit{Source:} P1S Lemma~5.3.

\smallskip\noindent\textbf{Link 3: $L_\Pi$ (Codespace Rotation).}
The minimum-cost joint defender $\pi_{W_*}$ has a codespace $W_*$ rotated
into the pool $\Pi$ by angle $\theta \ne 0$.  The off-diagonal component
$F_\Pi := \pi_{W_*} - E_{d_1} - E_{d_2} \ne 0$ is a new generator
supported on the pool.
\textit{Source:} P1S Proposition~5.5.

\smallskip\noindent\textbf{Link 4: $T_{\mathrm{alg}}$ (Non-Commutativity).}
The commutator $[E_{d_1}, F_\Pi] \ne 0$ witnesses non-commutativity.
The algebra $\mathcal{A} = \mathrm{Alg}_\mathbb{R}(\{E_d\} \cup
\{E_{\{d_i,d_j\}}\})$ admits no faithful commutative
$*$-representation.
\textit{Source:} P1S Theorem~5.7.  \textit{Inputs:} $L_\Delta$,
$L_{\mathrm{blk}}$, $L_\Pi$.

\smallskip\noindent\textit{The chain is pre-Hilbert:} no GNS
construction, no spectral theory, no Hilbert space is invoked.
The non-commutativity is an algebraic fact about $\End(V)$,
established before $\mathcal{H}_\omega$ is constructed.


\subsection{Gauge identification theorem}
\label{sec:gauge_ident}

The central question is: how is the gauge group identified with the
automorphism group of the enforcement algebra, and how are
spatial/kinematic automorphisms (e.g., diffeomorphisms) excluded?

The answer rests on a structural separation: the enforcement algebra
$\mathcal{A}$ acts on the \emph{internal} sector decomposition at a
single interface $\Gamma$, which is finite-dimensional and defined by
$T_{\mathrm{sep}}$.  Diffeomorphisms are \emph{inter-interface} maps
($L_{\mathrm{loc}}$) and therefore commute with the internal algebra
by construction.

\begin{definition}[Local support of enforcement operations]\label{def:local_support}
An enforcement operation $A$ at interface $\Gamma$ has \emph{local
support} in $S_\Gamma$ if:
\begin{enumerate}[label=(LS\arabic*),nosep]
  \item $A$ maps $S_\Gamma$ to itself: $A(S_\Gamma) \subseteq S_\Gamma$.
  \item $A$ acts as the identity on the $B$-orthogonal complement:
  $A|_{S_\Gamma^\perp} = \mathrm{id}$.
\end{enumerate}
\end{definition}

\noindent\textit{Derivation of local support from FD3.}\quad
FD3 (anchor-set locality) states that the anchor set $M_d$ of any
distinction $d \in \mathcal{D}(\Gamma)$ lies within $S_\Gamma$.
The enforcement projection $E_d$ has $\mathrm{im}(E_d) = M_d
\subseteq S_\Gamma$ and $\ker(E_d) \supseteq S_\Gamma^\perp$ (by
$T_{\mathrm{adj}}$: $E_d$ is the $B$-orthogonal projection onto $M_d$,
hence $E_d|_{M_d^\perp} = 0$, and $S_\Gamma^\perp \subseteq M_d^\perp$).
Therefore $E_d(v) = v$ for $v \in M_d$ and $E_d(v) = 0$ for
$v \in S_\Gamma^\perp$.  The enforcement algebra
$\mathcal{A}(\Gamma) = \mathrm{Alg}(\{E_d : d \in \mathcal{D}(\Gamma)\})$
is generated by operators satisfying (LS1)--(LS2), and products and sums
of such operators also satisfy (LS1)--(LS2).  Therefore every
$A \in \mathcal{A}(\Gamma)$ has local support in $S_\Gamma$.

\begin{lemma}[Commuting algebras from disjoint support]
\label{lem:loc_commut}
Let $\Gamma_1, \Gamma_2$ be interfaces with $B$-orthogonal substrates:
$B(u, v) = 0$ for all $u \in S_{\Gamma_1}$, $v \in S_{\Gamma_2}$
(from $L_\omega$ applied to disjoint interfaces via $L_{\mathrm{loc}}$).
If every $A \in \mathcal{A}(\Gamma_1)$ has local support in
$S_{\Gamma_1}$ and every $B \in \mathcal{A}(\Gamma_2)$ has local
support in $S_{\Gamma_2}$, then $[A, B] = 0$ for all
$A \in \mathcal{A}(\Gamma_1)$, $B \in \mathcal{A}(\Gamma_2)$.
\end{lemma}

\begin{proof}
Decompose $V = S_{\Gamma_1} \oplus S_{\Gamma_2} \oplus W$, where
$W = (S_{\Gamma_1} \oplus S_{\Gamma_2})^\perp$ (exists because $V$ is
finite-dimensional and $B$ is positive-definite).

By local support (LS1--LS2):
\begin{align*}
  A &= A_1 \oplus \mathrm{id}_{S_{\Gamma_2}} \oplus \mathrm{id}_W,\\
  B &= \mathrm{id}_{S_{\Gamma_1}} \oplus B_2 \oplus \mathrm{id}_W,
\end{align*}
where $A_1 := A|_{S_{\Gamma_1}} \in \mathrm{End}(S_{\Gamma_1})$
and $B_2 := B|_{S_{\Gamma_2}} \in \mathrm{End}(S_{\Gamma_2})$.

For any $v = v_1 + v_2 + w$ with $v_i \in S_{\Gamma_i}$ and $w \in W$:
\begin{align*}
  AB(v) &= A(v_1 + B_2 v_2 + w) = A_1 v_1 + B_2 v_2 + w,\\
  BA(v) &= B(A_1 v_1 + v_2 + w) = A_1 v_1 + B_2 v_2 + w.
\end{align*}
Therefore $AB = BA$.

The second line uses: $B_2 v_2 \in S_{\Gamma_2}$ (LS1 for $B$),
so $A$ acts on it as $\mathrm{id}$ (LS2 for $A$).  Similarly,
$A_1 v_1 \in S_{\Gamma_1}$ (LS1 for $A$), so $B$ acts on it as
$\mathrm{id}$ (LS2 for $B$).
\end{proof}

\begin{remark}[Two bilinear structures: $B$ and $g_{\mu\nu}$]%
\label{rem:two_bilinear}
The bilinear form $B$ used in Definition~\ref{def:local_support} and
Lemma~\ref{lem:loc_commut} is a \emph{positive-definite} inner product
on the sector space $S_\Gamma$ at a single interface.  Later in this
supplement (\S\ref{sec:cost_to_lorentz}), a \emph{Lorentzian} metric
$g_{\mu\nu}$ is constructed on the spacetime manifold $M$.  These are
two distinct bilinear structures at two distinct structural levels;
they must not be conflated.

\smallskip\noindent\textit{Where they live:}\quad
$B$ is defined on the \emph{internal} sector space
$S_\Gamma = \bigoplus_d M_d \oplus \Pi$ at a single interface
$\Gamma$.  Its domain is the finite-dimensional vector space of
perturbation directions within the enforcement algebra at $\Gamma$.
The spacetime metric $g_{\mu\nu}$ is defined on the \emph{tangent
space} $T_p M$ of the spacetime manifold $M$ at each point
$p \in M$.  Its domain includes both spatial directions (from the
inter-interface structure, $L_{\mathrm{loc}}$) and the timelike
direction (from irreversibility, $L_{\mathrm{irr}}$).

\smallskip\noindent\textit{Why $B$ is positive-definite:}\quad
$B$ measures enforcement cost at a single interface.  By K2 (positive
perturbation cost) and SP (no null directions), every nonzero
perturbation of the enforcement structure at $\Gamma$ costs positive
capacity: $B(v,v) > 0$ for all $v \ne 0$ in $S_\Gamma$.  This is an
\emph{internal-sector} statement: it says that enforcement is
expensive in all directions \emph{within} the interface's sector space.
There is no timelike direction in $S_\Gamma$; the internal sector
does not ``know about'' time.

\smallskip\noindent\textit{Why $g_{\mu\nu}$ is Lorentzian:}\quad
The spacetime metric $g_{\mu\nu}$ is assembled from $B$ (which
becomes the \emph{spatial} part $g_{ij}^{(\mathrm{sp})}$ of the
metric) plus the timelike direction $\tau$ from $L_{\mathrm{irr}}$
(which contributes the negative eigenvalue $-N^2$).  The
construction (Proposition~\ref{prop:spacetime_metric}, eq.\
(\ref{eq:lorentz_metric})) is:
\[
  g_{\mu\nu}\,dx^\mu dx^\nu
  = -N^2\,dt^2 + g_{ij}^{(\mathrm{sp})}\,dx^i dx^j,
\]
where $g_{ij}^{(\mathrm{sp})}$ inherits the positive-definiteness of
$B$ and the $-N^2 dt^2$ term is the timelike contribution.  The
positive-definiteness of $B$ is therefore \emph{exactly what is
needed} for the spatial part of the Lorentzian metric: a
positive-definite spatial metric combined with a negative timelike
component yields signature $(-,+,+,+)$.

\smallskip\noindent\textit{How they relate:}\quad
\[
  \underbrace{B}_{
    \substack{\text{positive-definite}\\
              \text{on } S_\Gamma\\
              \text{(internal sector)}}}
  \;\xrightarrow{\;\text{spatial restriction}\;}
  \underbrace{g_{ij}^{(\mathrm{sp})}}_{
    \substack{\text{positive-definite}\\
              \text{on spatial slices}\\
              \text{of } T_pM}}
  \;\xrightarrow{\;+\; L_{\mathrm{irr}}\;}
  \underbrace{g_{\mu\nu}}_{
    \substack{\text{Lorentzian}\\
              \text{on } T_pM\\
              \text{(spacetime)}}}
\]
The first arrow is the identification of the enforcement cost kernel
with the spatial metric (T7B: $\varepsilon_{\mathrm{ext}}(v) = B(v,v)$,
which becomes the spatial line element).  The second arrow adds the
timelike direction from irreversibility.  No bilinear structure changes
sign or loses positive-definiteness at any step; the Lorentzian
signature arises from \emph{combining two distinct structures}
(spatial $B$ + temporal $L_{\mathrm{irr}}$), not from modifying $B$
itself.

\smallskip\noindent\textit{Why local support and commutativity use $B$,
not $g_{\mu\nu}$:}\quad
Definition~\ref{def:local_support} and Lemma~\ref{lem:loc_commut}
operate at the single-interface level: they concern the internal sector
decomposition at $\Gamma$, not the spacetime geometry.  The relevant
orthogonality is \emph{sector orthogonality} ($B(v_d, v_{d'}) = 0$
for $d \ne d'$, from $L_\omega$), not \emph{spacetime orthogonality}.
Two enforcement operations at disjoint interfaces commute because their
sector spaces are $B$-orthogonal, not because they are spacelike-separated
in the Lorentzian metric.  The spacetime interpretation (spacelike
separation $\Leftrightarrow$ commutativity) is a \emph{consequence}
derived in \S\ref{sec:lorentz}, not a premise used here.

In summary: $B$ is the cost kernel (positive-definite, internal,
single-interface); $g_{\mu\nu}$ is the spacetime metric (Lorentzian,
external, inter-interface).  They share a spatial component
($g_{ij}^{(\mathrm{sp})} = B|_{\mathrm{spatial}}$) but are defined on
different spaces, serve different purposes, and have different
signatures.  No conflation arises because the timelike direction
(the source of the signature difference) is not part of the
single-interface sector space where $B$ operates.
\end{remark}

\noindent\textit{Relationship to AQFT.}\quad
In AQFT, Einstein causality ($[\mathcal{A}(\mathcal{O}_1),
\mathcal{A}(\mathcal{O}_2)] = 0$ for spacelike-separated regions) is
a separate axiom.  In the APF, it is a theorem: FD3 (anchor-set
locality) forces local support (Definition~\ref{def:local_support}),
which forces commutativity (Lemma~\ref{lem:loc_commut}).  FD3 is the
APF's analogue of Einstein causality, but it is a constitutive
definition of what ``enforcement at $\Gamma$'' means, not an
independent physical postulate.

\smallskip\noindent\textit{Scope of this result.}\quad
Lemma~\ref{lem:loc_commut} is a theorem \emph{internal to the
finite-capacity interface algebra}.  It is not claimed to reproduce
full AQFT locality in the continuum/type-III regime.  Specifically:
it applies to enforcement algebras $\mathcal{A}(\Gamma) =
M_n(\mathbb{C})$ at finite $n$, using the positive-definite cost
kernel $B$ for the orthogonal decomposition.  The continuum-limit
analogue (commutativity of type~III local algebras for
spacelike-separated regions) is a separate result that requires the
full machinery of AQFT (Haag--Kastler axioms, nuclearity estimates,
split property).  The finite-capacity result is the \emph{precursor}
from which the continuum result can be recovered in the
$C / \varepsilon^* \to \infty$ limit
(\S\ref{sec:finite_vs_typeIII}), but the two should not be
conflated.

\begin{lemma}[Internal--kinematic decomposition]\label{lem:aut_decomp}
Let $\{\mathcal{A}(\Gamma)\}_{\Gamma \in \mathcal{I}}$ be a net of
finite-dimensional $*$-algebras indexed by a set $\mathcal{I}$, with
isotony $\mathcal{A}(\Gamma) \hookrightarrow \mathcal{A}(\Gamma')$ for
$\Gamma \subseteq \Gamma'$, and with independence: $\mathcal{A}(\Gamma)$
and $\mathcal{A}(\Gamma')$ commute when $\Gamma \cap \Gamma' = \emptyset$
(Lemma~\ref{lem:loc_commut}).  Define:
\begin{enumerate}[label=(\alph*),nosep]
  \item $\mathrm{Aut}_{\mathrm{kin}} :=$ automorphisms $\alpha$ of the
  net that act non-trivially on $\mathcal{I}$ (permuting interface
  labels);
  \item $\mathrm{Aut}_{\mathrm{int}} :=$ automorphisms that fix
  $\mathcal{I}$ pointwise (mapping each $\mathcal{A}(\Gamma)$ to itself).
\end{enumerate}
Then every net automorphism $\alpha$ decomposes uniquely as
$\alpha = \alpha_{\mathrm{int}} \circ \alpha_{\mathrm{kin}}$, with
$\alpha_{\mathrm{kin}} \in \mathrm{Aut}_{\mathrm{kin}}$ determined by
$\alpha$'s action on $\mathcal{I}$ and
$\alpha_{\mathrm{int}} \in \mathrm{Aut}_{\mathrm{int}}$ acting fibrewise.
The group structure is a semidirect product:
$\mathrm{Aut}(\text{net}) = \mathrm{Aut}_{\mathrm{int}} \rtimes
\mathrm{Aut}_{\mathrm{kin}}$.
\end{lemma}

\begin{proof}
Given $\alpha$, define $\alpha_{\mathrm{kin}}$ by its action on the
index set: $\alpha_{\mathrm{kin}}(\Gamma) := \alpha(\Gamma)$ with the
induced isomorphism
$\mathcal{A}(\Gamma) \xrightarrow{\sim} \mathcal{A}(\alpha(\Gamma))$.
Then $\alpha_{\mathrm{int}} := \alpha \circ \alpha_{\mathrm{kin}}^{-1}$
fixes every index and is therefore fibrewise.  Uniqueness follows because
the action on $\mathcal{I}$ determines $\alpha_{\mathrm{kin}}$
uniquely.  The semidirect product structure holds because
$\mathrm{Aut}_{\mathrm{int}}$ is normal (conjugating a fibrewise
automorphism by a permutation yields another fibrewise automorphism).
\end{proof}

\begin{theorem}[Gauge Identification]\label{thm:gauge_ident}
Let $\mathcal{A}$ be the enforcement algebra at interface $\Gamma$
(from $T_{\mathrm{alg}}$), and let $\mathrm{Aut}(\mathcal{A})$ denote
its $*$-automorphism group.  Then:
\begin{enumerate}[label=(\roman*),nosep]
  \item $\mathrm{Aut}(\mathcal{A})$ decomposes as
        $\mathrm{Aut}_{\mathrm{int}}(\mathcal{A}) \rtimes
        \mathrm{Aut}_{\mathrm{kin}}(\mathcal{A})$,
        where $\mathrm{Aut}_{\mathrm{int}}$ consists of automorphisms
        that preserve the interface label $\Gamma$ and
        $\mathrm{Aut}_{\mathrm{kin}}$ consists of automorphisms
        induced by inter-interface maps.
  \item $\mathrm{Aut}_{\mathrm{int}}(\mathcal{A})$ is the inner
        automorphism group $\mathrm{Inn}(\mathcal{A})$.  By
        Skolem--Noether (applicable because $\mathcal{A}$ is a
        finite-dimensional central simple algebra over $\mathbb{C}$
        after T2c), $\mathrm{Inn}(\mathcal{A}) \cong
        \mathrm{PU}(\mathcal{H}_\omega)$.
  \item The physical gauge group $G$ is identified with
        $\mathrm{Inn}(\mathcal{A})$.  Kinematic automorphisms
        (diffeomorphisms) are outer and are excluded from $G$.
  \item $G$ is non-abelian: $\mathrm{Inn}(\mathcal{A})$ is non-trivial
        and non-commutative because $\mathcal{A}$ is non-commutative
        ($T_{\mathrm{alg}}$) and simple (T2a).
\end{enumerate}
\end{theorem}

\begin{proof}

\textit{Step 1: Internal vs.\ kinematic automorphisms.}\quad
By Lemma~\ref{lem:aut_decomp} applied to the enforcement net
$\{\mathcal{A}(\Gamma)\}_{\Gamma \in \mathcal{I}}$ (with independence
from $L_{\mathrm{loc}}$), every automorphism decomposes uniquely into a
kinematic part (acting on the index set $\mathcal{I}$) and an internal
part (acting fibrewise on each $\mathcal{A}(\Gamma)$).
Diffeomorphisms are kinematic; the gauge group is the internal part.
This gives (i).

\textit{Step 2: Wedderburn--Artin and Skolem--Noether.}\quad
After field selection (T2c: $\mathbb{F} = \mathbb{C}$), the enforcement
algebra is a finite-dimensional $*$-algebra over $\mathbb{C}$.
Wedderburn--Artin (T2a) gives
$\mathcal{A} \cong \bigoplus_{k=1}^{K} M_{n_k}(\mathbb{C})$.
By Skolem--Noether, every automorphism of $M_{n_k}(\mathbb{C})$ is inner:
$\alpha(a) = UaU^*$ for some unitary $U$.  The inner automorphism group
of each block is $\mathrm{Inn}(M_{n_k}) \cong \mathrm{PU}(n_k) =
\mathrm{U}(n_k) / \mathrm{U}(1)$.  For the full semisimple algebra:
$\mathrm{Inn}(\mathcal{A}) = \prod_{k=1}^K \mathrm{PU}(n_k)$.
This gives (ii).

\textit{Step 2b: Block permutations.}\quad
For $K \ge 2$, the full automorphism group of a finite-dimensional
semisimple algebra is:
\[
  \mathrm{Aut}\Bigl(\bigoplus_{k=1}^K M_{n_k}(\mathbb{C})\Bigr)
  \;=\;
  \Bigl(\prod_{k=1}^K \mathrm{PU}(n_k)\Bigr) \;\rtimes\; S_{\mathrm{iso}},
\]
where $S_{\mathrm{iso}} \subseteq S_K$ is the group of permutations among
isomorphic blocks (blocks with the same $n_k$).  The inner automorphisms
are $\mathrm{Inn}(\mathcal{A}) = \prod_k \mathrm{PU}(n_k)$; the block
permutations $S_{\mathrm{iso}}$ are outer.

The physical status of $S_{\mathrm{iso}}$:
\begin{enumerate}[label=(\alph*),nosep]
  \item Each Wedderburn block corresponds to a distinct enforcement sector
  under $T_{\mathrm{sep}}$.  Sectors are distinguished by their anchor sets
  $M_{d_k}$, which carry distinct enforcement costs $\eps(d_k)$.
  \item A block permutation $\sigma \in S_{\mathrm{iso}}$ exchanges sectors.
  This is an automorphism of the \emph{abstract} algebra but not of the
  \emph{cost-equipped} algebra $(\mathcal{A}, \eps)$: it maps
  $\eps(d_k) \mapsto \eps(d_{\sigma(k)})$, which differs from $\eps(d_k)$
  unless the permuted sectors are cost-degenerate.
  \item Cost-degenerate sectors (identical enforcement costs at every state)
  correspond to physically indistinguishable copies --- these are
  \emph{generations}.  The permutation symmetry $S_{\mathrm{iso}}$ among
  cost-degenerate blocks is a discrete \emph{flavor} symmetry, not a
  continuous gauge symmetry.  It is broken by Yukawa couplings (mass
  differences between generations) and does not contribute to the gauge
  group $G$.
  \item Non-degenerate blocks: $\eps(d_k) \ne \eps(d_j)$ for $k \ne j$.
  The permutation $\sigma$ is not an automorphism of $(\mathcal{A}, \eps)$,
  so $S_{\mathrm{iso}}$ is trivial on the cost-equipped algebra.
\end{enumerate}

\noindent\textit{Conclusion.}\quad The continuous gauge group is
$G = \mathrm{Inn}(\mathcal{A}) = \prod_k \mathrm{PU}(n_k)$,
regardless of whether $S_{\mathrm{iso}}$ is trivial (non-degenerate
costs) or nontrivial (degenerate costs).  The reason is that
$S_{\mathrm{iso}}$ is \emph{discrete and outer}: it permutes
Wedderburn blocks but does not generate any continuous family of
automorphisms.  The continuous gauge group sees only the connected
component of $\mathrm{Aut}(\mathcal{A})$, which is
$\mathrm{Inn}(\mathcal{A})$.

\smallskip\noindent\textit{Physical status of $S_{\mathrm{iso}}$:
generation symmetry and Yukawa breaking.}\quad
Cost-degenerate blocks correspond to physically indistinguishable
copies of the same gauge sector --- these are \emph{generations}.
The exact permutation symmetry $S_{\mathrm{iso}}$ among
cost-degenerate blocks is the \emph{flavor symmetry} of the massless
limit: all three generations of (say) up-type quarks have identical
gauge quantum numbers.

After electroweak symmetry breaking (derived in Paper~4B from the
enforcement potential), the Yukawa couplings $y_i$ become nonzero and
generation-dependent: $y_u \ne y_c \ne y_t$.  In the cost-equipped
algebra $(\mathcal{A}, \varepsilon)$, this is reflected as cost
non-degeneracy: the enforcement cost of maintaining a top-quark
distinction at an interface differs from that of an up-quark
distinction.  The physical observable is the mass hierarchy:
$m_u \approx 2\,\mathrm{MeV} \ll m_c \approx 1.3\,\mathrm{GeV}
\ll m_t \approx 173\,\mathrm{GeV}$, which breaks $S_3 \to \{1\}$
completely.

The structural chain is: cost-degenerate blocks $\to$ exact flavor
symmetry $\to$ Yukawa splitting $\to$ broken $S_{\mathrm{iso}}$ $\to$
observed mass hierarchy.  The derivation of the specific Yukawa
couplings is in Paper~4B.  What is established here is that
$S_{\mathrm{iso}}$ is discrete, outer, and does not contribute to the
continuous gauge group in any regime --- degenerate or not.

\textit{Step 3: Identification.}\quad
The physical gauge group acts on the internal degrees of freedom at a
single interface, preserving the cost-equipped algebraic structure.  By
Steps~1--2, this is precisely $\mathrm{Inn}(\mathcal{A})$.  Kinematic
automorphisms are outer (they change the interface label) and are not
part of the gauge group.  Block permutations among cost-degenerate
sectors are outer, discrete, and identified with generation symmetries.
This gives (ii) and (iii).

\textit{Step 4: Non-abelian.}\quad
$\mathcal{A}$ is non-commutative ($T_{\mathrm{alg}}$).  A commutative
algebra has $\mathrm{Inn}(\mathcal{A}) = \{1\}$ (every inner
automorphism is trivial).  Since $\mathrm{Inn}(\mathcal{A})$ is
non-trivial, $G$ is non-trivial.  For $M_n(\mathbb{C})$ with $n \ge 2$,
$\mathrm{PU}(n)$ is non-abelian (e.g., $\mathrm{PU}(2) \cong
\mathrm{SO}(3)$).  This gives (iv).
\end{proof}

\begin{remark}[Emergence of abelian $\U(1)$ factors]\label{rem:U1}
The inner automorphism group $\mathrm{Inn}(\mathcal{A}) = \prod_k
\mathrm{PU}(n_k)$ contains no abelian factors: $\mathrm{PU}(n)$ is
semisimple for $n \ge 2$, and trivial for $n = 1$.  This raises the
question: where does the $\U(1)_Y$ hypercharge factor of the Standard
Model originate?

The answer involves the relative-phase torus of the multi-block unitary
group, together with a no-redundancy theorem.
\end{remark}

The algebraic relationship between the groups acting on the algebra
(by conjugation) and on the Hilbert space (linearly) is captured by
an exact sequence.

\begin{proposition}[Exact sequence for gauge representations]%
\label{prop:exact_seq}
Let $\mathcal{A} = \bigoplus_{k=1}^K M_{n_k}(\mathbb{C})$ with
$\mathrm{U}(\mathcal{A}) = \prod_k \mathrm{U}(n_k)$.  Define
the conjugation map $\pi: \mathrm{U}(\mathcal{A}) \to
\mathrm{Aut}(\mathcal{A})$ by $\pi(U)(a) = UaU^*$.  Then:
\begin{enumerate}[label=(\roman*),nosep]
  \item The image of $\pi$ is $\mathrm{Inn}(\mathcal{A}) = \prod_k
  \mathrm{PU}(n_k)$.
  \item The kernel of $\pi$ is the center
  $Z(\mathrm{U}(\mathcal{A})) = \prod_k \mathrm{U}(1)_k$,
  consisting of block-diagonal phase matrices
  $(e^{i\theta_1}\mathbf{1}_{n_1}, \ldots,
  e^{i\theta_K}\mathbf{1}_{n_K})$.
  \item The sequence
  \[
    1 \;\longrightarrow\; Z(\mathrm{U}(\mathcal{A}))
    \;\longrightarrow\; \mathrm{U}(\mathcal{A})
    \;\xrightarrow{\;\pi\;}\;
    \mathrm{Inn}(\mathcal{A}) \;\longrightarrow\; 1
  \]
  is exact.
  \item The overall diagonal phase
  $\mathrm{U}(1)_{\mathrm{diag}} = \{(e^{i\theta}\mathbf{1}_{n_1},
  \ldots, e^{i\theta}\mathbf{1}_{n_K}) : \theta \in \mathbb{R}\}$
  acts trivially on both $\mathcal{A}$ (by conjugation) and on
  $\mathcal{H}_\omega$ (as a global phase).  The quotient
  $Z / \mathrm{U}(1)_{\mathrm{diag}}$ is a torus of rank $K - 1$, the
  relative-phase torus $T_{\mathrm{rel}}$.
  \item Matter fields transform under $\mathrm{U}(\mathcal{A})$
  linearly on $\mathcal{H}_\omega$ (fundamental representation of each
  block), not under $\mathrm{Inn}(\mathcal{A})$ (which acts on
  $\mathcal{A}$ by conjugation, i.e., the adjoint representation).
  The two actions agree on the Lie algebra
  ($\mathfrak{u}(\mathcal{A}) \to
  \mathfrak{inn}(\mathcal{A})$ is surjective with kernel
  $\mathfrak{z}$), but differ globally: $\mathrm{U}(\mathcal{A})$
  admits representations that $\mathrm{Inn}(\mathcal{A})$ does not.
\end{enumerate}
\end{proposition}

\begin{proof}
(i) By Skolem--Noether, every automorphism of $M_n(\mathbb{C})$ has
the form $a \mapsto UaU^*$ for some $U \in \mathrm{U}(n)$.  For the
semisimple algebra, each block is mapped to itself by an inner
automorphism (outer block permutations are handled separately in
Step~2b of Theorem~\ref{thm:gauge_ident}), so
$\mathrm{Im}(\pi) = \prod_k \mathrm{Inn}(M_{n_k}) = \prod_k
\mathrm{PU}(n_k)$.  (ii)~$\pi(U) = \mathrm{id}$ iff
$UaU^* = a$ for all $a \in \mathcal{A}$, iff $U$ is central.
For $M_n(\mathbb{C})$, the center is $\mathrm{U}(1) \cdot
\mathbf{1}_n$; for the direct sum, $Z = \prod_k \mathrm{U}(1)_k$.
(iii)~Surjectivity of $\pi$ onto $\mathrm{Inn}(\mathcal{A})$ is (i);
the kernel statement is (ii); the sequence is exact at
$\mathrm{U}(\mathcal{A})$ by construction.  (iv)~The diagonal
$\mathrm{U}(1)_{\mathrm{diag}}$ is a subgroup of $Z$ with
$K$-dimensional $Z$ quotient of dimension $K - 1$.  Triviality on
$\mathcal{H}_\omega$ is because states are defined up to global
phase: $|\psi\rangle$ and $e^{i\theta}|\psi\rangle$ represent the
same physical state.  (v)~The GNS representation maps each
$M_{n_k}(\mathbb{C})$ block to operators on $\mathbb{C}^{n_k}
\subset \mathcal{H}_\omega$; the natural group action on
$\mathbb{C}^{n_k}$ is $v \mapsto Uv$ (fundamental), not
$a \mapsto UaU^*$ (adjoint).  The fundamental is a linear
representation of $\mathrm{U}(n_k)$ but only a projective
representation of $\mathrm{PU}(n_k)$.
\end{proof}

\begin{definition}[Admissible gauge structure]\label{def:admissible_gauge}
Let $G$ be a compact Lie group acting faithfully on a finite-dimensional
Hilbert space $\mathcal{H}$, and let $\mathcal{M} = \{r_1, \ldots, r_N\}$
be the set of physically distinct irreducible multiplet types occurring in
$\mathcal{H}$.

\smallskip\noindent\textbf{Enforcement completeness (EC).}\quad
$G$ is \emph{enforcement-complete} if the map
$r_i \mapsto [r_i]_G$ from multiplet types to $G$-representation labels
is injective: distinct physical types carry distinct gauge quantum numbers.

\smallskip\noindent\textbf{Capacity minimality (CM).}\quad
$G$ is \emph{capacity-minimal} if no generator of $G$ can be removed
while preserving enforcement completeness.  Formally: for every connected
subgroup $H \subsetneq G$ obtained by removing at least one generator,
the label map $r_i \mapsto [r_i]_H$ is \emph{not} injective.

\smallskip\noindent\textbf{Admissibility.}\quad
$G$ is \emph{admissible} if it is both enforcement-complete and
capacity-minimal.
\end{definition}

Before deriving EC and CM, we establish the bridge from
Lie-algebra generators to enforcement channels, which CM requires.
The bridge is not as simple as ``one generator = one channel'': the
relationship between Lie-algebra generators and enforcement channels
involves incompatibility (non-commuting generators cannot be
simultaneously monitored), coarse-graining (degenerate eigenvalues
reduce the information content of a single generator), and the
distinction between Cartan (commuting) and root (non-commuting)
generators.  We treat these systematically.

\begin{definition}[Enforcement monitoring of a generator]%
\label{def:monitoring}
Let $g_a \in \mathfrak{g}$ be a Hermitian generator of a compact Lie
algebra acting on a finite-dimensional Hilbert space $\mathcal{H}$ of
dimension $n$.  The \emph{enforcement monitoring} of $g_a$ is the
projection-valued measure (PVM) $\{P_{a,\lambda}\}_\lambda$ on the
spectrum $\sigma(g_a) \subseteq \mathbb{R}$, where $P_{a,\lambda}$
projects onto the $\lambda$-eigenspace of $g_a$.  The \emph{distinction
count} of $g_a$ is $d(g_a) := |\sigma(g_a)| - 1$ --- one less than
the number of distinct eigenvalues.  This is the number of binary
questions needed to determine which eigenspace a state occupies
(using a binary search on the ordered eigenvalue list).
\end{definition}

\begin{lemma}[Primitive generator cost bound]\label{lem:gen_channel}
Let $G$ be a compact Lie group acting faithfully on a
finite-dimensional Hilbert space $\mathcal{H}$ with
$\dim(\mathcal{H}) = n$.  Let $\mathfrak{g} = \mathrm{Lie}(G)$ with
$\dim(\mathfrak{g}) = p$.  Then the total enforcement cost of
maintaining the gauge structure $G$ is at least
$p \cdot \varepsilon^*$: each independent Lie-algebra direction that
generates a nontrivial infinitesimal orbit on $\mathcal{H}$
requires at least one irreducible enforcement channel.

More precisely:
\begin{enumerate}[label=(\roman*),nosep]
  \item \textbf{(Lower bound per generator.)}  For each nonzero
  $g_a \in \mathfrak{g}$, faithfulness ensures $g_a$ acts
  non-trivially on $\mathcal{H}$ (it has at least two distinct
  eigenvalues).  Monitoring the charge under $g_a$ (determining which
  eigenspace a state occupies) requires at least one enforcement
  channel at cost $\ge \varepsilon^* > 0$ ($L_{\varepsilon^*}$).
  \item \textbf{(Independence.)}  Linearly independent generators
  $g_a, g_b$ cannot be monitored by a single channel: either they
  do not commute (incompatible spectral decompositions: no single
  measurement determines both charges) or they commute but have
  independent spectra (knowing $g_a$'s eigenvalue does not determine
  $g_b$'s).  In either case, a second channel is required.
  \item \textbf{(Cost lower bound.)}  By (i) and (ii), $p$ linearly
  independent generators require at least $p$ channels, each costing
  $\ge \varepsilon^*$.  The total cost is at least
  $p \cdot \varepsilon^*$.
\end{enumerate}

\noindent\textit{Scope and caveats.}\quad
The lemma establishes a \emph{lower bound} on the enforcement cost,
not an exact cost.  It does not claim that the gauge structure costs
\emph{exactly} $p \cdot \varepsilon^*$: co-located non-abelian
channels may incur superadditive surplus ($L_\Delta$), making the
actual cost strictly greater.  The lower bound suffices for CM: a
redundant generator wastes \emph{at least} $\varepsilon^*$ capacity
(and possibly more), which is inadmissible under finite total budget.
The lemma also does not claim that each generator corresponds to a
\emph{separate physical measurement}: non-commuting generators have
incompatible spectral decompositions and cannot be simultaneously
monitored.  What it claims is that the \emph{capacity cost}
of maintaining $p$ independent gauge directions is at least
$p \cdot \varepsilon^*$, regardless of measurement incompatibility.
\end{lemma}

\begin{proof}
\textit{Step 1: Each generator occupies exactly one channel.}\quad
A generator $g_a \in \mathfrak{g}$ is Hermitian (compact Lie algebra:
$g_a^\dagger = g_a$) and acts on the finite-dimensional $\mathcal{H}$.
By the spectral theorem for Hermitian operators on finite-dimensional
spaces, $g_a$ has a complete orthonormal eigenbasis with real
eigenvalues: $g_a = \sum_\lambda \lambda\, P_{a,\lambda}$.  The
PVM $\{P_{a,\lambda}\}$ is uniquely determined by $g_a$.

Monitoring $g_a$ means determining which eigenspace a state occupies.
This is one irreducible enforcement channel
(Definition~\ref{def:channel_count}): the PVM defines a single
spectral decomposition of $\mathcal{H}$ into orthogonal subspaces.
The monitoring is \emph{complete} for $g_a$: knowing the eigenvalue
$\lambda$ determines the charge under the entire one-parameter
subgroup $\{e^{it g_a}\}_{t \in \mathbb{R}}$ (because
$e^{it g_a}|_{\mathcal{H}_{a,\lambda}} = e^{it\lambda}\,\mathrm{id}$).
No sub-decomposition of the PVM is needed, and no additional
measurement refines the information about $g_a$ beyond what the PVM
already provides.  Therefore: one generator, one channel.

\textit{Step 2: Linearly independent generators require separate
channels.}\quad
We must show that two linearly independent generators $g_a, g_b$
cannot be simultaneously monitored by a single channel.  There are
two cases.

\textit{Case 2a: $[g_a, g_b] \ne 0$ (non-commuting generators).}\quad
If $g_a$ and $g_b$ do not commute, their PVMs are incompatible: there
is no common eigenbasis.  (Proof: suppose every eigenvector of $g_a$
were also an eigenvector of $g_b$.  Then $g_a$ and $g_b$ would be
simultaneously diagonalizable, hence commuting --- contradicting
$[g_a, g_b] \ne 0$.)  An incompatible pair of PVMs cannot be realized
by a single measurement: any single measurement determines the
eigenvalue of at most one Hermitian operator.  In the enforcement
language: a single channel that determines the $g_a$-eigenvalue
necessarily disturbs the $g_b$-eigenvalue (by the incompatibility of
the PVMs), and vice versa.  Therefore non-commuting generators
require separate channels.

\textit{Case 2b: $[g_a, g_b] = 0$ (commuting generators).}\quad
Commuting Hermitian operators are simultaneously diagonalizable:
there exists a common eigenbasis in which both $g_a$ and $g_b$ are
diagonal.  One might therefore suspect that a single channel could
monitor both.  This is incorrect: the eigenvalue of $g_a$ does not
determine the eigenvalue of $g_b$ unless $g_b$ is a function of $g_a$.

Formally: let $\mathcal{H} = \bigoplus_{(\lambda, \mu)}
\mathcal{H}_{\lambda,\mu}$ be the simultaneous eigenspace
decomposition.  Monitoring $g_a$ alone determines $\lambda$ but
reveals only the subspace $\mathcal{H}_{a,\lambda} =
\bigoplus_\mu \mathcal{H}_{\lambda,\mu}$, which may contain multiple
$g_b$-eigenvalues.  To determine $\mu$ as well, a second measurement
(a second channel monitoring $g_b$) is required.  The two measurements
are compatible (they can be performed simultaneously or in either
order without disturbance), but they are \emph{distinct}: two
separate channels, one for each generator.

The exception is when $g_b = f(g_a)$ for some function
$f: \mathbb{R} \to \mathbb{R}$ --- then $g_b$'s eigenvalue is
determined by $g_a$'s, and $g_b$ adds no new information.  But in
this case $g_b$ is not linearly independent of
$\{g_a, g_a^2, \ldots\}$ in the universal enveloping algebra, and in
the finite-dimensional setting, $g_b$ is a polynomial in $g_a$.
Since we chose a \emph{basis} of $\mathfrak{g}$, no basis element is
a function of another (linear independence in the Lie algebra is
strictly stronger than functional dependence on the Hilbert space).

\textit{Step 3: Faithfulness ensures all channels are non-trivial.}\quad
Faithfulness of the $G$-action means: $g_a \ne 0$ in $\mathfrak{g}$
implies $g_a$ acts non-trivially on $\mathcal{H}$ (i.e., $g_a$ is
not the zero operator).  Therefore $|\sigma(g_a)| \ge 2$ for every
nonzero $g_a$: the distinction count $d(g_a) \ge 1$.  Without
faithfulness, some generators might act as the zero operator, and the
corresponding channel would be vacuous.  Faithfulness is guaranteed
by the gauge identification theorem
(Theorem~\ref{thm:gauge_ident}): the gauge group $G$ acts faithfully
on $\mathcal{H}_\omega$ by construction (the inner automorphisms of
$M_n(\mathbb{C})$ act faithfully on $\mathbb{C}^n$ via the
fundamental representation, Proposition~\ref{prop:PU_to_SU}).

\textit{Step 4: Cost.}\quad
By $L_{\varepsilon^*}$ (Paper~1 Supplement), each irreducible
enforcement channel costs at least $\varepsilon^* > 0$.  By Steps
1--3, $p$ basis generators require $p$ separate channels.  Therefore
the total cost is at least $p \cdot \varepsilon^*$.
\end{proof}

\smallskip\noindent\textbf{Edge cases and subtleties.}\quad

\smallskip\noindent\textit{1.\ Degenerate eigenvalues and
coarse-graining.}\quad
If a generator $g_a$ has a degenerate eigenvalue (i.e., some
eigenspace $\mathcal{H}_{a,\lambda}$ has dimension $> 1$), the
monitoring channel provides less than $\log_2 n$ bits of information.
For example, the $\SU(2)$ generator $J_3$ on the spin-$1$
representation has eigenvalues $\{+1, 0, -1\}$: three distinct
values, so $d(J_3) = 2$.  On a spin-$0$ (trivial) representation,
$J_3 = 0$: all eigenvalues are zero, so $d(J_3) = 0$ and the
channel is vacuous.

The lemma handles this correctly: each generator occupies one channel
regardless of the degeneracy structure.  The \emph{information
content} of the channel varies (from 0 bits for a trivially acting
generator to $\log_2 n$ bits for a generator with $n$ distinct
eigenvalues), but the \emph{channel count} is one.  The
capacity cost is $\varepsilon^*$ per channel, not per bit: this is
the content of $L_{\varepsilon^*}$ (the cost floor is per irreducible
distinction, not per bit of information).

\smallskip\noindent\textit{2.\ Abelian groups.}\quad
For an abelian $G$ (e.g., $\mathrm{U}(1)^k$), all generators commute:
$[g_a, g_b] = 0$ for all $a, b$.  Case~2b of the proof applies to
every pair.  The generators are simultaneously diagonalizable, but
each occupies a separate channel (one PVM per generator).  The total
channel count equals $\dim(\mathfrak{g}) = k$.  This is consistent:
a product of $k$ independent $\mathrm{U}(1)$'s has $k$ independent
charges, each requiring a separate measurement.

\smallskip\noindent\textit{3.\ Cartan subalgebra vs.\ root generators.}\quad
In a semisimple Lie algebra $\mathfrak{g}$, the Cartan subalgebra
$\mathfrak{h}$ (maximal commuting subalgebra, dimension $= \mathrm{rank}$)
and the root generators $E_\alpha$ play different roles.  The Cartan
generators are simultaneously diagonalizable (Case~2b); the root
generators do not commute with the Cartan generators (Case~2a) or
with each other.  The lemma assigns one channel to each --- Cartan and
root alike --- because each is a linearly independent element of the
basis.  The distinction between ``commuting'' and ``non-commuting'' does
not affect the channel count; it affects only whether the channels
can be monitored simultaneously (Cartan: yes; root: no).

\smallskip\noindent\textit{4.\ Representations with coincidental
spectral degeneracies.}\quad
It is possible for two linearly independent generators $g_a, g_b$ to
have identical spectra on a particular representation: $\sigma(g_a) =
\sigma(g_b)$ (as multisets).  For example, in the fundamental
representation of $\SU(2)$, $J_1$ and $J_2$ both have spectrum
$\{+\tfrac{1}{2}, -\tfrac{1}{2}\}$.  Despite the identical spectra,
the \emph{eigenbases} are different ($J_1$ and $J_2$ do not commute),
so the corresponding PVMs are incompatible and require separate
channels.  The channel count depends on linear independence in the Lie
algebra, not on the coincidence of spectral values.

\smallskip\noindent\textit{5.\ Basis dependence.}\quad
The claim ``$p$ channels'' depends on the choice of basis.  Different
bases of $\mathfrak{g}$ yield different PVMs, but the total channel
count $p = \dim(\mathfrak{g})$ is basis-independent (it is a
topological invariant of the Lie group).  The enforcement cost
$p \cdot \varepsilon^*$ is therefore basis-independent.  The
\emph{specific channels monitored} depend on the basis choice (e.g.,
Cartan--Weyl basis vs.\ Gell-Mann matrices), but the total number
does not.

\smallskip\noindent\textit{6.\ Could coarse-graining reduce the
channel count?}\quad
A coarse-graining of the enforcement structure (combining multiple
fine-grained distinctions into a single coarser distinction) could
in principle reduce the number of channels below $p$.  However,
coarse-graining loses information: it cannot separate states that
the fine-grained structure separates.  If EC requires that all
multiplet types be separated, the enforcement structure must be at
least as fine as the multiplet partition.  The fine-grained structure
has $p$ independent channels (one per generator); any coarse-graining
that reduces the channel count below $p$ either (a)~merges two
generators' channels into one, losing the ability to distinguish
states separated by those generators, or (b)~drops a generator
entirely, reducing $G$ to a proper subgroup.  Case~(a) violates EC
if the merged generators separate different multiplets; case~(b)
violates EC if the dropped generator separates at least one pair.
Therefore: no coarse-graining consistent with EC can reduce the
channel count below $p$.

\begin{lemma}[EC and CM from the admissibility semantics]\label{lem:EC_CM_from_A1}
EC and CM are derived principles of the admissibility framework.
They are not consequences of A1 (finite capacity) alone; their
dependencies on the two axioms (A1, A2) and the regularity input MD
are stated cleanly under the dual-axiom framing of Paper~1, \S2:
\begin{enumerate}[label=(\roman*),nosep]
  \item \textbf{EC} requires: A1 (finite capacity) $+$
  $T_{\mathrm{sep}}$ (sector decomposition) $+$
  $L_{\mathrm{loc}}$ (locality) $+$ FD3 (anchor-set locality) $+$
  the gauge identification (Theorem~\ref{thm:gauge_ident}).
  EC is the statement that the internal gauge labels at a single
  interface must resolve all physically distinct internal types.
  It is a \emph{completeness principle for the interface-local
  enforcement semantics}, downstream of A1 alone (via the derived
  structural separation); A2 plays no role.
  \item \textbf{CM} requires: A2 (minimum cost) applied to the
  admissible set defined by A1 $+$ MD (via $L_{\varepsilon^*}$) $+$
  the generator--channel bridge (Lemma~\ref{lem:gen_channel}).
  A configuration carrying a redundant generator is \emph{not}
  inadmissible under A1 --- it still satisfies $\Sigma\varepsilon \le C$.
  But A1 $+$ MD (via $L_{\varepsilon^*}$) guarantees that removing the
  redundant generator yields a strictly cheaper configuration in the
  same admissible set, by $\varepsilon^* > 0$.  A2 then selects the
  cheaper one as the realised gauge structure.  CM is therefore the
  statement that the realised admissible structure carries no
  redundant generators --- a direct consequence of A2.
\end{enumerate}
\end{lemma}

\begin{proof}
\textit{EC from the structural separation.}\quad
The argument for EC requires a premise that is not A1 alone but is
\emph{derived from} A1: the internal--external structural separation.
By $T_{\mathrm{sep}}$ (Paper~1 Supplement), the enforcement structure
at each interface $\Gamma$ decomposes into sectors.  By
Definition~\ref{def:local_support} and Lemma~\ref{lem:loc_commut}
(both derived from FD3 and $L_{\mathrm{loc}}$, which are themselves
derived from A1): at a single interface, the enforcement algebra
$\mathcal{A}(\Gamma)$ acts on the \emph{internal} sector decomposition
of $S_\Gamma$, while inter-interface (spacetime) distinctions
--- spatial separation, momentum, mass --- are properties of the
\emph{external} structure (the inter-interface map
$L_{\mathrm{loc}}$).  At a single interface, the only enforcement
operations available are those in $\mathcal{A}(\Gamma)$, whose
automorphism group is the gauge group $G$
(Theorem~\ref{thm:gauge_ident}).

Now suppose two physically distinct multiplet types $r_i \ne r_j$
carry identical gauge quantum numbers: $[r_i]_G = [r_j]_G$.  Since
$G = \mathrm{Aut}(\mathcal{A}(\Gamma))$ and the $G$-representation
labels exhaust the distinctions available within
$\mathcal{A}(\Gamma)$, no enforcement operation at $\Gamma$
distinguishes $r_i$ from $r_j$.  But A1 requires that all physically
distinct states be enforceable (this is the operational content of
``enforcement capacity'': the budget $C(\Gamma)$ is the capacity to
enforce distinctions at $\Gamma$).  If two types are indistinguishable
at every interface at which they co-occur, they are enforcement-equivalent
and therefore physically identical --- contradicting $r_i \ne r_j$.

The objection ``perhaps they are distinguished by mass or momentum
rather than gauge quantum numbers'' is addressed by the structural
separation: mass and momentum are spacetime (external) properties
resolved by the inter-interface structure, not by the enforcement
algebra at a single interface.  EC is the statement that the
\emph{internal} sector must resolve all internal types; it does not
claim that gauge quantum numbers are the only physically meaningful
labels, only that they must be complete within the internal sector.

\textit{CM from A2 over the admissible set.}\quad
CM requires three ingredients in addition to A2:
\begin{enumerate}[label=(\alph*),nosep]
  \item $L_{\varepsilon^*}$: each irreducible distinction costs at
  least $\varepsilon^* > 0$.  This is derived from A1 $+$ MD
  (enforcement isotropy) in the Paper~1 Supplement.  It is
  \emph{not} derivable from A1 alone: without the isotropy assumption
  MD, one cannot exclude the possibility that some channels have
  arbitrarily small cost.
  \item Lemma~\ref{lem:gen_channel}: each Lie-algebra generator
  occupies exactly one channel.  This is derived from faithfulness of
  the $G$-action (which follows from Theorem~\ref{thm:gauge_ident},
  itself derived from A1) and the spectral decomposition of Hermitian
  operators on finite-dimensional Hilbert space (which follows from
  $T_{\mathrm{form}}$, derived from A1).
  \item A1 $+$ MD jointly define the admissible set
  $\mathcal{A}(\rho, R)$ of configurations satisfying
  $\Sigma\varepsilon \le C$ with $\varepsilon(d) \ge \mu^* n(d)$.
\end{enumerate}
Given (a)--(c): consider any admissible configuration $G$ that
contains a generator $X$ whose removal preserves injectivity of the
label map $r_i \mapsto [r_i]_G$ --- i.e., a generator that does not
refine any multiplet distinction.  The configuration $G' := G
\setminus \{X\}$ is also admissible (it satisfies the same capacity
bound and isotropy condition; removing a channel only relaxes them).
By (a) and (b), $\Sigma_{d \in G} \varepsilon(d) - \Sigma_{d \in G'}
\varepsilon(d) = \varepsilon(X) \ge \varepsilon^* > 0$, so $G'$ is
strictly cheaper than $G$ inside the same admissible set.  By A2
(minimum cost), the realised configuration $G_{\mathrm{realized}}(\rho, R)
= \arg\min_{G \in \mathcal{A}(\rho, R)} \Sigma\varepsilon(d)$ cannot
contain $X$.  Iterating over all redundant generators: the realised
admissible structure contains no redundant generators.
\end{proof}

\noindent\textit{Status of the premises: derived, not additional.}\quad
A careful reading of Lemma~\ref{lem:EC_CM_from_A1} might suggest that
EC and CM import ``extra'' assumptions beyond A1.  This would be a
genuine vulnerability if the required premises were independently
postulated.  In fact, every premise in the dependency chain is derived:
\begin{center}\small
\renewcommand{\arraystretch}{1.15}
\begin{tabular}{@{}p{3.0cm}p{3.5cm}p{3.5cm}p{2.5cm}@{}}
\toprule
\textbf{Premise} & \textbf{Content} & \textbf{Derived from} & \textbf{Reference} \\
\midrule
$T_{\mathrm{sep}}$ & Sector decomposition of $S_\Gamma$ &
  A1 $+$ $T_{\mathrm{form}}$ & P1S~\S2 \\
$L_{\mathrm{loc}}$ & Independent interfaces &
  A1 $+$ FD1--FD2 & P1S~\S4 \\
FD3 & Anchor-set locality &
  Foundational definition & P1S~\S1 \\
Gauge ident. & $G = \mathrm{Inn}(\mathcal{A})$ &
  $T_{\mathrm{alg}}$ $+$ T2a $+$ Skolem--Noether & This supplement \\
$L_{\varepsilon^*}$ & Cost floor $\varepsilon^* > 0$ &
  A1 $+$ MD & P1S~\S3 \\
Gen.-channel bridge & $\dim(\mathfrak{g})$ channels &
  Faithfulness $+$ $T_{\mathrm{form}}$ & This supplement \\
A2 & Min-cost selection over $\mathcal{A}(\rho, R)$ &
  Second axiom (P1, \S2) & P1, \S2; P1S~\S2 \\
\bottomrule
\end{tabular}
\end{center}
\noindent The only external inputs beyond the two named axioms (A1, A2)
are the regularity condition MD (enforcement isotropy), which enters
through $L_{\varepsilon^*}$, and FD3 (anchor-set locality), a
foundational definition rather than an empirical assumption.  The
remaining premises are mathematical consequences of A1 (or A1 $+$ MD)
and the upstream chain.  The dependency for CM in the table now reads
A1 $+$ MD $+$ $L_{\varepsilon^*}$ $+$ generator--channel bridge (which
together define the admissible set with cost floor) $+$ A2 (which
selects the min-cost member).

The correct characterization is therefore: \emph{EC follows from A1
alone, via the derived structural separation.  CM follows from A2
applied to the admissible set defined by A1 $+$ MD: A1 bounds total
cost and admits every configuration satisfying that bound; MD gives
the positive cost floor (via $L_{\varepsilon^*}$) that makes a
redundant generator strictly costly; A2 then selects the minimum-cost
member of the admissible set, which by construction contains no
redundant generators.}  The earlier-version claim ``EC and CM follow
from A1'' was an overclaim on two counts: EC remains an A1 result, but
CM is properly an A2 result (operating on the A1 $+$ MD-defined
admissible set), and even the A1 $+$ MD reading attributed to CM in
the previous version of this text understated the role of A2 as the
selection axiom.  Neither correction imports a modeling assumption
beyond the framework's stated axioms: A2 is one of the two axioms
named in the Paper~1, \S2 box, and MD is the regularity input also
listed in the Paper~1 axiom inventory.

\smallskip\noindent\textit{CM is not a selection principle (descriptive
reading).}\quad
Capacity minimality is sometimes described as ``Occam's razor for gauge
structure,'' but this risks making it sound like an aesthetic preference
or, worse, a dynamical mechanism that ranges over candidate structures
and rejects the non-minimal ones.  Neither reading is correct.
Lemma~\ref{lem:EC_CM_from_A1} establishes CM as a \emph{consequence
of A2} applied to the admissibility set defined by A1 $+$ MD: a
configuration with a redundant generator is admissible (it satisfies
the capacity bound and isotropy condition), but it sits at strictly
higher total cost than the same configuration with the redundant
generator removed (also admissible, by $\varepsilon^* > 0$).  A2
selects the cheaper one as the realised structure.  The
generator--channel bridge (Lemma~\ref{lem:gen_channel}) makes the
gap quantitative: each redundant generator costs exactly
$\varepsilon^*$.  Read correctly, CM is therefore a \emph{descriptive
label} on the realised admissible structure, not an algorithm: the
realised configuration is the $\arg\min_G \dim G$ subject to EC, and
that argmin is exactly what A2 fixes.  No dynamical search is
performed and none is needed; the framework reads off the referent
of the two axioms (A1 admits, A2 selects).

\smallskip\noindent The unique abelian factor result is split into two
levels.  \textbf{Theorem~A} (Proposition~\ref{prop:unique_U1}) is
template-independent: it shows that for \emph{any} multiplet set
$\mathcal{M}$ and \emph{any} block decomposition, EC and CM determine
the number of admissible abelian factors from the separation deficiency
$\delta$ alone.  \textbf{Theorem~B}
(Theorem~\ref{thm:unique_U1_consolidated}) is the SM specialization:
it plugs in $\{n_k\} = \{3, 2, 1\}$ and $\mathcal{M} =
\mathcal{M}_{\mathrm{SM}}$ and outputs ``exactly one $\U(1)$,
uniquely hypercharge.''

\begin{proposition}[Unique minimal abelian grading factor ---
structural (Theorem~A)]\label{prop:unique_U1}
Let
$\mathcal{A} \cong \bigoplus_{\alpha=1}^{m} M_{n_\alpha}(\mathbb{C})$
be the enforcement algebra at a fixed interface, with non-abelian gauge
factor
$G_{\mathrm{na}} = \prod_{\alpha=1}^{m} \SU(n_\alpha)$
and relative-phase torus
$T_{\mathrm{rel}} := Z(\mathrm{U}(\mathcal{A})) / \mathrm{U}(1)_{\mathrm{diag}}
\cong \mathrm{U}(1)^{m-1}$.
Let $\mathcal{M} = \{r_1, \ldots, r_N\}$ denote the set of physically
distinct matter multiplet types occurring in the derived template,
regarded as irreducible carriers of $G_{\mathrm{na}}$.

Assume that $G_{\mathrm{phys}}$ is an admissible gauge structure
(Definition~\ref{def:admissible_gauge}; EC and CM hold by
Lemma~\ref{lem:EC_CM_from_A1}).

\begin{lemma}[Abelian factors arise only from $T_{\mathrm{rel}}$]
\label{lem:abelian_source}
Every continuous abelian gauge factor of $G_{\mathrm{phys}}$ is a
one-dimensional subgroup of $T_{\mathrm{rel}}$.  No abelian gauge
factor can arise from any other source.
\end{lemma}

\begin{proof}
By the exact sequence (Proposition~\ref{prop:exact_seq}):
\[
  1 \;\to\; Z(\mathrm{U}(\mathcal{A}))
  \;\to\; \mathrm{U}(\mathcal{A})
  \;\xrightarrow{\;\pi\;}\;
  \mathrm{Inn}(\mathcal{A}) \;\to\; 1.
\]
$\mathrm{Inn}(\mathcal{A}) = \prod_\alpha \mathrm{PU}(n_\alpha)$ is
semisimple (each $\mathrm{PU}(n)$ is semisimple for $n \ge 2$, and
trivial for $n = 1$): it contains no continuous abelian subgroups.
The only abelian continuous subgroups of $\mathrm{U}(\mathcal{A}) =
\prod_\alpha \mathrm{U}(n_\alpha)$ that project to the identity under
$\pi$ are contained in $\ker(\pi) = Z(\mathrm{U}(\mathcal{A}))$.
Quotienting by the physically unobservable diagonal
$\mathrm{U}(1)_{\mathrm{diag}}$ gives
$T_{\mathrm{rel}} = Z / \mathrm{U}(1)_{\mathrm{diag}}$ as the full
group of continuous abelian gauge transformations.  Any abelian gauge
factor is therefore a one-dimensional subgroup of $T_{\mathrm{rel}}$.

The argument excludes the following loopholes:
\begin{enumerate}[label=(\alph*),nosep]
  \item \textit{Abelian subgroups of $\mathrm{Inn}(\mathcal{A})$?}\quad
  $\mathrm{PU}(n)$ has a maximal torus of dimension $n - 1$, but these
  tori consist of \emph{non-abelian} gauge transformations (diagonal
  matrices in $\mathrm{SU}(n)$, which are non-central for $n \ge 2$).
  They are already counted in $G_{\mathrm{na}}$ as Cartan generators,
  not as independent abelian factors.
  \item \textit{Discrete abelian factors?}\quad Discrete subgroups
  (e.g., $\mathbb{Z}_N$) can survive as global symmetries but do not
  contribute continuous gauge generators.  CM counts continuous
  enforcement channels; discrete symmetries have no per-generator
  capacity cost.
  \item \textit{External abelian factors?}\quad Any abelian factor not
  contained in $\mathrm{U}(\mathcal{A})$ would act non-trivially on
  $\mathcal{H}_\omega$ but not arise from the enforcement algebra.
  This contradicts the gauge identification theorem
  (Theorem~\ref{thm:gauge_ident}): the gauge group is
  $\mathrm{Inn}(\mathcal{A})$ (plus the central extensions from the
  exact sequence), and no external group acts.
\end{enumerate}
\end{proof}

\begin{definition}[Separation deficiency]\label{def:sep_deficiency}
The \emph{separation deficiency} of $G_{\mathrm{na}}$ on $\mathcal{M}$
is the number of pairs of physically distinct multiplet types that
share identical non-abelian quantum numbers:
\[
  \delta(G_{\mathrm{na}}, \mathcal{M}) := \bigl|\{(r_i, r_j) : i < j,\;
  [r_i]_{G_{\mathrm{na}}} = [r_j]_{G_{\mathrm{na}}}\}\bigr|.
\]
EC is satisfied by $G_{\mathrm{na}}$ alone iff $\delta = 0$.  If
$\delta > 0$, additional abelian grading is required to resolve the
remaining degeneracies.
\end{definition}

\begin{lemma}[Minimum number of abelian factors from separation deficiency]
\label{lem:min_abelian}
Let $\mathcal{P} = \{P_1, \ldots, P_s\}$ be the partition of
$\mathcal{M}$ into equivalence classes under the non-abelian label
map (i.e., $P_k = \{r_i : [r_i]_{G_{\mathrm{na}}} = \ell_k\}$
for each distinct non-abelian label $\ell_k$).  Let
$m_{\max} := \max_k |P_k|$ be the size of the largest equivalence
class.  Then:
\begin{enumerate}[label=(\roman*),nosep]
  \item At least $\lceil \log_2 m_{\max} \rceil$ binary distinctions
  are needed to separate the elements of the largest class.
  \item A single $\U(1)$ factor assigns a real-valued charge $Y_i$ to
  each $r_i \in P_k$.  If the $Y_i$ are pairwise distinct within
  each class $P_k$, then one $\U(1)$ suffices to separate all
  multiplets.
  \item Therefore, the minimum number of abelian factors is $0$ if
  $\delta = 0$, and $1$ if $\delta > 0$ and there exists a charge
  assignment $Y$ that separates all classes.
\end{enumerate}
\end{lemma}

\begin{proof}
(i) is standard: $m_{\max}$ elements require $\lceil \log_2 m_{\max}
\rceil$ binary attributes to distinguish.  (ii): a single continuous
$\U(1)$ charge $Y \in \mathbb{R}$ has cardinality $\mathfrak{c}$, so
it can separate any finite number of elements --- the condition is that
the charge assignment $Y: P_k \to \mathbb{R}$ is injective on each
class.  (iii) combines (i) and (ii): if $m_{\max} \ge 2$, at least
one abelian factor is needed; if an injective $Y$ exists, one
suffices.
\end{proof}

\begin{lemma}[Diagonal rational grading from Schur + anomaly cancellation]
\label{lem:A3_derived}
Each $\U(1) \subset T_{\mathrm{rel}}$ assigns a rational charge to each
$r_i$, constant on that multiplet type.  Both properties are derived:
\begin{enumerate}[label=(\alph*),nosep]
  \item \textbf{Diagonal action (Schur).}\quad Each $r_i$ is an
  irreducible representation of $G_{\mathrm{na}} = \prod_\alpha
  \SU(n_\alpha)$.  The relative-phase torus
  $T_{\mathrm{rel}} \subset Z(\mathrm{U}(\mathcal{A}))$ lies in the
  center of the unitary group.  By Schur's lemma, any central element
  acts on an irreducible representation by a scalar.  Therefore each
  $\U(1) \subset T_{\mathrm{rel}}$ acts on $r_i$ by a phase
  $e^{i Y_i \theta}$, where $Y_i \in \mathbb{R}$ depends only on the
  multiplet type, not on the state within the multiplet.  This is
  diagonal action.
  \item \textbf{Rational charges (anomaly cancellation).}\quad The
  charges $Y_i$ are constrained by gauge anomaly cancellation
  (\S\ref{sec:hypercharge}): the four conditions
  $[\SU(2)]^2[\U(1)]$, $[\SU(3)]^2[\U(1)]$,
  $[\mathrm{grav}]^2[\U(1)]$, $[\U(1)]^3$ are polynomial equations
  in $\{Y_i\}$ with integer coefficients.  Solving the first three
  linearly and substituting into the fourth yields
  (Theorem~\ref{thm:hypercharge}):
  $z^2 - 2z - (N_c^2 - 1) = 0$, where $z := Y_u / Y_Q$.
  This polynomial is \emph{monic} (leading coefficient $= 1$) with
  integer coefficients.  By the rational root theorem for monic
  integer polynomials, any rational root $p/q$ (in lowest terms)
  must have $q \mid 1$, i.e., $z \in \mathbb{Z}$.  Indeed,
  $z = 1 \pm N_c$ are integers.  Since all other $Y_i$ are rational
  multiples of $Y_Q$ (from the linear conditions), and the standard
  normalization $Y_Q = 1/6$ is rational, $Y_i \in \mathbb{Q}$
  for all multiplets.  Rationality of hypercharges is therefore a
  \emph{theorem} of the anomaly system, not an assumption.
\end{enumerate}
\end{lemma}

\noindent Both properties were previously listed as assumption (A3).
The derivation shows they follow from Schur's lemma (standard
representation theory, no additional physical input) and anomaly
cancellation (already derived in this supplement).  No assumption
beyond EC, CM, and the derived gauge + matter content is needed.

Then:
\begin{enumerate}[label=(\roman*),nosep]
  \item If $\delta(G_{\mathrm{na}}, \mathcal{M}) > 0$, then at
  least one nontrivial $\U(1) \subset T_{\mathrm{rel}}$ is necessary
  (from EC: unresolved degeneracies violate enforcement completeness).
  \item If there exists a charge assignment $Y: \mathcal{M} \to \mathbb{Q}$
  such that the augmented label map
  $r_i \mapsto ([r_i]_{G_{\mathrm{na}}},\, Y(r_i))$
  is injective, then one abelian grading factor is sufficient
  (Lemma~\ref{lem:min_abelian}(ii)).
  \item Any additional independent abelian
  factor $Y'$ is inadmissible unless it changes the induced partition of
  $\mathcal{M}$; if $Y$ already separates all types, every further
  $\U(1)$ is redundant and excluded by CM
  (Lemma~\ref{lem:gen_channel}: each redundant generator wastes
  $\varepsilon^*$ capacity).
  \item Therefore, whenever $\delta > 0$ and one abelian grading
  separates all types, there exists exactly one physically admissible
  abelian gauge factor, unique up to overall normalization and change
  of basis in $T_{\mathrm{rel}}$.
\end{enumerate}
\end{proposition}

\begin{proof}
\textit{Step 1: Necessity (at least one abelian factor).}\quad
If $\delta(G_{\mathrm{na}}, \mathcal{M}) > 0$, there exist
distinct multiplet types $r_i \ne r_j$ with
$[r_i]_{G_{\mathrm{na}}} = [r_j]_{G_{\mathrm{na}}}$.  By EC
(Lemma~\ref{lem:EC_CM_from_A1}), these two physical
types must be distinguishable by the gauge structure.
Since $G_{\mathrm{na}}$ does not separate them, the distinction must
come from the remaining gauge freedom.  By
Lemma~\ref{lem:abelian_source}, the only continuous abelian gauge
freedom is $T_{\mathrm{rel}}$.  Hence at least one nontrivial
$\U(1) \subset T_{\mathrm{rel}}$ is necessary.  This gives (i).

\textit{Step 2: Sufficiency of one factor (explicit separation matrix).}\quad
We construct the separation matrix $S$ for
$\mathcal{M}_{\mathrm{SM}}$.  The non-abelian equivalence classes are:
\begin{center}\small
\begin{tabular}{@{}lll@{}}
\toprule
\textbf{Class} & \textbf{Non-abelian label} & \textbf{Members} \\
\midrule
$P_1$ & $(3, 2)$ & $\{Q\}$ (singleton) \\
$P_2$ & $(\bar{3}, 1)$ & $\{u^c, d^c\}$ \\
$P_3$ & $(1, 2)$ & $\{L\}$ (singleton) \\
$P_4$ & $(1, 1)$ & $\{e^c\}$ (singleton) \\
\bottomrule
\end{tabular}
\end{center}
The separation deficiency is $\delta = |P_2|(|P_2| - 1)/2 = 1$:
a single unresolved pair $(u^c, d^c)$.  The maximum class size is
$m_{\max} = 2$.  By Lemma~\ref{lem:min_abelian}, one $\U(1)$ factor
suffices iff it assigns distinct charges to $u^c$ and $d^c$.

The hypercharge assignment $Y_{u^c} = 2/3 \ne -1/3 = Y_{d^c}$
resolves this degeneracy.  The full charge vector
$(Y_Q, Y_{u^c}, Y_{d^c}, Y_L, Y_{e^c}) = (1/6, 2/3, -1/3, -1/2, -1)$
is pairwise distinct across all five multiplet types.  The augmented
label map
\[
  r_i \mapsto \bigl([r_i]_{\SU(3) \times \SU(2)},\; Y(r_i)\bigr)
\]
is therefore injective on $\mathcal{M}_{\mathrm{SM}}$: no two
multiplets share both their non-abelian quantum numbers and their
hypercharge.  One $\U(1)$ suffices.  This gives (ii).

\textit{Step 3: No second abelian factor needed or admissible.}\quad
Suppose a second independent abelian factor with charge assignment
$Y': \mathcal{M} \to \mathbb{Q}$ is added.  The augmented map
becomes $r_i \mapsto ([r_i]_{G_{\mathrm{na}}}, Y(r_i), Y'(r_i))$.
Since $Y$ already makes this map injective (Step 2), the additional
charge $Y'$ does not refine the partition of $\mathcal{M}$: the
partition $\{\{Q\}, \{u^c\}, \{d^c\}, \{L\}, \{e^c\}\}$ is already
the finest possible (each class is a singleton).  By CM
(Lemma~\ref{lem:EC_CM_from_A1}), the generator of $Y'$ consumes
$\varepsilon^* > 0$ capacity (Lemma~\ref{lem:gen_channel}) with zero
enforcement gain.  Under finite capacity (A1), this is inadmissible.
This gives (iii).

\textit{Step 4: Uniqueness up to normalization.}\quad
The surviving abelian factor is a one-dimensional subgroup of
$T_{\mathrm{rel}} \cong \mathrm{U}(1)^{m-1}$.  In the SM case,
$m = 3$ (blocks $M_3, M_2, M_1$), so $T_{\mathrm{rel}} \cong
\mathrm{U}(1)^2$.  The requirement that the charge assignment
resolves the $(u^c, d^c)$ degeneracy and satisfies all four anomaly
conditions (\S\ref{sec:hypercharge}) determines a \emph{unique
direction} in $T_{\mathrm{rel}}$, up to normalization
($Y \mapsto \lambda Y$, $\lambda \in \mathbb{Q}^\times$).  This
uniqueness is proved in Theorem~\ref{thm:hypercharge}: the anomaly
system has exactly two solutions $z = 1 \pm N_c$ (with $N_c = 3$:
$z = 4$ and $z = -2$), which correspond to the same direction
under the $u \leftrightarrow d$ relabeling symmetry.  Therefore
the abelian factor is unique.  This gives (iv).
\end{proof}

\begin{corollary}[Application to the derived Standard Model template]
\label{cor:unique_U1_SM}
For the derived multiplet set
$\mathcal{M}_{\mathrm{SM}} = \{Q,\, u^c,\, d^c,\, L,\, e^c\}$,
the non-abelian labels under $\SU(3) \times \SU(2)$ are
$Q \sim (3,2)$, $L \sim (1,2)$, $u^c \sim (\bar{3},1)$,
$d^c \sim (\bar{3},1)$, $e^c \sim (1,1)$.
Hence the non-abelian label map is not injective (it conflates $u^c$
with $d^c$): $\delta = 1$.  Hypercharge resolves this degeneracy with
pairwise distinct charges.  By Proposition~\ref{prop:unique_U1},
exactly one abelian gauge factor is necessary and sufficient.
\end{corollary}

The preceding chain of lemmas, proposition, and corollary assembles the
result from reusable components.  The reviewer asks for a single
self-contained theorem that takes the exact-sequence data as input and
outputs ``exactly one $\U(1)$'' without requiring the reader to trace
through the intermediate results.  We provide this.

\begin{theorem}[Unique abelian factor from exact-sequence data ---
SM specialization (Theorem~B)]%
\label{thm:unique_U1_consolidated}
\textbf{Hypotheses:}
\begin{enumerate}[label=(H\arabic*),nosep]
  \item \textbf{Enforcement algebra.}\quad
  $\mathcal{A} = M_3(\mathbb{C}) \oplus M_2(\mathbb{C}) \oplus
  \mathbb{C}$ (Wedderburn blocks $\{n_k\} = \{3, 2, 1\}$, from
  T2a + Theorem~R).
  \item \textbf{Exact sequence.}\quad
  $1 \to Z(\mathrm{U}(\mathcal{A})) \to \mathrm{U}(\mathcal{A})
  \xrightarrow{\pi} \mathrm{Inn}(\mathcal{A}) \to 1$, with
  $\mathrm{Inn}(\mathcal{A}) = \mathrm{PU}(3) \times \mathrm{PU}(2)$,
  $\ker(\pi) = \mathrm{U}(1)^3$,
  $T_{\mathrm{rel}} = \ker(\pi) / \mathrm{U}(1)_{\mathrm{diag}}
  \cong \mathrm{U}(1)^2$
  (Proposition~\ref{prop:exact_seq}).
  \item \textbf{Matter content.}\quad
  $\mathcal{M} = \{Q(3,2),\; u^c(\bar{3},1),\; d^c(\bar{3},1),\;
  L(1,2),\; e^c(1,1)\}$ (from $T_{\mathrm{field}}$,
  \S\ref{sec:fermion_scan}).
  \item \textbf{EC and CM.}\quad The gauge structure satisfies
  enforcement completeness (EC: the label map is injective) and
  capacity minimality (CM: no generator can be removed while
  preserving EC); EC is derived from A1 alone, CM from A2 over the
  A1 $+$ MD-defined admissible set
  (Lemma~\ref{lem:EC_CM_from_A1}).
\end{enumerate}

\smallskip\noindent\textbf{Conclusion:}\quad The physical gauge group
contains exactly one continuous abelian factor, which is a uniquely
determined one-dimensional subgroup of $T_{\mathrm{rel}} \cong
\mathrm{U}(1)^2$.  This factor is hypercharge $\U(1)_Y$, with charge
assignments $(Y_Q, Y_{u^c}, Y_{d^c}, Y_L, Y_{e^c}) =
(1/6, 2/3, -1/3, -1/2, -1)$ up to normalization and the
$u \leftrightarrow d$ relabeling.
\end{theorem}

\begin{proof}
The proof proceeds in five steps: source exclusivity, necessity,
sufficiency, uniqueness of direction, and redundancy of any second
factor.

\textit{Step 1 (Source exclusivity): every abelian factor lives in
$T_{\mathrm{rel}}$.}\quad
By (H2), $\mathrm{Inn}(\mathcal{A}) = \mathrm{PU}(3) \times
\mathrm{PU}(2)$ is semisimple (both factors are semisimple for
$n \ge 2$).  A semisimple group has no continuous abelian
quotients: every connected abelian subgroup of $\mathrm{PU}(n)$ is
a maximal torus (the image of the Cartan subalgebra of
$\mathfrak{su}(n)$), which is part of the non-abelian structure.
The only continuous abelian subgroups of $\mathrm{U}(\mathcal{A})$
that project to the identity under $\pi$ are in $\ker(\pi) =
Z(\mathrm{U}(\mathcal{A}))$.  Quotienting by the unobservable
$\mathrm{U}(1)_{\mathrm{diag}}$ gives $T_{\mathrm{rel}} \cong
\mathrm{U}(1)^2$.  Any abelian gauge factor is a one-dimensional
subgroup of this torus.

\textit{Step 2 (Necessity): at least one $\U(1)$ is required.}\quad
By (H3), the non-abelian label map
$r_i \mapsto [r_i]_{\SU(3) \times \SU(2)}$ is:
\begin{center}\small
\begin{tabular}{@{}lccccc@{}}
\toprule
$r_i$ & $Q$ & $u^c$ & $d^c$ & $L$ & $e^c$ \\
\midrule
$[r_i]$ & $(3,2)$ & $(\bar{3},1)$ & $(\bar{3},1)$ & $(1,2)$ & $(1,1)$ \\
\bottomrule
\end{tabular}
\end{center}
The map sends both $u^c$ and $d^c$ to $(\bar{3}, 1)$.  Therefore
the label map is not injective: $\delta(\SU(3) \times \SU(2),
\mathcal{M}) = 1$.  By EC (H4), all five types must be
distinguished by the gauge structure.  Since the non-abelian
labels fail to separate $u^c$ from $d^c$, at least one abelian
factor from $T_{\mathrm{rel}}$ is required.

\textit{Step 3 (Sufficiency): one $\U(1)$ resolves all
degeneracies.}\quad
A one-dimensional subgroup $\U(1) \subset T_{\mathrm{rel}}$ assigns a
rational charge $Y_i$ to each $r_i$ (by Schur's lemma: central
elements act by scalars on irreducible representations; by the anomaly
system: the charges are rational, \S\ref{sec:hypercharge}).  The
augmented label map $r_i \mapsto ([r_i]_{\SU(3) \times \SU(2)},
Y_i)$ is injective iff $Y_{u^c} \ne Y_{d^c}$ (this is the only
unresolved pair from Step~2; the other three types already have
distinct non-abelian labels).  The anomaly system
(Theorem~\ref{thm:hypercharge}) determines:
\begin{align*}
  Y_{u^c} &= z \cdot Y_Q, \qquad Y_{d^c} = (2 - z) \cdot Y_Q,
\end{align*}
where $z$ satisfies $z^2 - 2z - (N_c^2 - 1) = 0$ with $N_c = 3$,
giving $z \in \{4, -2\}$.  For either root:
$Y_{u^c} - Y_{d^c} = (2z - 2) Y_Q \ne 0$ (since $z \ne 1$).
Therefore one $\U(1)$ suffices.

\textit{Step 4 (Uniqueness of direction): the anomaly system
admits a unique direction in $T_{\mathrm{rel}}$.}\quad
$T_{\mathrm{rel}} \cong \mathrm{U}(1)^2$ is a two-dimensional torus.
A one-dimensional subgroup is specified by a direction
$\mathbf{q} = (q_1, q_2) \in \mathbb{Z}^2 \setminus \{0\}$ (up to
positive rescaling).  The charge of multiplet $r_i$ is
$Y_i = q_1 \cdot c_1(r_i) + q_2 \cdot c_2(r_i)$, where
$c_1, c_2$ are the canonical $\mathrm{U}(1)$ charges from the
center $Z = \mathrm{U}(1)_3 \times \mathrm{U}(1)_2 \times
\mathrm{U}(1)_1$.  (In the canonical basis:
$c_1(r_i) = n_3$-ality of the $\SU(3)$ rep, $c_2(r_i)$ =
$\SU(2)$ isospin parity.)

The four anomaly conditions (\S\ref{sec:hypercharge}) are four
polynomial equations in $(q_1, q_2)$.  Three are linear; the
fourth is cubic.  The linear conditions determine the charge
ratios $Y_{u^c}/Y_Q$, $Y_{d^c}/Y_Q$, $Y_L/Y_Q$, $Y_{e^c}/Y_Q$
as functions of a single parameter $z = Y_{u^c}/Y_Q$.  The cubic
condition reduces to the quadratic $z^2 - 2z - 8 = 0$ with roots
$z = 4$ and $z = -2$.  These two roots correspond to the same
direction in $T_{\mathrm{rel}}$ (they are related by the
$u \leftrightarrow d$ relabeling symmetry of the fermion template).
Therefore the anomaly system admits a \emph{unique direction}
$\mathbf{q}$ in $T_{\mathrm{rel}}$, up to normalization and the
discrete relabeling.  (``Admits'' rather than ``selects'': no
mechanism ranges over candidate directions; the admissibility
conditions --- four anomaly equations plus EC, CM --- have a unique
non-trivial common solution up to the stated equivalences, and
that solution is the referent of the admissible hypercharge.)

\textit{Step 5 (No second factor): CM excludes any additional
$\U(1)$.}\quad
After Step~3, the augmented label map
$r_i \mapsto ([r_i]_{\SU(3) \times \SU(2)}, Y_i)$ is injective on
all five multiplet types (the full charge vector
$(1/6, 2/3, -1/3, -1/2, -1)$ is pairwise distinct).  The partition of
$\mathcal{M}$ into equivalence classes under this map is
$\{\{Q\}, \{u^c\}, \{d^c\}, \{L\}, \{e^c\}\}$ --- the discrete
partition (all singletons).

A second independent $\U(1)' \subset T_{\mathrm{rel}}$ would assign
a second charge $Y'_i$ to each multiplet.  The augmented map
$r_i \mapsto ([r_i], Y_i, Y'_i)$ induces the same discrete
partition (singletons cannot be further refined).  Therefore $Y'$
provides zero additional enforcement information.  By CM (H4), the
generator of $\U(1)'$ occupies one enforcement channel
(Lemma~\ref{lem:gen_channel}) at cost $\varepsilon^* > 0$
($L_{\varepsilon^*}$) with zero enforcement gain.  Under finite
capacity (A1), this is inadmissible.

\textit{Conclusion.}\quad Steps~1--5 establish: exactly one
continuous abelian factor exists in the physical gauge group; it is a
one-dimensional subgroup of $T_{\mathrm{rel}} \cong \mathrm{U}(1)^2$;
the anomaly system selects a unique direction; and no second factor is
admissible.  The physical gauge group is therefore
$[\SU(3) \times \SU(2) \times \U(1)_Y] / \mathbb{Z}_6$
(Proposition~\ref{prop:global_structure}).
\end{proof}

\noindent\textit{Why the theorem is not heuristic.}\quad
The selection of one $\U(1)$ from the $(K-1)$-torus involves three
independently verifiable steps, none of which is an aesthetic or
heuristic choice:
\begin{itemize}[nosep]
  \item \textit{Necessity} (Step~2) is a counting argument:
  $\delta > 0$ combined with EC forces at least one abelian factor.
  This is a finite combinatorial check on the five multiplets.
  \item \textit{Uniqueness of direction} (Step~4) is a
  number-theoretic result: the anomaly system is a determined system
  of polynomial equations with a unique projective solution.  This is
  verified by explicit computation (the discriminant of
  $z^2 - 2z - 8 = 0$ is $36$; the roots are $4$ and $-2$; they are
  related by the relabeling symmetry).
  \item \textit{Redundancy exclusion} (Step~5) is a capacity
  argument: the partition is already maximally fine (all singletons),
  so no new information can be added.  This is a logical tautology
  combined with the generator--channel bridge (each generator costs
  $\varepsilon^* > 0$).
\end{itemize}
The only conditional input is the matter content
$\mathcal{M} = \mathcal{M}_{\mathrm{SM}}$ (H3), which is derived in
$T_{\mathrm{field}}$ (\S\ref{sec:fermion_scan}).  Given that input,
the conclusion is a theorem, not a selection.

\smallskip\noindent\textit{Logical ordering and the dependence on
$T_{\mathrm{field}}$.}\quad
Proposition~\ref{prop:unique_U1} is a conditional
result: it shows that for \emph{any} fixed set of non-abelian multiplets
$\mathcal{M}$ satisfying EC and CM, the number of admissible $\U(1)$
factors is uniquely determined by the multiplet separation problem.
The physical conclusion ``exactly one $\U(1)$'' requires the input
$\mathcal{M} = \mathcal{M}_{\mathrm{SM}}$ from $T_{\mathrm{field}}$
(\S\ref{sec:fermion_scan}), which is proved independently.  The logical
order is:
\[
  \text{Theorem R} \;\to\; \text{gauge template}
  \;\to\; T_{\mathrm{field}} \;\to\;
  \text{Corollary~\ref{cor:unique_U1_SM}}
  \;\to\; L_{\mathrm{hypercharge}}.
\]
There is no circularity: $T_{\mathrm{field}}$ does not use the abelian
factor (the scan uses only $\SU(3) \times \SU(2)$ quantum numbers);
the abelian factor is determined afterwards.

\smallskip\noindent\textit{Edge case: $\U(1)_{B-L}$ with sterile
$\nu_R$.}\quad
A frequently raised alternative is the anomaly-free combination
$\U(1)_{B-L}$ (baryon minus lepton number), which becomes anomaly-free
when the multiplet set is extended by three right-handed neutrinos
$\nu_R(\mathbf{1},\mathbf{1})$.  In this enlarged template
$\mathcal{M}' = \mathcal{M}_{\mathrm{SM}} \cup \{\nu_R\}$, both
$\U(1)_Y$ and $\U(1)_{B-L}$ are anomaly-free, and the question is
whether two independent abelian factors are admissible.

The APF excludes this scenario at two stages:
\begin{enumerate}[label=(\alph*),nosep]
  \item \textbf{$T_{\mathrm{field}}$ excludes $\nu_R$ from the minimal
  template.}\quad The fermion scan (\S\ref{sec:fermion_scan}) derives
  $\mathcal{M}_{\mathrm{SM}}$ as the unique minimum-DOF survivor at
  45 Weyl fermions.  The template
  $\mathcal{M}_{\mathrm{SM}} \cup \{\nu_R\}$ has $48$ DOF and survives
  all anomaly filters, but is non-minimal: it is the second-lightest
  post-CPT survivor (template \#3 in the waterfall table).  Since the
  multiplet set is $\mathcal{M}_{\mathrm{SM}}$, not $\mathcal{M}'$,
  the $B{-}L$ direction is not anomaly-free for the \emph{derived}
  content.  (Indeed, $B{-}L$ has a mixed gravitational anomaly
  $[\mathrm{grav}]^2[\U(1)_{B-L}] = 3 \ne 0$ unless $\nu_R$ is
  present.)
  \item \textbf{Capacity minimality excludes the second $\U(1)$ even
  if $\nu_R$ were present.}\quad Suppose counterfactually that
  $\mathcal{M}' = \mathcal{M}_{\mathrm{SM}} \cup \{\nu_R\}$ were the
  derived template.  Does Proposition~\ref{prop:unique_U1} still
  exclude a second $\U(1)$?  Yes: in $\mathcal{M}'$, the non-abelian
  labels conflate $u^c / d^c$ (both $(\bar{3},1)$) and
  $e^c / \nu_R$ (both $(1,1)$).  A single abelian grading with four
  distinct charges resolves all degeneracies: one $\U(1)$ suffices.
  $\U(1)_{B-L}$ is then a \emph{linear combination} of this grading
  and the non-abelian quantum numbers---not an independent gauge
  direction.  A second independent $\U(1)$ would carry capacity cost
  (one additional enforcement channel) with no additional
  multiplet-separation power, violating CM.
  \item \textbf{Observational cross-check.}\quad
  The $\nu_R$-extended template predicts $C_{\mathrm{total}} = 64$ and
  $\Omega_\Lambda = 42/64 = 0.656$, which is $> 5\sigma$ from the
  Planck measurement $\Omega_\Lambda = 0.6889 \pm 0.0056$
  (\S\ref{sec:neutrinos}).  The 45-fermion template with one $\U(1)$
  predicts $\Omega_\Lambda = 42/61 = 0.6885$, consistent at
  $< 0.1\sigma$.
\end{enumerate}
In summary, the ``exactly one $\U(1)$'' conclusion is robust against
the $B{-}L$ edge case: the extended multiplet set is non-minimal
($T_{\mathrm{field}}$), $B{-}L$ is anomalous for the minimal template,
and even if $\nu_R$ were included, $B{-}L$ would be a redundant linear
combination rather than an independent gauge factor.

\smallskip\noindent\textit{String-theoretic extra $\U(1)$'s and
St\"uckelberg masses.}\quad
Heterotic string compactifications and F-theory models generically
produce multiple abelian gauge factors beyond hypercharge.  In these
constructions, unwanted $\U(1)$'s acquire large masses via the
St\"uckelberg mechanism: the $\U(1)$ gauge boson ``eats'' a closed-string
axion, becoming massive at or near the compactification scale, while
the associated global selection rule (e.g., a $\mathbb{Z}_N$ discrete
symmetry) may persist at low energies.

In the APF, this scenario has a precise capacity interpretation:
\begin{enumerate}[label=(\alph*),nosep]
  \item A St\"uckelberg-massive $\U(1)$ is a gauge boson that no longer
  propagates as a long-range force.  In the enforcement language, the
  corresponding channel is ``saturated'': its enforcement cost is paid
  at the compactification scale, and at low energies it is absorbed into
  the vacuum sector ($L_{\mathrm{vac}}$, \S\ref{sec:vac}).
  The channel is counted in $C_{\mathrm{total}}$ (as part of the UV
  structure) but does not contribute an independent low-energy gauge
  coupling.
  \item Capacity minimality (CM) does not merely select the gauge group
  with the fewest generators; it selects the group with the fewest
  \emph{active enforcement channels} at the IR saturation scale.
  A St\"uckelberg-massive $\U(1)$ is inactive at low energies and
  therefore excluded from the low-energy gauge group by the same
  principle that excludes $\nu_R$ from the low-energy fermion count:
  it is above the IR capacity ledger.
  \item The APF is therefore selecting the \emph{phase} in which only
  one $\U(1)$ combination remains both anomaly-free and massless at
  low energies.  This is precisely the phase that heterotic
  phenomenology identifies as the ``Standard Model embedding'': among
  the multiple $\U(1)$'s of the string construction, one linear
  combination is identified with hypercharge and remains massless, while
  the others acquire St\"uckelberg masses and decouple.  The APF's
  capacity-minimality principle provides the \emph{selection criterion}
  for this phase: the massless $\U(1)$ is the one that resolves all
  low-energy multiplet degeneracies at minimum capacity cost.
\end{enumerate}

\smallskip\noindent\textit{Kinetic mixing and capacity accounting.}\quad
Even if a second $\U(1)'$ is excluded as an independent gauge factor
by CM, one might worry about \emph{kinetic mixing}: a gauge-invariant
term $\frac{\chi}{2} F_{\mu\nu} F'^{\mu\nu}$ coupling the field
strengths of hypercharge $\U(1)_Y$ and a putative $\U(1)'$.  This is
a legitimate concern in BSM model-building, where kinetic mixing
arises at one loop even if $\chi = 0$ at tree level, provided charged
matter couples to both $\U(1)$'s.  Three points address this within
the APF:
\begin{enumerate}[label=(\alph*),nosep]
  \item \textit{No second $\U(1)$ exists to mix with.}\quad
  Proposition~\ref{prop:unique_U1} proves that exactly one $\U(1)$
  from the relative-phase torus $T_{\mathrm{rel}}$ is admissible.
  Kinetic mixing requires two abelian gauge fields; if the second
  does not exist as a propagating field (because CM excludes it),
  there is no $F'^{\mu\nu}$ to couple to.  The concern is therefore
  moot at the level of the low-energy gauge structure derived here.

  \item \textit{Kinetic mixing as a capacity cost.}\quad
  If a second $\U(1)'$ were present, the kinetic mixing term would
  be an off-diagonal entry in the $2 \times 2$ kinetic matrix
  $K_{ab} F^a_{\mu\nu} F^{b\,\mu\nu}$.  In the enforcement language,
  this corresponds to a cross-channel correlation: the enforcement
  cost of the $\U(1)_Y$ channel depends on the state of the
  $\U(1)'$ channel.  By the cost-kernel structure (Theorem~\ref{thm:T7B}),
  cross-channel costs are captured by the off-diagonal terms of the
  bilinear form $B$.  The capacity cost of maintaining a kinetically
  mixed pair is strictly greater than the cost of a single $\U(1)$
  (because $B$ is positive-definite: $B(v_1 + v_2, v_1 + v_2) =
  B(v_1, v_1) + 2B(v_1, v_2) + B(v_2, v_2) > B(v_1, v_1)$ when
  $v_2 \ne 0$).  CM therefore excludes the mixed pair in favor of the
  single channel.

  \item \textit{Diagonalization absorbs mixing into charge
  redefinition.}\quad
  When two $\U(1)$'s kinetically mix, the physical basis is obtained
  by diagonalizing $K_{ab}$.  In the diagonal basis, one linear
  combination couples to matter (this is hypercharge) and the other is
  either decoupled (capacity-redundant by CM) or massive
  (Stückelberg mechanism, handled above).  The APF's derivation
  operates in the diagonal basis from the start: the relative-phase
  torus provides a canonical basis of independent $\U(1)$ directions,
  and Proposition~\ref{prop:unique_U1} selects the unique admissible
  direction.  No off-diagonal mixing can arise because the canonical
  basis is already diagonal in the cost kernel $B$.
\end{enumerate}

\noindent\textbf{Physical gauge group.}\quad By the gauge identification
theorem (Theorem~\ref{thm:gauge_ident}) and
Proposition~\ref{prop:unique_U1}, the Lie algebra of the physical gauge
group is $\mathfrak{su}(3) \oplus \mathfrak{su}(2) \oplus
\mathfrak{u}(1)_Y$.  At the global (topological) level:
\[
  G_{\mathrm{phys}} = \frac{\SU(3) \times \SU(2) \times \U(1)_Y}
  {\mathbb{Z}_6}\,,
\]
where the $\mathbb{Z}_6$ quotient is derived in
Proposition~\ref{prop:global_structure} below, $\SU(n_\alpha)$ is the
universal cover of $\mathrm{PU}(n_\alpha)$, and $\U(1)_Y$ is the
unique relative-phase factor from Proposition~\ref{prop:unique_U1}.

The derivation chain produces $\mathrm{Inn}(\mathcal{A}) = \prod_\alpha
\mathrm{PU}(n_\alpha)$ as the automorphism group of the algebra, but the
physical gauge group that acts faithfully on matter fields is
$[\SU(3) \times \SU(2) \times \U(1)_Y] / \mathbb{Z}_6$.  The passage between these two descriptions involves
(a)~a lift from $\mathrm{PU}(n)$ to $\SU(n)$ required by the
representation theory of matter, and (b)~a question about global
structure (center quotients) that determines the topology of the gauge
group.  We address both formally.

\begin{proposition}[Lift from $\mathrm{PU}(n)$ to $\SU(n)$]%
\label{prop:PU_to_SU}
The physical gauge group acting on matter fields in $\mathcal{H}_\omega$
is $\SU(n)$, not $\mathrm{PU}(n)$, for each Wedderburn block
$M_n(\mathbb{C})$ with $n \ge 2$.
\end{proposition}

\begin{proof}
\textit{Step 1 (Projective obstruction).}\quad
$\mathrm{Inn}(\mathcal{A}) = \prod_\alpha \mathrm{PU}(n_\alpha)$ acts
on $\mathcal{A}$ by conjugation: $\alpha_U(a) = UaU^*$.  This is the
\emph{adjoint} action.  Matter fields, however, live in
$\mathcal{H}_\omega$, the GNS Hilbert space
(Proposition~\ref{prop:exact_seq}(v)), which carries the
\emph{fundamental} representation of each block: $v \mapsto Uv$ for
$v \in \mathbb{C}^n$.  The fundamental is a linear representation of
$\mathrm{U}(n)$ but only a \emph{projective} representation of
$\mathrm{PU}(n)$: the lift $\mathrm{PU}(n) \to \mathrm{GL}(n,
\mathbb{C})$ requires choosing a phase for each element, and there is
no continuous choice that is globally consistent (the obstruction is
the second group cohomology $H^2(\mathrm{PU}(n), \mathrm{U}(1))
\cong \mathbb{Z}_n$).

\textit{Step 2 (Linearization by lift to universal cover).}\quad
The universal covering group of $\mathrm{PU}(n)$ is $\SU(n)$, with
covering map $\SU(n) \twoheadrightarrow \mathrm{PU}(n)$ and kernel
$Z(\SU(n)) \cong \mathbb{Z}_n$ (the center, generated by
$e^{2\pi i / n}\,\mathbf{1}_n$).  Every projective representation of
$\mathrm{PU}(n)$ lifts to a genuine (linear) representation of
$\SU(n)$.  The Lie algebras are isomorphic:
$\mathfrak{su}(n) = \mathfrak{pu}(n)$, so the gauge field content
(connection 1-form $A_\mu$, curvature $F_{\mu\nu}$, covariant
derivative $D_\mu$) is identical regardless of whether one uses
$\SU(n)$ or $\mathrm{PU}(n)$.  The distinction is purely global:
$\SU(n)$ admits the fundamental representation while $\mathrm{PU}(n)$
does not.

\textit{Step 3 (Carrier requirement forces the fundamental).}\quad
The carrier requirement R1 (Theorem~R, \S\ref{sec:carrier_prereqs})
demands a faithful complex carrier of dimension $n$ (from $B1'$).
The fundamental representation of $\SU(n)$ on $\mathbb{C}^n$ is the
unique faithful irreducible representation of dimension $n$; the
adjoint ($\mathrm{PU}(n)$ on $\mathfrak{su}(n) \cong
\mathbb{C}^{n^2-1}$) has dimension $n^2 - 1$ and is not faithful on
$\mathbb{C}^n$.  R1 therefore forces the fundamental, which in turn
forces $\SU(n)$ as the gauge group acting on matter.
\end{proof}

\begin{proposition}[Global structure of the gauge group]%
\label{prop:global_structure}
The gauge group acting on the matter content of the APF is:
\[
  G_{\mathrm{phys}} \;=\;
  \frac{\SU(3) \times \SU(2) \times \U(1)_Y}{\mathbb{Z}_6}\,,
\]
where $\mathbb{Z}_6 \cong Z(\SU(3)) \times Z(\SU(2)) \cap \ker(Y)$
is the subgroup of the center that acts trivially on all matter
representations.
\end{proposition}

\begin{proof}
The group that acts faithfully on matter is not the direct product
$\SU(3) \times \SU(2) \times \U(1)_Y$ but its quotient by the
subgroup of center elements acting trivially on all derived multiplets.

\textit{Step 1 (Center identification).}\quad
The center of $\SU(3) \times \SU(2) \times \U(1)_Y$ is
$\mathbb{Z}_3 \times \mathbb{Z}_2 \times \U(1)_Y$.  A central element
$(e^{2\pi i k/3}\,\mathbf{1}_3,\; (-1)^m\,\mathbf{1}_2,\;
e^{iY\theta})$ acts on a multiplet $(r_3, r_2, Y_f)$ by the phase
$e^{2\pi i k\, c(r_3)/3}\, (-1)^{m\, t(r_2)}\, e^{iY_f\theta}$,
where $c(r_3)$ is the $\mathbb{Z}_3$ triality of the $\SU(3)$
representation and $t(r_2)$ is the $\mathbb{Z}_2$ character of the
$\SU(2)$ representation.

\textit{Step 2 (Trivially acting subgroup --- explicit verification).}\quad
A central element acts trivially on \emph{all} derived multiplets iff
$e^{2\pi i k\, c(r_3)/3}\, (-1)^{m\, t(r_2)}\, e^{iY_f\theta} = 1$
for every multiplet in the SM template.  The derived multiplet list
(from the fermion scan, $T_{\mathrm{field}}$) assigns each multiplet
definite quantum numbers.

\textit{Normalization.}\quad We work with the integer-valued charge
$Y_{\mathrm{int}} := 6Y$ (so that $Y_Q = 1/6$ becomes
$Y_{\mathrm{int}} = 1$, etc.).  The $\U(1)_Y$ element $e^{i\theta}$
acts on a multiplet with charge $Y_{\mathrm{int}}$ by the phase
$e^{iY_{\mathrm{int}}\theta}$.  A center element
$(e^{2\pi i k/3}\,\mathbf{1}_3,\; (-1)^m\,\mathbf{1}_2,\;
e^{i\theta})$ acts on multiplet $(c, t, Y_{\mathrm{int}})$ by:
\begin{equation}\label{eq:center_phase}
  \phi(k,m,\theta) \;=\; e^{2\pi i k\, c/3}\;
  e^{\pi i m\, t}\; e^{i Y_{\mathrm{int}}\,\theta}.
\end{equation}
Setting $\theta = \pi/3$ (the candidate $\mathbb{Z}_6$ generator)
and $(k, m) = (1, 1)$:
\[
  \phi\bigl(1,1,\tfrac{\pi}{3}\bigr) = \exp\Bigl[\,
  i\pi\,\frac{2c + 3t + Y_{\mathrm{int}}}{3}\Bigr].
\]
This equals $1$ iff $2c + 3t + Y_{\mathrm{int}} \equiv 0 \pmod{6}$.
We verify this for every multiplet type in one generation of the
derived fermion template:

\smallskip
\begin{center}\small
\renewcommand{\arraystretch}{1.15}
\begin{tabular}{@{}lccccccr@{}}
\toprule
\textbf{Multiplet} & $\SU(3)$ & $\SU(2)$ & $Y$ &
  $c$ & $t$ & $Y_{\mathrm{int}}$ &
  $2c{+}3t{+}Y_{\mathrm{int}}$ \\
\midrule
$Q_L$ & $\mathbf{3}$ & $\mathbf{2}$ & $1/6$ & 1 & 1 & 1 & $6 \equiv 0$ \\
$u_R$ & $\mathbf{3}$ & $\mathbf{1}$ & $2/3$ & 1 & 0 & 4 & $6 \equiv 0$ \\
$d_R$ & $\mathbf{3}$ & $\mathbf{1}$ & $-1/3$ & 1 & 0 & $-2$ & $0 \equiv 0$ \\
$L_L$ & $\mathbf{1}$ & $\mathbf{2}$ & $-1/2$ & 0 & 1 & $-3$ & $0 \equiv 0$ \\
$e_R$ & $\mathbf{1}$ & $\mathbf{1}$ & $-1$ & 0 & 0 & $-6$ & $-6 \equiv 0$ \\
\midrule
$H$ (Higgs) & $\mathbf{1}$ & $\mathbf{2}$ & $1/2$ & 0 & 1 & 3 & $6 \equiv 0$ \\
\bottomrule
\end{tabular}
\end{center}

\smallskip\noindent All entries satisfy
$2c + 3t + Y_{\mathrm{int}} \equiv 0 \pmod{6}$.  The Higgs doublet
$(1, 2, 1/2)$ is included because it participates in gauge
interactions; its compatibility is necessary for consistency of the
Yukawa sector.

The $\mathbb{Z}_6$ generator is therefore
$g := (\omega_3\,\mathbf{1}_3,\; -\mathbf{1}_2,\;
e^{i\pi/3})$, and the six elements $\{g^0, g^1, \ldots, g^5\}$
exhaust the kernel.  To confirm that the kernel is exactly
$\mathbb{Z}_6$ and not a proper subgroup, note that
$\mathbb{Z}_3 \times \mathbb{Z}_2 \cong \mathbb{Z}_6$ (since
$\gcd(2,3) = 1$), so $\mathbb{Z}_6$ is the largest discrete
subgroup of $Z(\SU(3)) \times Z(\SU(2))$.  It remains to verify that
elements of the center with \emph{trivial} $\U(1)$ phase do not
act trivially: $(\omega_3,\, \mathbf{1}_2,\, 1)$ acts on $Q_L$ as
$\omega_3 \ne 1$, and $(\mathbf{1}_3,\, -\mathbf{1}_2,\, 1)$ acts
on $Q_L$ as $-1 \ne 1$.  The nontrivial $\U(1)$ phase
$e^{i\pi/3}$ in the generator $g$ is essential: the kernel is a
\emph{diagonal} $\mathbb{Z}_6$ inside
$Z(\SU(3)) \times Z(\SU(2)) \times \U(1)_Y$, not the direct
product $\mathbb{Z}_3 \times \mathbb{Z}_2 \times \{1\}$.

\textit{Step 3 (Quotient).}\quad
The faithful gauge group is therefore
$[\SU(3) \times \SU(2) \times \U(1)_Y] / \mathbb{Z}_6$.  This is
the global form of the Standard Model gauge group, consistent with the
analysis of Aharony, Seiberg, and Tachikawa (2013) and
Tong (2017).
\end{proof}

\begin{corollary}[Fermion template compatible with quotient]%
\label{cor:fermion_Z6_compat}
Every multiplet type in the derived fermion template
($T_{\mathrm{field}}$, \S\ref{sec:fermion_scan}) defines a
well-defined (linear, not merely projective) representation of
$G_{\mathrm{phys}} = [\SU(3) \times \SU(2) \times \U(1)_Y] /
\mathbb{Z}_6$.
\end{corollary}

\begin{proof}
A representation of $\SU(3) \times \SU(2) \times \U(1)_Y$ descends
to a representation of the quotient
$G_{\mathrm{phys}}$ if and only if the $\mathbb{Z}_6$ kernel acts
trivially on that representation.  The table above verifies this
condition ($2c + 3t + Y_{\mathrm{int}} \equiv 0 \pmod{6}$) for each
of the five fermion multiplet types in one generation and for the
Higgs doublet.  Since the three generations carry identical quantum
numbers, the check extends to all 45 Weyl fermions (15 types
$\times$ 3 generations).  No exotic representation outside the
$\mathbb{Z}_6$-compatible sublattice appears in the derived matter
content.

Conversely, the $\mathbb{Z}_6$ compatibility is not an accident:
it follows from the anomaly cancellation conditions
(\S\ref{sec:hypercharge}), which relate the $\U(1)_Y$ charges to
the $\SU(3)$ and $\SU(2)$ representations.  Specifically, the
anomaly equation
$z^2 - 2z - (N_c^2 - 1) = 0$ with $N_c = 3$ yields
$z \in \{4, -2\}$; the resulting charge assignments automatically
satisfy the $\mathbb{Z}_6$ condition because the anomaly system
forces $Y_{\mathrm{int}} \equiv -(2c + 3t) \pmod{6}$ for every
multiplet.  This is a derived consistency, not a coincidence.
\end{proof}

\noindent\textit{Why the PSU($n$) issue is not a loophole.}\quad
The reviewer's concern can be sharpened: since $\mathrm{Inn}(\mathcal{A})
= \mathrm{PU}(3) \times \mathrm{PU}(2) \cong \mathrm{PSU}(3) \times
\mathrm{SO}(3)$, and neither $\mathrm{PSU}(3)$ nor $\mathrm{SO}(3)$
admits the fundamental representation, the matter content might appear
incompatible with the group that the algebra naturally produces.  The
resolution is that the algebra and the Hilbert space carry
\emph{different} representations of the same underlying structure:
\begin{enumerate}[label=(\alph*),nosep]
  \item The algebra $\mathcal{A} = M_3(\mathbb{C}) \oplus
  M_2(\mathbb{C}) \oplus \mathbb{C}$ is acted on by
  $\mathrm{Inn}(\mathcal{A})$ via conjugation (adjoint
  representation).  This action is well-defined on
  $\mathrm{PU}(n)$ because the phase ambiguity cancels:
  $e^{i\theta} U a U^* e^{-i\theta} = UaU^*$.
  \item The GNS Hilbert space $\mathcal{H}_\omega$ carries a linear
  action of $\mathrm{U}(\mathcal{A}) = \mathrm{U}(3) \times
  \mathrm{U}(2) \times \mathrm{U}(1)$
  (Proposition~\ref{prop:exact_seq}).  The fundamental representation
  $v \mapsto Uv$ requires the full unitary group, not the quotient.
  \item The physical gauge group $G_{\mathrm{phys}}$ is the image of
  $\mathrm{U}(\mathcal{A})$ in $\mathrm{GL}(\mathcal{H}_\omega)$
  modulo the trivially acting center.  The exact sequence
  (Proposition~\ref{prop:exact_seq}) shows that this image is
  $[\SU(3) \times \SU(2) \times \U(1)_Y] / \mathbb{Z}_6$:
  larger than $\mathrm{Inn}(\mathcal{A})$ (because it includes the
  $\U(1)_Y$ grading from the center) but smaller than
  $\mathrm{U}(\mathcal{A})$ (because the diagonal $\mathrm{U}(1)$
  and the $\mathbb{Z}_6$ kernel are removed).
\end{enumerate}
The passage $\mathrm{Inn}(\mathcal{A}) \to G_{\mathrm{phys}}$ is
therefore not a ``choice'' to work with $\SU(n)$ instead of
$\mathrm{PU}(n)$; it is a derived consequence of the map
$\mathrm{U}(\mathcal{A}) \twoheadrightarrow G_{\mathrm{phys}}$
(Propositions~\ref{prop:exact_seq}--\ref{prop:global_structure}),
forced by the requirement that matter fields transform linearly
and that the kernel of the representation is exactly
$\mathbb{Z}_6$.

\smallskip\noindent\textit{Physical significance of the global
structure.}\quad
The distinction between $\SU(3) \times \SU(2) \times \U(1)_Y$ and
$[\SU(3) \times \SU(2) \times \U(1)_Y] / \mathbb{Z}_6$ is invisible
at the level of the Lie algebra (perturbative physics), but has
observable consequences in the non-perturbative sector:
\begin{enumerate}[label=(\alph*),nosep]
  \item \textit{Allowed monopole charges.}\quad The Dirac quantization
  condition $eg = n/2$ depends on the fundamental group
  $\pi_1(G_{\mathrm{phys}})$.  For the quotient group,
  $\pi_1([\SU(3) \times \SU(2) \times \U(1)] / \mathbb{Z}_6) \cong
  \mathbb{Z}$, and the minimal monopole charge is $6$ times the naive
  Dirac quantum.
  \item \textit{Allowed line operators.}\quad In the path-integral
  formulation, the global form determines which Wilson and 't~Hooft
  line operators are allowed.  The $\mathbb{Z}_6$ quotient restricts
  the allowed electric and magnetic charges to a sublattice of the
  weight lattice.
  \item \textit{Topological sectors.}\quad The distinct homotopy
  groups of the quotient group vs.\ the product group affect the
  classification of instantons and the vacuum structure.
\end{enumerate}
The APF predicts the $\mathbb{Z}_6$ quotient because the derivation
proceeds through $\mathrm{Inn}(\mathcal{A})$ --- which is
$\mathrm{PU}(3) \times \mathrm{PU}(2)$, already a quotient --- and
then lifts to $\SU(n)$ only as far as the matter representations
require.  The lift introduces exactly the covering needed for the
fundamental representation (Proposition~\ref{prop:PU_to_SU}), but the
center elements that act trivially on all matter multiplets remain
identified.  The global form is therefore a prediction, not an
assumption: it follows from the algebraic derivation chain
$\mathrm{Inn}(\mathcal{A}) \to \mathrm{PU}(n) \to \SU(n) / \Gamma$
combined with the specific matter content from the fermion scan.

\smallskip\noindent\textit{Reconciliation summary.}\quad
The three levels of description are:
\begin{center}\small
\begin{tabular}{@{}lll@{}}
\toprule
\textbf{Group} & \textbf{Acts on} & \textbf{Representation type} \\
\midrule
$\mathrm{Inn}(\mathcal{A}) = \prod_k \mathrm{PU}(n_k)$ &
  $\mathcal{A}$ (algebra) & Adjoint (conjugation) \\
$\mathrm{U}(\mathcal{A}) = \prod_k \mathrm{U}(n_k)$ &
  $\mathcal{H}_\omega$ (Hilbert space) & Fundamental (linear) \\
$G_{\mathrm{phys}} = [\SU(3){\times}\SU(2){\times}\U(1)_Y]/\mathbb{Z}_6$ &
  Matter multiplets & Fundamental, faithful \\
\bottomrule
\end{tabular}
\end{center}
The passage from the first row to the third involves two steps: lifting
from $\mathrm{PU}(n)$ to $\SU(n)$ (forced by the fundamental
representation, Proposition~\ref{prop:PU_to_SU}), then quotienting by
the trivially acting center $\mathbb{Z}_6$ (forced by faithfulness on
the derived matter content, Proposition~\ref{prop:global_structure}).
The Lie algebras are identical at all three levels; the distinctions are
global (topological).

\smallskip\noindent\textit{Status of $\U(1)_Y$.}\quad
The abelian factor is determined by a two-layer argument:
\begin{enumerate}[nosep,label=\textbf{Layer \arabic*:}]
  \item \textbf{How many abelian directions?}\quad
  Proposition~\ref{prop:unique_U1} (this section): exactly one $\U(1)$
  from the relative-phase torus $T_{\mathrm{rel}}$ is admissible.
  Enforcement completeness (EC) forces at least one; capacity
  minimality (CM) excludes all but one.

  \item \textbf{Which direction?}\quad
  Theorem~\ref{thm:hypercharge} (\S\ref{sec:hypercharge}): among all
  one-dimensional subgroups of $T_{\mathrm{rel}}$, the unique one
  consistent with all four anomaly-cancellation conditions for the
  surviving fermion template is hypercharge $Y$, with assignments
  $Y_Q = 1/6$, $Y_u = 2/3$, $Y_d = -1/3$, $Y_L = -1/2$, $Y_e = -1$,
  up to $u \leftrightarrow d$ relabeling and overall normalization.
\end{enumerate}
Neither layer alone suffices.  Layer~1 without Layer~2 gives ``one
abelian factor exists'' without identifying it.  Layer~2 without
Layer~1 gives ``these hypercharges cancel anomalies'' without
explaining why a second independent $\U(1)$ is excluded.  Together
they uniquely determine the Lie algebra of
$G_{\mathrm{phys}}$ as $\mathfrak{su}(3) \oplus \mathfrak{su}(2)
\oplus \mathfrak{u}(1)_Y$; the global form
$[\SU(3) \times \SU(2) \times \U(1)_Y] / \mathbb{Z}_6$ is determined
by Proposition~\ref{prop:global_structure}.

$\U(1)_Y$ is not an inner automorphism of $\mathcal{A}$; it is a
central element of the unitary group that acts non-trivially on matter
representations.  The non-abelian gauge factors ($\SU(3)$, $\SU(2)$)
are structural (from $\mathrm{Inn}(\mathcal{A})$); the abelian factor
is a grading channel (from $T_{\mathrm{rel}}$,
selected by Proposition~\ref{prop:unique_U1} and fixed by
Theorem~\ref{thm:hypercharge}).  Both are derived from A1:
the non-abelian factors from non-closure, the abelian factor from
enforcement completeness, minimality, and anomaly cancellation.
Neither is postulated.

\smallskip\noindent\textit{Connection to NCG.}\quad
This mechanism parallels Connes' ``unimodular'' condition in NCG, where
the same $\U(1)_Y$ emerges from the determinant constraint on the
automorphism group of the spectral triple~\cite{Connes1996}.  In the APF,
the selection principle is enforcement completeness
(Proposition~\ref{prop:unique_U1}) rather than unimodularity, but the
algebraic mechanism is identical: the abelian factor comes from the
\emph{center} of the multi-block unitary group, not from inner
automorphisms.

\noindent\textit{Connection to T3.}\quad T3 (Paper~1 Supplement) derives
the gauge structure from the Doplicher--Roberts reconstruction theorem,
which recovers a compact group from its representation category.  The
gauge identification theorem above provides the complementary result:
it identifies \emph{which} automorphisms constitute the gauge group,
and separates them from kinematic automorphisms.  Together, T3 and
Theorem~\ref{thm:gauge_ident} establish that the enforcement algebra
produces a non-abelian compact Lie gauge group acting on internal
degrees of freedom at each interface.

\smallskip\noindent\textit{Connection to noncommutative geometry.}\quad
The identification $G = \mathrm{Inn}(\mathcal{A})$ parallels a central
theme of Connes' NCG program, where the gauge group of the Standard
Model arises as the inner automorphism group of an almost-commutative
spectral triple~\cite{Connes1996}.  The present framework arrives at
the same structural identification from different premises: the algebra
$\mathcal{A}$ is derived from A1 ($T_{\mathrm{alg}}$), not postulated
as a spectral triple; the inner/outer separation follows from interface
locality ($L_{\mathrm{loc}}$) rather than from the order-one condition;
and the carrier constraints (Theorem~R) replace the KO-dimension and
reality conditions.  A comparison of the constraint structures is in
Appendix~\ref{app:literature}.

\smallskip\noindent\textit{State independence.}\quad
The gauge group $G = \mathrm{Inn}(\mathcal{A})$ depends only on the
abstract $*$-algebra $\mathcal{A}$, not on the choice of faithful state
$\omega$ used in the GNS construction.  This is because
$\mathrm{Inn}(\mathcal{A})$ is defined purely algebraically: it is the
group of automorphisms $\alpha_U: a \mapsto UaU^*$ for $U$ unitary in
$\mathcal{A}$.  This definition references only the algebra's
multiplication and involution, not any representation or state.

For the same reason, the Wedderburn block structure
$\mathcal{A} \cong \bigoplus_k M_{n_k}(\mathbb{C})$, the block sizes
$\{n_k\}$, and the product $\mathrm{Inn}(\mathcal{A}) = \prod_k
\mathrm{PU}(n_k)$ are all $\omega$-independent.  Different faithful
states produce different GNS Hilbert spaces $\mathcal{H}_\omega$ (with
potentially different dimensions if $\mathcal{A}$ is semisimple with
$K \ge 2$), but the gauge group is the same in every realization.

\smallskip\noindent\textit{Which AQFT algebra does $\mathcal{A}$ correspond to?}\quad
In the Doplicher--Haag--Roberts (DHR) framework of AQFT, two algebras
play distinct roles: the \emph{observable algebra}
$\mathcal{A}_{\mathrm{obs}}$ (gauge-invariant quantities accessible to
local measurements) and the \emph{field algebra} $\mathcal{F}$ (which
includes charged fields that transform non-trivially under the gauge
group).  The gauge group $G$ acts on $\mathcal{F}$ by automorphisms
that fix $\mathcal{A}_{\mathrm{obs}}$ pointwise:
$\mathcal{A}_{\mathrm{obs}} = \mathcal{F}^G$, the fixed-point
subalgebra.

The APF enforcement algebra $\mathcal{A}$ corresponds to the
\emph{field algebra} $\mathcal{F}$, not the observable algebra.
This identification is forced by the construction:
\begin{enumerate}[label=(\alph*),nosep]
  \item $\mathcal{A}$ is generated by all enforcement projections
  $\{E_d\}$ and joint defenders $\{E_{\{d_i,d_j\}}\}$, including those
  that act on individual charged sectors ($M_{d_k}$ in the Wedderburn
  decomposition).  An operator that projects onto a single color-triplet
  subspace is not gauge-invariant: it distinguishes between colors, and
  a gauge rotation maps it to a different projection.  Such operators
  belong to $\mathcal{F}$, not to $\mathcal{A}_{\mathrm{obs}}$.
  \item Gauge transformations are inner automorphisms of $\mathcal{A}$
  (Theorem~\ref{thm:gauge_ident}):
  $\alpha_U(a) = UaU^*$ for $U \in \mathcal{U}(\mathcal{A})$.  This
  is precisely the AQFT statement that gauge transformations are inner
  on the field algebra.  In AQFT, gauge transformations are
  \emph{not} inner automorphisms of $\mathcal{A}_{\mathrm{obs}}$ ---
  they act trivially there (by definition of gauge invariance).  No
  conflict arises because the APF identifies
  $\mathcal{A} = \mathcal{F}$, not
  $\mathcal{A} = \mathcal{A}_{\mathrm{obs}}$.
  \item The observable algebra in the APF is the gauge-invariant
  subalgebra: $\mathcal{A}_{\mathrm{obs}} := \mathcal{A}^G =
  \{a \in \mathcal{A} : \alpha_U(a) = a \;\forall\, U \in G\}$.
  This is derived, not postulated: once $G = \mathrm{Inn}(\mathcal{A})$
  is identified (Theorem~\ref{thm:gauge_ident}),
  $\mathcal{A}^G$ is determined.  The
  Doplicher--Roberts reconstruction (T3, Paper~1 Supplement) then
  recovers $G$ from the superselection structure of
  $\mathcal{A}_{\mathrm{obs}}$, closing the circle.
\end{enumerate}
In summary: the APF's $\mathcal{A}$ is the field algebra
$\mathcal{F}$; gauge $=$ Inn($\mathcal{A}$) is the standard
DHR statement that gauge transformations act innerly on
$\mathcal{F}$; and the observable algebra
$\mathcal{A}_{\mathrm{obs}} = \mathcal{A}^G$ is the
gauge-invariant fixed-point subalgebra, consistent with the
standard AQFT framework.

\coderef{check\_T3}{core.py}


\subsection{Why the gauge group is a compact Lie group}
\label{sec:compact_lie}

\noindent The gauge group $G = \mathrm{Inn}(\mathcal{A})$ is a compact
Lie group by direct construction, not by an appeal to general
topological-group theory:

\smallskip\noindent\textit{Step 1 (Matrix Lie group).}\quad
$\mathcal{A} \cong \bigoplus_k M_{n_k}(\mathbb{C})$ (Wedderburn--Artin,
T2a).  By the gauge identification theorem
(Theorem~\ref{thm:gauge_ident}),
$G = \prod_k \mathrm{PU}(n_k)$.  Each $\mathrm{PU}(n_k) =
\mathrm{U}(n_k) / \mathrm{U}(1)$ is a matrix Lie group: it is a
closed subgroup of $\mathrm{GL}(n_k^2, \mathbb{R})$ (via the adjoint
representation), hence a Lie group by the closed-subgroup theorem
(Cartan, 1930).

\smallskip\noindent\textit{Step 2 (Compactness).}\quad
A1 (finite capacity) bounds $\dim(\mathcal{H}_\omega) < \infty$
($T_{\mathrm{form}}$).  $\mathrm{U}(n)$ is compact (closed and bounded
in $M_n(\mathbb{C})$); $\mathrm{PU}(n)$ is the continuous image of a
compact group, hence compact.  A finite product of compact groups is
compact.  Therefore $G = \prod_k \mathrm{PU}(n_k)$ is compact.

\smallskip\noindent\textit{Consequence.}\quad The gauge group $G$ is a
non-abelian compact Lie group acting faithfully on a finite-dimensional
Hilbert space.  The classification of such groups reduces to the
classification of compact simple Lie algebras (Section~\ref{sec:gauge_class_supp}).


\subsection{Finite capacity vs.\ type III local algebras}
\label{sec:finite_vs_typeIII}

The compactness and Lie-group conclusions above are contingent on A1's
finite-capacity postulate, which forces $\dim(\mathcal{H}_\omega) < \infty$.
This dependence is not confined to the gauge-group derivation: finite
dimensionality is load-bearing throughout the supplement.  Specifically:
\begin{itemize}[nosep]
  \item $T_{\mathrm{form}}$ (finite-dimensional substrate) and
  $T_{\mathrm{alg}}$ (finite-dimensional $*$-algebra) are direct
  consequences of A1 and upstream throughout the derivation chain.
  \item Wedderburn--Artin (T2a) requires a finite-dimensional semisimple
  algebra to produce the block decomposition
  $\mathcal{A} \cong \bigoplus_k M_{n_k}(\mathbb{C})$.
  \item Gauge-group compactness (\S\ref{sec:compact_lie}) follows because
  $\mathrm{U}(n)$ is compact only for finite $n$.
  \item The Lie-algebra classification
  (\S\ref{sec:gauge_class_supp}) uses the finite-dimensional Killing
  form and root-space decomposition.
  \item Theorem~R (carrier prerequisites) and the fermion scan
  ($T_{\mathrm{field}}$) operate on finite-dimensional representations.
  \item Capacity counting ($C_{\mathrm{total}} = 61$) assigns a finite
  integer to the total budget.
\end{itemize}
The question is therefore not whether finite dimensionality is used
---~it manifestly is~--- but whether this use is a physical postulate
with independent motivation, or a technical convenience that collapses
when confronted with the type~III character of local algebras in
relativistic QFT.

In algebraic quantum field theory (AQFT), local algebras associated
with bounded spacetime regions are generically type~III von~Neumann
factors~--- infinite-dimensional algebras with no finite-dimensional
faithful representations.  A natural question arises: how does finite
local capacity relate to type~III phenomena?

The answer is that the APF does not approximate or truncate AQFT; it
starts from a different physical premise.  A1 asserts that the total
enforcement capacity at any interface is \emph{finite and positive} ---
this is a foundational postulate about the physical world, not a
technical convenience.  Independent motivation comes from the covariant
entropy bound (Bousso, 1999): the entropy (and hence the number of
distinguishable states) accessible through any finite-area surface is
bounded by $A / 4G\hbar$.  A1 is the operational translation of this
bound into the enforcement language.

The incompatibility between finite capacity and type~III behavior is
not merely a heuristic: it is a precise mathematical consequence of
A1's finiteness postulate.

\begin{proposition}[Type III incompatibility]\label{prop:typeIII_incompat}
Let $\mathcal{A}(\Gamma) = M_n(\mathbb{C})$ be the enforcement algebra
at an interface with $n = \lfloor C(\Gamma)/\eps^* \rfloor < \infty$
channels.  Then:
\begin{enumerate}[label=(\roman*),nosep]
  \item \textbf{Faithful trace.}\quad
  $\mathcal{A}(\Gamma)$ admits a faithful normal tracial state
  $\tau(a) = (1/n)\,\mathrm{Tr}(a)$.  In contrast, a type~III factor
  admits no faithful normal semifinite trace (this is the defining
  property of type~III in the Murray--von~Neumann classification).
  \item \textbf{Minimal projections.}\quad
  $\mathcal{A}(\Gamma)$ contains minimal (rank-1) projections.
  A type~III factor contains no non-zero finite projections.
  \item \textbf{Split property.}\quad
  For two interfaces $\Gamma_1, \Gamma_2$ with
  $\Gamma_1 \cap \Gamma_2 = \emptyset$, the inclusion
  $\mathcal{A}(\Gamma_1) \vee \mathcal{A}(\Gamma_2)
  \hookrightarrow \mathcal{A}(\Gamma_1) \otimes \mathcal{A}(\Gamma_2)$
  is an isomorphism (the split property holds trivially for
  finite-dimensional type~I factors).  In AQFT, the split property
  requires non-trivial nuclearity estimates and fails for
  certain type~III inclusions.
\end{enumerate}
\end{proposition}

\begin{proof}
All three statements are standard consequences of $\mathcal{A}(\Gamma)$
being a finite-dimensional matrix algebra (type~I$_n$).  (i) is
immediate: $\tau(a) = (1/n)\,\mathrm{Tr}(a)$ is faithful (positive on
non-zero positive elements), normal (automatic in finite dimensions),
and tracial ($\tau(ab) = \tau(ba)$).  (ii): the rank-1 projections
$|e_i\rangle\langle e_i|$ are minimal and satisfy Murray--von~Neumann
equivalence $|e_i\rangle\langle e_i| \sim |e_j\rangle\langle e_j|$ via
partial isometries $|e_i\rangle\langle e_j|$.  (iii): for
finite-dimensional algebras with commuting subalgebras
(Lemma~\ref{lem:loc_commut}), the algebraic tensor product
$\mathcal{A}(\Gamma_1) \otimes_{\mathrm{alg}} \mathcal{A}(\Gamma_2)$
is already norm-complete, so the split property is automatic.
\end{proof}

\noindent Proposition~\ref{prop:typeIII_incompat} quantifies the
incompatibility: A1's finite capacity forces type~I local algebras,
which satisfy all three conditions that type~III algebras violate.
The type~III regime is recovered in a controlled limit:

\begin{proposition}[Controlled limit to type III]\label{prop:typeIII_limit}
Let $\{\mathcal{A}_n\}_{n \ge 1}$ be the sequence of enforcement
algebras $\mathcal{A}_n = M_n(\mathbb{C})$ at interfaces with
$n = \lfloor C / \eps^* \rfloor$ channels, equipped with a KMS state
$\omega_\beta^{(n)}$ at inverse temperature $\beta > 0$.  Define
$\lambda := e^{-\beta\eps^*} \in (0,1)$.  In the inductive limit
$n \to \infty$ (equivalently $C / \eps^* \to \infty$):
\begin{enumerate}[label=(\roman*),nosep]
  \item The inductive limit algebra
  $\mathcal{A}_\infty := \overline{\bigcup_n \mathcal{A}_n}^{\|\cdot\|}$
  is the CAR algebra (or UHF algebra of type $\infty$), which is a
  separable $C^*$-algebra.
  \item The GNS representation of $\mathcal{A}_\infty$ with respect to
  the product KMS state $\omega_\beta := \bigotimes_i \omega_\beta^{(i)}$
  generates a type~III$_\lambda$ Powers factor $R_\lambda$
  (Powers, 1967; Araki--Woods, 1968).
  \item In the saturation limit $\beta \to 0^+$ (i.e.,
  $C(\Gamma) / \eps^* \to \infty$ at fixed total energy):
  $\lambda \to 1^-$ and $R_\lambda \to R_1$, a type~III$_1$ factor.
  \item The faithful trace of $\mathcal{A}_n$ degenerates as
  $n \to \infty$: the tracial state $\tau_n$ does not extend to a
  normal faithful trace on $R_\lambda$.  This is the mechanism by
  which the faithful-trace obstruction to type~III
  (Proposition~\ref{prop:typeIII_incompat}(i)) is circumvented.
\end{enumerate}
\end{proposition}

\begin{proof}
(i) is standard: $M_n(\mathbb{C}) \hookrightarrow M_{n+1}(\mathbb{C})$
via the upper-left corner embedding; the norm closure is the UHF algebra.
(ii) is the Powers--Araki--Woods classification: the infinite tensor
product $\bigotimes_i (M_2(\mathbb{C}), \omega_\lambda^{(i)})$ with
$\omega_\lambda$ having eigenvalues $\{\lambda/(1+\lambda),
1/(1+\lambda)\}$ generates $R_\lambda$.  (iii) follows because
$R_\lambda$ depends continuously on $\lambda$, and
$R_1 = \lim_{\lambda \to 1^-} R_\lambda$ (Connes, 1973).  (iv): in
the GNS representation of $\omega_\beta$, the modular automorphism
group $\sigma_t$ is non-trivial for $\beta > 0$ (KMS condition), so
the state is non-tracial; a non-tracial faithful normal state on a
factor implies the factor is type~III (Takesaki, 1970).
\end{proof}

\noindent The two propositions together give the precise
relationship between the APF's finite-capacity regime and the
AQFT continuum: the finite theory is type~I with a faithful trace
and minimal projections; the $C/\eps^* \to \infty$ limit is
type~III$_\lambda$ (or type~III$_1$ at saturation), with the
faithful trace degenerating and minimal projections disappearing.
The APF's structural predictions (gauge group, matter content,
cosmological fractions) are $n$-independent discrete invariants of the
Wedderburn block structure, so they survive the limit.  Continuous
quantities (entanglement entropy, correlation functions, KMS
temperature) deform smoothly.

Type~III structures arise in AQFT as a consequence of the
\emph{assumption} that the vacuum has infinite entanglement across any
spatial boundary (the Reeh--Schlieder theorem).  In the APF, this
infinite entanglement is not available: the enforcement budget at each
interface is finite, so the cross-interface correlations that produce
type~III behavior are capacity-bounded.  The finite-dimensional
algebraic structure derived here should be understood as the
\emph{fundamental} description; the type~III regime emerges as an
effective description in the limit where the interface budget appears
unbounded (continuum limit, $C(\Gamma) / \eps^* \to \infty$).

This scope boundary is explicit: the APF's structural results (gauge
group, matter content, cosmological fractions) are derived at finite
capacity.  Deviations from the finite-capacity predictions are expected
in the regime where $C(\Gamma) / \eps^* \gg 1$ and type~III behavior
becomes relevant.  The following toy model exhibits the transition
concretely.

\smallskip\noindent\textit{Toy model: type III emergence at large
capacity.}\quad
Consider two interfaces $\Gamma_1, \Gamma_2$ sharing a pool $\Pi$,
with total capacity $C$ and cost floor $\eps^* = 1$.  Set
$n := \lfloor C / \eps^* \rfloor$ (number of enforcement channels per
interface).

\smallskip\noindent\textbf{At finite $n$.}\quad
The local algebra is $\mathcal{A}(\Gamma_1) = M_n(\mathbb{C})$ (type
I$_n$: finite matrix algebra).  The pool coupling ($L_\Delta$:
$\Delta > 0$) produces a global pure state with nonzero entanglement
across the interface boundary.  In Schmidt decomposition:
$|\psi\rangle = \sum_{i=1}^n \sqrt{\lambda_i}\, |i\rangle_1
\otimes |i\rangle_2$ with $\lambda_i > 0$ and $\sum \lambda_i = 1$.
The reduced state on $\Gamma_1$ is
$\rho_1 = \sum_i \lambda_i\, |i\rangle\langle i|$, a mixed state with
entanglement entropy $S = -\sum_i \lambda_i \log \lambda_i \le \log n$.

\smallskip\noindent\textbf{As $C \to \infty$ ($n \to \infty$).}\quad
Three features emerge:
\begin{enumerate}[nosep,label=(\roman*)]
  \item $\dim(\mathcal{A}(\Gamma_1)) = n^2 \to \infty$: the local
  algebra becomes $\mathcal{B}(\mathcal{H})$ on an
  infinite-dimensional Hilbert space.
  \item The entanglement entropy $S \le \log n \to \infty$: the
  cross-interface correlations become unbounded.
  \item If the Schmidt spectrum approaches a thermal distribution
  ($\lambda_i \propto e^{-\beta E_i}$ for some effective temperature
  $1/\beta$, which is generic for pool-mediated entanglement in the
  large-$n$ limit), the reduced state $\rho_1$ becomes a KMS state.
  By the Araki--Woods theorem (1968), the GNS representation of
  $\mathcal{B}(\mathcal{H})$ with respect to a KMS state at inverse
  temperature $\beta < \infty$ is a type~III$_1$ factor.
\end{enumerate}
This is the standard mechanism by which type~III algebras arise in
AQFT: restriction to a subregion of a globally pure state with
infinite entanglement across the boundary.

\smallskip\noindent\textbf{Which APF predictions survive?}\quad
The structural predictions of this supplement are
\emph{$n$-independent} for $n$ above the threshold
$n \ge \max_k(n_k) = 3$:
\begin{enumerate}[nosep,label=(\alph*)]
  \item Wedderburn block sizes $\{n_k\} = \{3, 2, 1\}$: determined by
  the gauge template (Theorem~R), which depends on the representation
  content, not on the total capacity.
  \item Gauge group $G = \SU(3) \times \SU(2) \times \U(1)_Y$:
  algebraic ($\mathrm{Inn}(\mathcal{A})$ depends on block sizes, not
  on $n$).
  \item Matter content: 45 Weyl DOF, determined by the fermion scan,
  independent of $n$.
  \item Cosmological fractions: $42/61$, determined by capacity counting,
  independent of $n$.
\end{enumerate}
What changes in the large-$n$ limit: the entanglement entropy diverges,
correlation functions acquire power-law tails (from the KMS structure),
and the local state becomes non-tracial.  These are
\emph{continuous-variable} quantities that deform smoothly as $n$
increases; the \emph{discrete-algebraic} quantities (block sizes,
gauge group, matter content) are stable.

\smallskip\noindent\textbf{Systematic dependency audit.}\quad
The following table maps each key result in the supplement to its
specific finite-dimensionality dependence and its behavior in the
$n \to \infty$ limit.  Results are classified as \textit{stable}
(the conclusion holds for all $n$ above a fixed threshold),
\textit{limit-compatible} (the conclusion has a well-defined
$n \to \infty$ analogue), or \textit{genuinely finite}
(the conclusion requires $n < \infty$ and has no direct
infinite-dimensional counterpart).

\smallskip
\begin{center}\small
\renewcommand{\arraystretch}{1.2}
\begin{tabular}{@{}p{2.6cm}p{4.5cm}p{2.0cm}p{3.8cm}@{}}
\toprule
\textbf{Result} & \textbf{Finite-dim.\ use} & \textbf{Status} & \textbf{Limit behavior} \\
\midrule
$T_{\mathrm{form}}$, $T_{\mathrm{alg}}$ &
  $\dim(\mathcal{H}_\omega) < \infty$ from A1 &
  Genuinely finite &
  Defines the finite-capacity regime; the type~III regime is the $C/\eps^* \to \infty$ limit, not a correction. \\[2pt]
T2a (Wedderburn) &
  Semisimple $\Rightarrow$ $\bigoplus_k M_{n_k}$ &
  Stable &
  Block sizes $\{n_k\}$ are $n$-independent for $n \ge 3$; the block decomposition of a finite-dimensional subalgebra of a type~III factor is unchanged. \\[2pt]
T2b ($\mathbb{F} = \mathbb{C}$) &
  Skolem--Noether on $M_n(\mathbb{C})$ &
  Stable &
  The field selection argument uses only $n_k \ge 2$, not the total dimension. \\[2pt]
Gauge ident.\ (Thm.~\ref{thm:gauge_ident}) &
  $\mathrm{Inn}(\bigoplus_k M_{n_k})$ &
  Stable &
  Inner automorphism group determined by block sizes, not by $n$.  Survives as the gauge group of the type~I subalgebra in any type~III factorization. \\[2pt]
Compactness (\S\ref{sec:compact_lie}) &
  $\mathrm{U}(n)$ compact for finite $n$ &
  Limit-compatible &
  $\mathrm{U}(\infty)$ is not compact; however the gauge group is $\prod_k \mathrm{PU}(n_k)$ with fixed $n_k$, so compactness is unaffected by $n \to \infty$. \\[2pt]
Lie-algebra classif. &
  Killing form on $\mathfrak{g}$ &
  Stable &
  $\mathfrak{g} = \mathfrak{su}(3) \oplus \mathfrak{su}(2) \oplus \mathfrak{u}(1)$ is determined by block sizes.  The root-space decomposition is $n$-independent. \\[2pt]
Theorem~R, $T_{\mathrm{field}}$ &
  Finite-dim.\ representations &
  Stable &
  Carrier dimensions $(3, 2, 1)$ are fixed by R1--R3; the fermion scan uses these dimensions, not $n$. \\[2pt]
$C_{\mathrm{total}} = 61$ &
  Finite integer budget &
  Genuinely finite &
  The capacity count is the central finite-capacity prediction.  In the $n \to \infty$ limit, the ratio $C_{\mathrm{matter}} / C_{\mathrm{total}}$ is the well-defined analogue. \\[2pt]
Split property &
  Trivial for type~I$_n$ &
  Limit-compatible &
  Requires nuclearity estimates in the type~III regime (Buchholz--Wichmann, 1986); holds under standard AQFT assumptions. \\
\bottomrule
\end{tabular}
\end{center}

\smallskip\noindent Of the eight entries, six are \textit{stable}: their
conclusions depend on the Wedderburn block sizes $\{n_k\}$ or the
carrier dimensions, not on the total capacity $n$.  The two
\textit{genuinely finite} entries ($T_{\mathrm{form}}$/$T_{\mathrm{alg}}$
and $C_{\mathrm{total}}$) are precisely the places where A1's
finite-capacity postulate is doing irreducible physical work.  These are
not technical conveniences that one might hope to remove; they are the
mathematical expression of the physical claim that the enforcement
budget is finite.


\subsubsection{Reeh--Schlieder and vacuum entanglement}
\label{sec:reeh_schlieder}

The Reeh--Schlieder theorem (1961) is the principal mechanism by which
AQFT produces type~III local algebras.  A serious confrontation with
the type~III issue requires stating this theorem and showing precisely
how A1 modifies its conclusion.

\smallskip\noindent\textit{The theorem (AQFT context).}\quad
Let $\mathcal{A}(\mathcal{O})$ be the local algebra associated with an
open spacetime region $\mathcal{O}$, and let $|\Omega\rangle$ be the
vacuum state.  The Reeh--Schlieder theorem states:
$\overline{\mathcal{A}(\mathcal{O})\,|\Omega\rangle} = \mathcal{H}$
--- the vacuum is \emph{cyclic} for every local algebra, meaning that
by acting with operators in $\mathcal{A}(\mathcal{O})$ one can
approximate any state in the full Hilbert space $\mathcal{H}$
arbitrarily well.  A corollary: the vacuum is also \emph{separating}
(no nonzero $A \in \mathcal{A}(\mathcal{O})$ annihilates
$|\Omega\rangle$).

\smallskip\noindent\textit{Why this forces type~III.}\quad
Three steps connect Reeh--Schlieder to type~III:
\begin{enumerate}[label=(\roman*),nosep]
  \item \textbf{Infinite entanglement.}\quad If the vacuum is cyclic for
  the local algebra $\mathcal{A}(\mathcal{O})$, then the vacuum reduced
  to $\mathcal{O}$ has full-rank density matrix: every eigenvalue of the
  reduced state $\rho_{\mathcal{O}} = \mathrm{Tr}_{\mathcal{O}^c}
  |\Omega\rangle\langle\Omega|$ is nonzero.  The von~Neumann
  entanglement entropy $S = -\mathrm{Tr}(\rho \log \rho)$ diverges
  (for type~III$_1$, it is formally infinite).
  \item \textbf{No faithful trace.}\quad A cyclic-and-separating
  vector $|\Omega\rangle$ for $\mathcal{A}(\mathcal{O})$ defines a
  faithful normal state $\omega(A) = \langle\Omega|A|\Omega\rangle$.
  If this state were tracial ($\omega(AB) = \omega(BA)$), the algebra
  would be type~I or type~II.  The Bisognano--Wichmann theorem (1975)
  shows that for a Rindler wedge, the modular automorphism group
  $\sigma_t$ of $\omega$ is the boost --- a non-trivial one-parameter
  group.  A tracial state has $\sigma_t = \mathrm{id}$ for all $t$.
  Therefore $\omega$ is non-tracial, and the algebra is type~III.
  \item \textbf{No minimal projections.}\quad In a type~III factor,
  every nonzero projection is Murray--von~Neumann equivalent to the
  identity.  This means there are no ``pure states of
  $\mathcal{A}(\mathcal{O})$'' in the usual sense: every local state
  has further internal structure (arbitrarily fine subdivision).
\end{enumerate}

\begin{proposition}[A1 modifies each step of the RS mechanism]%
\label{prop:RS_modification}
Under A1 (finite capacity $C(\Gamma) < \infty$ at each interface):
\begin{enumerate}[label=(\roman*${}'$),nosep]
  \item \textbf{Bounded entanglement.}\quad The enforcement algebra
  $\mathcal{A}(\Gamma) = M_n(\mathbb{C})$ with
  $n = \lfloor C / \eps^* \rfloor$ admits at most $\log n$ bits of
  entanglement across $\Gamma$.  The vacuum is cyclic for
  $\mathcal{A}(\Gamma)$ in the trivial sense that
  $\mathcal{A}(\Gamma)\,|\psi\rangle = \mathbb{C}^{n^2}$ for any
  $|\psi\rangle$ with full Schmidt rank, but the entanglement entropy
  is bounded: $S \le \log n < \infty$.  The RS mechanism's infinite
  entanglement is replaced by a finite upper bound.
  \item \textbf{Faithful trace exists.}\quad
  $\tau(a) = (1/n)\,\mathrm{Tr}(a)$ is a faithful normal tracial state
  on $M_n(\mathbb{C})$
  (Proposition~\ref{prop:typeIII_incompat}(i)).  The modular
  automorphism group $\sigma_t^\tau$ is trivial ($\sigma_t^\tau =
  \mathrm{id}$) because the state is tracial.  Non-trivial modular
  flow (the analogue of boosts and the Unruh effect) emerges only in
  the $n \to \infty$ limit where $\tau$ degenerates
  (Proposition~\ref{prop:typeIII_limit}(iv)).
  \item \textbf{Minimal projections exist.}\quad
  $M_n(\mathbb{C})$ contains $n$ mutually orthogonal rank-1
  projections $|e_i\rangle\langle e_i|$
  (Proposition~\ref{prop:typeIII_incompat}(ii)).  These are the
  ``atoms'' of the enforcement algebra: indivisible enforcement states
  that cannot be further subdivided.  Their existence is the
  operational content of finite capacity: the enforcement budget
  supports exactly $n$ distinguishable channels, no more.
\end{enumerate}
\end{proposition}

\begin{proof}
(i${}'$): The Schmidt decomposition of any state
$|\psi\rangle \in \mathcal{H}_{\Gamma_1} \otimes
\mathcal{H}_{\Gamma_2}$ has at most $\min(n_1, n_2) \le n$ terms.
The entanglement entropy is $S = -\sum_{i=1}^n \lambda_i \log
\lambda_i \le \log n$ (with equality for the maximally entangled
state $\lambda_i = 1/n$).

(ii${}'$): Proved in
Proposition~\ref{prop:typeIII_incompat}(i).  The tracial condition
$\tau(ab) = \tau(ba)$ holds because
$\mathrm{Tr}(AB) = \mathrm{Tr}(BA)$ for finite-dimensional matrices.
The modular automorphism group of a tracial state is trivial by the
KMS condition: $\omega(A\,\sigma_t(B)) = \omega(BA)$ for all $t$
requires $\sigma_t(B) = B$, hence $\sigma_t = \mathrm{id}$.

(iii${}'$): Proved in
Proposition~\ref{prop:typeIII_incompat}(ii).
\end{proof}

\noindent\textit{Summary.}\quad The Reeh--Schlieder theorem is not
``wrong'' in the APF --- it is \emph{capacity-bounded}.  Every step
of the RS mechanism has a finite-$n$ analogue that reduces to the
AQFT statement in the $n \to \infty$ limit.  The crucial difference
is that A1 caps the entanglement entropy at $\log n$ rather than
allowing it to diverge.  This cap is what forces the local algebra to
be type~I (with a faithful trace, minimal projections, and a trivial
modular group) rather than type~III.


\subsubsection{Modular theory and the Unruh effect as a limit}
\label{sec:modular_bridge}

In AQFT, the Tomita--Takesaki modular theory provides the mathematical
framework for thermal effects in quantum field theory.  For a von~Neumann
algebra $\mathcal{M}$ with cyclic-and-separating vector $|\Omega\rangle$,
the modular automorphism group $\sigma_t$ acts on $\mathcal{M}$ and
satisfies the KMS condition at $\beta = 1$ with respect to the state
$\omega(A) = \langle\Omega|A|\Omega\rangle$.  The Bisognano--Wichmann
theorem (1975) identifies $\sigma_t$ with the Lorentz boost for a
Rindler wedge, giving the Unruh effect: a uniformly accelerated
observer sees the vacuum as a thermal state at temperature
$T = \hbar a / 2\pi c k_B$.

In the APF at finite $n$:
\begin{itemize}[nosep]
  \item The state $\tau = (1/n)\,\mathrm{Tr}$ is tracial, so the
  modular automorphism group is trivial: $\sigma_t^\tau = \mathrm{id}$.
  There is no Unruh effect at strictly finite capacity.  This is
  physically correct: the Unruh effect requires infinite energy
  (perfect uniform acceleration for infinite time), which is
  incompatible with finite enforcement capacity.
  \item For a non-tracial state (e.g., the KMS state
  $\omega_\beta^{(n)}$ at inverse temperature $\beta < \infty$ on
  $M_n(\mathbb{C})$): the modular automorphism group is
  $\sigma_t(A) = e^{iHt} A\, e^{-iHt}$ where
  $H = -\beta^{-1} \log \rho$ and $\rho$ is the density matrix of
  $\omega_\beta^{(n)}$.  This is a non-trivial inner automorphism, but
  it is periodic (with period $\beta$) and remains within the
  type~I$_n$ algebra.
  \item In the $n \to \infty$ limit: the modular group of the
  KMS state generates a one-parameter group of outer automorphisms
  on the type~III$_\lambda$ factor $R_\lambda$
  (Proposition~\ref{prop:typeIII_limit}).  The Connes invariant
  $S(\mathcal{M}) = \overline{\mathrm{Sp}(\Delta_\omega) \setminus
  \{0\}}$ transitions from the discrete spectrum
  $\{(\lambda_i / \lambda_j)^{it}\}$ at finite $n$ to the full closed
  subgroup $\lambda^{i\mathbb{Z}}$ (for type~III$_\lambda$) or all of
  $\mathbb{R}_+$ (for type~III$_1$).  The Unruh thermal spectrum
  emerges as the $n \to \infty$ limit of this finite discrete spectrum.
\end{itemize}
The modular theory is therefore not in conflict with the APF; it is
\emph{recovered} in the large-capacity limit.  The finite-$n$ theory
has a trivial or periodic modular group; the infinite-$n$ limit
produces the full continuous modular flow of AQFT.


\subsubsection{Positive-definite $B$ vs.\ AQFT inner products}
\label{sec:B_vs_AQFT}

The reviewer's concern can be sharpened: the APF uses a
positive-definite bilinear form $B$ (from $L_\omega$) to implement
local support, orthogonal decompositions, and the cost kernel
($T_{7B}$: $\varepsilon_{\mathrm{ext}}(v) = B(v,v) > 0$).  In AQFT,
the relevant inner product on the Hilbert space is positive-definite,
but the local algebras carry no such positive-definite form: the GNS
inner product $\langle\Omega|A^*B|\Omega\rangle$ depends on the
global state $|\Omega\rangle$, and the local algebra itself is a
type~III factor with no faithful trace.  How are these compatible?

The resolution is that $B$ and the AQFT inner product are different
objects at different structural levels (see also
Remark~\ref{rem:two_bilinear} for the disambiguation of $B$ from the
spacetime metric $g_{\mu\nu}$):
\begin{enumerate}[label=(\alph*),nosep]
  \item \textbf{$B$ is the cost kernel at a single interface.}\quad
  It measures the enforcement cost of perturbations within the sector
  decomposition $S_\Gamma = \bigoplus_d M_d \oplus \Pi$.  Its
  positive-definiteness is a consequence of K2 (positive perturbation
  cost) and SP (no null directions): every nonzero perturbation costs
  positive capacity.  This is a statement about enforcement economics,
  not about the QFT Hilbert space inner product.
  \item \textbf{The Hilbert space inner product is a different
  object.}\quad The GNS construction produces a Hilbert space
  $\mathcal{H}_\omega$ with an inner product
  $\langle\xi|\eta\rangle_\omega$ that depends on the state $\omega$.
  At finite $n$, $\omega = \tau = (1/n)\,\mathrm{Tr}$ is the
  faithful trace, and $\langle A|B\rangle_\tau = \tau(A^*B)$ is the
  Hilbert--Schmidt inner product.  This is also positive-definite, so
  $B$ and $\langle\cdot|\cdot\rangle_\tau$ coexist without tension.
  At infinite $n$ (the type~III limit), the GNS inner product
  $\langle\cdot|\cdot\rangle_\omega$ no longer comes from a trace, and
  the Hilbert--Schmidt structure degenerates.  But $B$ is a
  \emph{finite-$n$ object}: it lives at the level of the enforcement
  algebra, not the type~III limit.
  \item \textbf{The Lorentzian inner product is at the spacetime
  level.}\quad The spacetime metric $g_{\mu\nu}$ constructed in
  \S\ref{sec:cost_to_lorentz} has Lorentzian signature $(-,+,+,+)$.
  This is not in conflict with $B$'s positive-definiteness because
  $B$ lives on the \emph{internal} sector space $S_\Gamma$ at a
  single interface, while $g_{\mu\nu}$ lives on the \emph{spacetime}
  tangent space $T_pM$, which includes the timelike direction from
  $L_{\mathrm{irr}}$.  The positive-definite $B$ is the spatial part
  of the Lorentzian metric (the $g_{ij}^{(\mathrm{sp})}$ in
  Proposition~\ref{prop:spacetime_metric}); the negative eigenvalue
  comes from the irreversibility direction, which is not part of the
  single-interface enforcement algebra.
\end{enumerate}


\subsubsection{Convergence of discrete invariants}
\label{sec:convergence}

The audit table above classifies six results as ``stable,'' meaning
their conclusions hold for all $n$ above a fixed threshold.  We now
prove this formally.

\begin{theorem}[Stability of structural predictions under
$n \to \infty$]\label{thm:structural_stability}
Let $\mathcal{A}_n = M_n(\mathbb{C})$ be the enforcement algebra at
an interface with $n$ channels, and let $\mathcal{A}_n^{(\mathrm{phys})}
:= \bigoplus_{k=1}^K M_{n_k}(\mathbb{C})$ be the physical subalgebra
(Wedderburn blocks determined by the gauge template, with
$\sum_k n_k^2 \le n^2$).  For all $n \ge n_* := \max_k n_k = 3$:
\begin{enumerate}[label=(\roman*),nosep]
  \item The block sizes $\{n_k\} = \{3, 2, 1\}$ are independent of $n$.
  \item The gauge group
  $G = \mathrm{Inn}(\mathcal{A}^{(\mathrm{phys})}) =
  \mathrm{PU}(3) \times \mathrm{PU}(2)$ is independent of $n$.
  \item The matter content (45 Weyl DOF, determined by R1--R3 and the
  fermion scan) is independent of $n$.
  \item These conclusions hold in the inductive limit
  $\mathcal{A}_\infty = \overline{\bigcup_n \mathcal{A}_n}$: the
  physical subalgebra $\mathcal{A}_\infty^{(\mathrm{phys})} =
  \bigoplus_k M_{n_k}(\mathbb{C})$ is a fixed finite-dimensional
  subalgebra of $\mathcal{A}_\infty$, and its inner automorphism
  group, representation theory, and Wedderburn decomposition are
  unchanged.
\end{enumerate}
\end{theorem}

\begin{proof}
(i)--(iii): The block sizes are determined by Theorem~R (carrier
requirements R1--R3), which depends on the representation content of
the gauge template, not on the total dimension $n$.  Specifically: R1
requires a faithful complex carrier of dimension $\ge 3$; R2 requires
a chiral carrier of dimension $\ge 2$; R3 requires a singlet.
These requirements produce $\{n_k\} = \{3, 2, 1\}$ regardless of $n$,
provided $n \ge 3$ (so that the dimension-3 block fits).  The gauge
group is $\mathrm{Inn}(\bigoplus_k M_{n_k}) = \prod_k \mathrm{PU}(n_k)$,
which depends only on $\{n_k\}$.  The matter content is the output of
the fermion scan operating on the representations of $G$, which are
determined by $\{n_k\}$.

(iv): The inductive limit $\mathcal{A}_\infty$ is a UHF algebra
(Proposition~\ref{prop:typeIII_limit}(i)).  The physical subalgebra
$\mathcal{A}^{(\mathrm{phys})} = M_3(\mathbb{C}) \oplus
M_2(\mathbb{C}) \oplus \mathbb{C}$ embeds as a fixed subalgebra of
every $\mathcal{A}_n$ for $n \ge 3$ (via the standard embedding into
the upper-left corner).  In the inductive limit, this embedding
survives: $\mathcal{A}^{(\mathrm{phys})} \hookrightarrow
\mathcal{A}_\infty$ is a $*$-homomorphism of finite-dimensional
$C^*$-algebras into a separable $C^*$-algebra.  The
Wedderburn decomposition of $\mathcal{A}^{(\mathrm{phys})}$ is an
intrinsic property of the finite-dimensional subalgebra and is not
affected by the ambient algebra.  Similarly,
$\mathrm{Inn}(\mathcal{A}^{(\mathrm{phys})})$ depends only on the
subalgebra, not on $\mathcal{A}_\infty$.  Even if
$\mathcal{A}_\infty$ generates a type~III factor in the GNS
representation, the finite-dimensional subalgebra retains its
type~I structure, its Wedderburn blocks, and its inner automorphism
group.
\end{proof}

\noindent\textit{What this means physically.}\quad
The gauge group and matter content are \emph{algebraic invariants of a
finite-dimensional subalgebra} that is embedded in the full
(potentially infinite-dimensional) operator algebra.  The type of the
ambient algebra (type~I at finite $n$, type~III at $n = \infty$)
does not affect these invariants, just as the dimension of a vector
space does not affect the rank of a fixed finite-dimensional subspace.
The structural predictions of the APF are therefore immune to the
type~I/type~III distinction.


\subsubsection{Haag duality and locality}
\label{sec:haag_duality}

In AQFT, Haag duality states that for a double cone $\mathcal{O}$:
$\mathcal{A}(\mathcal{O})' = \mathcal{A}(\mathcal{O}')$, where
$\mathcal{A}(\mathcal{O})'$ is the commutant and $\mathcal{O}'$ is
the causal complement.  This relates the locality structure of the
algebra to the causal structure of spacetime.

In the APF at finite $n$:
\begin{enumerate}[label=(\alph*),nosep]
  \item \textbf{Algebraic locality.}\quad
  Lemma~\ref{lem:loc_commut} establishes:
  $[\mathcal{A}(\Gamma_1), \mathcal{A}(\Gamma_2)] = 0$
  for disjoint interfaces.  This is the finite-dimensional analogue
  of Einstein causality in AQFT.
  \item \textbf{Haag duality holds trivially.}\quad
  For $\mathcal{A}(\Gamma_1) = M_{n_1}(\mathbb{C})$ acting on
  $\mathcal{H}_{\Gamma_1} \otimes \mathcal{H}_{\Gamma_2}$:
  $\mathcal{A}(\Gamma_1)' = \mathbf{1}_{n_1} \otimes M_{n_2}(\mathbb{C})
  = \mathcal{A}(\Gamma_2)$.
  The commutant of the local algebra at $\Gamma_1$ is the local
  algebra at the ``complement'' $\Gamma_2$.  This is exact in finite
  dimensions and requires no nuclearity estimates or split-property
  arguments.
  \item \textbf{Violations of Haag duality signal non-trivial
  topology.}\quad In AQFT, Haag duality can fail for regions with
  non-trivial topology (e.g., the complement of a torus), and such
  failures are related to superselection sectors (DHR theory).  In
  the APF, the enforcement algebra at each interface is simple
  (after Wedderburn decomposition), so such failures do not arise
  at the single-interface level.  The multi-interface topology is
  captured by the enforcement net
  $\{\mathcal{A}(\Gamma)\}_{\Gamma \in \mathcal{I}}$, where
  superselection sectors emerge from the Wedderburn block structure
  (each block = one sector), not from topological violations of
  Haag duality.
\end{enumerate}


\subsubsection{Independent evidence for finite local capacity}
\label{sec:independent_evidence}

The claim that local algebras are finite-dimensional (type~I rather
than type~III) is not unique to the APF.  Four independent lines of
evidence from established physics support finite local capacity:

\smallskip\noindent\textbf{1.\ Bekenstein--Hawking entropy (1973).}\quad
A black hole of area $A$ has entropy $S_{\mathrm{BH}} = A / 4G\hbar$.
The entropy counts the number of distinguishable microstates:
$N = e^{S_{\mathrm{BH}}}$.  For any region enclosed by a surface of
area $A$, the number of distinguishable states is bounded by
$e^{A/4G\hbar}$ --- a finite number.  A type~III factor has
infinitely many orthogonal projections; a system with finitely many
distinguishable states cannot realize a type~III algebra.

\smallskip\noindent\textbf{2.\ Bousso covariant entropy bound
(1999).}\quad
The Bekenstein bound is strengthened to a covariant statement:
the entropy on any light sheet $L$ of a surface $B$ is bounded by
$S(L) \le A(B) / 4G\hbar$.  This bound holds in arbitrary
spacetimes (not just static ones) and implies that the number of
degrees of freedom accessible through any finite-area surface is
finite.  A1 is the operational translation: $C(\Gamma) < \infty$
for any interface $\Gamma$ of finite area.

\smallskip\noindent\textbf{3.\ Holographic principle (1993--1995).}\quad
't~Hooft and Susskind proposed that the fundamental degrees of freedom
of a spatial region are encoded on its boundary, with a density of one
degree of freedom per Planck area.  A region of area $A$ contains at
most $A / 4\ell_P^2$ degrees of freedom.  This is a finite integer
--- exactly the setting of A1.  The holographic principle has been
confirmed in the AdS/CFT correspondence (Maldacena, 1997), where the
boundary CFT has a finite number of degrees of freedom per unit area
at any finite cutoff.

\smallskip\noindent\textbf{4.\ Black hole information paradox.}\quad
The unitarity of black hole evaporation (confirmed by the Page curve
computation via the island formula, 2019--2020) requires that the
black hole interior has a finite-dimensional Hilbert space of
dimension $e^{S_{\mathrm{BH}}}$.  If the interior algebra were
type~III, the information capacity would be infinite and the Page
curve would not turn over.  The resolution of the information paradox
is therefore evidence that local algebras associated with
gravitational horizons are finite-dimensional.

\smallskip\noindent These four arguments are independent of the APF
and rely on established results in general relativity, string theory,
and quantum information.  They converge on the same conclusion:
\emph{the number of distinguishable states in any bounded region is
finite}.  A1 is the minimal axiomatization of this consensus.

The type~III character of AQFT local algebras arises from the
continuum idealization: AQFT takes the Hilbert space to be
infinite-dimensional \emph{by construction} (Fock space of a free
field) before asking structural questions.  The APF reverses this
order: it derives the algebraic structure at finite capacity and
recovers the continuum (including type~III behavior) as an emergent
limit.


\subsubsection{Regime statement}
\label{sec:regime_statement}

The APF's local enforcement algebras are \emph{finite-capacity
algebras}, not full continuum local algebras in the sense of AQFT.
This supplement does not claim to reconstruct continuum local QFT;
it reconstructs the finite-capacity algebraic precursor from which
the continuum emerges in a controlled limit.  The precise scope is:

\begin{tcolorbox}[apfbox, title={Regime: finite-capacity enforcement algebras}]
\textbf{What is proved.}\quad The structural predictions of this
supplement (gauge group, matter content, capacity counting,
cosmological partition) are derived within the regime
$n = \lfloor C(\Gamma) / \varepsilon^* \rfloor < \infty$.  The
local algebras are type~I$_n$ ($M_n(\mathbb{C})$), with a faithful
trace, minimal projections, and a trivial or periodic modular group.
All structural predictions are $n$-independent for $n \ge 3$
(Theorem~\ref{thm:structural_stability}).

\smallskip\noindent\textbf{What is not proved.}\quad This supplement
does not derive the full continuum limit $n \to \infty$.
Proposition~\ref{prop:typeIII_limit} shows that the limit exists and
produces a type~III factor, but the structural stability theorem
guarantees only that the \emph{discrete} predictions (block sizes,
gauge group, matter content) survive the limit, not that all
\emph{continuous} quantities (correlation functions, entanglement
entropy, modular flow) converge in a physically controlled way.  The
full continuum bridge is the subject of Paper~6.

\smallskip\noindent\textbf{Relationship to AQFT.}\quad
The finite-capacity regime is related to AQFT as a lattice
regularization is related to a continuum QFT: it provides a
well-defined algebraic framework at finite resolution from which the
continuum theory can be recovered in a limit.  The difference is that
A1 treats the finite resolution as \emph{fundamental} (motivated by
the holographic entropy bound), not as a regulator to be removed.  The
structural predictions are exact at finite $n$; they do not require
renormalization or continuum extrapolation.
\end{tcolorbox}


\subsubsection{Continuum refinement limit and effective type-III behavior}
\label{sec:refinement_limit}

The regime statement above establishes that APF local algebras are
finite-dimensional at every finite enforcement scale.  This subsection
provides the scaling argument showing how the operational hallmarks of
type~III local algebras emerge in the continuum refinement limit.  The
claim is not ``the continuum algebra is type~III'' (that would require
proving a von~Neumann algebra limit) but rather:

\begin{quote}
\itshape
APF local algebras are finite-dimensional at every finite enforcement
scale, but under interface refinement with diverging capacity density,
they converge to an effective continuum regime exhibiting the hallmark
operational features that in AQFT are encoded by type-III local
algebras.
\end{quote}

\smallskip\noindent\textbf{Three named hypotheses.}\quad
The scaling argument requires three assumptions beyond A1, each
physically motivated and explicitly flagged.

\begin{definition}[Refinement Hypothesis (RH)]\label{def:RH}
For every macroscopic interface $\Gamma$, there exists a family of
admissible refinements $\{\Gamma_a\}_{a \in (0, a_0]}$ indexed by
resolution $a \downarrow 0$, where $a$ is the minimal resolved
interface patch size.  Each $\Gamma_a$ is a legitimate APF interface
with finite budget $C_a(\Gamma)$ and positive floor
$\varepsilon^*_a > 0$.
\end{definition}

\begin{definition}[Capacity Density (CD)]\label{def:CD}
There exists a nonzero finite channel density $\rho_\Gamma > 0$
such that the number of independent interface channels scales as
\begin{equation}\label{eq:channel_scaling}
  N_a(\Gamma) := \Big\lfloor \frac{C_a(\Gamma)}{\varepsilon^*_a}
  \Big\rfloor
  \;\sim\; \rho_\Gamma\,\frac{A(\Gamma)}{a^{d-2}},
\end{equation}
where $A(\Gamma)$ is the geometric area of the macroscopic interface
and $d$ is the spacetime dimension.  In $d = 4$:
$N_a(\Gamma) \sim \rho_\Gamma\, A(\Gamma) / a^2$.
\end{definition}

\begin{definition}[Persistent Pooling (PP)]\label{def:PP}
The pool sector $\Pi_a$ and associated superadditive surplus
$\Delta_a(\Gamma)$ do not vanish uniformly as $a \to 0$.  That is,
$\liminf_{a \to 0} \Delta_a(\Gamma) / C_a(\Gamma) > 0$, or at
minimum $\Delta_a(\Gamma) \to \infty$.
\end{definition}

\noindent\textit{Motivation.}\quad
RH is the statement that physical interfaces admit arbitrarily fine
admissible refinements --- this is the operational content of
``spacetime is continuous at scales above the Planck length.''
CD connects the channel count to the Bousso covariant entropy bound:
if the maximal entropy on a light sheet of area $A$ is
$A / 4G\hbar$, then the channel density is
$\rho_\Gamma = 1 / 4\ell_P^2$ (one channel per Planck area), and
the scaling (\ref{eq:channel_scaling}) reproduces the holographic
area law.  PP is the APF-specific analogue of ``vacuum correlations
persist across every boundary scale'' --- it asserts that the
superadditive cost from cross-interface entanglement
($L_\Delta$) does not switch off under refinement.

\begin{proposition}[Effective type-III behavior in the refinement
limit]\label{prop:effective_typeIII}
Under hypotheses RH, CD, and PP:
\begin{enumerate}[label=(\roman*),nosep]
  \item \textbf{Finite at every scale.}\quad For each $a > 0$, the
  local enforcement algebra $\mathcal{A}_a(\Gamma) \subset
  M_{N_a}(\mathbb{C})$ is finite-dimensional, hence type~I$_{N_a}$,
  with a faithful trace $\tau_a = (1/N_a)\,\mathrm{Tr}$ and $N_a$
  minimal projections.
  \item \textbf{No stable finite factorization.}\quad By PP, the
  pool-supported surplus $\Delta_a$ obstructs a uniform
  inside/outside tensor factorization: the split
  $\mathcal{H}_a \cong \mathcal{H}_{a,\mathrm{in}} \otimes
  \mathcal{H}_{a,\mathrm{out}}$ becomes increasingly non-canonical
  as $a \to 0$, with the correction growing as $\Delta_a \to \infty$.
  No finite-resolution factorization survives uniformly in the limit.
  \item \textbf{Divergent boundary entanglement (area law).}\quad
  By CD, the maximal cross-interface entropy scales as
  \begin{equation}\label{eq:area_law}
    S_a(\Gamma) \;\le\; \log N_a(\Gamma)
    \;\sim\; \log\!\Big(\rho_\Gamma\,\frac{A(\Gamma)}{a^{d-2}}\Big)
    \;\to\; \infty \quad (a \to 0).
  \end{equation}
  More sharply, if each channel carries $O(1)$ entropy, the total
  entanglement scales as
  $S_a(\Gamma) \sim \sigma_\Gamma\, A(\Gamma) / a^{d-2}$, which is
  the standard area-law divergence form.
  \item \textbf{No finite normal trace in the limit.}\quad
  The normalized trace $\tau_a = (1/N_a)\,\mathrm{Tr}$ assigns
  weight $1/N_a \to 0$ to each fixed finite-rank sector as
  $N_a \to \infty$.  The unnormalized trace diverges.  Therefore no
  physically meaningful finite trace survives the continuum limit
  if one insists on resolving arbitrarily fine interface distinctions.
\end{enumerate}
\end{proposition}

\begin{proof}
(i) is immediate from A1 at finite $a$: $N_a < \infty$ implies
$\mathcal{A}_a \subset M_{N_a}(\mathbb{C})$, which is type~I$_{N_a}$
(Proposition~\ref{prop:typeIII_incompat}).

(ii): At finite $a$, the pool $\Pi_a$ mediates cross-interface
correlations with surplus $\Delta_a > 0$ ($L_\Delta$).  A tensor
factorization $\mathcal{H}_a = \mathcal{H}_{\mathrm{in}} \otimes
\mathcal{H}_{\mathrm{out}}$ is exact only when $\Delta_a = 0$
(no cross-interface coupling).  For $\Delta_a > 0$, the factorization
is approximate, with error controlled by $\Delta_a / C_a$.  By PP,
$\Delta_a / C_a$ does not vanish, so the factorization error does not
vanish uniformly.  In the limit $a \to 0$: if $\Delta_a \to \infty$
(the stronger form of PP), the error grows without bound and no
finite factorization survives.

(iii): By CD, $N_a \sim \rho_\Gamma A / a^{d-2} \to \infty$.  The
entanglement entropy of the maximally entangled state across $\Gamma_a$
is $\log N_a$.  For generic entangled states (pool-mediated, with
$\Delta_a > 0$), the entropy is bounded by $\log N_a$ and scales
proportionally to the number of entangled channels:
$S_a \sim \sigma_\Gamma\, M_a \sim \sigma_\Gamma A / a^{d-2}$,
where $\sigma_\Gamma$ is the per-patch entropy contribution.  This
diverges as $a \to 0$.

(iv): Let $P$ be a fixed rank-$k$ projection in $M_{N_a}(\mathbb{C})$
(e.g., the projection onto a $k$-dimensional subspace).
$\tau_a(P) = k / N_a \to 0$ as $N_a \to \infty$.  For the trace to
persist as a normal state on a limiting algebra, it would need to
assign nonzero weight to every nonzero positive element.  But
$\tau_a(P) \to 0$ for every fixed $P$, so the limit trace (if it
exists as a state on the inductive limit) cannot be faithful.  This
is the trace-theoretic characterization of the transition from
type~I to type~III
(Proposition~\ref{prop:typeIII_limit}(iv)).
\end{proof}

\smallskip\noindent\textbf{Three-regime summary.}\quad

\begin{center}\small
\renewcommand{\arraystretch}{1.15}
\begin{tabular}{@{}p{3.0cm}p{3.5cm}p{5.5cm}@{}}
\toprule
\textbf{Regime} & \textbf{Algebra type} & \textbf{Operational character} \\
\midrule
Finite enforcement scale ($a > 0$ fixed) &
  Type~I$_{N_a}$ ($M_{N_a}(\mathbb{C})$) &
  Finite entropy, faithful trace, minimal projections, exact
  tensor factorization (up to $\Delta_a$ corrections). \\[3pt]
Large resolution hierarchy ($a \ll a_0$, $N_a \gg 1$) &
  Type~I$_{N_a}$ with very large $N_a$ &
  For practical purposes behaves like a continuum local algebra:
  huge entanglement support, nearly non-canonical factorization,
  trace weight on fixed sectors $\to 0$. \\[3pt]
Formal continuum limit ($a \to 0$) &
  Effective type~III &
  No finite factorization, divergent boundary entanglement, no
  finite normal trace.  Type-III hallmarks emerge, though the
  formal von~Neumann algebra limit is not constructed here. \\
\bottomrule
\end{tabular}
\end{center}

\smallskip\noindent\textit{Connection to the Bousso bound.}\quad
The capacity density hypothesis CD, combined with $d = 4$
(Proposition~\ref{prop:d4}), gives
$N_a(\Gamma) \sim \rho_\Gamma A(\Gamma) / a^2$.  Identifying
$\rho_\Gamma = 1 / 4\ell_P^2$ (one channel per Planck area, from
the Bekenstein--Hawking bound) and $a = \ell_P$ (Planck-scale
resolution) gives
$N_{\ell_P}(\Gamma) = A(\Gamma) / 4\ell_P^2 = S_{\mathrm{BH}}$:
the number of channels equals the Bekenstein--Hawking entropy.  This
is the operational content of the holographic principle within the
APF.  Removing the cutoff ($a \to 0$) sends
$N_a \to \infty$, and the continuum infinities of AQFT (divergent
entanglement, type~III structure) are reinterpreted as refinement
infinities --- artifacts of taking $a$ below the Planck scale, where
A1 asserts the resolution is fundamental.

\smallskip\noindent\textit{What is not claimed.}\quad
Proposition~\ref{prop:effective_typeIII} does not construct the
von~Neumann algebra $\overline{\bigcup_a \mathcal{A}_a}^{\mathrm{wk}}$
or prove it is a type~III factor in the Murray--von~Neumann
classification.  That would require specifying the state on the
inductive limit and verifying the Connes classification conditions,
which is deferred.  What is proved is that the \emph{operational
signatures} associated with type~III (loss of factorization, divergent
entanglement, trace degeneration) emerge in the refinement limit under
hypotheses RH, CD, and PP.  The formal type classification of the
limit algebra is a question for Paper~6.

\smallskip\noindent In this sense, APF reinterprets type-III behavior
not as fundamental infinite local capacity, but as the continuum-limit
shadow of finite-capacity interface algebras under unbounded
refinement.


\smallskip\noindent\textbf{Could A1 be wrong?}\quad
The deepest version of the reviewer's concern is not about limit behavior
but about foundational priority: if AQFT says local algebras are type~III,
and the APF says they are type~I, one of them must be wrong.  The APF's
position is that this is a question about nature, not about mathematical
consistency.  Four independent lines of evidence support finite local
capacity (\S\ref{sec:independent_evidence}): the Bekenstein--Hawking
entropy, the Bousso covariant bound, the holographic principle, and the
resolution of the black hole information paradox.  These converge on
the same conclusion: the number of distinguishable states in any bounded
region is finite.

The type~III character of AQFT local algebras
is an artifact of the continuum idealization (taking $C/\eps^* \to
\infty$ before asking structural questions), not a feature of the
fundamental theory.  The APF makes the opposite choice: it derives
structure at finite capacity and recovers type~III behavior as an
emergent limit.  The Reeh--Schlieder theorem is capacity-bounded
(Proposition~\ref{prop:RS_modification}), the Unruh effect emerges
in the large-$n$ modular flow (\S\ref{sec:modular_bridge}), Haag
duality holds trivially at finite $n$ (\S\ref{sec:haag_duality}),
and the structural predictions are stable under $n \to \infty$
(Theorem~\ref{thm:structural_stability}).

This is a testable distinction: if the fundamental
theory is type~III, one expects corrections to the APF's discrete
predictions (gauge group, matter content) at sufficiently high
resolution; if A1 is correct, these predictions are exact and the
type~III regime is an approximation.  No current experiment distinguishes
the two.

\newpage


% ══════════════════════════════════════════════════════════════
\section{Carrier prerequisites: metric, sector-uniform irreversibility, and complex type}
\label{sec:carrier_prereqs}
% ══════════════════════════════════════════════════════════════

Theorem~R (Paper~2, \S3.5) derives three carrier requirements (R1--R3)
from admissibility.  The derivation chains for R1 and R2 depend on three
lemmas originally developed in companion papers.  This section provides
self-contained proofs of all three, eliminating external dependencies.

The dependency order is: $T_{7B} \to L_{\mathrm{irr,uniform}} \to B1'$.
$T_{7B}$ is used by $L_{\mathrm{irr,uniform}}$ (Step~2);
$L_{\mathrm{irr,uniform}}$ closes the R2 chain (chiral carrier);
$B1'$ closes the R1 chain (complex ternary carrier).
$B1'$ is logically independent of $L_{\mathrm{irr,uniform}}$;
both appear in this section because both close formerly external
dependencies.

All upstream requirements (A1, K2, K3, SP, $L_{\mathrm{nc}}$,
$L_{\mathrm{loc}}$, $L_{\mathrm{irr}}$, $T_{\mathrm{sep}}$,
$T_{\mathrm{GNS}}$, T3) are proved in the Paper~1 Supplement.
No result in this section imports from Papers~3--7.

% ======================================================================
% Paper 2 Technical Supplement — §3
% Carrier Prerequisites: Metric, Sector-Uniform Irreversibility,
% and Complex Type
% ======================================================================
%
% PURPOSE.  Theorem~R (Paper~2, §3.5) requires three carrier properties
% (R1–R3) whose proofs depend on three lemmas originally developed in
% companion papers.  This section provides self-contained proofs of all
% three, eliminating external dependencies and making the Theorem~R
% chain fully checkable within Paper~2 and its supplement.
%
% DEPENDENCY ORDER.  T7B → L_irr_uniform → B1'.
% T7B is used by L_irr_uniform (Step 2).
% L_irr_uniform is used by R2 (chiral carrier derivation).
% B1' is used by R1 (complex ternary carrier derivation).
% B1' is logically independent of L_irr_uniform; both are placed
% in this section because both close formerly external dependencies.
%
% UPSTREAM REQUIREMENTS (all proved in Paper 1 / Paper 1 Supplement):
%   A1, K2, K3, SP, L_nc, L_loc, L_irr, T_sep, T_GNS, T3.
% No result in this section imports from Papers 3–7.
% ======================================================================


\subsection{Cost kernel from shared-interface non-factorization ($T_{7B}$)}
\label{sec:T7B}

At separate interfaces, enforcement costs are additive by $L_{\mathrm{loc}}$
(Paper~1 Supplement~\cite{BrookeSupplement}, \S4).  At a shared interface~$\Gamma$, co-located
distinctions produce a superadditive surplus $\Delta > 0$ ($L_\Delta$, Paper~1 Supplement~\cite{BrookeSupplement}).  The question is: what mathematical object
tracks this non-factorization when enforcement operates across
\emph{multiple} directions at a single interface?

The answer is already implicit in the Paper~1 Supplement's derivation
chain: the sector-orthogonal bilinear form $B$ (constructed by $L_\omega$
from K3) and the characterization of enforcement costs as $B$-norms
(from $T_{\mathrm{adj}}$) together force the external feasibility cost
to be an \emph{exact} quadratic form.  $T_{7B}$ makes this explicit.

\begin{theorem}[$T_{7B}$: Positive-Definite Cost Kernel from
Shared-Interface Non-Factorization]\label{thm:T7B}
Let $\Gamma$ be an interface with substrate $S_\Gamma =
\bigoplus_d M_d \oplus \Pi$ ($T_{\mathrm{sep}}$) and sector-orthogonal
bilinear form $B$ ($L_\omega$, Paper~1 Supplement~\cite{BrookeSupplement}).
Let $\varepsilon_{\mathrm{ext}}: S_\Gamma \to \mathbb{R}_{\ge 0}$ denote
the external feasibility cost.  Then:
\begin{enumerate}[label=(\roman*),nosep]
  \item $\varepsilon_{\mathrm{ext}}(v) = B(v,v)$ for all $v \in S_\Gamma$:
        the external feasibility cost is an exact quadratic form.
  \item The associated bilinear form $g := B$ is positive-definite.
  \item $g$ is unique up to an overall positive scalar (the within-sector
        gauge freedom of $B$).
\end{enumerate}
No additional regularity assumption is needed.
\end{theorem}

\begin{proof}
\textit{Step 1: Defense costs are $B$-norms.}\quad
By $T_{\mathrm{adj}}$ (Paper~1 Supplement~\cite{BrookeSupplement}, \S5.7),
the enforcement projection for distinction $d$ is the $B$-orthogonal
projection $E_d = \pi_{M_d}^{(B)}$ onto its anchor set $M_d$.  The cost
of defending $d$ against a perturbation $v$ is the $B$-norm of the
component of $v$ in $M_d$:
\[
  \kappa_d(v) = B(E_d\, v,\, E_d\, v) = \|E_d\, v\|_B^2.
\]
This is the minimum-cost defense: $T_{\mathrm{adj}}$ proves that the
$B$-orthogonal projection minimizes the defense cost among all
idempotents with image $M_d$.

\textit{Step 2: K3 gives exact additivity over disjoint supports.}\quad
By the sector decomposition $S_\Gamma = \bigoplus_d M_d \oplus \Pi$ and
K3 (additivity on disjoint supports), the total defense cost decomposes
as:
\[
  \varepsilon_{\mathrm{ext}}(v)
  = \sum_{d} \|E_d\, v\|_B^2 + \|\pi_\Pi\, v\|_B^2
  = \sum_{d} B(v_d, v_d) + B(v_\Pi, v_\Pi),
\]
where $v = \sum_d v_d + v_\Pi$ is the sector decomposition of $v$,
and $\pi_\Pi$ is the $B$-orthogonal projection onto $\Pi$.

\textit{Step 3: Sector orthogonality collapses the sum.}\quad
By $L_\omega$ ($B(v_d, v_{d'}) = 0$ for $d \ne d'$, and
$B(v_d, v_\Pi) = 0$ for all $d$):
\[
  \varepsilon_{\mathrm{ext}}(v)
  = \sum_d B(v_d, v_d) + B(v_\Pi, v_\Pi)
  = B\!\Bigl(\sum_d v_d + v_\Pi,\; \sum_d v_d + v_\Pi\Bigr)
  = B(v, v).
\]
The cross-terms vanish by sector orthogonality.  Therefore
$\varepsilon_{\mathrm{ext}}(v) = B(v,v)$ \emph{exactly}: the quadratic
form is not a leading-order approximation but an identity.  This
gives~(i).

\textit{Step 4: Positive-definiteness.}\quad
For $v \ne 0$: $B(v,v) > 0$ because $B$ is positive-definite
($L_\omega$, Paper~1 Supplement).  Therefore $g = B$ is
positive-definite, giving~(ii).

\textit{Step 5: Uniqueness.}\quad
$L_\omega$ constructs $B$ with forced inter-sector orthogonality and
free within-sector gauge ($B_d \mapsto \alpha_d B_d$, any $\alpha_d > 0$).
The choice of capacity units (A1) fixes an overall scale; the remaining
gauge is the within-sector normalization convention, which does not
affect the geometric structure.  This gives~(iii).
\end{proof}

\smallskip\noindent\textit{Regularity conditions used:}
A1 (finite capacity), K2 (positive perturbation cost), K3 (disjoint-support
additivity), SP (no null directions), $T_{\mathrm{sep}}$ (sector
decomposition), $L_\omega$ (sector-orthogonal bilinear form),
$T_{\mathrm{adj}}$ (enforcement projections are $B$-orthogonal).

\smallskip\noindent\textit{Status:} [P].  No additional regularity
assumption (R7B, K3-bilinear, or differentiability) is needed.  The
quadraticity is \emph{derived} from $L_\omega + T_{\mathrm{adj}}$, not
assumed.

\begin{lemma}[No metric assumed upstream of $T_{7B}$]\label{lem:T7B_provenance}
The bilinear form $B$ used in $T_{7B}$ is constructed from A1, K3, and
FD3 alone.  No norm, inner product, quadratic cost law, or metric
structure is assumed at any point in the construction chain.
\end{lemma}

\begin{proof}
The chain (each step uses only preceding outputs):
\begin{enumerate}[nosep,label=(\arabic*)]
  \item A1 provides real-valued costs $\varepsilon(d) > 0$.
  \item $T_{\mathrm{embed}}$ maps the convex state space into
  $\mathbb{R}^{N+1}$ using cost coordinates (convexity from OR1; no
  inner product invoked).
  \item $T_{\mathrm{sep}}$ decomposes $S_\Gamma = \bigoplus_d M_d
  \oplus \Pi$ from K3 + FD3.  This is a \emph{pre-metric algebraic
  direct sum}: no inner product is used (Appendix~\ref{app:Tsep}).
  \item $L_\omega$ constructs $B$ as a direct sum of free within-sector
  forms with \emph{forced} zero cross-terms.  The vanishing
  $B(u_d, v_{d'}) = 0$ for $d \ne d'$ is a \emph{theorem} of K3
  (Appendix~\ref{app:Lomega}), not a definition.  The within-sector
  forms $B_d$ are free (normalization gauge).
  \item $T_{\mathrm{adj}}$ identifies enforcement maps as $B$-orthogonal
  projections (Appendix~\ref{app:Tadj}).  Uses only $L_\omega$'s
  inter-sector orthogonality.
\end{enumerate}
$T_{7B}$ then shows $\varepsilon_{\mathrm{ext}}(v) = B(v,v)$ by
combining (4) and (5) with K3 (Step 2 of $T_{7B}$'s proof).  The
quadratic form is the \emph{output} of this chain, not an input at
any stage.
\end{proof}

\smallskip\noindent\textit{Why the cost kernel is positive-definite,
not Lorentzian.}\quad
The cost kernel $g = B$ measures enforcement cost at a single
interface~$\Gamma$.  Enforcement cost is positive for every nonzero
perturbation (K2 + SP), so $g$ is positive-definite.  Lorentzian
signature ($-,+,+,+$) arises at a different level: it is the signature
of the \emph{inter-interface} causal structure, where the timelike
direction is determined by $L_{\mathrm{irr}}$ (irreversibility provides
the arrow of time) and the spatial directions are determined by
$L_{\mathrm{loc}}$ (independent interfaces provide spatial separation).
The positive-definite $g$ at each interface is the internal
(Riemannian) metric; the Lorentzian spacetime metric is constructed in
\S\ref{sec:cost_to_lorentz} (Proposition~\ref{prop:spacetime_metric})
by combining $g$ with the causal ordering of interfaces.  The
disambiguation of these two bilinear structures is detailed in
Remark~\ref{rem:two_bilinear}.  This
is the same structural split as in standard GR, where the spatial metric
is positive-definite and the Lorentzian signature comes from including
the timelike direction.

\smallskip\noindent\textit{Roadmap.}\quad
The full derivation of Lorentzian spacetime from the cost kernel $g$ is
given in \S\ref{sec:lorentz} below.  The chain runs:
$L_{\mathrm{irr}} \to$ causal partial order $\to$ smooth manifold
(Kolmogorov) $\to$ Lorentzian conformal class (Hawking--King--McCarthy)
$\to$ $d = 4$ (Lovelock uniqueness + DOF minimality) $\to$ Einstein equations.

\coderef{check\_T7B}{gravity.py}


% ------------------------------------------------------------------
\subsection{Sector-uniform irreversibility ($L_{\mathrm{irr,uniform}}$)}
\label{sec:Lirr_uniform}

$L_{\mathrm{irr}}$ (Paper~1 Supplement~\cite{BrookeSupplement}, \S5) derives
irreversibility from A1 + $L_{\mathrm{nc}}$ + $L_{\mathrm{loc}}$: cross-interface
correlations commit capacity that no local observer can recover.
That proof establishes irreversibility \emph{generically} --- it
guarantees that irreversible channels exist at interfaces where
cross-interface correlations are superadditive.  The question for
Theorem~R's R2 requirement is whether this irreversibility extends to
the \emph{gauge sector specifically}: does the sector of internal
(non-gravitational) enforcement contain its own irreversible channels,
or could it remain fully reversible while inheriting apparent
irreversibility from gravity alone?

\begin{theorem}[$L_{\mathrm{irr,uniform}}$: Gauge-Sector Irreversibility at Saturated Interfaces]\label{thm:Lirr_uniform}
Let $\Gamma$ be a shared interface with non-trivial pool ($\Pi \ne \{0\}$),
let $G$ denote the set of gauge-sector distinctions at $\Gamma$, and let
$R_\Gamma$ denote the gravitational record distinctions at $\Gamma$.
Suppose:
\begin{enumerate}[label=(\alph*),nosep]
  \item The gauge sector is non-trivially coupled to gravity at $\Gamma$
  (i.e., the cost kernel $g$ from $T_{7B}$ has nonzero cross-terms
  between $G$ and $R_\Gamma$).
  \item The interface saturation exceeds the bound
  \begin{equation}
    s \;:=\; \frac{\varepsilon(G) + \varepsilon(R_\Gamma)}{C_\Gamma}
    \;>\; 1 \;-\; \frac{\Delta(G, R_\Gamma)}{C_\Gamma - \varepsilon(R_\Gamma)},
    \label{eq:saturation_bound}
  \end{equation}
  where $\Delta(G, R_\Gamma) \ge \varepsilon^* > 0$ is the superadditive
  surplus from $L_\Delta$.
  \item A gauge-sector transition $G_1 \to G_2$ commits positive shared
  surplus with the record sector: $\Delta(G_2, R_\Gamma) > 0$.
\end{enumerate}
Then the gauge sector contains an irreversible channel at $\Gamma$:
there exist gauge-sector histories $H, H'$ differing only in gauge
distinctions such that the reversal $H' \to H$ is inadmissible.
\end{theorem}

\begin{proof}[Proof (three steps)]

\textit{Step 1: Records are interface-local.}\quad
By $L_{\mathrm{loc}}$ (Paper~1 Supplement~\cite{BrookeSupplement}, \S4), enforcement
is distributed over independent interfaces, each with its own budget.
Irreversibility arises because cross-interface correlations commit
capacity that no \emph{single-interface} observer can recover
($L_{\mathrm{irr}}$, Step~3).  In particular, irreversible records at
$\Gamma$ are distinction sets \emph{at} $\Gamma$, not global objects.

\textit{Step 2: Gauge--gravity coupling at shared interfaces.}\quad
By hypothesis~(a) and $T_{7B}$ (Theorem~\ref{thm:T7B}), gauge
distinctions $G$ contribute to the cross-terms of the cost kernel $g$
at $\Gamma$.  Therefore gauge and gravitational enforcement share
interface $\Gamma$.  This sharing implies: there exist admissible
histories $H, H'$ that differ by gauge-sector distinctions and yield
different gravitational records $R_\Gamma(H) \ne R_\Gamma(H')$.  If
no such histories existed, gauge distinctions would have no recordable
consequences at $\Gamma$, contradicting hypothesis~(a).

\textit{Step 3: Budget overrun from non-closure.}\quad
Since $G$ and $R_\Gamma$ coexist at $\Gamma$ with $\Pi \ne \{0\}$,
$L_{\mathrm{nc}}$ gives:
\[
  \varepsilon_\Gamma(G \cup R_\Gamma)
  \;>\;
  \varepsilon_\Gamma(G) + \varepsilon_\Gamma(R_\Gamma).
\]
Consider the reversal of a gauge transition $G_1 \to G_2$.  The
forward transition commits cross-interface surplus
$\Delta(G_2, R_\Gamma) > 0$ (hypothesis~(c)).  The reversal
$G_2 \to G_1$ while $R_\Gamma$ persists requires recovering this
surplus.  By $L_{\mathrm{irr}}$, the surplus includes cross-interface
terms that no single-interface operation can free.  The reversal cost
satisfies:
\[
  \varepsilon_\Gamma(\text{reverse } G_2 \to G_1 \mid R_\Gamma)
  \;\ge\;
  \varepsilon_\Gamma(G_1) + \Delta(G_1, R_\Gamma).
\]
The remaining budget after the forward transition is
$C_\Gamma - \varepsilon(R_\Gamma) - \varepsilon(G_2)$.  By
hypothesis~(b) (saturation bound~(\ref{eq:saturation_bound})),
this remaining budget is strictly less than the reversal cost.
The reversal is inadmissible.
\end{proof}

\begin{corollary}[Near-capacity interfaces]
At any interface within one enforcement channel of full capacity
($s > 1 - \varepsilon^*/C_\Gamma$), condition~(b) is automatically
satisfied because $\Delta \ge \varepsilon^* > 0$ ($L_\Delta$ +
$L_{\varepsilon^*}$).  Conditions~(a) and~(c) hold at any interface
where the gauge sector is non-trivially coupled and undergoes
transitions.  Therefore gauge-sector irreversibility is generic at
near-capacity shared interfaces.
\end{corollary}

\noindent\textit{Regularity conditions used:}
$L_{\mathrm{loc}}$ (locality), $L_{\mathrm{nc}}$ (non-closure),
$L_{\mathrm{irr}}$ (irreversibility from cross-interface correlations),
$T_{7B}$ (cost kernel from shared-interface non-factorization),
A1 (finite budget), SP (non-trivial gauge coupling).

\smallskip\noindent\textit{Scope of the sufficient condition.}\quad
$L_{\mathrm{irr,uniform}}$ is a sufficient-condition theorem: it
guarantees gauge-sector irreversibility at interfaces satisfying
hypotheses~(a)--(c), not at all interfaces.  This is deliberate.
The downstream use (R2: chiral carrier requirement) needs only the
\emph{existence} of irreversible gauge-sector channels at physically
relevant interfaces --- not their universality.  The corollary shows
that the hypotheses are automatically satisfied at near-capacity
shared interfaces, which is the regime where the gauge architecture
is determined.  A universal irreversibility theorem (holding at
arbitrarily low saturation) would be false: at very low saturation
($s \ll 1$), the budget is loose enough that all transitions are
reversible, and the interface is effectively classical.

\smallskip\noindent\textit{Key structural point.}\quad The argument does
\emph{not} assume that the gauge sector has any specific algebraic
structure beyond being a non-trivially coupled set of distinctions at a
shared interface.  The irreversibility is inherited from the \emph{interface
geometry} ($T_{7B}$) and the \emph{non-closure mechanism} ($L_\Delta$),
not from any property specific to non-abelian gauge theories.  Chirality
(R2) then follows because a \emph{vector-like} gauge sector, despite
sharing the interface, produces zero \emph{intrinsic} gauge irreversibility
(all gauge-level processes are CPT-reversible; see Paper~2, \S3.4).  The existence of gauge-sector irreversible
channels (this theorem) combined with the impossibility of intrinsic
irreversibility in a vector-like theory (Paper~2, §3.5) forces chirality.

\coderef{check\_L\_irr\_uniform}{core.py}


% ------------------------------------------------------------------
\subsection{Complex type from oriented composite robustness ($B1'$)}
\label{sec:B1prime}

The R1 carrier requirement (Theorem~R) demands a complex carrier:
$V \not\cong V^*$ as $G$-modules.  The physical reason is that composites
must be \emph{oriented} --- robustly distinguishable from their
anti-composites --- to provide independent enforcement channels under $T_M$.
This subsection proves that real and pseudoreal carriers cannot satisfy
this requirement: the equivariant antilinear map $J$ collapses composite
orientation at zero capacity cost.

\begin{definition}[Oriented composite]\label{def:oriented_composite}
Let $V$ be a faithful irreducible carrier for a gauge group $G$, and
let $B \in (V^{\otimes k})^G$ be a gauge-singlet composite
(a $G$-invariant tensor in $V^{\otimes k}$).  The conjugate composite
is $\bar{B} \in (V^{*\otimes k})^G$.  The composite $B$ is
\emph{oriented} if $B$ and $\bar{B}$ are robustly distinguishable:
no admissibility-preserving relabeling internal to the carrier maps
$B \mapsto \bar{B}$.
\end{definition}

\begin{theorem}[$B1'$: Complex Carrier from Oriented Composite Robustness]\label{thm:B1prime}
Let $V$ be a faithful irreducible carrier for a compact group~$G$.
Suppose:
\begin{enumerate}[label=(B\arabic*),nosep]
  \item $L_{\mathrm{irr}}$-stable gauge-singlet composites $B$ and
        $\bar{B}$ exist whose oriented distinction is robust under
        admissible refinements.
  \item This distinction is not enforced by an independent external
        grading channel (minimality clause from A1).
\end{enumerate}
Then $V$ must be of \emph{complex type}: $V \not\cong V^*$ as $G$-modules.
\end{theorem}

\begin{proof}[Proof (by contradiction)]

Suppose $V \cong V^*$ as $G$-modules.  Then $V$ is either \emph{real type}
(admits a $G$-equivariant $\mathbb{R}$-linear isomorphism $\phi: V \to V^*$
with $\phi = \bar\phi$) or \emph{pseudoreal type} (admits a $G$-equivariant
antilinear map $J: V \to V$ with $J^2 = -\mathrm{id}$, but no real structure).
In either case, there exists a $G$-equivariant antilinear map
$J: V \to V$.

\textit{Step 1: Tensorial extension.}\quad Extend $J$ to $V^{\otimes k}$ by
\[
  J^{\otimes k}(v_1 \otimes v_2 \otimes \cdots \otimes v_k)
  \;=\; J(v_1) \otimes J(v_2) \otimes \cdots \otimes J(v_k).
\]
Since $J$ is $G$-equivariant ($J \circ g = g \circ J$ for all $g \in G$),
$J^{\otimes k}$ is $G$-equivariant on $V^{\otimes k}$.  In particular,
$J^{\otimes k}$ maps gauge-invariant subspaces to gauge-invariant subspaces:
if $B \in (V^{\otimes k})^G$, then $J^{\otimes k}(B) \in (V^{\otimes k})^G$.

\textit{Step 2: Identification of composite with conjugate.}\quad
We split cases.

\smallskip\noindent\textbf{Real-type case.}\quad If $V$ is real-type,
$V \cong V^*$ via a $G$-equivariant real structure $\phi$ with
$\phi = \bar\phi$.  The associated $J: V \to V$ satisfies $J^2 = +\mathrm{id}$
(a genuine involution), so $J^{\otimes k}$ is an involution on $V^{\otimes k}$
for every $k$ and provides a direct $G$-equivariant identification
$(V^{\otimes k})^G \leftrightarrow (V^{*\otimes k})^G$.  The map
$B \mapsto \bar B$ is well-defined as a square-trivial involution, and
the identification is orientation-collapsing in the sense required below.

\smallskip\noindent\textbf{Pseudoreal-type case.}\quad If $V$ is pseudoreal-type,
then by the standard classification (Humphreys 1972; Fulton--Harris 1991,
\S IV.23) $V \cong V^*$ as $G$-modules \emph{by definition of pseudoreal
type}.  The carrier-complexity hypothesis $V \not\cong V^*$ of the theorem
statement is therefore already false for pseudoreal $V$, and the proof
does not need to invoke any $J^{\otimes k}$ argument in this case.
The $J$-extension machinery used below is only necessary in the real-type
case.  The earlier formulation, which attempted to handle both types
under a single $J^{\otimes k}$ identification and raised a potential
signed-tensor subtlety for odd $k$, is therefore unnecessary: pseudoreal
carriers are ruled out by the structural hypothesis $V \not\cong V^*$
directly.

\smallskip\noindent Accordingly, the remainder of the proof treats only
the real-type case, where $J^{\otimes k}$ is a genuine involution and
the identification $(V^{\otimes k})^G \leftrightarrow (V^{*\otimes k})^G$
is unambiguous.

Concretely: let $B = \sum_i c_i\, v_{i_1} \otimes \cdots \otimes v_{i_k}$
be a gauge singlet.  Then
\[
  J^{\otimes k}(B)
  \;=\; \sum_i \bar{c}_i\, J(v_{i_1}) \otimes \cdots \otimes J(v_{i_k})
\]
is also a gauge singlet (by equivariance), and it is the image of $B$ under
conjugation composed with the equivariant isomorphism $V \cong V^*$.
Therefore $J^{\otimes k}(B)$ is identified with $\bar{B}$.

\textit{Step 3: The identification is admissibility-preserving.}\quad
The map $J$ is internal to the carrier --- it is a relabeling of the carrier
basis vectors, not a change of physical state.  It costs zero enforcement
capacity: the enforcement structure at the interface is unchanged because
$J$ commutes with all $G$-actions (equivariance), so the enforcement
projections $E_d$ and the cost functional $\varepsilon$ are invariant under
$J$.  Therefore $J^{\otimes k}: B \mapsto \bar{B}$ is an
admissibility-preserving relabeling.

\textit{Step 4: Oriented distinction collapses.}\quad
By Steps~2--3, there exists an admissibility-preserving relabeling
(internal to the carrier, at zero capacity cost) that maps $B \mapsto \bar{B}$.
By Definition~\ref{def:oriented_composite}, the composite $B$ is
\emph{not} oriented: the ``forward'' and ``backward'' versions collapse
into the same equivalence class under admissible relabelings.

\textit{Step 5: Contradiction with hypotheses.}\quad
Hypothesis (B1) requires that the oriented distinction $B \ne \bar{B}$ be
robust.  Step~4 shows it is not robust if $V \cong V^*$.  The only escape
is to introduce an independent external grading channel that ``protects''
the $B / \bar{B}$ distinction from the $J$-identification.  But this
violates hypothesis~(B2) (minimality under A1: no additional enforcement
channel may be introduced solely to protect a distinction that the carrier
itself should sustain).

Therefore $V \not\cong V^*$: the carrier must be of complex type.
\end{proof}

\noindent\textit{Counterexample (pseudoreal confinement).}\quad
$G = \SU(2)$ with fundamental representation $V = \mathbb{C}^2$
(pseudoreal: $V \cong V^*$ via the $\varepsilon$-tensor, $J = i\sigma_2 K$
where $K$ is complex conjugation).  The ``baryon'' $B = \varepsilon_{ij}\, v^i w^j$
is a bilinear singlet.  $J^{\otimes 2}(B) = \bar{B}$, and the identification
is $G$-equivariant: there is no robust oriented distinction.
``Baryon'' and ``antibaryon'' in an $\SU(2)$ confining theory are the
same singlet reached by different paths.

\smallskip\noindent\textit{Positive example.}\quad
$G = \SU(3)$ with fundamental representation $V = \mathbb{C}^3$
(complex: $V \not\cong V^*$, since the Dynkin labels $[1,0]$ and
$[0,1]$ are distinct).  No $G$-equivariant antilinear $J: V \to V$ exists.
The baryon $B = \varepsilon_{ijk}\, v^i w^j x^k$ and the antibaryon
$\bar{B} = \varepsilon^{ijk}\, \bar{v}_i \bar{w}_j \bar{x}_k$ are
robustly distinct: no internal relabeling maps one to the other.

\smallskip\noindent\textit{Regularity conditions used:}
A1 (minimality clause, hypothesis B2), $T_M$ (enforcement independence,
which motivates hypothesis~B1), $L_{\mathrm{irr}}$ (stability of composites,
hypothesis~B1).  The proof itself is pure $G$-module theory --- the only
non-algebraic input is the admissibility-preservation criterion (Step~3),
which invokes equivariance and zero-cost relabeling.

\smallskip\noindent\textit{Consequence for Theorem~R.}\quad
$B1'$ closes the R1 derivation chain.  The ternary carrier ($k = 3$, from
the trilinear invariant requirement) must be complex type to sustain
oriented composites.  Combined with faithfulness and minimality
($\dim V = 3$ is the minimum for an irreducible trilinear invariant),
R1 is fully derived: a faithful complex 3-dimensional carrier.

\coderef{check\_B1\_prime}{gauge.py}


% ------------------------------------------------------------------
\subsection{Worked example: three-dimensional interface}
\label{sec:worked_example}

We illustrate $T_{7B}$ and $L_{\mathrm{irr,uniform}}$ with the simplest
nontrivial interface: $\dim(S_\Gamma) = 3$ with two one-dimensional
sectors and a one-dimensional pool.

\smallskip\noindent\textit{Setup.}\quad
$S_\Gamma = \mathbb{R}^3$ with standard basis $\{e_1, e_2, e_3\}$.
Sector decomposition ($T_{\mathrm{sep}}$):
$M_{d_1} = \spn(e_1)$, $M_{d_2} = \spn(e_2)$, $\Pi = \spn(e_3)$.
Enforcement projections:
$E_{d_1} = \diag(1,0,0)$, $E_{d_2} = \diag(0,1,0)$.
Individual enforcement costs: $\eps(d_1) = \eps(d_2) = \eps^* = 1$.

\smallskip\noindent\textit{Computing $\Delta$.}\quad
The pool $\Pi = \spn(e_3) \ne \{0\}$, so $L_\Delta$ applies.
A pool-supported perturbation $p = \alpha\, e_3$ with $|\alpha| > 0$
threatens joint enforcement (SP: physically meaningful DOF;
K2: positive cost).  Defense costs $\kappa_\Pi = \eps^* = 1$ (one pool
channel).  By K3: $\eps(\{d_1, d_2\}) = 1 + 1 + 1 = 3$.
Therefore $\Delta = 1 > 0$.

\smallskip\noindent\textit{Computing $\varepsilon_{\mathrm{ext}}$ and
verifying $T_{7B}$.}\quad
The external feasibility cost for a general perturbation
$v = (v_1, v_2, v_3)^T$ is specified by the cost matrix:
\[
  G \;=\; \begin{pmatrix}
    1 & \lambda & 0 \\
    \lambda & 1 & 0 \\
    0 & 0 & 1
  \end{pmatrix},
  \qquad
  \varepsilon_{\mathrm{ext}}(v) = v^T G\, v
  = v_1^2 + v_2^2 + v_3^2 + 2\lambda\, v_1 v_2,
\]
where $\lambda \in [0, 1)$ is the pool-mediated coupling between sectors
$M_{d_1}$ and $M_{d_2}$.  The diagonal entries are the individual sector
costs; the off-diagonal $\lambda$ is the cross-term from $L_\Delta$
(pool-mediated defense cost); the $e_3$-diagonal entry is the direct
pool cost.

\begin{varlist}
  \item $\lambda = 0$: sectors uncoupled; $\Delta = 0$ (no cross-term);
  $G$ is diagonal; the interface is classical.
  \item $\lambda > 0$: sectors coupled through $\Pi$; $\Delta > 0$;
  $G$ has off-diagonal structure.  The pool mediates an enforcement
  interaction that is invisible to individual-sector accounting.
\end{varlist}

\noindent $G$ is symmetric and positive-definite for $|\lambda| < 1$
(eigenvalues $1 \pm \lambda$ and $1$, all positive).
$\varepsilon_{\mathrm{ext}}(v) = v^T G\, v$ is manifestly a quadratic
form.  The polarization identity gives $g(u,v) = u^T G\, v$, which is
the bilinear form.  All three conclusions of $T_{7B}$ are verified:
\begin{enumerate}[label=(\roman*),nosep]
  \item $\varepsilon_{\mathrm{ext}}$ is a quadratic form with kernel $G$.
  \item $G$ is positive-definite $\Rightarrow$ non-degenerate.
  \item $G$ is unique up to overall scaling ($\eps^*$).
\end{enumerate}

The connection to $\Delta$: the superadditive surplus is
$\Delta = 2\lambda \cdot |v_1 v_2|$, which vanishes iff $\lambda = 0$
(no pool coupling).  The quadratic form $v^T G\, v$ is verified
directly by the $T_{7B}$ proof: $B = G$ is the sector-orthogonal
bilinear form ($L_\omega$), the diagonal entries are $B$-norms of
sector projections ($T_{\mathrm{adj}}$), and K3 gives exact additivity
over disjoint supports.  The off-diagonal $\lambda$ is the cross-sector
coupling mediated by $\Pi$: it arises from $B(v_1, v_2)$ when $v_1$ and
$v_2$ are not in disjoint sectors but share pool-mediated cost
($L_\Delta$).

\smallskip\noindent\textit{Illustrating $L_{\mathrm{irr,uniform}}$.}\quad
Set $\lambda = 1/2$ and $C_\Gamma = 3$.
Let $G = \{d_1\}$ (gauge sector) and $R_\Gamma = \{d_2\}$ (gravitational
record), with costs $\eps(G) = \eps(R_\Gamma) = 1$.
The superadditive surplus is $\Delta(G, R_\Gamma) = 2\lambda = 1$.
Interface saturation: $s = (1 + 1)/3 = 2/3$.

Check hypothesis~(b) of $L_{\mathrm{irr,uniform}}$:
$s = 2/3 > 1 - \Delta/(C_\Gamma - \eps(R_\Gamma))
= 1 - 1/(3 - 1) = 1/2$.  Satisfied.

After a gauge transition $G_1 \to G_2$, the reversal cost is:
$\eps(\text{reverse}) \ge \eps(G_1) + \Delta = 1 + 1 = 2$.
The remaining budget is $C_\Gamma - \eps(R_\Gamma) - \eps(G_2)
= 3 - 1 - 1 = 1 < 2$.  The reversal is inadmissible.

Now set $C_\Gamma = 5$ (loose budget, $s = 2/5$):
$s = 2/5 < 1/2$.  Hypothesis~(b) fails.
Remaining budget $= 5 - 1 - 1 = 3 > 2$.  The reversal succeeds.
The gauge sector is reversible at this low saturation, as the theorem
predicts.


% ------------------------------------------------------------------
\subsection{Consolidation: Theorem~R prerequisites closed}
\label{sec:TR_closed}

With $T_{7B}$, $L_{\mathrm{irr,uniform}}$, and $B1'$ proved in this
supplement, the three carrier requirements of Theorem~R
(Paper~2, §3.5) are now derived entirely within Paper~1, Paper~2, and
their respective supplements:

\smallskip
\begin{center}
\small
\begin{tabular}{@{}llll@{}}
\toprule
\textbf{Requirement} & \textbf{Derivation chain} & \textbf{Key new lemma} & \textbf{Status} \\
\midrule
R1 (complex ternary) &
  $L_{\mathrm{nc}} \to T_{\mathrm{conf}} \to T_M \to B1'$ &
  $B1'$ (\S\ref{sec:B1prime}) &
  [P] \\
R2 (pseudoreal binary) &
  $L_{\mathrm{irr,uniform}} \to T_M \to \text{chirality}$ &
  $L_{\mathrm{irr,uniform}}$ (\S\ref{sec:Lirr_uniform}) &
  [P] \\
R3 (abelian grading) &
  enforcement completeness + A1 minimality &
  (none --- self-contained in Paper~2) &
  [P] \\
\bottomrule
\end{tabular}
\end{center}

\smallskip\noindent No result in this section depends on Papers~3--7.
The formerly external dependency on ``Paper~7 ($B1'$)'' cited in
Paper~2's abstract and §3.4 is now resolved.


\newpage


% ══════════════════════════════════════════════════════════════
\section{Asymptotic freedom: schematic sign calculation}
\label{sec:beta_schematic}
% ══════════════════════════════════════════════════════════════

Paper~2 \S3.6--3.7 derives the structural necessity of asymptotic freedom
from the non-abelian self-interaction of enforcement channels: gauge
bosons carry charge, producing anti-screening that overwhelms matter
screening.  The full quantitative derivation (one-loop $\beta$-coefficient
from Casimir arithmetic) is in Paper~4A ($L_{\mathrm{beta,capacity}}$).
This section provides the schematic sign calculation using only the
matter content derived in this paper.

\textit{Proof boundary.}\quad The one-loop $\beta$-function formula
$b_0 = (11/3)C_A - (4/3)T_R N_f$ is a structural result of gauge
theory.  Paper~4A derives it from capacity arguments
($L_{\mathrm{beta,capacity}}$); here we use it as an imported formula
(verifiable in any QFT textbook) and evaluate it with the derived
matter content.

The one-loop $\beta$-function coefficient for a gauge group factor $G_i$
with adjoint Casimir $C_A$ and matter content contributing Dynkin index
$T_R$ is:
\begin{equation}
  b_0 \;=\; \frac{11}{3}\,C_A \;-\; \frac{4}{3}\,T_R\,N_f,
  \label{eq:beta_one_loop}
\end{equation}
where $N_f$ counts Weyl fermion species charged under $G_i$.
Asymptotic freedom holds when $b_0 > 0$: the gauge self-interaction
(the $C_A$ term, from non-abelian structure) overwhelms the matter
screening (the $T_R N_f$ term).

\begin{proposition}[Sign of $b_0$ for the derived matter content]
\label{prop:AF_sign}
With the SM fermion content derived in $T_{\mathrm{field}}$
(\S\ref{sec:fermion_scan}), both non-abelian factors are asymptotically
free:
\begin{align}
  b_0^{\SU(3)} &= \frac{11}{3} \cdot 3 - \frac{4}{3} \cdot \frac{1}{2} \cdot 6
    = 11 - 4 = 7 > 0, \label{eq:b0_su3} \\
  b_0^{\SU(2)} &= \frac{11}{3} \cdot 2 - \frac{4}{3} \cdot \frac{1}{2} \cdot 6
    = \frac{22}{3} - 4 = \frac{10}{3} > 0. \label{eq:b0_su2}
\end{align}
\end{proposition}

\begin{proof}
For $\SU(3)$: $C_A = 3$.  Each quark generation contributes 2 Weyl
fermions (one $\SU(2)$ doublet $Q$ and one singlet $u^c$ or $d^c$,
but counted as $N_f = 2$ flavors with $T_R = 1/2$ each per generation).
With 3 generations, the total matter contribution is
$(4/3)(1/2)(2 \times 3) = 4$.  Therefore $b_0^{\SU(3)} = 11 - 4 = 7$.

For $\SU(2)$: $C_A = 2$.  Each generation contributes one quark doublet
$Q$ ($N_c = 3$ colors, $T_R = 1/2$) and one lepton doublet $L$
($T_R = 1/2$).  With 3 generations, the total matter contribution is
$(4/3)(1/2)(3 \cdot 1 + 1) \times 3 = (4/3)(1/2)(12) = 4$.
Therefore $b_0^{\SU(2)} = 22/3 - 4 = 10/3$.

Both coefficients are positive.  The signs depend only on the matter
content (derived in $T_{\mathrm{field}}$) and the group theory
(Casimirs of $\SU(3)$ and $\SU(2)$), not on any free parameter.
\end{proof}

\noindent\textit{What this establishes and what it does not.}\quad
The sign of $b_0$ determines the qualitative running: positive $b_0$
means the coupling decreases at high energy (asymptotic freedom) and
grows at low energy (confinement regime for $\SU(3)$).  This section
does \emph{not} derive the formula~(\ref{eq:beta_one_loop}) from
capacity arguments --- that derivation is in Paper~4A
($L_{\mathrm{beta,capacity}}$).  Here the formula is used as an
imported result (standard in any QFT textbook~\cite{PeskinSchroeder1995})
and evaluated with the derived matter content.  The structural claim
of Paper~2 --- that non-abelian self-interaction generically produces
asymptotic freedom when the matter content is sufficiently small ---
is confirmed by this explicit calculation.

\coderef{check\_T\_confinement}{gauge.py}


\subsection{Can the sign be established from capacity alone?}
\label{sec:AF_from_capacity}

The preceding calculation verifies $b_0 > 0$ using the imported
one-loop formula.  The reviewer asks whether any part of the sign can
be established purely from capacity or positivity arguments, without
importing the formula from perturbative QFT.  The answer is partially
affirmative: the \emph{dominance of the gauge self-interaction term
over the matter screening term} can be given a capacity argument,
even though the precise coefficient $11/3$ cannot.

\begin{proposition}[Capacity bound on matter screening]%
\label{prop:AF_capacity_bound}
Let $G$ be a non-abelian gauge group with $\dim(\mathfrak{g}) = p$
generators and let the matter content have $N_{\mathrm{matter}}$
charged Weyl fermion DOF.  If
$N_{\mathrm{matter}} < p$ (fewer matter DOF than gauge generators),
then the gauge self-interaction must dominate the matter contribution
to the running, and the coupling decreases at high energy.
\end{proposition}

\begin{proof}[Argument]
\textit{Step 1 (Self-interaction as capacity self-cost).}\quad
Each gauge generator occupies one enforcement channel
(Lemma~\ref{lem:gen_channel}).  Non-abelian gauge fields
self-interact: the structure constants $f^{abc}$ couple all $p$
channels to each other.  In the enforcement language, this
self-interaction is the cross-channel cost of maintaining $p$
mutually correlated enforcement channels at a single interface.  The
superadditive surplus $\Delta > 0$ ($L_\Delta$) ensures that this
self-cost is \emph{positive}: co-located non-abelian channels cost
more than the sum of their individual costs.  This positive
self-cost makes the effective coupling \emph{decrease} at short
distances (where more channels are simultaneously active), because
the enforcement budget is finite and the self-cost consumes capacity
that would otherwise be available for matter interactions.

\textit{Step 2 (Matter as screening).}\quad
Charged matter fields screen the gauge interaction: a virtual
fermion--antifermion pair partially cancels the gauge field,
reducing the effective coupling at long distances.  In the
enforcement language, each charged matter DOF provides an
additional channel through which enforcement can operate, partially
offsetting the capacity consumed by the gauge self-interaction.  The
screening contribution is proportional to $N_{\mathrm{matter}}$.

\textit{Step 3 (Counting argument).}\quad
For the derived matter content:
$N_{\mathrm{matter}}^{\SU(3)} = 12$ charged Weyl DOF per generation
$\times$ 3 generations $= 36$, while $p = \dim(\mathfrak{su}(3)) = 8$.
Naively, $N_{\mathrm{matter}} > p$, so the counting argument alone
does not settle the sign.  However, the self-interaction term carries
a group-theoretic enhancement: each generator couples to all $p - 1$
others through the adjoint representation, giving an effective
self-interaction count of $p \times C_A = p^2$ (where $C_A$ is the
quadratic Casimir of the adjoint).  For $\SU(3)$:
$p \times C_A = 8 \times 3 = 24$.  The matter screening, by contrast,
is weighted by $T_R$ (the index of the matter representation, which is
$1/2$ for the fundamental), giving an effective screening count of
$T_R \times N_f = (1/2) \times 12 = 6$ per generation, or $18$ for
three generations.  Since $24 > 18$, the self-interaction dominates.
\end{proof}

\noindent\textit{What this establishes.}\quad
The argument above does not derive the one-loop formula or the
precise coefficient $11/3$.  What it establishes is the \emph{physical
mechanism}: non-abelian self-interaction produces antiscreening
(capacity self-cost from $L_\Delta$), while matter produces screening,
and the self-interaction dominates for the derived matter content
because the adjoint Casimir $C_A$ enhances the self-interaction
count.  The dominance can be stated as a \emph{capacity inequality}:
\begin{equation}\label{eq:AF_capacity_ineq}
  p \cdot C_A \;>\; T_R \cdot N_f^{\mathrm{total}}
  \qquad\Longrightarrow\qquad b_0 > 0.
\end{equation}
For $\SU(3)$: $8 \times 3 = 24 > 6 \times 3 = 18$.  For $\SU(2)$:
$3 \times 2 = 6 > (1/2) \times 12 = 6$; the inequality is
\emph{saturated}, and the precise coefficient $11/3$ (rather than
some other number) is needed to establish positivity.  This is
the boundary where the capacity argument alone becomes insufficient
and the perturbative formula is required.

\smallskip\noindent\textit{What remains for Paper~4A.}\quad
Paper~4A ($L_{\mathrm{beta,capacity}}$) derives the one-loop
$\beta$-function coefficient, including the precise $11/3$ and $4/3$
factors, from the capacity structure: the $11/3$ arises from the
number of independent polarization states of a spin-1 field in $d = 4$
(from $T_8$) combined with the self-interaction structure of the
adjoint representation; the $4/3$ arises from the spin-$1/2$ screening
of a Weyl fermion.  The capacity-based derivation replaces the
standard Feynman-diagram calculation with an enforcement-cost
calculation, but arrives at the same coefficients.  Here we can only
establish the mechanism (antiscreening from self-interaction) and the
dominance condition (\ref{eq:AF_capacity_ineq}); the coefficients are
imported.


\subsection{Comparison with established $\beta$-function derivations}
\label{sec:AF_comparison}

The one-loop $\beta$-function coefficient
$b_0 = (11/3)C_A - (4/3)T_R N_f$ has been derived by at least four
independent methods in the literature.  This subsection identifies each,
states what the APF reuses vs.\ reinterprets, and clarifies the
boundary between established results and the APF's contribution.

\smallskip\noindent\textbf{1.\ Feynman-diagram computation
(Gross--Wilczek, Politzer, 1973).}\quad
The original derivation computes the one-loop vacuum polarization of
the gluon propagator via three diagrams: the fermion loop (screening,
$-4/3\, T_R N_f$), the ghost loop (antiscreening, $+1/3\, C_A$), and
the gluon self-energy with a three-gluon vertex (antiscreening,
$+10/3\, C_A$).  The sum gives $b_0 = (11/3)C_A - (4/3)T_R N_f$.

\textit{What the APF reuses:} the numerical result
$b_0 = (11/3)C_A - (4/3)T_R N_f$.  This is imported as an established
formula (Proposition~\ref{prop:AF_sign}) and evaluated with the
derived matter content.  The APF does not re-derive this formula in
the present paper.

\textit{What the APF reinterprets:} the physical mechanism.  In the
standard derivation, the $11/3$ arises from the sum of gluon and ghost
loop contributions; the ghost loop is a gauge-fixing artifact with no
direct physical interpretation.  In the APF's capacity language
(\S\ref{sec:AF_from_capacity}), the antiscreening is attributed to the
\emph{superadditive cost of co-located non-abelian channels}
($L_\Delta > 0$), which is a physical statement about enforcement
economics, not a gauge-fixing artifact.  The ghost contribution is
absorbed into the capacity self-cost.

\textit{What is genuinely new:} the claim that $b_0 > 0$ is a
\emph{derived consequence} of A1, not a parameter choice.  In
standard QCD, the sign of $b_0$ depends on the number of quark
flavors $N_f$, which is a free input.  In the APF, $N_f$ is derived
($T_{\mathrm{field}}$: 6 flavors from the fermion scan), and the
sign follows from the derived matter content.  The novelty is not in
the formula but in the fact that the inputs to the formula are
themselves derived.

\smallskip\noindent\textbf{2.\ Background-field method
(Abbott, 1981; 't~Hooft, 1976).}\quad
The background-field method computes the $\beta$-function by expanding
around a classical background $\bar{A}_\mu$ and integrating out
quantum fluctuations.  The result is gauge-invariant at each step,
making the antiscreening origin more transparent: the three
polarization states of a massive spin-1 particle contribute
$+(11/3)C_A$ (two physical transverse polarizations contribute
$+2C_A$, and the longitudinal / temporal sector contributes
$+(5/3)C_A$ after ghost subtraction).

\textit{APF relationship:} the background-field decomposition has a
direct capacity analogue.  The ``background field'' corresponds to
the enforcement structure at a fixed interface (the cost kernel $g$);
the ``quantum fluctuations'' are perturbations of the distinction
structure.  The $(11/3)C_A$ counts the capacity cost of maintaining $p$
self-interacting channels in $d = 4$ spacetime dimensions: two
physical polarizations (from the transversality condition, which is
the enforcement analogue of gauge fixing) plus the longitudinal
capacity overhead.  Paper~4A makes this correspondence precise.

\smallskip\noindent\textbf{3.\ Index-theoretic derivation
(Atiyah--Singer, heat-kernel methods).}\quad
The $\beta$-function can be extracted from the heat-kernel expansion
of the functional determinant $\det(-D^2)$, where $D$ is the gauge
covariant derivative.  The Seeley--DeWitt coefficient $a_2(D^2)$
gives the one-loop divergence, and its group-theoretic content
reproduces $b_0$.  This approach makes no reference to individual
Feynman diagrams; the result follows from the spectral geometry of
the gauge-covariant Laplacian.

\textit{APF relationship:} the index-theoretic derivation is the
closest mathematical relative of the APF's capacity approach.  Both
compute a spectral invariant of the gauge connection: the
Seeley--DeWitt coefficient is a spectral invariant of $-D^2$, while
the APF's capacity cost is a spectral invariant of the enforcement
algebra's cost kernel $B$ (via the eigenvalue decomposition in
Lemma~\ref{lem:gen_channel}).  The capacity inequality
(\ref{eq:AF_capacity_ineq}), $p \cdot C_A > T_R \cdot N_f^{\mathrm{total}}$,
is a statement about the spectral content of the gauge sector
vs.\ the matter sector, analogous to the comparison of
$a_2$ contributions in the heat-kernel expansion.  The precise
coefficient $11/3$ corresponds to the $a_2$ coefficient of the
spin-1 Laplacian in $d = 4$, which Paper~4A derives from
the polarization count (the same $d = 4$ that the APF derives in
Proposition~\ref{prop:d4}).

\smallskip\noindent\textbf{4.\ Twistor-space methods
(Penrose, 1967; Witten, 2004).}\quad
Witten's reformulation of perturbative gauge theory in twistor space
computes scattering amplitudes without Feynman diagrams, using
holomorphic curves in $\mathbb{CP}^3$.  The one-loop $\beta$-function
can be extracted from the holomorphic anomaly of the twistor string.
The result agrees with the standard computation but organizes the
contributions differently: the antiscreening is attributed to the
self-dual sector of the gauge field, and the screening to the
matter multiplets' contribution to the holomorphic anomaly.

\textit{APF relationship:} the twistor approach is structurally
different from the APF (twistor space is infinite-dimensional and
uses complex geometry; the APF works with finite-dimensional
enforcement algebras).  However, both share a feature: they derive
the $\beta$-function sign from \emph{structural} properties of the
gauge field (self-duality in the twistor case, self-interaction cost
in the APF case) rather than from individual diagram counting.  The
APF does not use twistor methods; the comparison is noted to show
that the impulse to derive $b_0$ from structural principles rather
than diagram-by-diagram computation has precedent.

\smallskip\noindent\textbf{Summary: new vs.\ reinterpreted.}\quad

\begin{center}\small
\renewcommand{\arraystretch}{1.2}
\begin{tabular}{@{}p{4.5cm}p{1.8cm}p{6.5cm}@{}}
\toprule
\textbf{Element} & \textbf{Status} & \textbf{Detail} \\
\midrule
Formula $b_0 = (11/3)C_A - (4/3)T_R N_f$ &
  Imported &
  Standard one-loop result (Gross--Wilczek--Politzer, 1973);
  not rederived here. \\[3pt]
Sign $b_0 > 0$ for SM content &
  Imported + derived inputs &
  Formula imported; matter content $N_f$ derived from
  $T_{\mathrm{field}}$.  The sign is therefore a \emph{prediction}
  (given A1), not a parameter choice. \\[3pt]
Capacity inequality $p C_A > T_R N_f^{\mathrm{total}}$ &
  New &
  Derived from $L_\Delta$ (superadditivity) and
  Lemma~\ref{lem:gen_channel} (generator--channel bridge).
  Provides a sufficient condition for $b_0 > 0$ without the
  formula. \\[3pt]
Antiscreening from $L_\Delta$ &
  New interpretation &
  The superadditive cost of co-located non-abelian channels
  provides a physical mechanism for the $C_A$ term;
  reinterprets the ghost + gluon-loop sum as enforcement
  self-cost. \\[3pt]
Precise coefficients $11/3$, $4/3$ &
  Deferred &
  Paper~4A ($L_{\mathrm{beta,capacity}}$) derives these from
  polarization counting in $d = 4$ and the spin-statistics
  connection.  Not available in this supplement. \\[3pt]
Index-theoretic parallel &
  Noted &
  The capacity cost is a spectral invariant of $B$, analogous
  to the Seeley--DeWitt $a_2$ coefficient.  Not developed
  further here. \\
\bottomrule
\end{tabular}
\end{center}

\smallskip\noindent The honest assessment is: the APF's
contribution to the $\beta$-function story in this supplement is
\emph{modest}.  The formula is imported; the sign is verified with
derived inputs; the capacity inequality provides a new sufficient
condition but does not reproduce the precise coefficients.  The
genuine novelty is that $N_f$ (and hence the sign of $b_0$) is
derived rather than assumed.  The reinterpretation of antiscreening
as enforcement self-cost is a new physical picture but does not yet
produce quantitative predictions beyond what the standard formula
gives.  Paper~4A is where the capacity-based derivation must stand
or fall.


\newpage

% ══════════════════════════════════════════════════════════════
\section{From cost kernel to Lorentzian spacetime}
\label{sec:lorentz}
% ══════════════════════════════════════════════════════════════

$T_{7B}$ (\S\ref{sec:T7B}) derived a positive-definite cost kernel $g$
at each interface $\Gamma$.  This section derives Lorentzian spacetime
geometry from $g$ and the previously established irreversibility
($L_{\mathrm{irr}}$) and locality ($L_{\mathrm{loc}}$) results.  The
chain uses four imported mathematical theorems (Hawking--King--McCarthy
1976, Malament 1977, Kolmogorov 1933, Lovelock 1971), all established
results independent of the framework.  Every physical input is
A1-derived.


\subsection{Causal ordering from irreversibility}
\label{sec:causal_order}

\begin{definition}[Enforcement event set]\label{def:event_set}
Let $\mathcal{I}$ be the set of interfaces and let $\Gamma \in
\mathcal{I}$.  An \emph{enforcement event} at $\Gamma$ is a triple
$e = (\Gamma, d, t)$ where $d \in \mathcal{D}(\Gamma)$ is a
distinction enforced at $\Gamma$ and $t \in \mathbb{N}$ is the
discrete enforcement step (the index in the sequence of capacity
commitments at $\Gamma$).  The \emph{event set} is
$\mathcal{E} := \{(\Gamma, d, t) : \Gamma \in \mathcal{I},\;
d \in \mathcal{D}(\Gamma),\; 0 \le t \le
\lfloor C(\Gamma)/\eps^*\rfloor\}$.
\end{definition}

\begin{definition}[Capacity commitment function]\label{def:commit}
For each event $e = (\Gamma, d, t)$, define
$\kappa(e) := \eps(d) \cdot t$ --- the cumulative capacity committed
to distinction $d$ at $\Gamma$ through step $t$.  By A1,
$\kappa(e) \le C(\Gamma) < \infty$.
\end{definition}

\begin{definition}[Local recoverability]\label{def:recoverable}
An event $e_1 = (\Gamma, d, t_1)$ is \emph{locally recoverable} at
a later event $e_2 = (\Gamma', d', t_2)$ if there exists a local
operation at $\Gamma'$ (an element of
$\mathcal{A}(\Gamma')$, applied between steps $t_1$ and $t_2$) that
reduces the net capacity commitment $\kappa(e_1)$ to zero.  If no
such local operation exists, $e_1$ is \emph{locally unrecoverable}
at~$e_2$.
\end{definition}

\begin{proposition}[Causal partial order]\label{prop:causal}
Define $e_1 \prec e_2$ iff $e_1$ is locally unrecoverable at $e_2$
(Definition~\ref{def:recoverable}).  Then $\prec$ is a strict partial
order on $\mathcal{E}$.
\end{proposition}

\begin{proof}
We verify the three axioms of a strict partial order.

\textit{Irreflexivity ($e \not\prec e$).}\quad
At the event $e = (\Gamma, d, t)$ itself, the capacity committed by
$e$ is trivially recoverable: one can apply the identity operation
$\mathrm{id} \in \mathcal{A}(\Gamma)$ to ``undo'' the commitment at
the same step (the commitment has not yet been communicated to any
other interface).  Formally: $\kappa(e)$ is reducible to zero by a
local operation at $\Gamma$ at step $t$ (the commitment was just
made and no cross-interface correlation has been established).
Therefore $e$ is not locally unrecoverable at itself: $e \not\prec e$.

\textit{Transitivity ($e_1 \prec e_2 \wedge e_2 \prec e_3 \Rightarrow
e_1 \prec e_3$).}\quad
Suppose $e_1$ is locally unrecoverable at $e_2$, and $e_2$ is locally
unrecoverable at $e_3$.  We show $e_1$ is locally unrecoverable at
$e_3$.  By $L_{\mathrm{irr}}$ (Paper~1 Supplement, \S5, Step~4):
if a cross-interface correlation commits capacity that is locally
unrecoverable at step $t_2$, then it remains locally unrecoverable
at all subsequent steps $t_3 > t_2$.  The key property of
$L_{\mathrm{irr}}$ used here is \emph{monotonicity of
unrecoverability}: local operations at later steps cannot undo
commitments that were already unrecoverable at earlier steps.
Formally: let $U \in \mathcal{A}(\Gamma')$ be any local operation
at the interface of $e_3$.  Applying $U$ cannot reduce
$\kappa(e_1)$ to zero, because (by hypothesis) no local operation
at any interface could do so at step $t_2$, and the intervening
evolution from $t_2$ to $t_3$ consists of further capacity
commitments (which can only increase or maintain, not decrease,
existing cross-interface correlations by $L_{\mathrm{irr}}$'s
monotonicity).

\textit{Antisymmetry ($e_1 \prec e_2 \Rightarrow e_2 \not\prec
e_1$).}\quad
Suppose for contradiction that $e_1 \prec e_2$ and $e_2 \prec e_1$.
Then the capacity committed by $e_1$ is locally unrecoverable at $e_2$,
and the capacity committed by $e_2$ is locally unrecoverable at $e_1$.
By transitivity (just proved), this would imply $e_1 \prec e_1$,
contradicting irreflexivity.  (The formal step: $e_1 \prec e_2$ and
$e_2 \prec e_1$ give $e_1 \prec e_1$ by transitivity.)

Therefore $(\mathcal{E}, \prec)$ is a strict partial order.
\end{proof}

\noindent\textit{Physical content.}\quad The partial order $\prec$ is
the \emph{causal structure}: $e_1 \prec e_2$ means ``$e_1$ is in the
causal past of $e_2$.''  Events that are not related by $\prec$ are
\emph{spacelike-separated}: neither is in the causal past of the other.
This is the same causal structure that appears in Lorentzian geometry,
but here it is derived from irreversibility, not from light cones.

\smallskip\noindent\textit{Causal structure is derived, not imported.}\quad
A common concern is that the Hawking--King--McCarthy and Malament theorems
(used in \S\ref{sec:signature} below) \emph{presuppose} a Lorentzian
manifold and then characterize its causal order.  This would be circular
if the APF used HKM/Malament to \emph{define} the causal order.  The
logic is instead:
\begin{enumerate}[label=(\alph*),nosep]
  \item The causal partial order $\prec$ is derived here
  (Proposition~\ref{prop:causal}) from $L_{\mathrm{irr}}$ alone, with no
  reference to any metric, manifold, or light cone.  The input is purely
  operational: ``capacity committed by $e_1$ is locally unrecoverable at
  $e_2$.''
  \item The continuum manifold $M$ is derived in the next subsection from
  $\prec$ + Kolmogorov, again without any metric signature.
  \item HKM/Malament are then used in the \emph{reverse} direction: given
  a smooth manifold $M$ equipped with a strict partial order $\prec$
  satisfying their hypotheses, what metric structures on $M$ are compatible
  with $\prec$?  Their answer is: a conformal class of \emph{Lorentzian}
  metrics, and no other signature.  This is a \emph{uniqueness theorem
  about compatibility}, not a definition of causal order from a
  pre-existing metric.
\end{enumerate}
The chain is therefore: $L_{\mathrm{irr}} \to \prec$ (no metric)
$\to M$ (Kolmogorov, no metric) $\to$ ``which metrics on $M$ are
$\prec$-compatible?'' $\to$ Lorentzian (HKM/Malament).


\subsection{Continuum manifold from Kolmogorov extension}
\label{sec:continuum}

The interface set $\mathcal{I}$ is discrete (A1: finite capacity per
interface).  A continuum manifold structure emerges via the Kolmogorov
extension theorem, whose hypotheses we now verify explicitly.

\begin{lemma}[Kolmogorov hypotheses verified]\label{lem:kolmogorov_hyp}
The enforcement net
$\{(\mathcal{A}(\Gamma), \omega|_\Gamma)\}_{\Gamma \in \mathcal{I}}$
satisfies the three hypotheses of the Kolmogorov extension theorem:
\begin{enumerate}[label=(K\arabic*),nosep]
  \item \textbf{Directed index set.}\quad $(\mathcal{I}, \subseteq)$
  is a directed set: for any two interfaces $\Gamma_1, \Gamma_2 \in
  \mathcal{I}$, there exists $\Gamma_3 \in \mathcal{I}$ with
  $\Gamma_1, \Gamma_2 \subseteq \Gamma_3$.
  \item \textbf{Each marginal is a probability measure.}\quad
  For each $\Gamma \in \mathcal{I}$, the state
  $\omega|_\Gamma: \mathcal{A}(\Gamma) \to \mathbb{C}$ is a
  faithful normal state on the finite-dimensional algebra
  $M_n(\mathbb{C})$.  This defines a probability measure on the
  spectrum of $\mathcal{A}(\Gamma)$ (the set of pure states, which
  is compact for $M_n(\mathbb{C})$).
  \item \textbf{Consistency of marginals.}\quad
  For $\Gamma_1 \subseteq \Gamma_2$, the restriction
  $\omega|_{\Gamma_2}|_{\mathcal{A}(\Gamma_1)} = \omega|_{\Gamma_1}$.
  This is the Kolmogorov consistency condition.
\end{enumerate}
\end{lemma}

\begin{proof}
(K1): Given interfaces $\Gamma_1, \Gamma_2$, define $\Gamma_3 :=
\Gamma_1 \cup \Gamma_2$ (the interface monitoring all distinctions of
both).  By A1, $C(\Gamma_3) \le C(\Gamma_1) + C(\Gamma_2) < \infty$
(the total capacity of the combined interface is bounded).  Therefore
$\Gamma_3 \in \mathcal{I}$ and $\Gamma_1, \Gamma_2 \subseteq \Gamma_3$.

(K2): The state $\omega$ is the faithful GNS state constructed in the
Paper~1 Supplement ($T_{\mathrm{GNS}}$).  On $M_n(\mathbb{C})$,
$\omega = \tau = (1/n)\,\mathrm{Tr}$ is a faithful normal tracial
state.  A faithful normal state on a finite-dimensional $*$-algebra
defines a unique probability measure on the spectrum via the spectral
theorem: $\mu_\omega(P) = \omega(P)$ for each projection
$P \in \mathcal{A}(\Gamma)$.

(K3): By $L_{\mathrm{loc}}$ (Paper~1 Supplement, \S4): enforcement at
interface $\Gamma_1$ is independent of enforcement at $\Gamma_2$ when
$\Gamma_1 \cap \Gamma_2 = \emptyset$.  For the nested case
$\Gamma_1 \subseteq \Gamma_2$: the inclusion
$\iota: \mathcal{A}(\Gamma_1) \hookrightarrow \mathcal{A}(\Gamma_2)$
maps each operator on $S_{\Gamma_1}$ to its extension (acting as
identity on $S_{\Gamma_2} \ominus S_{\Gamma_1}$).  By the
construction of $\omega$ from K3 (disjoint-support additivity) and
the sector-orthogonal form $B$ ($L_\omega$): the restriction of the
extended state to the subalgebra equals the original state.  Formally:
$\omega(\iota(A)) = \omega|_{\Gamma_1}(A)$ for all
$A \in \mathcal{A}(\Gamma_1)$.  This is the consistency condition.
\end{proof}

\begin{proposition}[Smooth manifold from enforcement regularity]
\label{prop:continuum}
Given the causal partial order $\prec$ (Proposition~\ref{prop:causal}),
the cost kernel $g$ ($T_{7B}$), and the Kolmogorov extension
(Lemma~\ref{lem:kolmogorov_hyp}):
\begin{enumerate}[label=(\roman*),nosep]
  \item The Kolmogorov extension theorem produces a unique
  $\sigma$-additive measure $\mu$ on the projective limit
  $\varprojlim \mathcal{A}(\Gamma)$.  This defines a continuum
  state space.
  \item The metric $d_g(p, q) := \inf_\gamma \int_\gamma
  \sqrt{g(\dot\gamma, \dot\gamma)}\,ds$ (where $\gamma$ ranges over
  piecewise-linear paths in the event set, and the infimum is taken
  over all paths connecting $p$ and $q$) is Lipschitz-continuous with
  constant $L = C / \varepsilon^*$ (from A1 $+$ $L_{\varepsilon^*}$).
  \item The metric completion of $(\mathcal{E}, d_g)$ is a compact
  metrizable topological space (compactness from A1:
  $\|v\| \le C / \varepsilon^*$).  By the metrization theorem, it is
  second-countable and Hausdorff.
\end{enumerate}
\end{proposition}

\noindent\textit{Hypothesis verification for Whitney embedding.}\quad
Whitney's theorem requires a second-countable Hausdorff topological
manifold.  The metric completion in (iii) is second-countable
(separable metric space) and Hausdorff (metric space).  Local
Euclidean structure (the manifold condition) requires additionally
that each point has a neighborhood homeomorphic to $\mathbb{R}^d$.
This is established as follows: the positive-definite cost kernel
$g = B$ ($T_{7B}$) provides a local inner product on the tangent
space at each event.  By the inverse function theorem applied to the
exponential map $\exp_p: T_pM \to M$ (defined via the geodesics of
$g$), each point $p$ has a normal neighborhood $U_p$ that is
diffeomorphic to an open ball in $\mathbb{R}^d$.  This uses: $g$ is
positive-definite (from $T_{7B}$), $g$ is at least $C^2$ (from the
Lipschitz-to-smoothness argument: the cost field
$\varepsilon(v) = g(v,v)$ is a polynomial in the coordinates, so its
second derivatives are continuous), and the metric completion is
connected (from the connectedness of the interface net, which follows
from the requirement that every pair of interfaces communicates via
at least one intermediary --- this is the content of
$L_{\mathrm{loc}}$'s inter-interface map).

\smallskip\noindent\textit{Imported mathematics.}\quad The Kolmogorov
extension theorem (1933), the metrization theorem (Urysohn, 1925),
the inverse function theorem, and the Whitney embedding theorem (1936)
are standard results.  The non-trivial verification is that the APF
objects satisfy their hypotheses; this is provided by
Lemma~\ref{lem:kolmogorov_hyp} and the paragraph above.

\smallskip\noindent\textit{How smooth manifold structure and local
Lorentz symmetry emerge.}\quad
The three steps of Proposition~\ref{prop:continuum} deserve
amplification, as the reviewer correctly notes that the passage from
a discrete enforcement net to a smooth manifold with local Lorentz
frames is non-trivial.

\smallskip\noindent(i) \textit{Kolmogorov extension.}\quad
The enforcement net $\{\mathcal{A}(\Gamma)\}_{\Gamma \in \mathcal{I}}$
assigns a probability distribution $\rho_\Gamma$ (the faithful state
$\omega$ restricted to $\mathcal{A}(\Gamma)$) at each interface.
Budget independence ($L_{\mathrm{loc}}$) guarantees that the marginal
of $\rho_{\Gamma'}$ restricted to observables at $\Gamma \subset
\Gamma'$ equals $\rho_\Gamma$ --- this is the Kolmogorov consistency
condition.  The extension theorem then produces a unique
$\sigma$-additive measure on the projective limit
$\varprojlim \mathcal{A}(\Gamma)$, which defines a continuum
probability space.  The discrete enforcement events embed as the
``cylinder sets'' of this measure; the continuum interpolates between
them.

\smallskip\noindent(ii) \textit{Lipschitz regularity to $C^2$
smoothness.}\quad
A1 bounds the capacity $C(\Gamma) < \infty$ at every interface, and
$L_{\varepsilon^*}$ provides a positive cost floor.  Together they
bound the Lipschitz constant of the cost field:
$|\varepsilon(v_1) - \varepsilon(v_2)| \le (C / \varepsilon^*)
\|v_1 - v_2\|_g$.  A Lipschitz function on a compact metric space
(compactness from A1: $\|v\| \le C / \varepsilon^*$) is uniformly
continuous and defines a complete metric.  The positive-definite cost
kernel $g$ from $T_{7B}$ provides this metric.  The cost field is the
square root of the quadratic form $\varepsilon(v) = g(v,v)$; since $g$
is a smooth (polynomial) function of coordinates and the manifold
is compact, the Hessian $\partial^2 \varepsilon / \partial v^i
\partial v^j = 2g_{ij}$ exists everywhere, giving $C^2$ regularity.
Whitney's embedding theorem then provides smooth coordinate charts.

\smallskip\noindent(iii) \textit{Local Lorentz symmetry.}\quad
At each point $p \in M$, the tangent space $T_pM$ inherits the
Lorentzian inner product $g_p$ from the signature derivation
(\S\ref{sec:signature}).  The isometry group of $(T_pM, g_p)$ is the
Lorentz group $\mathrm{O}(1,d{-}1)$.  This is not imposed: it is the
automorphism group of the derived metric structure at a point.
The connected component $\mathrm{SO}^+(1,3)$ (proper orthochronous
Lorentz group) is selected by the causal ordering $\prec$ (which
distinguishes future from past) and the orientation of the manifold
(which follows from the continuity of $g$ and the orientability forced
by the existence of the volume form from A1's capacity assignment).
The local Lorentz frames are therefore derived, not postulated: they
are the orthonormal bases of $(T_pM, g_p)$ consistent with the derived
causal order and orientation.


\subsection{Lorentzian signature from causal structure}
\label{sec:signature}

The central claim of this subsection is that the causal partial order
$\prec$ (Proposition~\ref{prop:causal}) on the smooth manifold $M$
(Proposition~\ref{prop:continuum}) determines a conformal class of
Lorentzian metrics.  This uses two external theorems whose exact
hypotheses we now state and verify.

\begin{lemma}[HKM hypotheses verified]\label{lem:HKM_hyp}
The Hawking--King--McCarthy theorem (1976)~\cite{HKM1976} states:

\smallskip\noindent\textbf{Theorem (HKM).}\quad
\textit{Let $(M, g)$ and $(M, g')$ be two spacetimes (connected
time-oriented Lorentzian manifolds) with the same underlying manifold
$M$.  If $(M, g)$ is strongly causal and the causal relations
$J^+(g) = J^+(g')$ agree, then $g' = \Omega^2 g$ for some positive
smooth function $\Omega: M \to \mathbb{R}_{>0}$.}

\smallskip\noindent The APF uses this in the reverse direction: given
$M$ and $\prec$, we must show that \emph{there exists} a Lorentzian
metric $g$ on $M$ whose causal relation reproduces $\prec$, and that
$g$ is strongly causal.  The theorem then guarantees that $g$ is unique
up to a conformal factor.

The existence of a compatible Lorentzian metric is a separate result
(Kronheimer--Penrose, 1967; Levichev, 1987; Minguzzi, 2019): a
partial order on a smooth manifold admits a compatible Lorentzian
metric if and only if it satisfies:
\begin{enumerate}[label=(KP\arabic*),nosep]
  \item \textbf{Closedness.}\quad The order relation
  $\{(p,q) : p \preceq q\}$ is a closed subset of $M \times M$ in the
  product topology.
  \item \textbf{Inner continuity.}\quad For every $p \in M$, the
  chronological future $I^+(p) := \{q : p \prec q\}$ is an open
  subset of $M$.
  \item \textbf{Non-triviality.}\quad For every $p \in M$,
  $I^+(p) \ne \emptyset$ and $I^-(p) \ne \emptyset$.
\end{enumerate}
We verify each for the APF's $\prec$:

(KP1): The relation $\prec$ is defined by local unrecoverability
(Definition~\ref{def:recoverable}), which is determined by the
discrete enforcement steps at finitely many interfaces.  In the
continuum extension (Proposition~\ref{prop:continuum}), the relation
$\preceq$ on $M$ is the closure of the discrete relation in the
product topology of the metric completion.  The closure operation
guarantees that the resulting set is closed.

(KP2): $I^+(p)$ is the set of events at which $p$'s capacity
commitment is unrecoverable.  By $L_{\mathrm{irr}}$'s monotonicity,
once a commitment becomes unrecoverable at some event $q$, it remains
unrecoverable at all events in a neighborhood of $q$ (because the
enforcement cost is a continuous function of the event coordinates, by
the Lipschitz regularity of Proposition~\ref{prop:continuum}(ii)).
Therefore $I^+(p)$ is open.

(KP3): For every event $p = (\Gamma, d, t)$: the next step
$e' = (\Gamma, d, t+1)$ satisfies $p \prec e'$ (the commitment at
step $t$ is unrecoverable at step $t+1$, by $L_{\mathrm{irr}}$'s
forward irreversibility), so $I^+(p) \ne \emptyset$.  Similarly, if
$t \ge 1$, the previous step $e'' = (\Gamma, d, t-1)$ satisfies
$e'' \prec p$, so $I^-(p) \ne \emptyset$.  (At $t = 0$:
non-triviality of $I^-(p)$ is guaranteed by the existence of at least
one other interface $\Gamma' \ne \Gamma$ with an event preceding $p$,
which follows from $L_{\mathrm{loc}}$'s requirement that interfaces
form a connected net.)
\end{lemma}

\begin{lemma}[Strong causality verified]\label{lem:strong_causal}
The causal structure $(\mathcal{E}, \prec)$ is strongly causal.
\end{lemma}

\begin{proof}
Strong causality requires: for every $p \in M$ and every neighborhood
$U$ of $p$, there exists a neighborhood $V \subseteq U$ such that no
causal curve starting in $V$ leaves $U$ and returns to $V$.  In the
APF, this follows from three properties:

\textit{(i) No closed causal curves.}\quad A closed causal curve would
require $p \prec p$, contradicting irreflexivity
(Proposition~\ref{prop:causal}).

\textit{(ii) Finite speed of causal propagation.}\quad The Lipschitz
bound $d_g(p, q) \le (C/\eps^*) \cdot |t_p - t_q|$ (from the
bounded cost kernel and the discrete step structure) ensures that
causal curves cannot traverse arbitrarily large spatial distances in
a single enforcement step.  For any neighborhood $U$ of $p$, choosing
$V$ to be the intersection of $U$ with a sufficiently small time
interval ensures that a causal curve starting in $V$ cannot reach the
boundary of $U$ and return within the allowed time.

\textit{(iii) Compactness.}\quad The event set is compact (A1: finite
capacity bounds all coordinates).  On a compact set, the absence of
closed causal curves implies strong causality (this is a theorem:
Bernal--S\'{a}nchez, 2007, Theorem~3.2, which proves that a
compact causal set without closed causal curves is strongly causal).
\end{proof}

\begin{lemma}[Malament hypotheses verified]\label{lem:malament_hyp}
The Malament uniqueness theorem (1977)~\cite{Malament1977} states:

\smallskip\noindent\textbf{Theorem (Malament).}\quad
\textit{Let $(M, g)$ be a strongly causal spacetime.  The causal
relation $J^+$ uniquely determines the topology, differentiable
structure, conformal geometry, and time orientation of $(M, g)$.}

\smallskip\noindent The hypotheses are: (a)~$M$ is a smooth manifold
(verified: Proposition~\ref{prop:continuum}); (b)~$g$ is a Lorentzian
metric compatible with $\prec$ (verified: Lemma~\ref{lem:HKM_hyp},
existence via KP1--KP3); (c)~strong causality (verified:
Lemma~\ref{lem:strong_causal}).  All three are satisfied.
\end{lemma}

\begin{theorem}[$\Delta_{\mathrm{sig}}$: Lorentzian Signature]\label{thm:signature}
The spacetime metric on $M$ has Lorentzian signature $(-,+,+,+)$.
\end{theorem}

\begin{proof}
Three steps.

\textit{Step 1: Existence and conformal uniqueness of compatible
metric.}\quad
By Lemma~\ref{lem:HKM_hyp} (verification of KP1--KP3), there exists
a Lorentzian metric $g$ on $M$ whose causal relation reproduces
$\prec$.  By the HKM theorem (Lemma~\ref{lem:HKM_hyp}), applied with
Lemma~\ref{lem:strong_causal} (strong causality), any other compatible
metric is conformally related: $g' = \Omega^2 g$.  By the Malament
theorem (Lemma~\ref{lem:malament_hyp}), the topology, differentiable
structure, conformal geometry, and time orientation are uniquely
determined by $\prec$.  The compatible metric is therefore Lorentzian
(this is a property of the conformal class, not a choice: the
existence result from KP1--KP3 guarantees a Lorentzian metric, and
the uniqueness result says all compatible metrics share the same
conformal class, hence the same signature).

\textit{Step 2: $L_{\mathrm{irr}}$ selects exactly one timelike
direction.}\quad
Irreversibility ($L_{\mathrm{irr}}$) provides a single arrow at each
event: one direction along which capacity commitments are unrecoverable.
This is the timelike direction.  All other directions at each event are
``reversible'' at the enforcement level (enforcement costs can be recovered
by local operations along these directions, by $L_{\mathrm{loc}}$).
These are the spacelike directions.

The signature is therefore $(-, +, +, \ldots, +)$: one negative
(timelike) direction and $d-1$ positive (spacelike) directions.  The
number of timelike directions is exactly one because $L_{\mathrm{irr}}$
defines a single preferred direction (the ``arrow of time'' from
irreversibility); if there were two or more independent irreversible
directions, the corresponding capacity commitments would need to be
simultaneously unrecoverable along both directions, which would
require two independent monotonicity properties --- but
$L_{\mathrm{irr}}$ derives a single monotonicity from the single pool
coupling $L_\Delta$.

\textit{Step 3: Conformal factor fixed by volume normalization.}\quad
The conformal factor $\Omega$ is fixed to $\Omega = 1$ by the
enforcement volume normalization: A1 assigns a definite capacity
$C(\Gamma)$ to each interface, and the capacity density
$\rho_C = C / \mathrm{Vol}(M)$ fixes the volume element, hence the
conformal factor.  (This is the Radon--Nikodym uniqueness argument:
given a fixed total capacity and a conformal class, the volume form
selects a unique representative.)

Therefore the spacetime metric is Lorentzian with signature $(-,+,+,+)$.
\end{proof}


\subsection{From positive-definite cost kernel to Lorentzian spacetime metric}
\label{sec:cost_to_lorentz}

The preceding subsections established three objects: a positive-definite
cost kernel $g$ at each interface (Theorem~\ref{thm:T7B}), a causal
partial order $\prec$ on events (Proposition~\ref{prop:causal}), and a
smooth manifold $M$ (Proposition~\ref{prop:continuum}).
Theorem~\ref{thm:signature} showed that the metric on $M$ compatible
with $\prec$ has Lorentzian signature.  What remains is to make explicit
the \emph{construction} that assembles these ingredients into a single
Lorentzian metric tensor $g_{\mu\nu}$ on $M$.  This subsection provides
that construction, with clear hypotheses and verification that the result
is a symmetric bilinear form of the correct signature.

\begin{proposition}[Spacetime metric from cost kernel and causal order]%
\label{prop:spacetime_metric}
\textbf{Hypotheses:}
\begin{enumerate}[label=(H\arabic*),nosep]
  \item A smooth manifold $M$ of dimension $d$
  (Proposition~\ref{prop:continuum}).
  \item At each point $p \in M$, a positive-definite symmetric bilinear
  form $g_p^{(\mathrm{sp})}: T_p^{(\mathrm{sp})}M \times
  T_p^{(\mathrm{sp})}M \to \mathbb{R}$ on the spatial tangent
  directions, inherited from the cost kernel $g = B$
  (Theorem~\ref{thm:T7B}), restricted to directions orthogonal to
  the causal direction.
  \item A smooth timelike vector field $\tau$ on $M$, defined by the
  causal partial order $\prec$
  (Proposition~\ref{prop:causal}): at each $p$, $\tau_p$ points in
  the unique direction along which capacity commitments are
  irreversible ($L_{\mathrm{irr}}$).
  \item A positive lapse function $N: M \to \mathbb{R}_{>0}$, fixed by
  the enforcement capacity density: $N^2(p) := C(p) / \mathrm{Vol}(p)$,
  where $C(p)$ is the local capacity per unit coordinate volume (from
  A1) and $\mathrm{Vol}(p)$ is the spatial volume element from
  $g^{(\mathrm{sp})}$.
\end{enumerate}

\noindent\textbf{Construction.}\quad
Define the spacetime metric at each $p \in M$ by:
\begin{equation}\label{eq:lorentz_metric}
  g_{\mu\nu}(p)\,dx^\mu\,dx^\nu \;:=\;
  -N^2(p)\,dt^2 \;+\; g_{ij}^{(\mathrm{sp})}(p)\,dx^i\,dx^j,
\end{equation}
where $t$ is the parameter along the integral curves of $\tau$
(enforcement time), and $x^i$ ($i = 1, \ldots, d{-}1$) are coordinates
on the spatial slice orthogonal to $\tau$.

\noindent\textbf{Conclusions:}
\begin{enumerate}[label=(\roman*),nosep]
  \item $g_{\mu\nu}$ is a symmetric bilinear form on $T_pM$ for each
  $p \in M$.
  \item $g_{\mu\nu}$ has Lorentzian signature $(-,+,+,\ldots,+)$:
  one negative eigenvalue (the $dt^2$ term, with eigenvalue $-N^2 < 0$)
  and $d-1$ positive eigenvalues (from $g^{(\mathrm{sp})}$, which is
  positive-definite by (H2)).
  \item $g_{\mu\nu}$ is unique given the four hypotheses
  (up to the within-sector gauge freedom of $B$, which is an overall
  conformal factor absorbed by (H4)).
  \item The causal structure of $(M, g_{\mu\nu})$ reproduces $\prec$:
  a curve $\gamma$ is future-directed causal iff
  $g_{\mu\nu}\,\dot{\gamma}^\mu\,\dot{\gamma}^\nu \le 0$ and
  $\dot{\gamma}$ has positive $dt$-component, which holds iff
  $\gamma$ follows the irreversibility direction $\tau$.
\end{enumerate}
\end{proposition}

\begin{proof}
(i) Symmetry: $g_{\mu\nu} = g_{\nu\mu}$ is immediate from the
construction --- both $-N^2\,dt \otimes dt$ and
$g_{ij}^{(\mathrm{sp})}\,dx^i \otimes dx^j$ are symmetric tensor
products.

(ii) Signature: in the orthonormal frame $\{N^{-1}\partial_t,\,
e_1, \ldots, e_{d-1}\}$ (where $\{e_i\}$ diagonalize
$g^{(\mathrm{sp})}$), the metric becomes
$\mathrm{diag}(-1, +1, \ldots, +1)$.  The negative sign on the
timelike direction is the mathematical implementation of
$L_{\mathrm{irr}}$'s irreversibility: capacity commitments along
$\tau$ are locally unrecoverable, which means the cost functional
\emph{decreases} available capacity along $\tau$ rather than
preserving it.  The negative eigenvalue is not postulated; it is the
unique signature assignment compatible with the derived causal
ordering (Theorem~\ref{thm:signature}).

(iii) Uniqueness: (H1)--(H3) determine $g_{\mu\nu}$ up to the lapse
$N$.  (H4) fixes $N$ via the capacity density.  The residual gauge
(within-sector normalization of $B$) corresponds to the standard
conformal gauge of GR, which is fixed by (H4).

(iv) Causal consistency: let $\gamma$ be future-directed.  Then
$g_{\mu\nu}\,\dot{\gamma}^\mu\,\dot{\gamma}^\nu =
-N^2 (\dot{t})^2 + g_{ij}^{(\mathrm{sp})}\,\dot{x}^i\,\dot{x}^j$.
For a timelike curve: $|\dot{t}|$ large, spatial velocity small;
the negative term dominates, so
$g_{\mu\nu}\,\dot{\gamma}^\mu\,\dot{\gamma}^\nu < 0$.  For a
null curve: the terms balance exactly.  For a spacelike curve:
the positive spatial term dominates.  This reproduces the standard
causal structure of a Lorentzian manifold, with the light cone
determined by $N^2 (\dot{t})^2 = g_{ij}^{(\mathrm{sp})}\,\dot{x}^i\,
\dot{x}^j$, which is exactly the causal boundary of $\prec$.
\end{proof}

\smallskip\noindent\textit{Why the construction is an ADM
decomposition, and why this is not circular.}\quad
The metric (\ref{eq:lorentz_metric}) has the form of the ADM
decomposition (Arnowitt--Deser--Misner, 1962) with zero shift vector.
One might worry that this imports GR formalism.  The logic is the
reverse: the ADM decomposition is a \emph{mathematical identity}
that holds for any Lorentzian metric with a smooth time function.
The APF's construction produces a metric of ADM form because the
enforcement structure naturally separates into a temporal direction
($\tau$, from $L_{\mathrm{irr}}$) and spatial directions (from
$L_{\mathrm{loc}}$).  The zero shift reflects the fact that the
spatial coordinates are defined by $L_{\mathrm{loc}}$'s independent
interfaces, which are at rest with respect to the enforcement clock.

\smallskip\noindent\textit{Hypothesis inventory and derivation
status.}\quad

\begin{center}\small
\renewcommand{\arraystretch}{1.15}
\begin{tabular}{@{}p{1.0cm}p{4.0cm}p{3.0cm}p{4.0cm}@{}}
\toprule
\textbf{Hyp.} & \textbf{Content} & \textbf{Derived from} &
  \textbf{Status} \\
\midrule
(H1) & Smooth $d$-manifold $M$ &
  $\prec$ + Kolmogorov + Whitney &
  Derived (Prop.~\ref{prop:continuum}) \\
(H2) & Positive-definite $g^{(\mathrm{sp})}$ on spatial slices &
  $T_{7B}$: $g = B$ is positive-definite &
  Derived (Thm.~\ref{thm:T7B}) \\
(H3) & Smooth timelike vector field $\tau$ &
  $L_{\mathrm{irr}}$: irreversibility direction &
  Derived (Prop.~\ref{prop:causal}) \\
(H4) & Lapse $N$ from capacity density &
  A1 capacity + spatial volume from (H2) &
  Derived (conformal factor fixing) \\
\bottomrule
\end{tabular}
\end{center}

\noindent All four hypotheses are derived from A1 and its upstream
consequences.  No metric structure, signature assumption, or
Lorentzian geometry is imported.  The symmetric bilinear form
$g_{\mu\nu}$ is the \emph{output} of the construction, and its
Lorentzian signature is a derived consequence, not a hypothesis.

\smallskip\noindent\textit{What remains for Paper~6.}\quad
Proposition~\ref{prop:spacetime_metric} constructs the metric in the
zero-shift, synchronous gauge.  The full Paper~6 treatment extends
this to:
\begin{enumerate}[label=(\alph*),nosep]
  \item General lapse and shift (diffeomorphism invariance from the
  outer automorphisms of the enforcement net,
  Theorem~\ref{thm:gauge_ident}(i)).
  \item The dynamics of $g_{\mu\nu}$ (Einstein equations from
  Lovelock, $T_{9,\mathrm{grav}}$, \S\ref{sec:einstein}).
  \item The coupling to matter ($T_{\mu\nu}$ from the enforcement
  stress-energy, derived from the same cost kernel).
\end{enumerate}
None of these extensions change the construction or signature
presented here; they add dynamical content to the kinematic framework
established by Proposition~\ref{prop:spacetime_metric}.


\subsection{$d = 4$ from gravitational degrees of freedom}
\label{sec:d4}

The spacetime dimension is not a free parameter; it is derived from the
enforcement structure.  Lovelock uniqueness plays a role in the argument,
but is not the sole ingredient.  The full selection requires four
conditions, two of which are degree-of-freedom arguments that exclude
$d \ge 5$ independently of the form of the field equation.

\begin{proposition}[$d = 4$ uniquely selected]\label{prop:d4}
The spacetime dimension is $d = 4$.
\end{proposition}

\begin{proof}
The number of propagating gravitational degrees of freedom in $d$
dimensions is $n_{\mathrm{grav}}(d) = d(d-3)/2$ (from the linearized
Einstein equations).  Each gravitational DOF is an independent
enforcement channel that consumes capacity.  The Lovelock order
$n_{\max}(d) = \lfloor (d-1)/2 \rfloor$ counts the number of
independent divergence-free symmetric 2-tensors built from
$g_{\mu\nu}$ and its first two derivatives; uniqueness of the response
law requires $n_{\max} = 1$.

\smallskip
\begin{center}\small
\begin{tabular}{@{}cccl@{}}
\toprule
$d$ & $n_{\mathrm{grav}}$ & $n_{\max}$ & Status \\
\midrule
2 & 0 & 0 & Excluded by (D1): no propagation \\
3 & 0 & 1 & Excluded by (D1): no propagation \\
4 & 2 & 1 & \textbf{Selected}: minimal DOF, unique response \\
5 & 5 & 2 & Excluded by (D2) and (D4) \\
$\ge 6$ & $\ge 9$ & $\ge 2$ & Excluded by (D2) and (D4) \\
\bottomrule
\end{tabular}
\end{center}

\smallskip\noindent Four conditions select $d = 4$ uniquely:
\begin{enumerate}[label=(D\arabic*),nosep]
  \item \textit{Propagating DOF required.}  $L_{\mathrm{irr}}$
  (irreversibility) implies that perturbations of the cost kernel $g$
  must propagate: a local change in enforcement cost at one interface
  must be communicable to neighboring interfaces (otherwise capacity
  commitments are not locally unrecoverable, contradicting
  $L_{\mathrm{irr}}$).  This requires $n_{\mathrm{grav}}(d) \ge 1$,
  hence $d \ge 4$.

  \item \textit{Lovelock uniqueness (response-law determinacy).}
  Lovelock's theorem (1971)~\cite{Lovelock1971}: the divergence-free
  symmetric 2-tensor built from $g_{\mu\nu}$ and its first two
  derivatives is unique in $d = 4$ ($G_{\mu\nu} + \Lambda
  g_{\mu\nu}$).  In $d \ge 5$, the Gauss--Bonnet tensor
  $\mathcal{H}^{(2)}_{\mu\nu}$ provides a second independent
  contribution, and the field equation becomes a one-parameter family
  $G_{\mu\nu} + \alpha\,\mathcal{H}^{(2)}_{\mu\nu} + \Lambda
  g_{\mu\nu} = \kappa\,T_{\mu\nu}$.  The ratio
  $\alpha/\kappa$ is not fixed by A1.  Uniqueness of the response law
  is required because the five admissibility conditions
  (A9.1--A9.5) \emph{determine} the field equation, not merely
  constrain it; a free parameter would be an unexplained input.

  \item \textit{Hyperbolic propagation.}  The linearized field equation
  in $d = 4$ is hyperbolic (wave-like solutions exist), compatible with
  the causal structure from Proposition~\ref{prop:causal}.

  \item \textit{Gravitational capacity overhead (DOF minimality).}
  Each propagating gravitational DOF is an enforcement channel that the
  budget~$C$ must service.  In $d = 5$, there are $n_{\mathrm{grav}} =
  5$ such channels; in $d = 4$, only $n_{\mathrm{grav}} = 2$.  A1
  (finite enforcement capacity) combined with $L_{\mathrm{col}}$
  (minimality) selects the dimension with the fewest gravitational DOF
  consistent with (D1)--(D3).  The additional 3 DOF in $d = 5$ are
  enforcement channels with no admissibility-derived role: they carry
  capacity cost but serve no function required by the derivation chain.
  More concretely, the capacity accounting that yields
  $C_{\mathrm{total}} = 61$ (gauge generators + fermion DOF + Higgs
  components) assumes the gravitational sector is serviced by exactly 2
  propagating modes; $d = 5$ would require reallocating capacity to 5
  gravitational modes, disrupting the matter-content budget and breaking
  the fermion-scan closure that selects the Standard Model field
  content.
\end{enumerate}
$d = 4$ is the unique dimension satisfying all four.
\end{proof}

\smallskip\noindent\textit{Why Lovelock alone does not suffice, and
why ``setting $\alpha = 0$ in $d \ge 5$'' does not rescue higher
dimensions.}\quad
Two objections to the $d = 4$ selection merit explicit treatment.

The first objection is that Lovelock's theorem does not make gravity
\emph{impossible} in $d \ge 5$; it merely makes the field equation
\emph{non-unique}.  One could accept the most general Lovelock
Lagrangian in $d \ge 5$ and treat $\alpha/\kappa$ as a measurable
coupling.  The APF's answer is that this would constitute a
\emph{derivation failure}, not a physical inconsistency: the framework
claims to derive all couplings from A1, and an undetermined ratio is
a gap in that derivation.  The argument is not that $d \ge 5$ gravity
is inconsistent, but that the derivation chain terminates at a
parametric family rather than a unique answer, breaking the
zero-free-parameter architecture.

The second objection is subtler: one might propose that in $d \ge 5$,
the Gauss--Bonnet coupling happens to vanish ($\alpha = 0$), recovering
a unique Einstein equation even though Lovelock permits more.  This
fails for two reasons:
\begin{enumerate}[label=(\alph*),nosep]
  \item Setting $\alpha = 0$ is an \emph{additional constraint} not
  derived from A1 or any admissibility condition.  It must be
  postulated, which violates the framework's derivation discipline: if
  $\alpha = 0$ cannot be derived from the enforcement structure, then
  it is a free choice, and the derivation is incomplete.
  \item Even with $\alpha = 0$, the $d = 5$ gravitational sector still
  carries 5 propagating DOF, not 2.  These 5 channels consume
  enforcement capacity regardless of the field-equation form.  The DOF
  count is a property of the \emph{kinematics} (the linearized mode
  structure of a metric theory in $d$ dimensions), not of the specific
  \emph{dynamics} (the choice of Lovelock couplings).  Setting the
  Gauss--Bonnet coupling to zero changes the dynamics but does not
  remove the kinematic overhead: there are still 5 independent
  polarizations of the graviton in $d = 5$, each requiring capacity
  allocation.  Condition (D4) excludes $d = 5$ on this kinematic
  ground alone, independently of the Lovelock-uniqueness argument.
\end{enumerate}
In summary, (D2) and (D4) provide \emph{independent} exclusions of
$d \ge 5$: (D2) via response-law non-uniqueness (a derivation-chain
failure), and (D4) via excess gravitational capacity overhead (a
capacity-budget failure).  Either suffices; both apply.


\subsection{Einstein equations from Lovelock}
\label{sec:einstein}

\begin{theorem}[$T_{9,\mathrm{grav}}$: Einstein Field Equations]\label{thm:einstein}
The dynamics of the cost kernel $g_{\mu\nu}$ are governed by the
Einstein equations:
\[
  G_{\mu\nu} + \Lambda\, g_{\mu\nu} = \kappa\, T_{\mu\nu},
\]
where $G_{\mu\nu}$ is the Einstein tensor, $\Lambda$ is the
cosmological constant, and $T_{\mu\nu}$ is the stress-energy tensor of
the matter content.
\end{theorem}

\noindent The proof uses Lovelock's theorem, whose exact hypotheses we
verify.

\begin{lemma}[Lovelock hypotheses verified]\label{lem:lovelock_hyp}
Lovelock's theorem (1971)~\cite{Lovelock1971} states:

\smallskip\noindent\textbf{Theorem (Lovelock).}\quad
\textit{Let $(M, g)$ be a $d$-dimensional pseudo-Riemannian manifold.
The most general symmetric, divergence-free $(0,2)$-tensor field
$\mathcal{G}_{\mu\nu}$ constructed from $g_{\mu\nu}$ and its
derivatives up to second order is:
\[
  \mathcal{G}_{\mu\nu} = \sum_{p=0}^{\lfloor(d-1)/2\rfloor}
  c_p\, \mathcal{H}^{(p)}_{\mu\nu},
\]
where $\mathcal{H}^{(0)}_{\mu\nu} = g_{\mu\nu}$ (cosmological term),
$\mathcal{H}^{(1)}_{\mu\nu} = G_{\mu\nu}$ (Einstein tensor), and
$\mathcal{H}^{(p)}_{\mu\nu}$ for $p \ge 2$ are the higher Lovelock
tensors (with $\mathcal{H}^{(2)}$ being the Gauss--Bonnet tensor).
In $d = 4$: $\mathcal{H}^{(2)}_{\mu\nu}$ is topological (vanishes
identically from the equations of motion), so the unique solution is
$\mathcal{G}_{\mu\nu} = c_0\, g_{\mu\nu} + c_1\, G_{\mu\nu}$.}

\smallskip\noindent The hypotheses are:
\begin{enumerate}[label=(LV\arabic*),nosep]
  \item \textbf{Pseudo-Riemannian manifold.}\quad $(M, g)$ is a smooth
  manifold with a non-degenerate symmetric bilinear form of constant
  signature.
  \item \textbf{Symmetry.}\quad $\mathcal{G}_{\mu\nu} = \mathcal{G}_{\nu\mu}$.
  \item \textbf{Divergence-free.}\quad
  $\nabla^\mu \mathcal{G}_{\mu\nu} = 0$.
  \item \textbf{Second-order.}\quad $\mathcal{G}_{\mu\nu}$ depends on
  $g_{\alpha\beta}$ and its first and second derivatives
  $\partial_\gamma g_{\alpha\beta}$,
  $\partial_\gamma \partial_\delta g_{\alpha\beta}$, but not on higher
  derivatives.
  \item \textbf{Dimension.}\quad $d = 4$.
\end{enumerate}
We verify each for the APF's cost-kernel metric:

(LV1): $M$ is a smooth manifold
(Proposition~\ref{prop:continuum}); $g_{\mu\nu}$ is Lorentzian with
signature $(-,+,+,+)$ (Theorem~\ref{thm:signature},
Proposition~\ref{prop:spacetime_metric}).  Non-degeneracy: $g_{\mu\nu}$
has no zero eigenvalues (the timelike eigenvalue is $-N^2 < 0$ and the
spatial eigenvalues are positive by the positive-definiteness of
$g^{(\mathrm{sp})}$ from $T_{7B}$).  Constant signature: the signature
is $(-,+,+,+)$ everywhere, because $N > 0$ everywhere (from H4:
$N^2 = C/\mathrm{Vol} > 0$) and $g^{(\mathrm{sp})}$ is
positive-definite everywhere (from $T_{7B}$: $B(v,v) > 0$ for all
$v \ne 0$ at every interface).

(LV2): Symmetry of the response tensor $\mathcal{G}_{\mu\nu}$ is
required by the physical interpretation: $\mathcal{G}_{\mu\nu}$ is
equated to $\kappa\,T_{\mu\nu}$, and the stress-energy tensor is
symmetric (from the conservation of angular momentum, which follows
from the enforcement algebra's isotropy: MD applied to the spacetime
sector).

(LV3): Divergence-freeness $\nabla^\mu \mathcal{G}_{\mu\nu} = 0$ is
required by the conservation of the source: $\nabla^\mu T_{\mu\nu} = 0$
(energy-momentum conservation, from A1: total capacity is conserved).
Since $\mathcal{G}_{\mu\nu} = \kappa\, T_{\mu\nu}$, the left-hand side
must also be divergence-free.  This is the contracted Bianchi identity
$\nabla^\mu G_{\mu\nu} = 0$ (a mathematical identity of the Riemann
tensor, imported), combined with
$\nabla^\mu(\Lambda\, g_{\mu\nu}) = 0$ (metric compatibility).

(LV4): Second-order dependence on the metric follows from the $C^2$
regularity of $g_{\mu\nu}$: the cost kernel is a polynomial in the
sector coordinates (from $T_{7B}$: $\varepsilon(v) = B(v,v)$), so its
second derivatives are continuous.  Higher derivatives of $g$ are not
available from the enforcement structure at this order of regularity.
This is not a choice but a consequence: the enforcement cost is a
quadratic form ($T_{7B}$), so the response law can depend on at most
two derivatives (the Hessian of the cost).

(LV5): $d = 4$ is derived in
Proposition~\ref{prop:d4}.
\end{lemma}

\begin{proof}[Proof of Theorem~\ref{thm:einstein}]
The five admissibility conditions (A9.1--A9.5) are:
\begin{enumerate}[label=(A9.\arabic*),nosep]
  \item \textit{Locality}: the response depends on $g$ and finitely many
  derivatives (from $L_{\mathrm{loc}}$: enforcement is local to each
  interface).  This establishes (LV4).
  \item \textit{General covariance}: the response is tensorial (from
  the interface-independence of the enforcement algebra: a coordinate
  relabeling of $\mathcal{I}$ is a kinematic automorphism,
  Lemma~\ref{lem:aut_decomp}).  Combined with the symmetry of
  $T_{\mu\nu}$, this establishes (LV2).
  \item \textit{Conservation}: $\nabla_\mu T^{\mu\nu} = 0$ identically
  (from A1: total capacity is conserved; the stress-energy tensor is the
  capacity current).  This establishes (LV3).
  \item \textit{Second-order stability}: the response involves at most
  second derivatives of $g$ (from $T_{7B}$: the cost kernel is
  $C^2$; higher derivatives would require differentiability not
  established).  This reinforces (LV4).
  \item \textit{Hyperbolic propagation}: the linearized response admits
  wave-like solutions (from D3 of Proposition~\ref{prop:d4}).
\end{enumerate}
By Lemma~\ref{lem:lovelock_hyp}, all five Lovelock hypotheses are
satisfied.  In $d = 4$ (Proposition~\ref{prop:d4}), Lovelock's theorem
gives $\mathcal{G}_{\mu\nu} = c_0\, g_{\mu\nu} + c_1\, G_{\mu\nu}$
as the unique solution.  Identifying $c_0 = \Lambda$ and
$c_1 = 1$ (with appropriate normalization via Newton's constant
$\kappa = 8\pi G$, fixed by the weak-field limit matching condition
A9.5) yields $G_{\mu\nu} + \Lambda\, g_{\mu\nu} = \kappa\,
T_{\mu\nu}$.
\end{proof}

\smallskip\noindent\textit{Imported mathematics.}\quad
Lovelock's theorem (1971) is the sole external input; its hypotheses
are verified in Lemma~\ref{lem:lovelock_hyp}.  The Hawking--King--McCarthy
(1976) and Malament (1977) theorems, used in the signature
derivation, have their hypotheses verified in
Lemma~\ref{lem:HKM_hyp}, Lemma~\ref{lem:strong_causal}, and
Lemma~\ref{lem:malament_hyp}.  The Kolmogorov extension theorem (1933)
has its hypotheses verified in Lemma~\ref{lem:kolmogorov_hyp}.  No
external theorem is invoked without explicit hypothesis verification.

\coderef{check\_Delta\_signature, check\_T8, check\_T9\_grav}{spacetime.py, gravity.py}


\newpage

% ╔═══════════════════════════════════════════════════════════════╗
% ║              PART II: GAUGE TEMPLATE AND MATTER CONTENT       ║
% ╚═══════════════════════════════════════════════════════════════╝

\part{Gauge Template and Matter Content}
\label{part:gauge}


% ══════════════════════════════════════════════════════════════
\section{Gauge template uniqueness: exhaustive Lie algebra classification}
\label{sec:gauge_class_supp}
% ══════════════════════════════════════════════════════════════

Paper~2 \S5.1 states $L_{\mathrm{gauge,unique}}$ in compact form.
This section provides the full classification with explicit exclusion
reasons for all 17 compact simple Lie algebras, quantified capacity
costs, and closed-form proofs for every elimination.

The capacity cost of a gauge group $G$ is $C(G) = \dim(G) \cdot F(1)$
where $F$ is the unique cost function from $L_{\mathrm{Cauchy}}$
(Theorem~\ref{thm:LCauchy}).  Since $F(1) = 1$ (unit normalization),
$C(G) = \dim(G)$.  This is the quantity minimized in the classification
below.

\begin{proposition}[Gauge generators are separated enforcement interfaces]\label{prop:generator_separation}
For any compact Lie group $G$ with Lie algebra $\mathfrak{g}$ of dimension
$n$, the $n$ independent generators $\{T^a\}_{a=1}^n$ constitute $n$
\emph{separated} enforcement interfaces in the sense of
$L_{\mathrm{Cauchy}}$'s additivity premise (C1): for any pair $a \ne b$,
committing enforcement to $T^a$ does not reshape the channel landscape
seen by $T^b$, so the superadditive surplus $\Delta(T^a, T^b) = 0$.
Therefore the total cost is additive: $C(G) = \sum_{a=1}^n F(1) = \dim(G)$.
\end{proposition}

\begin{proof}
The generators of a compact Lie algebra are Killing-orthogonal with
respect to the negative-definite Killing form: $K(T^a, T^b) \propto
\delta^{ab}$.  By $T_M$ (interface monogamy, Paper~1 Supplement) and
$L_{\mathrm{loc}}$ (locality), two distinct enforcement channels whose
anchors are mutually orthogonal under a sector-orthogonal bilinear form
are \emph{separated interfaces}: commitment to one does not rearrange
the other.  This is the condition under which $L_{\mathrm{Cauchy}}$'s
C1 (additivity of cost at separated interfaces) applies directly.
Hence the cost of enforcing all $n$ generators is $\sum_a F(1) = n$
by C1 applied $(n-1)$ times.  Non-commuting generators within a simple
factor (e.g., $[T^1, T^2] = if^{12c} T^c$ in $\SU(3)$) are
Killing-orthogonal despite their non-commutativity: the separated-interface
condition is stated in terms of the $B$-orthogonality of the enforcement
projections, not in terms of Lie-algebra commutativity.  Therefore the
cost is linear in $\dim(G)$ for every compact simple Lie group,
including non-abelian simple groups such as $E_6$ whose 78 generators
form an irreducible Killing-orthogonal basis.
\end{proof}

\smallskip\noindent\textit{Remark.}\quad This proposition closes an otherwise
tacit step: without it, the cost formula $C(G) = \dim(G)$ would implicitly
assume that gauge generators are C1-admissible without derivation.
The capacity-minimality comparisons throughout this section
($E_6 \colon 78 \text{ vs.\ } \SU(3) \colon 8$, etc.) all rely on
Proposition~\ref{prop:generator_separation}; it is the bridge from
$L_{\mathrm{Cauchy}}$'s $\mathbb{N}$-domain additivity to the linearity
of cost in $\dim(G)$.

% ======================================================================
% Paper 2 Technical Supplement — §5
% Gauge Template Uniqueness: Exhaustive Lie Algebra Classification
% ======================================================================
%
% PURPOSE.  Paper 2 §5.1 states L_gauge,unique in a compact table
% covering the 17 compact simple Lie algebras.  The reviewer asks (Q10):
% "Can you provide a theorem excluding other groups with complex 3D
% irreps or alternative minimal faithful carriers (e.g., embeddings
% into SU(N>3)) on 'capacity minimality' grounds, with a quantified
% cost function?"  This section provides the full classification with
% explicit exclusion reasons, quantified capacity costs, and
% closed-form proofs for every elimination.
%
% UPSTREAM: Theorem R (§3), L_Cauchy (§1), Lie algebra classification
% (imported mathematics: Killing 1888, Cartan 1894).
% ======================================================================


\subsection{Setup: the classification problem}
\label{sec:gauge_setup}

Theorem~R (\S\ref{sec:TR_closed}) provides three abstract carrier
requirements:
\begin{enumerate}[label=\textbf{R\arabic*},nosep,leftmargin=3em]
  \item A faithful complex carrier of dimension $\ge 3$ with irreducible
        totally antisymmetric invariant of rank $k \ge 3$.
  \item A faithful pseudoreal carrier of dimension exactly~$2$.
  \item A single abelian factor ($\U(1)$) with distinct charge
        assignments for all physical multiplets.
\end{enumerate}
R1--R3 are derived from A1 without reference to any specific Lie group.
The problem is now purely mathematical: \emph{which compact Lie group(s)
can host all three carriers simultaneously, at minimum total capacity cost?}

The capacity cost of a gauge group $G$ is $\dim(G) \cdot \varepsilon^*$
($L_{\mathrm{Cauchy}}$, Theorem~\ref{thm:LCauchy}): the number of
independent gauge generators, each costing one enforcement channel.
This cost functional is unique (no free parameter) and is the quantity
to be minimized.


\subsection{Exhaustive classification of the confining factor (R1)}
\label{sec:R1_classification}

\noindent\textit{Input.}  The classification of compact simple Lie algebras
is a completed theorem of pure mathematics.  There are exactly four infinite
families ($A_n$, $B_n$, $C_n$, $D_n$ for $n \ge 1, 2, 3, 4$ respectively)
and five exceptional algebras ($G_2$, $F_4$, $E_6$, $E_7$, $E_8$).
For each, the fundamental (lowest-dimensional faithful) representation
has a known dimension and reality type (real, pseudoreal, or complex).
This data is imported from standard references
\cite{Humphreys1972,FultonHarris1991} and is not framework-dependent.

\noindent\textit{Filters.}  R1 imposes three necessary conditions on the
confining factor $G_{\mathrm{conf}}$:
\begin{enumerate}[label=(F\arabic*),nosep]
  \item The fundamental representation must be \emph{complex}
        ($V \not\cong V^*$ as $G$-modules).  This is forced by $B1'$
        (Theorem~\ref{thm:B1prime}): real and pseudoreal carriers
        cannot sustain oriented composites.
  \item The representation must admit a totally antisymmetric invariant
        of rank $k \ge 3$.  This is forced by the oriented-composite
        requirement: bilinear ($k=2$) invariants produce self-conjugate
        composites ($B = \bar{B}$).
  \item The representation must be \emph{faithful}: every non-identity
        group element acts non-trivially.  This is forced by enforcement
        completeness (A1): if a gauge transformation acts trivially on
        all matter, it is physically undetectable and should not be
        counted.
\end{enumerate}

\begin{theorem}[$L_{\mathrm{gauge,unique}}$ --- Step 2: Confining Factor Classification]\label{thm:R1_class}
The only compact simple Lie algebra whose fundamental representation
satisfies (F1)--(F3) is $A_n = \mathfrak{su}(n+1)$ with $n \ge 2$
(equivalently, $\SU(N_c)$ with $N_c \ge 3$).
\end{theorem}

\begin{proof}
We test each family and each exceptional algebra against (F1)--(F3).
The full data is:

\smallskip
\begin{center}
\small
\renewcommand{\arraystretch}{1.15}
\begin{tabular}{@{}llccccl@{}}
\toprule
\textbf{Algebra} & \textbf{Group} & $\dim_{\mathrm{fund}}$ &
\textbf{Type} & \textbf{(F1)}  & \textbf{(F2)} & \textbf{Verdict} \\
\midrule
\multicolumn{7}{@{}l}{\textit{Family $A_n = \mathfrak{su}(n+1)$}} \\
$A_1$ & $\SU(2)$  &  2 & Pseudoreal & Fail & ---  & Excluded (F1) \\
$A_2$ & $\SU(3)$  &  3 & Complex    & Pass & $k{=}3$: Pass & \textbf{Pass} \\
$A_3$ & $\SU(4)$  &  4 & Complex    & Pass & $k{=}4$: Pass & Pass$^*$ \\
$A_4$ & $\SU(5)$  &  5 & Complex    & Pass & $k{=}5$: Pass & Pass$^*$ \\
$A_{n\ge5}$ & $\SU(n{+}1)$ & $n{+}1$ & Complex & Pass & $k{=}n{+}1$: Pass & Pass$^*$ \\[4pt]
\multicolumn{7}{@{}l}{\textit{Family $B_n = \mathfrak{so}(2n+1)$, $n \ge 2$}} \\
$B_2$ & $\mathrm{SO}(5)$  &  5 & Real & Fail & --- & Excluded (F1) \\
$B_3$ & $\mathrm{SO}(7)$  &  7 & Real & Fail & --- & Excluded (F1) \\
$B_{n\ge2}$ & $\mathrm{SO}(2n{+}1)$ & $2n{+}1$ & Real & Fail & --- & Excluded (F1) \\[4pt]
\multicolumn{7}{@{}l}{\textit{Family $C_n = \mathfrak{sp}(2n)$, $n \ge 3$}} \\
$C_3$ & $\mathrm{Sp}(6)$  &  6 & Pseudoreal & Fail & --- & Excluded (F1) \\
$C_{n\ge3}$ & $\mathrm{Sp}(2n)$ & $2n$ & Pseudoreal & Fail & --- & Excluded (F1) \\[4pt]
\multicolumn{7}{@{}l}{\textit{Family $D_n = \mathfrak{so}(2n)$, $n \ge 4$}} \\
$D_4$ & $\mathrm{SO}(8)$  &  8 & Real & Fail & --- & Excluded (F1) \\
$D_{n\ge4}$ & $\mathrm{SO}(2n)$ & $2n$ & Real & Fail & --- & Excluded (F1) \\[4pt]
\multicolumn{7}{@{}l}{\textit{Exceptional algebras}} \\
$G_2$ &          &  7 & Real       & Fail & --- & Excluded (F1) \\
$F_4$ &          & 26 & Real       & Fail & --- & Excluded (F1) \\
$E_6$ &          & 27 & Complex    & Pass & $k{=}3$: Pass & Excluded$^\dagger$ \\
$E_7$ &          & 56 & Pseudoreal & Fail & --- & Excluded (F1) \\
$E_8$ &          &248 & Real       & Fail & --- & Excluded (F1) \\
\bottomrule
\end{tabular}
\end{center}

\smallskip\noindent ${}^*$Pass$^*$: satisfies (F1)--(F3) but is
non-minimal (higher $\dim(G)$ than $\SU(3)$; eliminated in Step~3
below).

\smallskip\noindent ${}^\dagger$$E_6$ satisfies (F1) and (F2) but is
excluded by capacity minimality: $\dim(E_6) = 78$ vs.\
$\dim(\SU(3)) = 8$.  The fundamental representation of $E_6$ has
dimension 27, so even if one ignored cost, the matter field
degrees of freedom per generation would be vastly larger than the SM.

\medskip\noindent\textit{Mechanism of exclusion.}\quad
The $B_n$, $C_n$, $D_n$ families and the exceptionals $G_2$, $F_4$,
$E_7$, $E_8$ are all excluded by (F1) alone: their fundamental
representations are real or pseudoreal, and $B1'$
(Theorem~\ref{thm:B1prime}) proves that real/pseudoreal carriers
cannot sustain oriented composites.  This is not a numerical
coincidence --- it reflects the deep mathematical fact that only $A_n$
($n \ge 2$) and $E_6$ have complex fundamental representations among
compact simple Lie algebras.

$E_6$ passes (F1) and (F2) but is excluded on capacity grounds:
$\dim(E_6) = 78$ requires 78 gauge generators, each costing one
enforcement channel.  The $\SU(3)$ alternative requires 8 generators.
By $L_{\mathrm{Cauchy}}$ (Theorem~\ref{thm:LCauchy}), the cost
functional is linear in generator count, so $E_6$ costs
$78/8 = 9.75$ times as much as $\SU(3)$.  A1 (capacity minimality)
selects the cheapest adequate structure.

The only surviving family is $A_n = \SU(N_c)$ with $N_c \ge 3$.
\end{proof}

\smallskip\noindent\textit{Where the $E_6$ cut actually happens.}\quad
$E_6$ satisfies every representation-theoretic filter: it has complex
fundamental representations (F1) and antisymmetric invariants of rank
$\ge 3$ (F2).  The difference between ``$\SU(N_c)$ survives'' and
``$E_6$ is rejected'' is not a representation-theory result; it is a
capacity-minimality result under A1.  A reader who accepts the carrier
requirements R1--R3 but doubts the capacity-minimality clause of A1
has not yet been given a reason to prefer $\SU(3)$ over $E_6$.  The
minimality clause is what closes the gap, and it is derived from A1
in $L_{\mathrm{col}}$ and the enforcement-cost ledger.  The $E_6$
exclusion is therefore a direct consequence of the minimum-cost
ontology: $E_6$ is admissible (its cost fits below the capacity bound)
but is not minimum-cost; $\SU(3)$ is.  ``$E_6$ is rejected'' is
shorthand for ``$E_6$ has higher cost than $\SU(3)$ and both are
admissible,'' not for an act of rejection.


\subsection{Capacity optimization: $N_c = 3$}
\label{sec:Nc_optimization}

Within the surviving $\SU(N_c)$ family, all $N_c \ge 3$ satisfy
(F1)--(F3).  The antisymmetric invariant has rank $k = N_c$
(the $\varepsilon$-tensor $\varepsilon_{i_1 \cdots i_{N_c}}$).
The oriented-composite requirement forces $k \ge 3$
($B1'$ excludes $k = 2$), so $N_c \ge 3$ is already established.

Two further filters apply:

\begin{proposition}[Witten anomaly excludes even $N_c$]\label{prop:witten}
For gauge group $\SU(N_c) \times \SU(2) \times \U(1)$, the number of
$\SU(2)$ doublets per generation is $N_c + 1$ (one quark doublet
contributing $N_c$ colors, plus one lepton doublet).  The Witten anomaly
\cite{Witten1982} requires this number to be even.  Therefore $N_c$ must
be odd.
\end{proposition}

\noindent\textit{Joint consistency with $T_{\mathrm{field}}$.}\quad
The statement ``one quark doublet plus one lepton doublet per generation''
is the SM-like template structure that emerges from $T_{\mathrm{field}}$
(\S\ref{sec:fermion_scan}) rather than an input independent of it.  The
Witten-anomaly cut applied in the current section therefore functions as
a \emph{joint consistency} check with the fermion scan: for any odd
$N_c \ge 3$, the minimum-DOF anomaly-free template has the (quark-doublet,
lepton-doublet) structure, and the Witten filter is satisfied; for any
even $N_c$, no admissible template of this structure survives.  The
logical ordering is not ``N$_c$ is fixed in \S5 then templates are scanned
in \S6.''  Rather, the admissible N$_c$'s and admissible fermion templates
are a joint optimum, and the supplement exhibits that joint optimum by
first narrowing N$_c$ to odd values under the generic template ansatz
(this section) and then closing N$_c = 3$ by the full scan (\S\ref{sec:fermion_scan}).
Readers who prefer a single-step derivation may take Propositions~\ref{prop:witten}
and \ref{prop:Nc3} together with Theorem~\ref{thm:Tfield} as a simultaneous
result: $(N_c = 3, \text{SM template})$ is the unique joint
minimum-cost admissible pair.

\begin{proposition}[Capacity minimality selects $N_c = 3$]\label{prop:Nc3}
The total gauge group dimension is
$\dim(G) = (N_c^2 - 1) + 3 + 1 = N_c^2 + 3$.  For odd $N_c \ge 3$:

\smallskip
\begin{center}
\small
\begin{tabular}{@{}ccccl@{}}
\toprule
$N_c$ & Witten & $\dim(G)$ & $C_{\mathrm{total}}$ & Verdict \\
\midrule
3 & Pass ($4$ doublets) & 12 & 61 & \textbf{Minimum cost} \\
5 & Pass ($6$ doublets) & 28 & 97 & Non-minimal ($2.3\times$) \\
7 & Pass ($8$ doublets) & 52 & 141 & Non-minimal ($4.3\times$) \\
9 & Pass ($10$ doublets) & 84 & 193 & Non-minimal ($7\times$) \\
\bottomrule
\end{tabular}
\end{center}

\smallskip\noindent $C_{\mathrm{total}} = (4N_c + 3) \cdot 3 + 4 + (N_c^2 + 3)$:
fermions $+$ Higgs $+$ gauge bosons.  The cost is strictly increasing
and unbounded.  $N_c = 3$ is the unique minimum.
\end{proposition}


\subsection{Chiral factor classification (R2)}
\label{sec:R2_classification}

\begin{proposition}[$\SU(2)$ is the unique chiral factor]\label{prop:R2_unique}
R2 requires a compact simple Lie group with a faithful pseudoreal
representation of dimension exactly~$2$.  Among all compact simple Lie
algebras:
\begin{enumerate}[label=(\roman*),nosep]
  \item $\mathfrak{su}(2)$ has a 2-dimensional irreducible representation
        (the fundamental doublet) which is pseudoreal
        ($J = i\sigma_2 K$, $J^2 = -1$).
  \item All other compact simple Lie algebras have minimum faithful
        dimension $\ge 3$: $\mathfrak{su}(n+1)$ for $n \ge 2$ has
        $\dim_{\mathrm{fund}} = n + 1 \ge 3$; $\mathfrak{so}(2n+1)$
        has $\dim_{\mathrm{fund}} \ge 5$; $\mathfrak{sp}(2n)$ for $n \ge 2$
        has $\dim_{\mathrm{fund}} \ge 4$; all exceptionals have
        $\dim_{\mathrm{fund}} \ge 7$.
\end{enumerate}
Therefore $\SU(2)$ is the unique compact simple Lie group satisfying R2.

At the group level, the gauge factor must be $\SU(2)$ rather than its
quotient $\mathrm{SO}(3) = \SU(2)/\mathbb{Z}_2$: the doublet
representation does not descend to $\mathrm{SO}(3)$ (it is not
single-valued under $2\pi$ rotations), so faithfulness requires
the simply connected cover $\SU(2)$.
\end{proposition}


\subsection{Abelian factor and simple-group alternatives (R3)}
\label{sec:R3_simple}

\begin{proposition}[$\U(1)$ is the unique abelian factor]\label{prop:R3_unique}
R3 requires exactly one abelian factor.  $\U(1) \cong \mathbb{R}/\mathbb{Z}$
is the unique connected compact 1-dimensional abelian Lie group.  A
second $\U(1)$ would be non-minimal under A1: one $\U(1)$ with distinct
hypercharge assignments already resolves all multiplet degeneracies
(Theorem~R, R3 derivation in \S\ref{sec:carrier_prereqs}).  Additional
$\U(1)$ factors add $\ge 1$ gauge generator each with no enforcement gain.
\end{proposition}

\begin{proposition}[Simple-group alternatives are excluded]\label{prop:no_simple}
Any simple compact Lie group $G_{\mathrm{simple}}$ containing
$\SU(3) \times \SU(2) \times \U(1)$ as a subgroup has
$\dim(G_{\mathrm{simple}}) \ge 24$:

\smallskip
\begin{center}
\small
\begin{tabular}{@{}lccc@{}}
\toprule
\textbf{Group} & $\dim$ & \textbf{Cost ratio} & \textbf{Contains SM?} \\
\midrule
$\SU(3) \times \SU(2) \times \U(1)$ & 12 & $1.0\times$ & (Product) \\
$\SU(5)$   & 24 & $2.0\times$ & Yes (Georgi--Glashow) \\
$\mathrm{SO}(10)$ & 45 & $3.75\times$ & Yes \\
$E_6$ & 78 & $6.5\times$ & Yes \\
\bottomrule
\end{tabular}
\end{center}

\smallskip\noindent Every simple envelope costs at least $2\times$ the
product structure.  A1 (capacity minimality) selects the product.
Additionally, enforcement independence (Step~1 of
$L_{\mathrm{gauge,unique}}$) requires factorization: $T_M$ (monogamy)
and $L_{\mathrm{loc}}$ (locality) force independent enforcement
channels to live in commuting subgroups with independent couplings.
A simple group has a single coupling constant, contradicting the
requirement that the confining factor (which confines at low energy)
and the chiral factor (which must \emph{not} confine, to preserve
chirality) run independently.
\end{proposition}


\subsection{Consolidation: unique gauge template}
\label{sec:gauge_consolidation}

Combining Steps~1--6:

\begin{theorem}[$L_{\mathrm{gauge,unique}}$: Gauge Template Uniqueness]\label{thm:gauge_unique}
The unique gauge group satisfying R1--R3 at minimum capacity cost is
\[
  G \;=\; \SU(3) \times \SU(2) \times \U(1),
  \qquad \dim(G) = 12.
\]
No alternative compact Lie group structure satisfies all three carrier
requirements.  No simple-group envelope is capacity-competitive.  No
additional abelian factor is needed.

\smallskip\noindent\textit{Derivation chain:}\quad
$\text{A1} \to \text{Theorem R (R1+R2+R3)} \to
\text{Lie classification} \to
\SU(N_c) \times \SU(2) \times \U(1) \to
\text{Witten} + \text{minimality} \to N_c = 3$.

\smallskip\noindent\textit{Imported mathematics:}\quad
The classification of compact simple Lie algebras (Killing 1888, Cartan 1894)
and the Witten anomaly theorem (Witten 1982) are used as external
mathematical inputs.  Both are established results with independent
proofs; neither is framework-specific.  All physical inputs (R1--R3)
are derived from A1 in \S\ref{sec:carrier_prereqs}.
\end{theorem}

\coderef{check\_L\_gauge\_template\_uniqueness}{gauge.py}



\newpage


% ══════════════════════════════════════════════════════════════
\section{Fermion content: exhaustive scan}
\label{sec:fermion_scan}
% ══════════════════════════════════════════════════════════════

\smallskip\noindent\textit{Descriptive reading of the scan.}\quad
The enumeration in this section is \emph{not} a physical process by
which nature ``chooses'' the Standard Model fermion content from a
space of candidates.  The admissibility conditions (five exclusion
proofs bounding the search space, plus seven sequential filters
encoding anomaly cancellation, asymptotic freedom, and chirality)
define a predicate
$\mathrm{Adm}(\mathcal{T})$ on chiral fermion templates.  The scan
is an \emph{algorithmic verification} that the set
$\{\mathcal{T} : \mathrm{Adm}(\mathcal{T})\}$ has a unique member up
to the stated equivalences --- namely the 45-Weyl-fermion SM
template.  The 1{,}680 distinct templates (4{,}680 ordered
iterations) are artefacts of the verification procedure, not
candidates in a dynamical competition.  The claim of this section is
therefore descriptive: ``the admissibility predicate is satisfied by
exactly one template,'' not ``the framework selects one template
from among 1{,}680.''

This section specifies the 1{,}680-template search space, defines the
seven sequential filters, provides the survivor count after each filter,
and proves that no candidates exist outside the explicitly bounded set.

% ── §6: Fermion content: exhaustive scan ──

\subsection{Template definition}
\label{sec:template_def}

A \emph{chiral fermion template} is a finite multiset
$\mathcal{T} = \{(r_i^{(3)}, r_i^{(2)})\}_{i=1}^{|\mathcal{T}|}$
where $r_i^{(3)}$ is an irreducible representation of $\SU(3)$ and
$r_i^{(2)}$ is an irreducible representation of $\SU(2)$.  Each entry
represents a left-handed Weyl fermion field.  The total number of Weyl
degrees of freedom per generation is
$\mathrm{DOF} = \sum_i \dim(r_i^{(3)}) \cdot \dim(r_i^{(2)})$.
The template describes one generation; the full content is
$N_{\mathrm{gen}} \times \mathrm{DOF}$ with $N_{\mathrm{gen}} = 3$
(derived in Paper~2, \S9).

The search space is bounded by five closed-form exclusion proofs that
restrict both the allowed representations and the number of field types.


\subsection{Search space bounds: five exclusion proofs}
\label{sec:exclusion_proofs}

\begin{proposition}[P1: Large $\SU(3)$ representations are AF-excluded]
\label{prop:P1}
Any $\SU(3)$ representation with dimension $\ge 10$ (i.e.,
$\mathbf{10}$, $\mathbf{15}$, $\ldots$) exceeds the $\SU(3)$
asymptotic freedom budget as a single field.  For the $\mathbf{10}$:
the Dynkin index is $T(\mathbf{10}) = 15/2$; the AF contribution of
a single field in the $\mathbf{10}$ with $\SU(2)$ singlet and 3
generations is $(4/3)(15/2)(1)(3) = 30 > 11 = (11/3) \cdot 3$.
Therefore no template containing any $\SU(3)$ rep of dimension
$\ge 10$ can be asymptotically free.  The allowed $\SU(3)$
representations are: $\mathbf{1}, \mathbf{3}, \bar{\mathbf{3}},
\mathbf{6}, \bar{\mathbf{6}}, \mathbf{8}$.
\end{proposition}

\begin{proposition}[P2: Colored $\SU(2)$ triplets and higher are AF-excluded]
\label{prop:P2}
Any colored field ($\dim(r^{(3)}) \ge 3$) in an $\SU(2)$
representation of dimension $\ge 3$ exceeds the $\SU(2)$ AF budget.
For a $(\mathbf{3}, \mathbf{3})$ field: $T(\mathbf{3}_{\SU(2)}) = 2$,
$\dim(\mathbf{3}_{\SU(3)}) = 3$, contribution $= (4/3)(2)(3)(3) = 24
> 22/3 \cdot 2 = 44/3 \approx 14.7$.  Therefore colored fields are
restricted to $\SU(2)$ representations $\{\mathbf{1}, \mathbf{2}\}$.
\end{proposition}

\begin{proposition}[P3: Colorless $\SU(2)$ triplets are DOF-excluded]
\label{prop:P3}
Replacing a colorless $\SU(2)$ doublet $(\mathbf{1}, \mathbf{2})$ with
a triplet $(\mathbf{1}, \mathbf{3})$ adds $(3-2) \times 3 = 3$ DOF per
3 generations.  Since the SM template already has 45 DOF and is the
minimum survivor (as verified by the scan), any template with colorless
$\SU(2)$ triplets has DOF $\ge 48 > 45$ and is non-minimal.
\end{proposition}

\begin{proposition}[P4: Multiple colored doublets are DOF-excluded]
\label{prop:P4}
A template with two colored $\SU(2)$ doublets requires at least two
colored singlets for $[\SU(3)]^3$ anomaly cancellation plus lepton content.
The tightest such template is
$(\mathbf{3},\mathbf{2}) + (\bar{\mathbf{3}},\mathbf{2}) +
(\mathbf{3},\mathbf{1}) + (\bar{\mathbf{3}},\mathbf{1}) +
(\mathbf{1},\mathbf{2}) + (\mathbf{1},\mathbf{1})$, with
DOF $= (6 + 6 + 3 + 3 + 2 + 1) \times 3 = 63 > 45$.
\end{proposition}

\begin{proposition}[P5: More than 5 field types are DOF-excluded]
\label{prop:P5}
Each field type contributes at least $1 \times 1 \times 3 = 3$ DOF
(the smallest possible: a colorless $\SU(2)$ singlet with 3 generations).
A template with $> 5$ field types has DOF $\ge 45 + 3 > 45$ and is
non-minimal.
\end{proposition}

\noindent These five proofs establish that the search space is
\emph{finite and bounded}: $\SU(3)$ reps $\in \{3, \bar{3}, 6,
\bar{6}, 8\}$ (excluding singlets, which are colorless),
$\SU(2)$ reps $\in \{1, 2\}$, and $|\mathcal{T}| \le 5$.
Combined with multiplicity choices (0--3 colored singlets drawn
with replacement, 0--2 lepton singlets, doublet present or absent),
this yields 1{,}680 distinct candidate templates.\footnote{Colored
singlets are unordered within a template, so the correct enumeration
uses combinations with replacement: $5 \times (1 + 5 + 15 + 35)
\times 2 \times 3 = 1{,}680$.  A standalone verification script
with exact rational arithmetic and four built-in red-team challenges
is archived at \url{https://doi.org/10.5281/zenodo.19154197}; an
interactive Colab walkthrough is at
\url{https://colab.research.google.com/drive/1otO1Qwp_Lr3OsK6ykkeMRATanecPkKg6};
source code is at
\url{https://github.com/Ethan-Brooke/APF-Paper-2-The-Structure-of-Admissible-Physics}.}


\subsection{Seven filters}
\label{sec:seven_filters}

Each filter is derived from the framework.  They are applied
sequentially; each eliminates candidates that passed all previous
filters.

\begin{enumerate}[label=\textbf{F\arabic*},nosep,leftmargin=3em]
  \item \textbf{$\SU(3)$ asymptotic freedom.}\quad
  $b_0^{(3)} = 11 - (4/3) \sum_i T(r_i^{(3)}) \dim(r_i^{(2)}) \cdot N_{\mathrm{gen}} > 0$.
  Derived from non-abelian self-interaction dominance (\S\ref{sec:beta_schematic}).

  \item \textbf{$\SU(2)$ asymptotic freedom.}\quad
  $b_0^{(2)} = 22/3 - (4/3) \sum_i T(r_i^{(2)}) \dim(r_i^{(3)}) \cdot N_{\mathrm{gen}} > 0$.

  \item \textbf{Chirality.}\quad
  The template must contain both colored doublets and colored singlets.
  A template with only colored doublets or only colored singlets is
  vector-like and has zero intrinsic gauge irreversibility (R2).

  \item \textbf{$[\SU(3)]^3$ anomaly cancellation.}\quad
  $\sum_i A(r_i^{(3)}) \cdot \dim(r_i^{(2)}) = 0$, where $A$ is the
  cubic anomaly coefficient ($A(\mathbf{3}) = 1/2$,
  $A(\bar{\mathbf{3}}) = -1/2$, $A(\mathbf{6}) = 5/2$,
  $A(\bar{\mathbf{6}}) = -5/2$, $A(\mathbf{8}) = 0$).

  \item \textbf{Witten anomaly freedom.}\quad
  The total number of $\SU(2)$ doublets (counted with $\SU(3)$ multiplicity)
  must be even: $\sum_i \dim(r_i^{(3)}) \cdot [r_i^{(2)} = \mathbf{2}]
  \equiv 0 \pmod{2}$.

  \item \textbf{Full anomaly system solvable.}\quad
  The combined system of $[\SU(2)]^2[\U(1)]$, $[\SU(3)]^2[\U(1)]$,
  $[\mathrm{grav}]^2[\U(1)]$, and $[\U(1)]^3$ anomaly cancellation
  conditions must admit rational hypercharge solutions.  Templates
  containing colored singlets with $A(R) = 0$ (e.g.\ the $\mathbf{8}$)
  are handled by spectator reduction: such fields carry $Y = 0$ and
  decouple from all anomaly conditions; the reduced template is then
  tested analytically.  For the spectator-free core, the system reduces
  to a quadratic in $z = Y_u/Y_Q$ having rational roots
  (\S\ref{sec:hypercharge}).

  \item \textbf{CPT quotient.}\quad
  Templates related by charge conjugation (replacing every representation
  with its conjugate) describe the same physics.  The CPT equivalence
  class is identified and only the canonical representative is counted.
\end{enumerate}


\subsection{Waterfall table}
\label{sec:waterfall}

\begin{center}
\small
\begin{tabular}{@{}lrrl@{}}
\toprule
\textbf{Filter} & \textbf{Survivors} & \textbf{Eliminated} & \textbf{Mechanism} \\
\midrule
Search space (P1--P5) & 1{,}680 & --- & Bounded by exclusion proofs \\
F1+F2: Asymptotic freedom & 348 & 1{,}332 & $b_0^{(3)} > 0$ and $b_0^{(2)} > 0$ \\
F3: Chirality & 324 & 24 & Colored doublets + singlets required \\
F4: $[\SU(3)]^3$ anomaly & 24 & 300 & $\sum A_3 d_2 = 0$ \\
F5: Witten anomaly & 12 & 12 & $\#$doublets even \\
F6: Full anomaly system & 8 & 4 & Rational hypercharges exist \\
F7: CPT quotient & 4 & 4 & Conjugate pairs identified \\
Minimality & \textbf{1} & 3 & Min DOF = 45 selected \\
\bottomrule
\end{tabular}
\end{center}

\smallskip\noindent\textit{Counting convention (narrative vs.\ enumeration).}\quad
The counts in the table above are the number of \emph{distinct}
templates surviving each filter, with unordered colored-singlet positions.
The reference implementation (\texttt{check\_T\_field} in
\texttt{gauge.py} and \texttt{fermion\_scan\_standalone.py}) uses an
ordered-tuple enumeration (\texttt{itertools.product}) and therefore
tests each distinct template multiple times, producing a raw iteration
count of $4{,}680$ rather than $1{,}680$, and downstream stage counts
$48 \to 24 \to 4 \to 2 \to 1$ after quotienting by the ordering
redundancy $48/2 = 24$, etc.  The survivor \emph{set} (as an unordered
collection of templates) is identical under either convention, and the
unique minimum-DOF winner is the Standard Model in both cases.
A reviewer running the code will see the enumeration multiplicities;
the table reports the unordered count because that is what the
five exclusion proofs (P1--P5) actually bound.

\smallskip\noindent\textit{Spectator octet handling.}\quad
Templates~\#5--\#8 in the post-F6 listing below include the adjoint
$\mathbf{8}$ as a colored $\SU(2)$ singlet.  Because $A(\mathbf{8}) = 0$,
spectator-reduction analysis (removing the anomaly-free spectator field
and testing the reduced template) keeps these templates through F6 in
the prose derivation, where they fail only at the minimality stage
(DOF $\ge 69 > 45$).  The reference implementation takes a shorter
path: its F6 predicate (\texttt{passes\_full\_anomaly}) requires every
colored singlet to have SU(3) dimension equal to $N_c = 3$, which
eliminates the $(\mathbf{8},\mathbf{1})$ at F6 rather than at minimality.
Both treatments converge on the same survivor set (the SM at DOF $= 45$),
and the two are equivalent for admissibility bookkeeping: spectator
octets are eliminated either by analytical reduction plus minimality,
or by the uniform-dimension test, at no change to the final result.
The table's F6 count of 8 is the analytical count before spectator
reduction; the code's F6 count of 4 is the count after uniform-dimension
filtering.  Both routes agree after CPT quotient and minimality.

\smallskip\noindent The final winner is exact and reproduced by
\texttt{check\_T\_field} in \texttt{gauge.py} and independently by
\texttt{fermion\_scan\_standalone.py}~\cite{BrookeFermionScan}.

\begin{proposition}[Filter-order independence]\label{prop:filter_order}
The survivor set $\mathcal{S} = \{\mathcal{T} : \mathcal{T} \text{
passes F1--F7}\}$ is independent of the order in which the seven
filters are applied.
\end{proposition}

\begin{proof}
Each filter $F_k$ defines a predicate $P_k: \mathcal{T} \to
\{0, 1\}$ on templates.  A template survives all filters iff it
satisfies the conjunction $P_1 \wedge P_2 \wedge \cdots \wedge P_7$.
Logical conjunction is commutative and associative:
$P_{\sigma(1)} \wedge \cdots \wedge P_{\sigma(7)} = P_1 \wedge
\cdots \wedge P_7$ for any permutation $\sigma \in S_7$.  Therefore
the survivor set $\mathcal{S} = \{\mathcal{T} : \bigwedge_k
P_k(\mathcal{T}) = 1\}$ is identical regardless of evaluation order.

What changes with filter order is the \emph{waterfall count} (number
of templates eliminated at each stage), because a template eliminated
by $F_k$ in the standard order might instead be eliminated by
$F_j$ with $j > k$ if the order is reversed.  The final survivor
set, however, is order-invariant.  The standalone verification
script confirms this: running all $7! = 5{,}040$ permutations on
the 1{,}680-template set yields the same 4 post-CPT survivors (and
the same unique DOF-45 winner) in every case.
\end{proof}

\smallskip\noindent\textit{Search-space bias and completeness.}\quad
A distinct concern is whether the search space itself (the 1{,}680
templates bounded by P1--P5) introduces bias by excluding candidates
that might survive the filters.  This is addressed by the five
exclusion proofs, each of which proves an \emph{unconditional} bound:
P1 excludes $\SU(3)$ reps of dimension $\ge 10$ because a single
such field exceeds the AF budget regardless of what other fields are
present; P2--P5 similarly establish bounds that depend only on the
AF constraints and the DOF lower bound, not on the other filters.
Crucially, the exclusion proofs use only AF and DOF bounds, which are
\emph{monotone}: adding fields to a template can only increase the
DOF and decrease the AF margin.  Therefore any template outside the
search space fails either AF or minimality, and would be eliminated
by F1/F2 or the DOF bound regardless of the remaining filters.
The search space is the exact set of templates that \emph{could}
survive; no candidates are hidden outside it.

\smallskip\noindent\textit{Explicit listing of all post-F6 survivors.}\quad
The 8 templates surviving F1--F6 (before CPT quotient) are:

\smallskip
\begin{center}\small
\renewcommand{\arraystretch}{1.15}
\begin{tabular}{@{}clcl@{}}
\toprule
\# & Template & DOF & Status \\
\midrule
1 & $(3,2) + (\bar{3},1) + (\bar{3},1) + (1,2) + (1,1)$ & 45 & SM \\
2 & $(\bar{3},2) + (3,1) + (3,1) + (1,2) + (1,1)$ & 45 & CPT of \#1 \\
3 & $(3,2) + (\bar{3},1) + (\bar{3},1) + (1,2) + (1,1) + (1,1)$ & 48 & SM + extra singlet \\
4 & $(\bar{3},2) + (3,1) + (3,1) + (1,2) + (1,1) + (1,1)$ & 48 & CPT of \#3 \\
5 & $(3,2) + (\bar{3},1) + (\bar{3},1) + (8,1) + (1,2) + (1,1)$ & 69 & SM + spectator octet \\
6 & $(\bar{3},2) + (3,1) + (3,1) + (8,1) + (1,2) + (1,1)$ & 69 & CPT of \#5 \\
7 & $(3,2) + (\bar{3},1) + (\bar{3},1) + (8,1) + (1,2) + (1,1) + (1,1)$ & 72 & SM + octet + singlet \\
8 & $(\bar{3},2) + (3,1) + (3,1) + (8,1) + (1,2) + (1,1) + (1,1)$ & 72 & CPT of \#7 \\
\bottomrule
\end{tabular}
\end{center}

\smallskip\noindent Templates \#5--\#8 contain the adjoint $\mathbf{8}$
as a colored $\SU(2)$ singlet.  Since $A(\mathbf{8}) = 0$, this field
carries $Y = 0$ and decouples from all anomaly conditions---it is an
anomaly-free spectator.  Its DOF cost ($8 \times 3 = 24$ per generation)
is far above the minimum.

\smallskip\noindent After the CPT quotient (F7), templates \#2, \#4,
\#6, and \#8 are identified with \#1, \#3, \#5, and \#7 respectively,
leaving 4 CPT-inequivalent survivors.  Minimality selects DOF $= 45$,
giving exactly one winner: the Standard Model.  The second-lightest
survivor (\#3, DOF $= 48$) differs from the SM by one additional
right-handed neutrino-like singlet per generation;
its exclusion by minimality is the $C = 64 \to 5\sigma$ result of
\S\ref{sec:neutrinos}.


\subsection{The unique survivor}
\label{sec:unique_survivor}

\begin{theorem}[$T_{\mathrm{field}}$: SM Fermion Content]\label{thm:Tfield}
The unique minimal chiral fermion template consistent with
$\SU(3) \times \SU(2) \times \U(1)$, three generations, and all
admissibility constraints is:
\[
  \{Q(\mathbf{3},\mathbf{2}),\;
    u^c(\bar{\mathbf{3}},\mathbf{1}),\;
    d^c(\bar{\mathbf{3}},\mathbf{1}),\;
    L(\mathbf{1},\mathbf{2}),\;
    e^c(\mathbf{1},\mathbf{1})\}
  \times 3 \text{ generations},
\]
with $\mathrm{DOF} = (6+3+3+2+1) \times 3 = 45$ Weyl fermions.
\end{theorem}

\noindent No other template among the 1{,}680 candidates has
$\mathrm{DOF} \le 45$ and passes all seven filters.
The five exclusion proofs (P1--P5) guarantee that no candidate outside
the search space can have $\mathrm{DOF} < 45$.  The derivation is
therefore exhaustive: the SM fermion content is the unique solution.

\coderef{check\_T\_field}{gauge.py}; standalone verification: \texttt{fermion\_scan\_standalone.py}~\cite{BrookeFermionScan}


\newpage


% ══════════════════════════════════════════════════════════════
\section{Anomaly cancellation and hypercharge uniqueness}
\label{sec:hypercharge}
% ══════════════════════════════════════════════════════════════

% ── §7: Anomaly cancellation and hypercharge uniqueness ──

\subsection{The four anomaly conditions}
\label{sec:anomaly_conditions}

The quantum consistency of a chiral gauge theory requires the
cancellation of all gauge anomalies.  For
$\SU(3) \times \SU(2) \times \U(1)$ with the surviving template, there
are four independent conditions on the five hypercharges
$\{Y_Q, Y_u, Y_d, Y_L, Y_e\}$.  Three are linear:
\begin{align}
  [\SU(2)]^2[\U(1)] &= 0: &
    N_c\, Y_Q + Y_L &= 0, \label{eq:anom1} \\
  [\SU(3)]^2[\U(1)] &= 0: &
    2Y_Q - Y_u - Y_d &= 0, \label{eq:anom2} \\
  [\mathrm{grav}]^2[\U(1)] &= 0: &
    2N_c\, Y_Q + 2Y_L - N_c\, Y_u - N_c\, Y_d - Y_e &= 0. \label{eq:anom3}
\end{align}
The fourth is the cubic anomaly:
\begin{equation}
  [\U(1)]^3 = 0: \quad
  2N_c\, Y_Q^3 + 2Y_L^3 - N_c\, Y_u^3 - N_c\, Y_d^3 - Y_e^3 = 0.
  \label{eq:anom4}
\end{equation}

\subsection{Reduction to a quadratic}
\label{sec:reduction_quadratic}

Solving (\ref{eq:anom1})--(\ref{eq:anom3}) gives
$Y_L = -N_c\, Y_Q$, $Y_d = 2Y_Q - Y_u$, $Y_e = -2N_c\, Y_Q$.
Substituting into (\ref{eq:anom4}) and dividing by $Y_Q^3 \ne 0$
(non-trivial grading), define $z := Y_u / Y_Q$:
\begin{equation}
  z^2 - 2z - (N_c^2 - 1) = 0.
  \label{eq:z_quadratic}
\end{equation}

\noindent\textit{Why the roots are rational.}\quad
Equation~(\ref{eq:z_quadratic}) is a \emph{monic} polynomial (leading
coefficient $= 1$) with integer coefficients ($-2$ and $-(N_c^2-1)$,
where $N_c \in \mathbb{Z}$).  By the rational root theorem for monic
integer polynomials: if $z = p/q$ in lowest terms is a root, then
$q \mid 1$, so $z \in \mathbb{Z}$.  The discriminant
$\Delta = 4 + 4(N_c^2 - 1) = 4N_c^2$ is a perfect square, confirming
integer roots:
\begin{equation}
  z = 1 \pm N_c.
\end{equation}
Since $Y_L = -N_c Y_Q$, $Y_d = 2Y_Q - Y_u = (2-z)Y_Q$, and
$Y_e = -2N_c Y_Q$ are all integer multiples of $Y_Q$, and $Y_Q$ itself
is a free normalization (any nonzero rational value is admissible, since
it sets the overall scale of the $\U(1)$ generator), choosing
$Y_Q = 1/6$ gives $Y_i \in \mathbb{Q}$ for all five multiplets.
Charge quantization --- the fact that all hypercharges are rational
multiples of a single unit --- is thus a \emph{consequence} of anomaly
cancellation, not an input.
For $N_c = 3$: $z \in \{4, -2\}$.  These are related by
$z \to 2 - z$ (the $u \leftrightarrow d$ relabeling), so they describe
the same physics.  Choosing $z = 4$ and normalizing $Y_Q = 1/6$:

\begin{theorem}[$L_{\mathrm{hypercharge}}$: Unique Hypercharge Assignment]
\label{thm:hypercharge}
\[
  Y_Q = \tfrac{1}{6}, \quad
  Y_u = \tfrac{2}{3}, \quad
  Y_d = -\tfrac{1}{3}, \quad
  Y_L = -\tfrac{1}{2}, \quad
  Y_e = -1.
\]
These are the unique hypercharges (up to normalization and $u \leftrightarrow d$)
consistent with anomaly cancellation for the derived template with $N_c = 3$.
\end{theorem}

\noindent\textit{Consequences.}\quad Electric charges $Q = T_3 + Y$ give
$Q(u) = 2/3$, $Q(d) = -1/3$, $Q(e) = -1$, $Q(\nu) = 0$.  Charge
quantization (all charges are multiples of $1/3$) is derived from the
anomaly structure, not assumed.  The quark--lepton relation
$Y_L = -N_c\, Y_Q$ ties lepton hypercharge to the number of colors.

\coderef{check\_T\_field (hypercharge section)}{gauge.py}


\newpage


% ══════════════════════════════════════════════════════════════
\section{Neutrino masses and the 45-fermion prediction}
\label{sec:neutrinos}
% ══════════════════════════════════════════════════════════════

The derived fermion content (45 Weyl fermions per 3 generations) contains
no right-handed neutrino.  This is a prediction, not a gap.

% ── §8: Neutrino masses and the 45-fermion prediction ──

The 45-fermion count is the unique minimum-DOF survivor of the 7-filter
scan.  It contains no right-handed neutrino $\nu_R$.  This is a
\emph{prediction}: the framework derives the minimal fermion content
and finds that $\nu_R$ is not among it.

\begin{proposition}[Neutrino mass mechanism]\label{prop:neutrino_mass}
Neutrino masses are accommodated within the 45-fermion framework via
the dimension-5 Weinberg operator:
\[
  \mathcal{O}_5 \;=\; \frac{c_{ij}}{\Lambda}\,(L_i H)(L_j H),
\]
where $L_i$ is the lepton doublet of generation $i$, $H$ is the Higgs
doublet, and $\Lambda$ is a high scale.  After electroweak symmetry
breaking ($\langle H \rangle = v$), this generates Majorana masses
$m_\nu \sim c\, v^2 / \Lambda$ for left-handed neutrinos.
\end{proposition}

\noindent\textit{Why $\nu_R$ is not counted.}\quad The Weinberg operator
can be UV-completed by a heavy Majorana fermion at scale $\Lambda$
(the type-I seesaw).  This heavy state is above the IR saturation scale
($T_{\mathrm{confinement}}$) and is not a distinct enforcement channel
at low energies.  The capacity count ($C_{\mathrm{total}} = 61$)
reflects the \emph{low-energy} enforcement ledger; states above the
saturation scale contribute to the UV completion but not to the
structural partition.

\smallskip\noindent\textit{Prediction.}\quad Neutrino masses are
Majorana (lepton number violated by 2 units).  Neutrinoless double-beta
decay ($0\nu\beta\beta$) is predicted to occur at a rate determined by
the Majorana mass matrix.

\smallskip\noindent\textit{Observational constraint.}\quad If $\nu_R$
were added as a low-energy enforcement channel (3 species for 3
generations), $C_{\mathrm{total}} = 64$ and
$\Omega_\Lambda = 42/64 = 0.656$.  Planck measures
$\Omega_\Lambda = 0.6889 \pm 0.0056$~\cite{Planck2020}: the 64-type
prediction is $>5\sigma$ from observation.  The 45-fermion count is
therefore observationally constrained, not merely a theoretical
preference for minimality.

\coderef{check\_T\_field}{gauge.py}


\newpage

% ╔═══════════════════════════════════════════════════════════════╗
% ║              PART III: CAPACITY COUNTING AND COSMOLOGY        ║
% ╚═══════════════════════════════════════════════════════════════╝

\part{Capacity Counting and Cosmology}
\label{part:cosmo}


% ══════════════════════════════════════════════════════════════
\section{Capacity counting and the BRST question}
\label{sec:capacity_count}
% ══════════════════════════════════════════════════════════════

% ── §9: Capacity counting and the BRST question ──

\subsection{Operational definition of capacity types}

A \emph{capacity type} is a structurally distinct enforcement role
in the low-energy ledger.  The count is:
\begin{equation}
  C_{\mathrm{total}} = \underbrace{45}_{\text{fermion types}}
    + \underbrace{4}_{\text{Higgs real components}}
    + \underbrace{12}_{\text{gauge generators}} = 61.
\end{equation}

\noindent\textit{Fermion types} (45): the 15 Weyl fermions per
generation ($Q = 6$ colored-doublet components $+$ $u^c = 3$
colored-singlet $+$ $d^c = 3$ $+$ $L = 2$ weak-doublet components
$+$ $e^c = 1$) times 3 generations.  Color and weak-isospin
multiplicity \emph{are} counted here as independent matter-sector
enforcement roles, because each distinct quantum-number configuration
is a separately enforceable distinction in the capacity ledger.  The
arithmetic is $(6+3+3+2+1) \times 3 = 15 \times 3 = 45$, consistent
with \texttt{check\_L\_count} in \texttt{gauge.py} (\texttt{per\_gen =
6+3+3+2+1 = 15}).  Internal-structure gauge channels (the 27
fermion-internal channels counted in $L_{\mathrm{vac}}$, Paper~2
\S7.1) are a distinct bookkeeping allocated to the vacuum sector;
they are not the matter-sector fermion types counted here.  This
resolves an earlier prose formulation that read ``the 5 multiplet
types'' and counted $5 \times 3 = 15$ rather than 45; the correct
count is $(6+3+3+2+1) \times 3 = 45$.

\noindent\textit{Higgs components} (4): the 4 real degrees of freedom
of the Higgs doublet.  After SSB, 3 become Goldstone bosons (absorbed
into $W^\pm, Z$) and 1 is the physical Higgs.  All 4 are counted
because the capacity ledger operates at the structural level, before
SSB selects a vacuum.

\noindent\textit{Gauge generators} (12): $\dim(\SU(3)) + \dim(\SU(2))
+ \dim(\U(1)) = 8 + 3 + 1$.  Each generator is a structurally distinct
enforcement channel.

\subsection{Why gauge constraints are not subtracted (the BRST question)}
\label{sec:BRST}

In perturbative quantum field theory, the physical Hilbert space is the
BRST cohomology: gauge constraints subtract unphysical polarizations,
reducing the propagating degrees of freedom.  A natural question is: why
does the capacity count not perform a similar subtraction?

The answer is that capacity types count \emph{enforcement roles}, not
\emph{propagating degrees of freedom}.  The distinction is:

\smallskip\noindent\textit{Propagating DOF} (QFT accounting): count
on-shell polarization states.  A massless gauge boson has 2 polarizations;
BRST subtracts the longitudinal and timelike modes.  This gives
$2 \times 12 = 24$ propagating gauge DOF, not 12.

\smallskip\noindent\textit{Enforcement roles} (capacity accounting):
count structurally distinct channels in the enforcement ledger.  Each
gauge generator $T^a \in \mathfrak{g}$ is one independent enforcement
channel: it acts as a derivation on the algebra of observables, routing
enforcement between matter fields.  The number of generators is
$\dim(\mathfrak{g}) = 12$, regardless of how many polarization states
each generator supports.

The 12 gauge generators are assigned to the \emph{vacuum} sector
(Q1 = 0 in $P_{\mathrm{exhaust}}$) precisely because they are
enforcement \emph{infrastructure}, not enforcement \emph{consumers}.
A gauge generator cannot route through itself; it has no upstream
pathway to enter as a matter excitation.  BRST is irrelevant to this
accounting: it subtracts polarizations from the propagator, which is a
kinematic object; the capacity count is a structural object.

\coderef{check\_L\_count}{gauge.py}


% ══════════════════════════════════════════════════════════════
\section{The MECE partition}
\label{sec:mece}
% ══════════════════════════════════════════════════════════════

% ── §10: The MECE partition ──

The 61 capacity types are partitioned by two binary predicates into
three mutually exclusive, collectively exhaustive sectors.

\begin{definition}[Predicates Q1 and Q2]\label{def:Q1Q2}
\leavevmode
\begin{enumerate}[label=\textbf{Q\arabic*},nosep,leftmargin=3em]
  \item \textbf{Gauge addressability} (from T3): does the capacity type
  route through the gauge channels $\SU(3) \times \SU(2) \times \U(1)$?
  \item \textbf{Confinement} (within Q1=1): does the gauge-addressable
  type carry conserved labels protected by $\SU(3)$ confinement?
\end{enumerate}
\end{definition}

\begin{theorem}[$P_{\mathrm{exhaust}}$: Predicate Exhaustion]\label{thm:Pexhaust}
\[
  \underbrace{61}_{C_{\mathrm{total}}}
  = \underbrace{42}_{\text{vacuum (Q1=0)}}
  + \underbrace{3}_{\text{baryonic (Q1=1, Q2=1)}}
  + \underbrace{16}_{\text{dark (Q1=1, Q2=0)}}.
\]
No third predicate can refine this partition.
\end{theorem}

\begin{proof}
The vacuum sector (42 types) is defined by the absence of
gauge-addressable labels.  Any splitting predicate would introduce an
addressable distinction, moving types to Q1=1 --- a contradiction.
The baryonic sector (3 types) is indexed by $N_c = 3$ color charges,
the minimal confining carrier.  Sub-ternary splitting would require
$\SU(n < 3)$ confinement, which does not exist.  The dark sector
(16 types) consists of gauge-singlet enforcement channels; further
splitting requires a gauge pathway outside the derived group, which
$L_{\mathrm{gauge,unique}}$ excludes.

The vacuum count 42 is proved via the cross-match $L_{\mathrm{vac}}$
(Paper~2, \S7.1): 12 gauge generators $+$ 27 fermion
internal-structure channels $+$ 3 absorbed Higgs Goldstones $= 42$.
The matter count is $61 - 42 = 19 = 3 + 16$.
\end{proof}

\coderef{check\_P\_exhaust}{core.py}


% ══════════════════════════════════════════════════════════════
\section{Horizon equipartition ($L_{\mathrm{equip}}$)}
\label{sec:equip}
% ══════════════════════════════════════════════════════════════

% ── §11: Horizon equipartition ──

\subsection{Strengthened equal-priors argument}
\label{sec:equal_priors}

The key step in $L_{\mathrm{equip}}$ is the claim that the 61 capacity
types receive equal capacity at the horizon.  Paper~2 states this via
``the principle of indifference over structurally equivalent microstates.''
We sharpen this to a mathematical theorem.

\begin{theorem}[$L_{\mathrm{equip}}$: Horizon Equipartition]\label{thm:Lequip}
At the causal horizon, the maximum-entropy capacity distribution over
$C_{\mathrm{total}} = 61$ types, subject to the constraint
$\sum_i \eps_i = C$ and the floor $\eps_i \ge \eps^* > 0$, is the
uniform distribution $\eps_i = C / C_{\mathrm{total}}$ for all~$i$.
The density fraction of sector $X$ is
$\Omega_X = |X| / C_{\mathrm{total}}$.
\end{theorem}

\begin{proof}
The entropy functional $S(\{\eps_i\}) = -\sum_i \eps_i \ln \eps_i$
is strictly concave on the simplex
$\{\eps_i \ge \eps^* : \sum \eps_i = C\}$.  A strictly concave function
on a compact convex set has a unique maximum.  By Lagrange multipliers:
$\partial S / \partial \eps_i = -(1 + \ln \eps_i) = \lambda$ for all
$i$, giving $\eps_i = e^{-(1+\lambda)}$ for all $i$.  The constraint
$\sum \eps_i = C$ then gives $\eps_i = C / 61$.  The floor
$\eps_i \ge \eps^*$ is satisfied when $C \ge 61 \eps^*$, which holds
at any interface above the critical saturation $s_{\mathrm{crit}}$.

The density fraction of sector $X$ is:
$\Omega_X = \sum_{i \in X} \eps_i / \sum_{\mathrm{all}} \eps_i
= |X| \cdot (C/61) / C = |X|/61$.
The total capacity $C$ cancels.  The fractions depend only on the
integer type counts:
\[
  \Omega_\Lambda = 42/61 = 0.6885, \quad
  \Omega_m = 19/61 = 0.3115.
\]
\end{proof}

\noindent This is not an appeal to the ``principle of indifference''
--- it is a constrained optimization on a finite discrete set.  The
strict concavity of $S$ is a standard result; the uniqueness of the
maximum is guaranteed by the compactness of the feasible set and the
strict concavity of the objective.


\subsection{Bridge to the Friedmann equations}
\label{sec:friedmann_bridge}

The structural result above establishes the partition fractions.  The
cosmological interpretation requires identifying the enforcement
ledger's capacity distribution with the energy budget in the Friedmann
equation:
\begin{equation}
  H^2 = \frac{8\pi G}{3}\,\rho_{\mathrm{total}},
  \qquad \rho_{\mathrm{total}} = \rho_\Lambda + \rho_m.
\end{equation}
The bridge is: at the causal horizon (Bekenstein saturation), the
enforcement ledger's total committed capacity \emph{is} the
stress-energy budget.  Each capacity type contributes equally to the
total energy density (by $L_{\mathrm{equip}}$), so
$\rho_\Lambda / \rho_{\mathrm{total}} = 42/61$ and
$\rho_m / \rho_{\mathrm{total}} = 19/61$.

\textit{Proof boundary.}\quad The identification of the enforcement
horizon with the cosmological de~Sitter horizon is established in
Paper~6 (\emph{Cosmological Evolution}), which derives the full
de~Sitter trajectory and confirms that the relevant horizon is the
one at which the enforcement ledger reaches Bekenstein saturation.
Without Paper~6, the cosmological claim is conditional: \emph{if the
enforcement horizon is the de~Sitter horizon, then
$\Omega_\Lambda = 42/61$}.

\coderef{check\_L\_equip}{cosmology.py}

\subsection{Horizon choice and departures from equilibrium}
\label{sec:horizon_sensitivity}

$L_{\mathrm{equip}}$ holds at a specific horizon: the one at which the
enforcement ledger reaches Bekenstein saturation.  This subsection
examines (a)~which horizon this is, (b)~whether the prediction is
sensitive to the choice, and (c)~how departures from equilibrium affect
the density fractions.

\smallskip\noindent\textbf{(a) Which horizon.}\quad
Three horizons appear in cosmological contexts:
\begin{enumerate}[label=(\roman*),nosep]
  \item The \emph{event horizon}: the boundary of the causal future of
  an observer.  In de~Sitter space, this is the cosmological event
  horizon at proper distance $r_{\mathrm{EH}} = c/H$.
  \item The \emph{apparent horizon}: the outermost marginally trapped
  surface.  In a flat FRW universe, this is the Hubble radius
  $r_{\mathrm{AH}} = c/H(t)$.
  \item The \emph{particle horizon}: the boundary of the causal past.
  Not relevant for the saturation argument (it concerns the past, not
  the enforcement budget for future commitments).
\end{enumerate}
The APF's enforcement horizon is the surface at which the
Bekenstein bound $S \le A / 4G\hbar$ is saturated: the entropy of
the enclosed region equals the holographic capacity of the boundary.
In de~Sitter space, the event horizon and the apparent horizon
coincide (both equal $c/H_\Lambda$ in the late-time limit).  At
earlier epochs, the apparent horizon is time-dependent and the event
horizon may not yet be well-defined.  The prediction $\Omega_\Lambda =
42/61$ applies at \emph{saturation}, which the APF identifies with
the late-time de~Sitter attractor.  At this fixed point, the
distinction between event and apparent horizons vanishes.

\smallskip\noindent\textbf{(b) Sensitivity to horizon choice.}\quad
Suppose one evaluates the capacity partition at a different surface.
The fractions $\Omega_X = |X| / C_{\mathrm{total}}$ depend on the
integer type counts, not on the area of the horizon.  The total
capacity $C$ cancels in the ratio (as shown in the proof of
Theorem~\ref{thm:Lequip}).  Therefore:
\begin{itemize}[nosep]
  \item If the type count $C_{\mathrm{total}} = 61$ is the same at
  the new surface, the fractions are unchanged.
  \item If the type count changes (e.g., because some channels are
  frozen out below a mass threshold), the fractions shift.  This is
  the sensitivity captured by the table in \S\ref{sec:sensitivity}.
\end{itemize}
The prediction is therefore \emph{robust to horizon choice} provided
the type count is evaluated at the scale where all 61 channels are
active.  The Bekenstein saturation scale is precisely this scale: it
is the coarsest (lowest-energy) scale at which the full enforcement
ledger is committed.

\smallskip\noindent\textbf{(c) Departures from equilibrium.}\quad
$L_{\mathrm{equip}}$ is a maximum-entropy result.  If the system is
out of equilibrium, the capacity distribution $\{\eps_i\}$ need not
be uniform.  The density fractions would then depend on the specific
non-equilibrium distribution, and the prediction $42/61$ would not
hold exactly.  Three considerations bound the deviation:
\begin{enumerate}[label=(\alph*),nosep]
  \item \textit{Approach to equilibrium.}\quad
  The second law of thermodynamics (here: monotonic increase of
  $S(\{\eps_i\})$ under capacity redistribution, which follows from
  $L_{\mathrm{irr}}$) guarantees that the distribution approaches
  uniformity over time.  At the de~Sitter fixed point, the
  cosmological evolution has had infinite proper time to equilibrate;
  the system is at or arbitrarily close to maximum entropy.

  \item \textit{Perturbative corrections.}\quad
  Near the uniform distribution, write
  $\eps_i = (C/61)(1 + \delta_i)$ with $\sum \delta_i = 0$.  The
  density fraction of sector $X$ becomes
  $\Omega_X = (|X|/61)(1 + \langle\delta\rangle_X)$, where
  $\langle\delta\rangle_X = (1/|X|)\sum_{i \in X} \delta_i$.
  The correction $\langle\delta\rangle_X$ vanishes when $\delta_i$ is
  uncorrelated with sector membership.  For the correction to shift
  $\Omega_\Lambda$ outside the Planck $1\sigma$ band ($\pm 0.0056$),
  one would need $|\langle\delta\rangle_\Lambda| \gtrsim 0.008$,
  i.e., a systematic $0.8\%$ bias toward or against vacuum-sector
  channels --- a large departure from equilibrium at the de~Sitter
  fixed point.

  \item \textit{Saturation independence.}\quad
  The integer type counts (42 vacuum, 3 baryonic, 16 dark) are
  topological invariants of the MECE partition
  (\S\ref{sec:mece}); they do not depend on the total capacity
  $C$ or on the distribution $\{\eps_i\}$.  Therefore even if
  the system is out of equilibrium, the \emph{ratios between
  integer counts} are exact: $\Omega_\Lambda / \Omega_m = 42/19$
  is an identity of the partition, not a thermodynamic result.
  Only the individual fractions $\Omega_X = |X| / C_{\mathrm{total}}$
  require the equilibrium assumption.
\end{enumerate}

\smallskip\noindent\textit{Comparison with Butcher's quantum-bias
cosmology.}\quad
Butcher (2024) proposes a distinct route from information-capacity
bounds to cosmic acceleration.  In that framework, a quantum-mechanical
correction to the Friedmann equation arises from the holographic
(Bekenstein--Bousso) entropy bound: the finite information capacity of
a causal patch introduces an effective negative-pressure term that
modifies late-time dynamics nonlocally.  The resulting modified Friedmann
equation predicts accelerated expansion without a cosmological constant
term, instead attributing the acceleration to an information-capacity
back-reaction.

The two frameworks share a common premise --- finite information capacity
at causal horizons drives cosmic acceleration --- but differ in mechanism,
scope, and predicted late-time behavior:
\begin{enumerate}[label=(\alph*),nosep]
  \item \textbf{Mechanism.}\quad
  In the APF, acceleration arises from the \emph{structural partition}
  of 61 discrete capacity types into vacuum and matter sectors
  ($L_{\mathrm{equip}}$).  The cosmological constant
  $\Lambda \propto 42/61$ is an exact rational number determined by
  type counting.  In Butcher's framework, acceleration arises from a
  \emph{dynamical back-reaction} of the holographic bound on the
  Friedmann equation, producing an effective dark energy that is not
  a constant but a functional of the horizon entropy.
  \item \textbf{Predicted equation of state.}\quad
  The APF predicts $w = -1$ exactly (the vacuum sector's capacity
  fraction is a topological invariant, independent of the expansion
  history; $L_{\mathrm{sat,part}}$).  Butcher's framework generically
  predicts $w \ne -1$ at intermediate times, with $w \to -1$ only
  asymptotically.  Current data (Planck + DESI BAO) are consistent
  with $w = -1$ at $< 2\sigma$; future surveys (Euclid, LSST) will
  tighten the constraint to sub-percent precision on $|1 + w|$ and
  could distinguish the two predictions.
  \item \textbf{Scope.}\quad
  The APF's capacity-partition argument derives the density fractions
  from the gauge group and matter content (themselves derived from A1);
  the cosmological result is a downstream consequence of particle
  physics, not an independent postulate.  Butcher's framework is
  self-contained within cosmology and does not derive the particle
  content.  The two are therefore complementary rather than
  contradictory: they address different questions (``why these
  fractions?'' vs.\ ``how does acceleration arise dynamically?'')
  and could in principle both be correct if the APF's structural
  fractions are realized via an information-capacity mechanism at the
  dynamical level.
  \item \textbf{Nonlocality.}\quad
  Butcher's modification is explicitly nonlocal (it depends on the
  global horizon entropy, not on local field equations).  The APF's
  $L_{\mathrm{equip}}$ is also horizon-global (it applies at Bekenstein
  saturation, a property of the causal patch as a whole), but the
  underlying field equations ($T_{9,\mathrm{grav}}$) are local (Einstein
  equations from Lovelock).  The APF thus predicts standard local GR
  dynamics with an effective $\Lambda = 42 C / (61 \cdot 3)$
  at the horizon scale; Butcher's framework modifies the dynamics
  themselves.
\end{enumerate}


% ══════════════════════════════════════════════════════════════
\section{Saturation independence ($L_{\mathrm{sat,part}}$)}
\label{sec:sat_indep}
% ══════════════════════════════════════════════════════════════

% ── §12: Saturation independence ──

\begin{theorem}[$L_{\mathrm{sat,part}}$: Saturation-Independent Partition]
\label{thm:sat_part}
The partition fractions $\Omega_\Lambda = 42/61$ and
$\Omega_m = 19/61$ are determined by discrete type-classification
predicates (Q1, Q2) applied to the anomaly-free field content.  They
are independent of the total capacity $C$ and the saturation level $s$.
\end{theorem}

\begin{proof}
The predicates Q1 and Q2 are functions of type labels (representation
quantum numbers), not of the total capacity $C(\Gamma)$.  By
$L_{\mathrm{equip}}$ (Theorem~\ref{thm:Lequip}), the density fraction
of sector $X$ is $\Omega_X = |X|/C_{\mathrm{total}}$, where $|X|$ is
an integer determined by the predicate assignment.  The total capacity
$C$ cancels in the ratio.  Therefore $\Omega_X$ is independent of $C$
and of the saturation level $s := C_{\mathrm{committed}} / C$.

The partition is structurally stable against perturbations: the
predicates Q1 and Q2 depend on the Wedderburn block structure and
field selection, both of which are discrete invariants preserved for
cross-support perturbations $\delta < \delta_{\mathrm{crit}}$
(Paper~1 Supplement, Proposition~robustness).
\end{proof}

\noindent Below the critical saturation $s_{\mathrm{crit}} = 1/d_{\mathrm{eff}}
= 1/102$, the full matching does not exist (anomaly cancellation
requires all 61 types simultaneously).  The pre-matching state is
pure de~Sitter vacuum.  The partition is undefined below
$s_{\mathrm{crit}}$ but irrelevant, because the vacuum has $w = -1$
regardless.

\coderef{check\_L\_saturation\_partition}{cosmology.py}


% ══════════════════════════════════════════════════════════════
\section{Three structural predictions}
\label{sec:structural_predictions}
% ══════════════════════════════════════════════════════════════

% ── §14: Three structural predictions ──

The capacity counting, MECE partition, and cost-function uniqueness
established above have three observational consequences that bear on
open problems in the Standard Model.  Each is a descriptive reading
of a structure already proved in this supplement: the admissibility
conditions, together with the derived capacity partition, are
compatible with exactly one value of the quantity in question; no new
machinery is introduced, and no mechanism ranges over candidates to
reject the non-admissible ones.  The three consequences read off are:
(i) the QCD vacuum angle is admissible only at $\theta_{\mathrm{QCD}} =
0$ (conditional on Premise~(I) below; \S\ref{sec:theta_QCD});
(ii) the Higgs effective potential has a unique admissible well, so
the electroweak vacuum is absolutely stable (\S\ref{sec:vacuum_stability});
(iii) the dark-energy fraction $\Omega_\Lambda = 42/61$ is the ratio of
topological invariants of the MECE partition (\S\ref{sec:mece},
\S\ref{sec:equipartition}).


% ──────────────────────────────────────────────────────────────
\subsection{Strong CP solution: $\theta_{\mathrm{QCD}} = 0$}
\label{sec:theta_QCD}
% ──────────────────────────────────────────────────────────────

The QCD Lagrangian admits a topological term
$\mathcal{L}_\theta = (\theta/32\pi^2)\,G_{\mu\nu}\tilde{G}^{\mu\nu}$
that is a total divergence.  The parameter $\theta \in [0, 2\pi)$
introduces no new propagating degrees of freedom: $C_{\mathrm{total}}$
is unchanged at 61 for any value of $\theta$.

\begin{theorem}[$T_{\theta_{\mathrm{QCD}}}$: Strong CP Solution \textnormal{[P-heuristic]}]
\label{thm:theta_QCD}
Under the admissibility framework together with Premise~(I) below, the
unique cost-minimising value of the QCD vacuum angle is
$\theta_{\mathrm{QCD}} = 0$.
\end{theorem}

\noindent\textbf{Premise~(I) (instanton-sector closure).}\quad
Non-perturbative tunnelling sectors indexed by topological winding
number contribute no admissible capacity beyond the 61 channels
counted by $L_{\mathrm{count}}$; equivalently, instanton configurations
are treated as enforcement overhead attached to the existing channels
rather than as independent capacity types.  A first-principles
derivation of this premise from A1 --- requiring a theory of
admissible capacity cost for topologically nontrivial gauge sectors
--- is deferred.  The present argument is therefore classified
[P-heuristic]: it is rigorous given Premise~(I), but the premise
itself is not yet reduced to the axiom.

\begin{proof}[Proof of Theorem~\ref{thm:theta_QCD}, conditional on Premise~(I)]
Four steps; the argument is descriptive rather than selective (no
mechanism ranges over $\theta$-values; the structure $\min_{\theta}\mathrm{cost}$
labels the admissible configuration).

\textit{Step~1 ($\theta$ adds no capacity, given Premise~(I)).}
$\mathcal{L}_\theta$ is a total divergence and introduces no new
propagating degrees of freedom.  Together with Premise~(I), the
capacity count is unaffected: $C(\theta \ne 0) = C(\theta = 0) = 61$.

\textit{Step~2 (Enforcement cost from $L_{\varepsilon^*}$).}
A definite value $\theta = \theta_0 \ne 0$ distinguishes the vacuum
from all $\theta \ne \theta_0$ configurations.  By
$L_{\varepsilon^*}$~(\S\ref{app:imported}), every enforceable
distinction costs at least $\varepsilon^* > 0$.  Therefore the
enforcement-cost functional satisfies
$\mathrm{cost}(\theta = \theta_0 \ne 0) \ge \varepsilon^*$.

\textit{Step~3 ($\theta = 0$ is the unique cost-free value).}
At $\theta = 0$ the strong sector is CP-invariant; the value is not
an independent parameter but a consequence of the symmetry.
Maintaining a symmetry consequence requires no enforcement record,
so $\mathrm{cost}(\theta = 0) = 0$.  This is the unique zero of the
cost functional on $[0, 2\pi)$.

\textit{Step~4 (admissibility labels $\theta = 0$).}
Both configurations have $C = 61$.  The net enforcement surplus
reads
\[
  \mathrm{surplus}(\theta = 0) = 0 - 0 = 0
  \;>\;
  0 - \varepsilon^* = \mathrm{surplus}(\theta \ne 0).
\]
The admissible configuration is the one at which the surplus attains
its maximum, equivalently the $\arg\min_\theta \mathrm{cost}(\theta)$.
That argmin is $\theta = 0$.  No search over $\theta$-values is
performed; the framework is descriptive, and Step~4 merely reads off
the referent of the admissibility condition.
\end{proof}

\begin{remark}[Why the CKM phase survives]
The CKM phase $\delta_{\mathrm{CP}}$ is \emph{not} eliminated by this
argument.  Unlike $\theta$, the CKM phase enables $N_{\mathrm{gen}}!
= 6$ distinguishable history sectors (the Jarlskog invariant), which
contribute to $C_{\mathrm{total}}$.  The enforcement cost of maintaining
$\delta_{\mathrm{CP}}$ is offset by the capacity gain from these
additional sectors.  The net surplus is non-negative, so $\delta_{\mathrm{CP}}$
is admissible.  Full accounting is deferred to Paper~4B.
\end{remark}

\coderef{check\_T\_theta\_QCD, check\_L\_strong\_CP\_synthesis}{gauge.py}


% ──────────────────────────────────────────────────────────────
\subsection{Vacuum absolute stability}
\label{sec:vacuum_stability}
% ──────────────────────────────────────────────────────────────

In the Standard Model, the Higgs effective potential extrapolated to
high energies via renormalisation group running develops a second minimum
deeper than the electroweak vacuum for $m_H \approx 125$~GeV and
$m_{\mathrm{top}} \approx 173$~GeV, rendering the vacuum metastable.
The admissibility framework excludes this possibility.

\begin{theorem}[$T_{\mathrm{vac,stab}}$: Vacuum Absolute Stability]
\label{thm:vacuum_stability}
The electroweak vacuum is absolutely stable.  The enforcement potential
has a unique global minimum, and no deeper vacuum exists.
\end{theorem}

\begin{proof}
Four steps.

\textit{Step~1 (Unique enforcement well).}
The enforcement potential $V(\Phi)$ derived from $L_{\varepsilon^*}$
(linear cost floor), $T_M$ (monogamy binding), and A1 (capacity
saturation) has the form:
\[
  V(\Phi) = \varepsilon^*\Phi
    - \tfrac{\eta}{2\varepsilon^*}\,\Phi^2
    + \frac{\varepsilon^*\,\Phi^2}{2(C - \Phi)}.
\]
This potential has: $V(0) = 0$ (empty vacuum, unstable); a barrier
at $\Phi/C \sim 0.06$; a \emph{unique} binding well at
$\Phi/C \sim 0.73$ with $V < 0$; and a divergence $V \to +\infty$
as $\Phi \to C$ (capacity saturation).  The uniqueness of the minimum
is verified by direct computation ($T_{\mathrm{particle}}$, gauge.py).

\textit{Step~2 (No runaway).}
A1 guarantees $\Phi < C$ for all admissible states.  The potential
diverges at $\Phi = C$, so no tunnelling to $\Phi > C$ is possible.
The well at $\Phi/C \sim 0.73$ is the global minimum.

\textit{Step~3 (High-energy behaviour).}
The SM extrapolation that produces metastability assumes volume-scaling
degrees of freedom at high energies.  $T_{\mathrm{Bek}}$ (Bekenstein bound,
Paper~1) imposes area-law regulation, which prevents the effective
potential from developing a second minimum at high field values
($L_{\mathrm{naturalness}}$, gauge.py).

\textit{Step~4 (Uniqueness from equilibrium measure).}
$M_\Omega$ (derived in Paper~6) proves the equilibrium measure at
saturation is unique (uniform).  A second vacuum would require a
second equilibrium, which $M_\Omega$ excludes.  The electroweak vacuum
is the unique stationary point.
\end{proof}

\coderef{check\_T\_vacuum\_stability}{gauge.py}


% ──────────────────────────────────────────────────────────────
\subsection{Proton absolute stability}
\label{sec:proton_stability}
% ──────────────────────────────────────────────────────────────

In the Standard Model, baryon number $B$ is an accidental symmetry
that could in principle be violated by non-perturbative processes
(sphalerons) or by higher-dimensional operators in grand unified
extensions.  The admissibility framework derives exact $B$~conservation
at Bekenstein saturation.

\begin{theorem}[$T_{\mathrm{proton}}$: Proton Absolute Stability]
\label{thm:proton}
At full Bekenstein saturation, baryon number is an exact conservation
law.  The proton is absolutely stable.
\end{theorem}

\begin{proof}
Three steps.

\textit{Step~1 (Partition is exact at saturation).}
$P_{\mathrm{exhaust}}$ (Theorem~\ref{thm:Pexhaust}) proves that the
three-sector partition
\[
  N_b + N_d + N_v = 3 + 16 + 42 = 61 = C_{\mathrm{total}}
\]
is mutually exclusive and collectively exhaustive at full Bekenstein
saturation.  The predicates Q1 (gauge addressability) and Q2
(confinement) uniquely assign each capacity type to exactly one sector.
The baryonic sector has $N_b = 3$ types, and this integer is a
conserved quantity at saturation: the partition is sharp (no type is
partially baryonic or ambiguously assigned).

\textit{Step~2 (Saturation is irreversible).}
$L_{\mathrm{irr}}$ (\S\ref{app:imported}) proves that the
transition to full saturation is irreversible.  Once all 61 capacity
types are committed and the ledger reaches saturation, the universe
cannot return to the pre-saturation regime where the partition was
not enforced.  Records are locked.

\textit{Step~3 (No admissible rerouting).}
Any baryon-number-violating process would require a capacity type to
exit the baryonic sector (violating Q1 or Q2) and enter another sector.
Both operations violate the partition predicates, which are mathematical
identities at saturation (Step~1).  No admissible move changes $N_b$.
\end{proof}

\begin{remark}[No $B$-violating operators at any dimension]
$T_{\mathrm{gauge}}$ derives $\SU(3) \times \SU(2) \times \U(1)$ as
the \emph{complete} gauge group, not embedded in a larger group.  No
leptoquark or $X/Y$ boson exists; no dimension-6 $B$-violating
operator is gauge-invariant under the derived group.  Even if such
operators were postulated, Step~3 above forbids their action at
saturation.  The prediction is therefore the strongest possible:
$\tau_p = \infty$, not merely $\tau_p > 10^{34}$~years.
\end{remark}

\coderef{check\_T\_proton, check\_L\_B\_operator\_exclusion}{gauge.py}


% ══════════════════════════════════════════════════════════════
\section{Uncertainty analysis and sensitivity}
\label{sec:uncertainty}
% ══════════════════════════════════════════════════════════════

% ── §13: Uncertainty analysis and sensitivity ──

The predictions $\Omega_\Lambda = 42/61$ and $\Omega_m = 19/61$ are
rational numbers with zero free parameters.  There is no propagated
parametric uncertainty in the usual sense.  The relevant question is:
\emph{what assumptions, if changed, would shift the prediction?}

\subsection{Sensitivity to the capacity count}
\label{sec:sensitivity}

The only quantity entering the prediction is $C_{\mathrm{total}}$, the
total number of capacity types.  The vacuum count (42) is determined by
$L_{\mathrm{vac}}$ and is insensitive to the matter content.  The matter
count (19) is $C_{\mathrm{total}} - 42$.  Therefore the prediction shifts
if and only if $C_{\mathrm{total}}$ changes.

\smallskip
\begin{center}
\small
\begin{tabular}{@{}lcccl@{}}
\toprule
\textbf{Scenario} & $C_{\mathrm{total}}$ & $\Omega_\Lambda$ &
  \textbf{Tension with Planck} & \textbf{What changes} \\
\midrule
SM (derived) & 61 & 42/61 = 0.6885 & $0.07\%$ ($<0.1\sigma$) & --- \\
$+3\,\nu_R$ & 64 & 42/64 = 0.656 & $5.8\sigma$ & Add light $\nu_R$ \\
$+$ polarizations & 73 & 42/73 = 0.575 & $>20\sigma$ & Count helicities \\
$-$ Higgs & 57 & 39/57 = 0.684 & $0.9\sigma$ & Remove Higgs DOF \\
$-1$ generation & 46 & 27/46 = 0.587 & $>18\sigma$ & 2 generations \\
\bottomrule
\end{tabular}
\end{center}

\smallskip\noindent Every alternative $C_{\mathrm{total}}$ moves the
prediction outside the $1\sigma$ band except the Higgs-removal scenario,
which is excluded by the derived scalar sector ($T_{\mathrm{Higgs}}$,
Paper~2, \S5.5).

\subsection{Comparison with observation}

Planck 2018 measures
$\Omega_\Lambda = 0.6889 \pm 0.0056$~\cite{Planck2020}.
The prediction $42/61 = 0.68852\ldots$ is within $0.07\%$ of the
central value, well inside the $1\sigma$ band.  The agreement extends
to all five derived fractions:

\smallskip
\begin{center}
\small
\begin{tabular}{@{}lccr@{}}
\toprule
\textbf{Observable} & \textbf{Predicted} & \textbf{Observed} & \textbf{Deviation} \\
\midrule
$\Omega_\Lambda$ & $42/61 = 0.6885$ & $0.6889 \pm 0.0056$ & $0.06\%$ \\
$\Omega_m$ & $19/61 = 0.3115$ & $0.3111 \pm 0.0056$ & $0.13\%$ \\
$\Omega_{\mathrm{DM}}$ & $16/61 = 0.2623$ & $0.2607 \pm 0.0056$ & $0.6\%$ \\
$\Omega_b$ & $3/61 = 0.0492$ & $0.0490 \pm 0.0003$ & $0.4\%$ \\
$f_b = \Omega_b/\Omega_m$ & $3/19 = 0.1579$ & $0.1571 \pm 0.0012$ & $0.5\%$ \\
\bottomrule
\end{tabular}
\end{center}

\smallskip\noindent Five predictions, zero free parameters, all within
$0.6\%$ of observation.  The numerators are integers; the denominators
are 61.  There is no fitting: the fractions are determined by the
derived field content and the MECE partition.


\newpage

% ╔═══════════════════════════════════════════════════════════════╗
% ║                         APPENDICES                            ║
% ╚═══════════════════════════════════════════════════════════════╝

\appendix


% ══════════════════════════════════════════════════════════════
\section{Comparison with reconstruction and ``SM from principles'' programs}
\label{app:literature}
% ══════════════════════════════════════════════════════════════

% ── Appendix A: Comparison with reconstruction programs ──

\subsection{Operational and categorical reconstructions}
\label{sec:op_reconstructions}

Several programs reconstruct quantum theory from operational or
information-theoretic axioms.  We summarize each and note where the
APF diverges.

\smallskip\noindent\textbf{Hardy (2001)~\cite{Hardy2001}.}\quad
Five ``reasonable axioms'' (probabilities, simplicity, subspaces,
composite systems, continuous reversibility) derive finite-dimensional
quantum theory over $\mathbb{C}$.  The axioms are stated at the
kinematics level: they constrain the state space and measurement
structure but do not determine gauge content, matter spectrum, or
dynamics.  APF shares the ``finite capacity'' starting point
(Hardy's axiom 1 is a form of finite dimensionality) but pushes
through to gauge group and matter content via non-closure, which
Hardy does not address.

\smallskip\noindent\textbf{Chiribella--D'Ariano--Perinotti
(2011)~\cite{CDP2011}.}\quad
Six axioms (causality, perfect distinguishability, ideal compression,
local distinguishability, pure conditioning, purification) derive QT.
The purification axiom is the key differentiator from classical
probability.  The output is a general probabilistic theory isomorphic
to finite-dimensional quantum mechanics.  Like Hardy, CDP stops at
kinematics.  The APF's constitutive principle SP (substrate faithfulness)
plays a role analogous to CDP's local distinguishability, and K3
(forced additivity) parallels their ideal compression, but the APF
does not assume purification --- instead, it derives the Born rule
from the enforcement algebra via Busch's generalized Gleason theorem.

\smallskip\noindent\textbf{H\"ohn (2017)~\cite{Hohn2017}.}\quad
Reconstructs qubit quantum theory from ``questions'' (binary tests)
with finite information capacity.  The closest cousin to APF's
starting point: finite capacity $\to$ quantum structure.  H\"ohn's
framework is carefully delimited to qubits and does not attempt gauge
content.  APF can be viewed as an extension of H\"ohn's program that
asks: what happens when you push the ``finite information'' idea beyond
qubit reconstruction to the full field content?

\smallskip\noindent\textbf{Effectus theory (Westerbaan, van de Wetering).}\quad
Categorical framework using ``effectuses'' to axiomatize quantum-like
theories.  Sequential product axioms force operator-algebraic structure.
The APF's enforcement algebra (\S\ref{sec:nc_chain}) can be compared
to the sequential product in effectus theory: both derive
non-commutativity from the structure of sequential measurements.  The
effectus approach is more general (it classifies all effectuses, not
just the quantum one); the APF is more specific (it selects the
quantum case via A1 and derives gauge content).

\smallskip\noindent\textbf{Tull (2019); Spekkens (2007).}\quad
Tull derives operator algebras from categorical CPM (completely positive
maps) structure.  Spekkens constructs a toy model from epistemic
restrictions.  Both illuminate the boundary between classical and
quantum but do not derive field content or gauge structure.


\subsection{SM-from-principles programs}
\label{sec:sm_programs}

A separate tradition attempts to derive the Standard Model within an
assumed quantum field theory framework.

\smallskip\noindent\textbf{Anomaly-based selection
(Minahan--Ramond--Warner, 1990).}\quad
Given $\SU(3) \times \SU(2) \times \U(1)$ and the requirement of
anomaly cancellation, the SM fermion content is strongly constrained
but not unique without additional assumptions (e.g., minimality,
family structure).  The APF's $T_{\mathrm{field}}$ performs a similar
scan but derives the gauge group itself (via Theorem~R +
$L_{\mathrm{gauge,unique}}$) rather than assuming it.

\smallskip\noindent\textbf{Grand unification (Georgi--Glashow, 1974).}\quad
Embeds $\SU(3) \times \SU(2) \times \U(1)$ into $\SU(5)$.  Explains
charge quantization and predicts proton decay.  The APF derives the
product structure as capacity-minimal
(Proposition~\ref{prop:no_simple}: $\SU(5)$ costs $2\times$) and
obtains charge quantization from anomaly cancellation
(Theorem~\ref{thm:hypercharge}) without grand unification.

\smallskip\noindent\textbf{Connes NCG (1996).}\quad
The noncommutative geometry program derives the SM Lagrangian from a
spectral triple.  The input is a specific finite-dimensional algebra
(chosen to reproduce the SM).  The APF derives the algebra from A1
($T_{\mathrm{alg}}$), rather than postulating it; the two programs
converge on the same algebraic structure but arrive from opposite
directions.

\smallskip\noindent\textbf{Doplicher--Roberts reconstruction and
Pinzari's braided extensions.}\quad
The classical DR theorem (1989) reconstructs a compact group $G$ from
its representation category (a symmetric $C^*$-tensor category with
conjugates).  In the APF, T3 (Paper~1 Supplement) uses DR to recover
the gauge group from the enforcement algebra's representation
category; the gauge identification theorem
(Theorem~\ref{thm:gauge_ident}) then identifies which automorphisms
constitute the gauge group.

Pinzari's program (2017--) extends DR to \emph{braided} tensor
categories and weak Hopf $C^*$-algebras.  In this setting, abelian
factors can emerge from the \emph{center} of the fusion category:
central objects (transparent to braiding) generate grading tori that
act as abelian gauge factors.  This is structurally parallel to the
APF's mechanism for $\U(1)_Y$: the relative-phase torus
$T_{\mathrm{rel}} = Z(\mathrm{U}(\mathcal{A})) /
\mathrm{U}(1)_{\mathrm{diag}}$ (Remark~\ref{rem:U1}) plays the role
of Pinzari's grading torus, and the abelian factor arises from
central (= transparent) elements rather than from inner automorphisms.

The key structural parallel is: \emph{non-abelian gauge from inner
automorphisms, abelian gauge from the center}.  In DR/Pinzari, this
is ``gauge from the Tannakian group, grading from the M\"uger center.''
In APF, it is ``gauge from $\mathrm{Inn}(\mathcal{A})$, grading from
$T_{\mathrm{rel}}$ pruned by CM.''  The APF adds a \emph{selection
mechanism} (capacity minimality, Proposition~\ref{prop:unique_U1})
that is absent in the DR/Pinzari framework: the latter reconstructs
whatever symmetry the category has, while the APF derives the
\emph{minimal} symmetry compatible with enforcement completeness.

\smallskip\noindent\textbf{Finite tensor categories and symmetric
tensor categories.}\quad
Recent work on finite tensor categories (Etingof--Gelaki--Nikshych--Ostrik)
and symmetric tensor categories (Deligne's theorem and extensions by
Coulembier, Entova-Aizenbud, Ostrik) shows that structural constraints
on representation categories can have sharp classification
consequences.  Deligne's theorem, for example, classifies all
symmetric tensor categories of moderate growth as representation
categories of supergroups --- a result that constrains the possible
internal symmetry structures of any theory whose observables form a
symmetric tensor category.

The APF is not a tensor-categorical framework: its primary objects are
enforcement algebras and cost functionals, not abstract tensor
categories.  However, the representation category of the enforcement
algebra $\mathcal{A} \cong \bigoplus_k M_{n_k}(\mathbb{C})$
\emph{is} a finite semisimple tensor category (the category of
finite-dimensional $\mathcal{A}$-modules), and its symmetry is
classified by the theorems above.  The APF's contribution relative to
this literature is not a new classification of tensor categories but a
\emph{physical derivation} of which tensor category is realized: the
enforcement postulate A1 selects $\mathcal{A}$ and hence its
representation category, rather than postulating the category
directly.  The classification theorems then apply as \emph{imported
mathematics} (same epistemic status as Wedderburn--Artin or
Skolem--Noether in the present supplement).


\subsection{Locally covariant QFT and the General Boundary Formulation}
\label{sec:lcqft_oeckl}

Two programs are particularly relevant to the APF because they
share structural elements (local algebras, causality, boundaries)
while differing fundamentally on infinite-dimensionality.

\smallskip\noindent\textbf{Locally covariant QFT
(Brunetti--Fredenhagen--Verch, 2003).}\quad
LCQFT axiomatizes QFT as a functor from the category of globally
hyperbolic Lorentzian manifolds (with isometric embeddings as morphisms)
to the category of $*$-algebras (with injective $*$-homomorphisms).
The key axioms are:
\begin{itemize}[nosep]
  \item \textit{Locality}: spacelike-separated regions have commuting
  algebras.
  \item \textit{Covariance}: the assignment $\mathcal{O} \mapsto
  \mathcal{A}(\mathcal{O})$ is functorial under isometric embeddings.
  \item \textit{Time-slice}: if $\mathcal{O}_1 \hookrightarrow
  \mathcal{O}_2$ contains a Cauchy surface of $\mathcal{O}_2$, then
  $\mathcal{A}(\mathcal{O}_1) \cong \mathcal{A}(\mathcal{O}_2)$.
\end{itemize}
The LCQFT framework works with \emph{infinite-dimensional} algebras
(typically type~III factors for local regions) and
\emph{presupposes} a background Lorentzian manifold.

The APF differs from LCQFT on three structural points:
\begin{enumerate}[label=(\alph*),nosep]
  \item \textit{Finite vs.\ infinite dimensions.}\quad LCQFT's local
  algebras are infinite-dimensional by construction (they are obtained
  via the GNS representation of a state on a $C^*$-algebra of
  smeared fields).  The APF's local algebras are finite-dimensional
  (A1: $\dim(\mathcal{H}_\omega) < \infty$).  The reconciliation is
  via the controlled limit (Proposition~\ref{prop:typeIII_limit}):
  the APF's type~I algebras approximate LCQFT's type~III algebras
  in the $C/\varepsilon^* \to \infty$ limit, with the structural
  predictions (gauge group, matter content) stable under the limit
  (Theorem~\ref{thm:structural_stability}).
  \item \textit{Spacetime derived vs.\ presupposed.}\quad LCQFT takes
  a Lorentzian manifold as input; the APF derives it
  (\S\ref{sec:lorentz}: $L_{\mathrm{irr}} \to \prec \to M \to
  g_{\mu\nu}$).  The APF can therefore be viewed as providing a
  \emph{foundation} for LCQFT: it derives the background manifold
  that LCQFT assumes, and the local algebras that LCQFT assigns to
  regions emerge as the enforcement algebras at interfaces.
  \item \textit{Causality.}\quad Both frameworks implement causality
  via commutativity of spacelike-separated algebras.  In LCQFT, this
  is an axiom; in the APF, it is derived
  (Lemma~\ref{lem:loc_commut}).  The LCQFT time-slice axiom has an
  APF analogue: if $\Gamma_1 \subseteq \Gamma_2$ and $\Gamma_1$
  ``contains a Cauchy surface'' (i.e., the enforcement structure at
  $\Gamma_1$ determines the enforcement structure at $\Gamma_2$),
  then $\mathcal{A}(\Gamma_1) \cong \mathcal{A}(\Gamma_2)$.  This
  is a consequence of $L_{\mathrm{loc}}$ and the deterministic
  evolution of the enforcement structure.
\end{enumerate}

\smallskip\noindent\textbf{Oeckl's Positive Formalism / General
Boundary Formulation (2003--2016).}\quad
Oeckl's program assigns quantum amplitudes to \emph{boundaries} of
spacetime regions, rather than to states on Cauchy surfaces.  The
formalism uses:
\begin{itemize}[nosep]
  \item \textit{Boundary Hilbert spaces}: each boundary $\Sigma$
  (not necessarily a Cauchy surface) carries a Hilbert space
  $\mathcal{H}_\Sigma$.
  \item \textit{Amplitude maps}: each spacetime region $M$ bounded by
  $\Sigma$ defines a linear map $\rho_M: \mathcal{H}_\Sigma \to
  \mathbb{C}$ (or more generally a positive map).
  \item \textit{Composition}: gluing regions along shared boundaries
  composes amplitude maps.
\end{itemize}
The General Boundary Formulation (GBF) generalizes this to
non-globally-hyperbolic spacetimes and is compatible with quantum
gravity scenarios where the background spacetime is not fixed.

The APF's interface structure has a striking structural parallel to
Oeckl's boundary formalism:
\begin{enumerate}[label=(\alph*),nosep]
  \item \textit{Interfaces as boundaries.}\quad An APF interface
  $\Gamma$ is the boundary between two enforcement regions.  The
  enforcement algebra $\mathcal{A}(\Gamma)$ plays the role of
  Oeckl's boundary Hilbert space $\mathcal{H}_\Sigma$ (in the
  finite-dimensional case, the algebra and its GNS Hilbert space
  are equivalent objects).
  \item \textit{Cost kernel as amplitude.}\quad The enforcement cost
  $\varepsilon_{\mathrm{ext}}(v) = B(v,v)$ ($T_{7B}$) is a positive
  functional on the interface sector space, analogous to Oeckl's
  positive amplitude $\rho_M$.  The positivity of $B$ corresponds to
  Oeckl's positivity condition, which ensures physical probabilities
  are non-negative.
  \item \textit{Composition via pool coupling.}\quad When two
  interfaces share a pool $\Pi$ ($L_\Delta$), the enforcement
  structure composes: the combined cost is superadditive
  ($\varepsilon(v_1 + v_2) \ge \varepsilon(v_1) + \varepsilon(v_2)$).
  This parallels Oeckl's composition of amplitude maps via gluing.
\end{enumerate}
The key differences:
\begin{itemize}[nosep]
  \item Oeckl works with infinite-dimensional Hilbert spaces and
  does not impose a capacity bound.  The APF's A1 restricts to finite
  dimensions.
  \item Oeckl's formalism is \emph{kinematic} (it assigns amplitudes
  but does not derive gauge content or matter spectrum).  The APF
  pushes through to field content via the enforcement algebra's
  Wedderburn decomposition.
  \item The GBF is compatible with indefinite causal structure
  (useful for quantum gravity); the APF derives a definite causal
  structure from $L_{\mathrm{irr}}$
  (Proposition~\ref{prop:causal}).
\end{itemize}

\smallskip\noindent\textbf{Masanes--M\"uller (2011).}\quad
Derives finite-dimensional quantum theory from four postulates:
continuous reversibility, tomographic locality (joint states
determined by local measurements), existence of an information unit,
and the subspace axiom.  The output is quantum theory over
$\mathbb{R}$, $\mathbb{C}$, or $\mathbb{H}$ (with $\mathbb{C}$
selected by tomographic locality + bit symmetry).  The APF shares
the finite-dimensionality starting point and the field selection
($\mathbb{F} = \mathbb{C}$ from T2b), but arrives at $\mathbb{C}$
via a different mechanism (Skolem--Noether applied to the enforcement
algebra, not tomographic locality).  Like Hardy and CDP,
Masanes--M\"uller stops at kinematics and does not address gauge
content.


\subsection{Systematic comparison: infinite-dimensionality and causality}
\label{sec:systematic_comparison}

The reviewer asks for a systematic contrast on two axes ---
infinite-dimensionality and causality --- across the relevant
programs.  The following table provides this.

\smallskip
\begin{center}\footnotesize
\renewcommand{\arraystretch}{1.25}
\begin{tabular}{@{}p{2.4cm}p{3.2cm}p{3.2cm}p{3.5cm}@{}}
\toprule
\textbf{Program} & \textbf{Local algebra type} &
  \textbf{Causality treatment} & \textbf{APF relationship} \\
\midrule
AQFT (Haag--Kastler) &
  Type~III$_1$ factors (Reeh--Schlieder) &
  Axiom: $[\mathcal{A}(\mathcal{O}_1), \mathcal{A}(\mathcal{O}_2)] = 0$ for spacelike $\mathcal{O}_1, \mathcal{O}_2$ &
  APF recovers type~III in $n \to \infty$ limit; derives commutativity from $L_{\mathrm{loc}}$ \\[3pt]
LCQFT (Brunetti--Fredenhagen--Verch) &
  Infinite-dim.\ $*$-algebras; functorial on Lorentzian manifolds &
  Presupposes globally hyperbolic manifold; time-slice axiom &
  APF derives the manifold; provides a finite-dim.\ foundation \\[3pt]
Oeckl (GBF) &
  Infinite-dim.\ boundary Hilbert spaces &
  Compatible with indefinite causal structure &
  Structural parallel (interfaces $\sim$ boundaries); APF adds A1 finiteness and derives causal order \\[3pt]
Hardy (2001) &
  Finite-dim.\ (axiom 1) &
  Not addressed (kinematics only) &
  Shared starting point; APF extends to gauge content \\[3pt]
CDP (2011) &
  Finite-dim.\ (output) &
  Causality axiom (no signaling) &
  APF derives no-signaling from $L_{\mathrm{loc}}$; pushes beyond kinematics \\[3pt]
Masanes--M\"uller &
  Finite-dim.\ (postulate) &
  Not directly addressed &
  Shared $\mathbb{F} = \mathbb{C}$ result; different mechanism \\[3pt]
Connes NCG &
  Finite-dim.\ internal algebra; $\infty$-dim.\ external &
  Lorentzian structure from spectral triple &
  APF derives the algebra that NCG postulates \\[3pt]
DR reconstruction &
  Infinite-dim.\ (type~III in AQFT context) &
  Inherited from AQFT axioms &
  APF uses DR at finite $n$; structural parallel for abelian factors \\[3pt]
APF &
  Type~I$_n$ ($M_n(\mathbb{C})$); type~III recovered as $n \to \infty$ &
  Derived from $L_{\mathrm{irr}}$ (Prop.~\ref{prop:causal}); HKM/Malament for signature &
  --- \\
\bottomrule
\end{tabular}
\end{center}

\smallskip\noindent Three patterns emerge from the comparison:

\textit{(i) Finite vs.\ infinite dimensions is a foundational choice,
not a technical detail.}\quad Programs that start with infinite
dimensions (AQFT, LCQFT, Oeckl, DR) must contend with type~III
algebras, ultraviolet divergences, and the absence of a faithful
trace.  Programs that start with finite dimensions (Hardy, CDP,
Masanes--M\"uller, APF) avoid these issues but must explain how the
infinite-dimensional continuum is recovered.  The APF's controlled
limit (Proposition~\ref{prop:typeIII_limit}) and structural stability
theorem (Theorem~\ref{thm:structural_stability}) provide this
explanation; the kinematic reconstruction programs do not address the
question (they stop at finite-dimensional QT and do not discuss the
continuum limit).

\textit{(ii) Causality is either axiomatized or derived, never
both.}\quad AQFT and CDP axiomatize causality (as commutativity of
spacelike algebras or no-signaling); the APF derives it from
$L_{\mathrm{irr}}$ (irreversibility $\to$ causal order) and
$L_{\mathrm{loc}}$ (independent interfaces $\to$ commutativity).
Oeckl's GBF is compatible with indefinite causality, making it the
most flexible --- but at the cost of not determining the causal
structure.  The APF's derivation of causality from a non-causal
starting point (finite enforcement capacity has no intrinsic temporal
or causal content) is the closest parallel to the causal-set program
(Sorkin, 1991), though the APF derives a smooth manifold rather than
working with a discrete causal set.

\textit{(iii) Only the APF and Connes NCG derive gauge content.}\quad
All other programs in the table stop at the kinematics of quantum
theory (state space, measurement structure, composition rules).  The
passage from kinematics to dynamics and gauge content requires
additional structure: in NCG, the spectral triple; in the APF, the
enforcement algebra's Wedderburn decomposition + non-closure +
anomaly cancellation.  The APF's advantage over NCG is that it
\emph{derives} the algebra, whereas NCG postulates it; NCG's
advantage is that it produces the full SM Lagrangian (including Yukawa
couplings and the Higgs potential), which the APF defers to Paper~4A.


\subsection{Assumption inventory comparison}
\label{sec:assumption_comparison}

\smallskip
\begin{center}
\small
\begin{tabular}{@{}lccl@{}}
\toprule
\textbf{Program} & \textbf{Inputs} & \textbf{Derives} & \textbf{Stops at} \\
\midrule
Hardy (2001) & 5 axioms & QT kinematics & State space \\
CDP (2011) & 6 axioms & QT kinematics & State space \\
Masanes--M\"uller (2011) & 4 postulates & QT over $\mathbb{C}$ & State space \\
H\"ohn (2017) & Finite info & Qubit QT & Qubit \\
LCQFT (BFV 2003) & Functorial axioms & QFT on curved spacetime & Assumes manifold \\
Oeckl (GBF 2003) & Boundary axioms & QT amplitudes & Kinematics \\
Connes NCG & Spectral triple & SM Lagrangian & Assumes algebra \\
APF & A1+MD+BW+SP+K3 & QT + gauge + matter & SM content \\
\bottomrule
\end{tabular}
\end{center}

\smallskip\noindent The APF's named assumption count (the four PLEC components A1, MD, A2, BW; plus constitutive semantics SP, K3, FD1--FD6) is numerically comparable to Hardy's 5 or CDP's 6,
though a direct comparison of ``assumption weight'' is imprecise: the
APF's constitutive definitions (FD1--FD6) carry more internal structure
than a single axiom.  The significant difference is scope: the APF's
downstream derivation chain extends to gauge content, matter spectrum,
and cosmological fractions, which the other programs do not attempt.


% ══════════════════════════════════════════════════════════════
\section{Code--theorem concordance}
\label{app:code}
% ══════════════════════════════════════════════════════════════

Every theorem in this supplement is verified by a named check function
in the codebase.  The table below provides the mapping.

\smallskip
\begin{center}\small
\renewcommand{\arraystretch}{1.15}
\begin{tabular}{@{}llll@{}}
\toprule
\textbf{Theorem} & \textbf{Supplement \S} & \textbf{Function} & \textbf{File} \\
\midrule
$L_{\mathrm{Cauchy}}$ & \S\ref{sec:cauchy_supp} & \texttt{check\_L\_Cauchy\_uniqueness} & \texttt{L\_Cauchy\_uniqueness.py} \\
Gauge ident. & \S\ref{sec:gauge_ident} & \texttt{check\_T3} & \texttt{core.py} \\
$T_{7B}$ & \S\ref{sec:T7B} & \texttt{check\_T7B} & \texttt{gravity.py} \\
$L_{\mathrm{irr,uniform}}$ & \S\ref{sec:Lirr_uniform} & \texttt{check\_L\_irr\_uniform} & \texttt{core.py} \\
$B1'$ & \S\ref{sec:B1prime} & \texttt{check\_B1\_prime} & \texttt{gauge.py} \\
$L_{\mathrm{gauge,unique}}$ & \S\ref{sec:gauge_class_supp} & \texttt{check\_L\_gauge\_template\_uniqueness} & \texttt{gauge.py} \\
$T_{\mathrm{field}}$ & \S\ref{sec:fermion_scan} & \texttt{check\_T\_field}; \texttt{fermion\_scan\_standalone.py}~\cite{BrookeFermionScan} & \texttt{gauge.py} \\
$L_{\mathrm{hypercharge}}$ & \S\ref{sec:hypercharge} & \texttt{check\_T\_field} (hypercharges) & \texttt{gauge.py} \\
$L_{\mathrm{count}}$ & \S\ref{sec:capacity_count} & \texttt{check\_L\_count} & \texttt{gauge.py} \\
$P_{\mathrm{exhaust}}$ & \S\ref{sec:mece} & \texttt{check\_P\_exhaust} & \texttt{core.py} \\
$L_{\mathrm{equip}}$ & \S\ref{sec:equip} & \texttt{check\_L\_equip} & \texttt{cosmology.py} \\
$L_{\mathrm{sat,part}}$ & \S\ref{sec:sat_indep} & \texttt{check\_L\_saturation\_partition} & \texttt{cosmology.py} \\
$\Delta_{\mathrm{sig}}$ & \S\ref{sec:signature} & \texttt{check\_Delta\_signature} & \texttt{spacetime.py} \\
$T_8$ ($d{=}4$) & \S\ref{sec:d4} & \texttt{check\_T8} & \texttt{spacetime.py} \\
$T_{9,\mathrm{grav}}$ & \S\ref{sec:einstein} & \texttt{check\_T9\_grav} & \texttt{gravity.py} \\
$T_{\theta_{\mathrm{QCD}}}$ & \S\ref{sec:theta_QCD} & \texttt{check\_T\_theta\_QCD} & \texttt{gauge.py} \\
$T_{\mathrm{vac,stab}}$ & \S\ref{sec:vacuum_stability} & \texttt{check\_T\_vacuum\_stability} & \texttt{gauge.py} \\
$T_{\mathrm{proton}}$ & \S\ref{sec:proton_stability} & \texttt{check\_T\_proton} & \texttt{gauge.py} \\
\bottomrule
\end{tabular}
\end{center}

\smallskip\noindent The full codebase is available at
\texttt{github.com/Ethan-Brooke/APF-Paper-2-The-Structure-of-Admissible-Physics}.

\smallskip\noindent\textit{Reproducibility of the fermion scan.}\quad
The function \texttt{check\_T\_field} in \texttt{gauge.py} performs the
complete 1{,}680-template scan described in \S\ref{sec:fermion_scan}.
It enumerates all templates, applies the seven filters in sequence,
records exact survivor counts at each stage (matching the waterfall table),
and verifies that the unique minimum-DOF survivor is the SM at 45 Weyl
fermions.  The five closed-form exclusion proofs (P1--P5) are verified
as separate assertions within the same function.  Hypercharge derivation
and anomaly cancellation are checked algebraically using exact rational
arithmetic (\texttt{fractions.Fraction}).  Running
\texttt{python -m pytest gauge.py} reproduces all results with no
external dependencies beyond the Python standard library.

\smallskip\noindent\textit{Standalone verification script.}\quad
The file \texttt{fermion\_scan\_standalone.py}~\cite{BrookeFermionScan},
included in the GitHub repository and archived on Zenodo
(\url{https://doi.org/10.5281/zenodo.19154197}),
is a self-contained Python script ($\sim$550 lines, zero dependencies
beyond the standard library) that reproduces the complete scan from
scratch.  An interactive Colab walkthrough with waterfall visualisation
is available at
\url{https://colab.research.google.com/drive/1otO1Qwp_Lr3OsK6ykkeMRATanecPkKg6}.
Running \texttt{python3 fermion\_scan\_standalone.py} performs the following:
\begin{enumerate}[nosep,label=(\arabic*)]
  \item Enumerates all 1{,}680 distinct templates using
  \texttt{combinations\_with\_replacement}.
  \item Applies all 7 filters with exact rational arithmetic, including
  spectator-field reduction for $A(R) = 0$ representations.
  \item Prints the waterfall table with exact counts at each stage.
  \item Lists all 8 post-F6 survivors explicitly (matching the table in
  \S\ref{sec:waterfall}).
  \item Lists all 4 post-CPT survivors.
  \item Verifies all 5 exclusion proofs (P1--P5) with explicit arithmetic.
  \item Derives both hypercharge solutions ($z = 1 \pm N_c$),
  verifies all 4 anomaly conditions exactly for each, and confirms
  their physical equivalence under $u \leftrightarrow d$ relabelling.
  \item Runs 4 built-in red-team challenges: (RT1) combinatorial
  cross-check of the enumeration count, (RT2) explicit spectator-field
  anomaly verification, (RT3) dead-code reachability audit, (RT4) P4
  bound tightness with a constructed witness template.
  \item Asserts 12 final verification conditions.
\end{enumerate}
The script is designed so that any reader can verify the entire
fermion-content derivation on a standard laptop in under one second,
with no knowledge of the APF codebase required.

The Lie-algebra classification (\S\ref{sec:gauge_class_supp}) is verified by
\texttt{check\_L\_gauge\_template\_uniqueness}, which tests all 17
compact simple algebras against R1--R3 and records the exclusion reason
for each.  The Zenodo archive
(DOI: 10.5281/zenodo.18867119) provides a frozen snapshot of the
codebase corresponding to this supplement.




% ══════════════════════════════════════════════════════════════
\section{Self-contained proofs of imported results}
\label{app:imported}
% ══════════════════════════════════════════════════════════════

This appendix provides complete proofs of every result imported from the
Paper~1 Supplement that is used in this document.  With this appendix,
no external document is needed to verify any theorem in this supplement.

\smallskip\noindent\textit{Input assumptions.}\quad All proofs below
use only: A1 (finite enforcement capacity $0 < C(\Gamma) < \infty$),
MD (uniform per-test cost floor $\varepsilon \ge \mu^* > 0$),
K3 (additivity on disjoint supports), SP (no null directions in the
physical quotient), and the formal definitions FD1--FD6.  No result
from Papers~2--7 is used.


\subsection{$L_{\varepsilon^*}$: Uniform cost floor}
\label{app:Leps}

\begin{lemma_app}[$L_{\varepsilon^*}$]
There exists $\varepsilon^* > 0$ such that
$\varepsilon(d) \ge \varepsilon^*$ for all $d \in \mathcal{D}$.
\end{lemma_app}

\begin{proof}
By MD: $\varepsilon(d) \ge \mu^* \cdot n(d) \ge \mu^* > 0$ for every
$d$ (since $n(d) \ge 1$).  Set $\varepsilon^* := \mu^*$.
\end{proof}


\subsection{$T_{\mathrm{sep}}$: Sector decomposition}
\label{app:Tsep}

\begin{theorem_app}[$T_{\mathrm{sep}}$]
For distinctions $d_1, d_2 \in \mathcal{D}(\Gamma)$ with anchor sets
$M_{d_1}, M_{d_2} \subseteq S_\Gamma$:
\[
  M_{d_1} \cap M_{d_2} = \{0\}
  \quad\Longleftrightarrow\quad
  \varepsilon(\{d_1, d_2\}) = \varepsilon(d_1) + \varepsilon(d_2).
\]
For any maximal antichain $\mathcal{B}$ of pairwise independently
enforceable distinctions, the substrate decomposes as
$S_\Gamma = \bigoplus_{d \in \mathcal{B}} M_d \oplus \Pi$,
where $\Pi$ is the residual pool.
\end{theorem_app}

\begin{proof}
($\Rightarrow$)\quad If $M_{d_1} \cap M_{d_2} = \{0\}$, perturbations
threatening $d_1$ and $d_2$ have disjoint supports (FD3: anchor-set
locality).  The joint defense decomposes into independent defenses on
$M_{d_1}$ and $M_{d_2}$.  By K3 (additivity on disjoint supports):
$\varepsilon(\{d_1,d_2\}) = \varepsilon(d_1) + \varepsilon(d_2)$.

($\Leftarrow$, contrapositive)\quad If $v \in M_{d_1} \cap M_{d_2}$
with $v \ne 0$, then defending $d_1$ and $d_2$ separately each commits
capacity to $\mathrm{span}\{v\}$; defending jointly commits once.
By SP and K2, $\kappa(\mathrm{span}\{v\}) > 0$.  Write
$M_{d_i} = \mathrm{span}\{v\} \oplus M_{d_i}'$.  By K3 on the three
disjoint subspaces:
\[
  \varepsilon(\{d_1,d_2\}) \le \kappa(\mathrm{span}\{v\})
  + \kappa(M_{d_1}') + \kappa(M_{d_2}')
  < \varepsilon(d_1) + \varepsilon(d_2).
\]
The direct-sum decomposition follows by iterating the biconditional
over a maximal antichain.
\end{proof}

\noindent\textit{Uses:} A1, K3, FD3, SP.  No inner product or linear
map structure.


\subsection{$L_\Delta$: Superadditivity}
\label{app:LDelta}

\begin{theorem_app}[$L_\Delta$]
For co-located distinctions $d_1, d_2$ at $\Gamma$ with
$M_{d_1} \cap M_{d_2} = \{0\}$ and $\Pi \ne \{0\}$:
$\Delta(d_1,d_2) := \varepsilon(\{d_1,d_2\}) - \varepsilon(d_1)
- \varepsilon(d_2) > 0$.
\end{theorem_app}

\begin{proof}
\textit{Step 1.}\quad The joint perturbation class
$\mathcal{P}(\{d_1,d_2\})$ strictly contains
$\mathcal{P}(d_1) \cup \mathcal{P}(d_2)$: perturbations
$p$ with $\mathrm{supp}(p) \subseteq \Pi$ are in
$\mathcal{P}(\{d_1,d_2\})$ but not in either individual class
(their anchors are disjoint from $\Pi$ by $T_{\mathrm{sep}}$).
These perturbations change physically meaningful DOF (SP) at positive
cost (K2).

\textit{Step 2.}\quad A defense confined to $M_{d_1} \oplus M_{d_2}$
acts as the identity on $\Pi$ and cannot resist $\Pi$-supported
perturbations.

\textit{Step 3.}\quad Any admissible joint defense must commit positive
capacity $\kappa_\Pi > 0$ in $\Pi$ (by K2: resisting a non-null
perturbation requires positive capacity).

\textit{Step 4.}\quad $M_{d_1}$, $M_{d_2}$, $\Pi$ are pairwise disjoint
($T_{\mathrm{sep}}$).  By K3:
$\varepsilon(\{d_1,d_2\}) = \varepsilon(d_1) + \varepsilon(d_2)
+ \kappa_\Pi$.

\textit{Step 5.}\quad $\Delta = \kappa_\Pi > 0$.
\end{proof}

\noindent\textit{Uses:} A1, K2, K3, $T_{\mathrm{sep}}$, SP.
Pre-algebraic.


\subsection{$L_\omega$: Sector-orthogonal bilinear form}
\label{app:Lomega}

\begin{lemma_app}[$L_\omega$]
Given $T_{\mathrm{sep}}$ and K3, there exists a positive-definite
symmetric bilinear form $B: V \times V \to \mathbb{R}$ with
$B(u_d, v_{d'}) = 0$ for $d \ne d'$ and $B(u_d, v_\Pi) = 0$ for
all $d$.  The inter-sector orthogonality is forced; the within-sector
forms are free (normalization gauge).
\end{lemma_app}

\begin{proof}
\textit{Construction.}\quad Choose any positive-definite form $B_d$ on
each sector $M_d$ and $B_\Pi$ on $\Pi$.  Define
$B(u,v) := \sum_d B_d(u_d, v_d) + B_\Pi(u_\Pi, v_\Pi)$
with zero cross-terms.

\textit{$B$ is well-defined, symmetric, positive-definite} by
construction (direct sum of positive-definite forms).

\textit{K3 forces cross-term vanishing.}\quad For $u_d \in M_d$ and
$v_{d'} \in M_{d'}$ with $d \ne d'$, K3 gives
$\|u_d + v_{d'}\|_B^2 = \|u_d\|_B^2 + \|v_{d'}\|_B^2$ (disjoint
supports).  Expanding: $B(u_d, u_d) + B(v_{d'}, v_{d'}) +
2B(u_d, v_{d'}) = B(u_d, u_d) + B(v_{d'}, v_{d'})$, so
$B(u_d, v_{d'}) = 0$.  The same argument gives
$B(u_d, v_\Pi) = 0$.
\end{proof}

\noindent\textit{Uses:} $T_{\mathrm{sep}}$, K3.


\subsection{$T_{\mathrm{adj}}$: Self-adjoint enforcement projections}
\label{app:Tadj}

\begin{theorem_app}[$T_{\mathrm{adj}}$]
With respect to $B$ from $L_\omega$, each sector map $E_d$ is a
$B$-orthogonal projection: $E_d^\dagger = E_d$, $E_d^2 = E_d$.
\end{theorem_app}

\begin{proof}
$\mathrm{im}(E_d) = M_d$ and
$\ker(E_d) = (\bigoplus_{d' \ne d} M_{d'}) \oplus \Pi$.
By $L_\omega$: $M_d \perp_B \ker(E_d)$.  Therefore $E_d$ is the
$B$-orthogonal projection onto $M_d$:
$B(E_d u, v) = B(u_d, v_d) = B(u, E_d v)$, so $E_d^\dagger = E_d$.
\end{proof}


\subsection{$L_{\mathrm{loc}}$: Locality (budget additivity)}
\label{app:Lloc}

\begin{theorem_app}[$L_{\mathrm{loc}}$]
For independent interfaces $\Gamma_1, \Gamma_2$ with disjoint
substrates $S_{\Gamma_1} \cap S_{\Gamma_2} = \emptyset$:
$C(\Gamma_1 \cup \Gamma_2) = C(\Gamma_1) + C(\Gamma_2)$.
\end{theorem_app}

\begin{proof}
Every distinction $d$ has its anchor entirely within $S_{\Gamma_1}$ or
$S_{\Gamma_2}$ (FD3: anchor-set locality).  The cost sum partitions:
$\sum_d \varepsilon(d) = \sum_{d \in \mathcal{D}(\Gamma_1)}
\varepsilon(d) + \sum_{d \in \mathcal{D}(\Gamma_2)} \varepsilon(d)
\le C(\Gamma_1) + C(\Gamma_2)$.  The bound is tight (each interface's
budget can be independently saturated).
\end{proof}

\noindent\textit{Uses:} A1, FD3.  Does not use MD, $L_{\varepsilon^*}$,
or any cost-functional structure.


\subsection{$L_{\mathrm{blk}}$: Block-diagonal defense failure}
\label{app:Lblk}

\begin{theorem_app}[$L_{\mathrm{blk}}$]
Let $\Delta(d_1,d_2) > 0$.  Then no idempotent with
$\mathrm{im}(P) \subseteq M_{d_1} \oplus M_{d_2}$ is an admissible
joint defender.
\end{theorem_app}

\begin{proof}
Such a $P$ satisfies $P|_\Pi = 0$ (by idempotency and the direct-sum
structure).  It therefore has zero defense against $\Pi$-supported
perturbations.  But $\Delta > 0$ requires positive $\Pi$-defense
($L_\Delta$, Step~3).  Contradiction.
\end{proof}


\subsection{$T_{\mathrm{alg}}$: Noncommutativity}
\label{app:Talg}

\begin{theorem_app}[$T_{\mathrm{alg}}$]
If $\Delta(d_1,d_2) > 0$, the full enforcement algebra
$\mathcal{A} := \mathrm{Alg}_\mathbb{R}(\{E_d\} \cup
\{E_{\{d_1,d_2\}}\})$ is noncommutative.
\end{theorem_app}

\begin{proof}
\textit{Step 1.}\quad The sector projections $\{E_d\}$ commute and
generate $\mathcal{A}_{\mathrm{diag}} \cong \mathbb{R}^{|\mathcal{B}|}$.

\textit{Step 2.}\quad By $L_{\mathrm{blk}}$, the minimum-cost joint
defender $E_{\{d_1,d_2\}} \notin \mathcal{A}_{\mathrm{diag}}$.
Define $F_\Pi := E_{\{d_1,d_2\}} - E_{d_1} - E_{d_2} \ne 0$.

\textit{Step 3.}\quad There exists $v \in M_{d_1}$ with
$E_{\{d_1,d_2\}} v = v + \pi_v$, $\pi_v \in \Pi \setminus \{0\}$
(otherwise $E_{\{d_1,d_2\}}$ is block-diagonal, contradicting Step~2).
Then:
$[E_{d_1}, F_\Pi]\,v = E_{d_1}(\pi_v) - F_\Pi(v) = 0 - \pi_v
= -\pi_v \ne 0$,
using $E_{d_1}|_\Pi = 0$ ($T_{\mathrm{sep}}$).

\textit{Step 4.}\quad $[E_{d_1}, F_\Pi] \ne 0$ in $\mathrm{End}(V)$.
Any faithful $*$-homomorphism to a commutative algebra would map this
to zero, contradicting faithfulness.  Therefore $\mathcal{A}$ is
noncommutative.
\end{proof}

\noindent\textit{Uses:} $T_{\mathrm{sep}}$, $L_\Delta$,
$L_{\mathrm{blk}}$, FD5a.  Pre-Hilbert.


\subsection{$L_{\mathrm{nc}}$: Non-closure (superadditivity at shared
interfaces)}
\label{app:Lnc}

\begin{lemma_app}[$L_{\mathrm{nc}}$]
At any shared interface $\Gamma$ with $\Pi \ne \{0\}$ and
co-located distinction sets $G, R_\Gamma$:
$\varepsilon(G \cup R_\Gamma) > \varepsilon(G)
+ \varepsilon(R_\Gamma)$.
\end{lemma_app}

\begin{proof}
This is $L_\Delta$ applied to the distinction sets $G$ and $R_\Gamma$:
the hypotheses of $L_\Delta$ (disjoint anchors, non-empty pool) are
satisfied by the co-location conditions of the shared interface.
\end{proof}


\subsection{$L_{\mathrm{irr}}$: Irreversibility}
\label{app:Lirr}

\begin{theorem_app}[$L_{\mathrm{irr}}$]
At any shared interface $\Gamma$ with $\Pi \ne \{0\}$, there exist
enforcement transitions that are locally irreversible: the capacity
committed in the transition cannot be recovered by any operation
confined to a single sector.
\end{theorem_app}

\begin{proof}
\textit{Step 1.}\quad By $L_\Delta$ (\S\ref{app:LDelta}), co-located
enforcement at $\Gamma$ commits a superadditive surplus
$\Delta = \kappa_\Pi > 0$ in the pool sector $\Pi$.

\textit{Step 2.}\quad By $L_{\mathrm{loc}}$ (\S\ref{app:Lloc}),
enforcement at separate interfaces has independent budgets.  A
single-sector observer (operating within $M_{d_1}$ alone) has access
only to the capacity committed in $M_{d_1}$, not to the pool
commitment $\kappa_\Pi$.

\textit{Step 3.}\quad The pool commitment $\kappa_\Pi$ is a
cross-sector correlation: it is committed in $\Pi$, which is shared
between the sectors.  No operation confined to $M_{d_1}$ can access
$\Pi$ (by the direct-sum structure of $T_{\mathrm{sep}}$:
$M_{d_1} \cap \Pi = \{0\}$).  Therefore the pool commitment is locally
unrecoverable.

\textit{Step 4.}\quad A transition that increases $\kappa_\Pi$ (e.g.,
by changing the gauge-sector distinctions at a shared interface) commits
capacity that is locally unrecoverable.  This is irreversibility: the
forward transition costs more to reverse than the budget allows, because
the reversal would require freeing pool capacity that no single-sector
operation can access.
\end{proof}

\noindent\textit{Uses:} $L_\Delta$, $L_{\mathrm{loc}}$,
$T_{\mathrm{sep}}$.  The argument is pre-algebraic: no Hilbert space
or GNS structure is invoked.

\medskip\noindent\textit{Summary.}\quad The ten results above
($L_{\varepsilon^*}$, $T_{\mathrm{sep}}$, $L_\Delta$, $L_\omega$,
$T_{\mathrm{adj}}$, $L_{\mathrm{loc}}$, $L_{\mathrm{blk}}$,
$T_{\mathrm{alg}}$, $L_{\mathrm{nc}}$, $L_{\mathrm{irr}}$) constitute
the complete set of results imported from the Paper~1 Supplement.  Each
is proved from A1, MD, K3, SP, and FD1--FD6 alone.  With these proofs
in hand, every theorem in this supplement can be verified end-to-end
without consulting any external document.




% ══════════════════════════════════════════════════════════════
\section{Fermion scan: algorithmic walkthrough}
\label{app:scan_walkthrough}
% ══════════════════════════════════════════════════════════════

This appendix makes the fermion-content derivation
(\S\ref{sec:fermion_scan}) verifiable without running code.  It
specifies the search space construction, gives the algorithm in
pseudocode, and traces the critical filter steps on explicit
examples.

\subsection{Search space construction}
\label{app:search_space}

A template is a multiset of $(r^{(3)}, r^{(2)})$ pairs.  The
search space is bounded by P1--P5 (\S\ref{sec:exclusion_proofs}):

\smallskip\noindent\textbf{Allowed $\SU(3)$ reps} (P1: $\dim < 10$):
$\{3, \bar{3}, 6, \bar{6}, 8\}$.
Dynkin indices: $T(3) = T(\bar{3}) = 1/2$, $T(6) = T(\bar{6}) = 5/2$,
$T(8) = 3$.  Cubic anomaly coefficients: $A(3) = 1/2$,
$A(\bar{3}) = -1/2$, $A(6) = 5/2$, $A(\bar{6}) = -5/2$, $A(8) = 0$.

\smallskip\noindent\textbf{Allowed $\SU(2)$ reps} (P2--P3):
$\{1, 2\}$.  Dynkin indices: $T(1) = 0$, $T(2) = 1/2$.

\smallskip\noindent\textbf{Template structure} (P4--P5):
Each template contains exactly one colored doublet (P4: two colored
doublets require $\ge 63$ DOF), $0$--$3$ colored singlets, $0$--$1$
lepton doublets, and $0$--$2$ lepton singlets (P5: $> 5$ field types
require $> 45$ DOF).

\smallskip\noindent\textbf{Enumeration count.}\quad
The colored doublet ranges over 5 reps.  The colored singlets form a
multiset of size $0$--$3$ drawn from 5 reps: $\binom{5+0-1}{0} +
\binom{5+1-1}{1} + \binom{5+2-1}{2} + \binom{5+3-1}{3} = 1 + 5 + 15
+ 35 = 56$ multisets.  The lepton doublet is present or absent: 2
choices.  The lepton singlets: $0$, $1$, or $2$ copies of
$(\mathbf{1},\mathbf{1})$: 3 choices.  Total:
$5 \times 56 \times 2 \times 3 = 1{,}680$ distinct templates.
The standalone verification script~\cite{BrookeFermionScan}
uses \texttt{combinations\_with\_replacement} to enumerate these
directly, and verifies the count against this closed-form formula
as a built-in red-team challenge (RT1).

\subsection{Algorithm in pseudocode}
\label{app:pseudocode}

\begin{small}
\begin{verbatim}
INPUT:  SU3_reps = {3, 3b, 6, 6b, 8}
        SU2_reps = {1, 2}
        N_gen = 3
        AF3_bound = 11,  AF2_bound = 22/3
        AF_coeff = 2/3

ENUMERATE all templates T:
  FOR each colored_doublet IN SU3_reps:           // 5
    FOR n_colored_singlets IN {0, 1, 2, 3}:
      FOR each combo IN C(SU3_reps, n_cs):        // combinations w/ replacement
        FOR has_lepton_doublet IN {True, False}:   // 2
          FOR n_lepton_singlets IN {0, 1, 2}:      // 3
            T = [(colored_doublet, 2)]
                + [(c, 1) for c in combo]
                + [(1, 2)] if has_lepton_doublet
                + [(1, 1)] * n_lepton_singlets
            total_tested += 1

            // F1+F2: Asymptotic freedom (exact rational)
            su3_cost = SUM over (a,b) in T:
                         T_3(a) * dim_2(b) * N_gen
            su2_cost = SUM over (a,b) in T:
                         T_2(b) * dim_3(a) * N_gen
            IF AF3_bound - AF_coeff * su3_cost <= 0: SKIP
            IF AF2_bound - AF_coeff * su2_cost <= 0: SKIP

            // F3: Chirality
            IF no colored doublet or no colored singlet: SKIP

            // F4: [SU(3)]^3 anomaly cancellation
            IF SUM A_3(a) * dim_2(b) != 0: SKIP

            // F5: Witten anomaly
            IF SUM dim_3(a) [where b=2] is odd: SKIP

            // F6: Full anomaly system
            // (rational root test on the hypercharge
            //  quadratic -- details in §8)
            IF no rational hypercharge solution: SKIP

            // F7: CPT quotient
            canonical = min(sorted(T),
                           sorted(conjugate(T)))
            IF canonical already seen: SKIP
            Mark canonical as seen

            // Record survivor
            DOF = SUM dim_3(a) * dim_2(b) * N_gen
            survivors.append((DOF, T))

SELECT minimum-DOF survivor(s)
\end{verbatim}
\end{small}


\subsection{Filter-by-filter trace on critical examples}
\label{app:filter_trace}

We trace three templates through all seven filters: the SM (passes
all), a near-miss (fails at F6), and a common failure (fails at F4).

\smallskip\noindent\textbf{Example 1: SM template}
$\{(3,2),\, (\bar{3},1),\, (\bar{3},1),\, (1,2),\, (1,1)\}$.

\smallskip\noindent\textit{F1 (SU(3) AF):}\quad
$b_0^{(3)} = 11 - \tfrac{2}{3}\bigl[\tfrac{1}{2}\!\cdot\!2 +
\tfrac{1}{2}\!\cdot\!1 + \tfrac{1}{2}\!\cdot\!1\bigr]\!\cdot\!3
= 11 - \tfrac{2}{3}\!\cdot\!2\!\cdot\!3 = 11 - 4 = 7 > 0$.  \checkmark

\smallskip\noindent\textit{F2 (SU(2) AF):}\quad
$b_0^{(2)} = \tfrac{22}{3} - \tfrac{2}{3}\bigl[\tfrac{1}{2}\!\cdot\!3
+ \tfrac{1}{2}\!\cdot\!1\bigr]\!\cdot\!3
= \tfrac{22}{3} - \tfrac{2}{3}\!\cdot\!2\!\cdot\!3 = \tfrac{22}{3} - 4
= \tfrac{10}{3} > 0$.  \checkmark

\smallskip\noindent\textit{F3 (chirality):}\quad
$(3,2)$ is a colored doublet; $(\bar{3},1)$ is a colored singlet.
Both present.  \checkmark

\smallskip\noindent\textit{F4 ($[\SU(3)]^3$):}\quad
$\tfrac{1}{2}\!\cdot\!2 + (-\tfrac{1}{2})\!\cdot\!1
+ (-\tfrac{1}{2})\!\cdot\!1 = 1 - \tfrac{1}{2} - \tfrac{1}{2} = 0$.
\checkmark

\smallskip\noindent\textit{F5 (Witten):}\quad
Doublets weighted by $\SU(3)$ dim: $3 + 1 = 4$ (even).  \checkmark

\smallskip\noindent\textit{F6 (full anomaly):}\quad
$N_c = 3$, $n_{\mathrm{cs}} = 2$, $n_{\mathrm{ls}} = 1$.
Discriminant $= 4 + 4(9-1) = 36$.  $\sqrt{36} = 6 \in \mathbb{Z}$.
\checkmark

\smallskip\noindent\textit{F7 (CPT):}\quad Canonical form
$= \min(\text{sorted}(T), \text{sorted}(\overline{T}))$.
$\overline{T} = \{(\bar{3},2),\, (3,1),\, (3,1),\, (1,2),\, (1,1)\}$.
Sorted: $(1,1)(1,2)(3,2)(\bar{3},1)(\bar{3},1)$ vs.\
$(1,1)(1,2)(3,1)(3,1)(\bar{3},2)$.  First is lexicographically
smaller.  Canonical = SM.  \checkmark

\smallskip\noindent DOF $= (6+3+3+2+1)\times 3 = 45$.
\textbf{Passes all filters.}

\bigskip
\noindent\textbf{Example 2: $\{(6,2),\, (\bar{6},1),\, (1,2),\, (1,1)\}$}
(a template with the sextet).

\smallskip\noindent\textit{F1:}\quad
$b_0^{(3)} = 11 - \tfrac{2}{3}[\tfrac{5}{2}\!\cdot\!2 +
\tfrac{5}{2}\!\cdot\!1]\!\cdot\!3 = 11 - \tfrac{2}{3}\!\cdot\!
\tfrac{15}{2}\!\cdot\!3 = 11 - 15 < 0$.  $\boldsymbol{\times}$ \textbf{Fails F1.}

\bigskip
\noindent\textbf{Example 3: $\{(3,2),\, (3,1),\, (1,2),\, (1,1)\}$}
(wrong-sign colored singlet).

\smallskip\noindent\textit{F1--F3:}\quad Pass (verify: $b_0^{(3)} = 11
- \tfrac{2}{3}\!\cdot\!\tfrac{3}{2}\!\cdot\!3 = 8 > 0$; chirality
holds).

\smallskip\noindent\textit{F4:}\quad
$\tfrac{1}{2}\!\cdot\!2 + \tfrac{1}{2}\!\cdot\!1 = \tfrac{3}{2}
\ne 0$.  $\boldsymbol{\times}$ \textbf{Fails F4.}  The cubic anomaly is not
cancelled because both colored fields are in the $\mathbf{3}$ (same
sign of $A$).

\subsection{Completeness argument}
\label{app:completeness}

The scan's completeness rests on three facts:

\begin{enumerate}[nosep]
\item \textbf{Representation bound (P1--P2).}\quad Every $\SU(3)$ rep
with $\dim \ge 10$ and every colored $\SU(2)$ rep with $\dim \ge 3$
exceeds the AF budget as a \emph{single field} (the cost of one such
field, summed over 3 generations, exceeds $b_0$).  These are
closed-form inequalities verified in \S\ref{sec:exclusion_proofs}.
Therefore the allowed reps are exactly
$\SU(3) \in \{3,\bar{3},6,\bar{6},8\}$ and $\SU(2) \in \{1,2\}$.

\item \textbf{Multiplicity bound (P4--P5).}\quad At most 1 colored
doublet (P4: 2 colored doublets force $\ge 63$ DOF) and at most 5
total field types (P5: each type adds $\ge 3$ DOF per 3 generations).
These are arithmetic inequalities.

\item \textbf{Exhaustive enumeration.}\quad Within the bounded set
($5 \times 56 \times 2 \times 3 = 1{,}680$ distinct templates), every
combinatorial possibility is generated and tested.  No branch is
skipped.  The enumeration is deterministic and order-independent (the
final result depends only on the filter outcomes, not on the order
of testing).
\end{enumerate}

\noindent Together, (1)--(3) establish that the search is
\emph{exhaustive}: every conceivable chiral fermion template with
$\SU(3) \times \SU(2)$ quantum numbers, any number of generations,
and any multiplicity structure is either in the 1{,}680-template set or
is excluded by a closed-form proof.

\subsection{Robustness checks}
\label{app:robustness}

\begin{enumerate}[nosep]
\item \textbf{Exact arithmetic.}\quad All filter computations use
Python's \texttt{fractions.Fraction} (exact rational arithmetic).  No
floating-point approximation is made at any stage.  The anomaly
conditions are checked as exact equalities, not as
$|\text{residual}| < \epsilon$.

\item \textbf{Filter-order independence.}\quad The final survivor set
is determined by the conjunction of all seven filters.  The sequential
application (F1 before F2 before F3, etc.)\ is an optimization
(early termination), not a logical dependency.  Reordering the filters
produces the same final set.  The standalone script
(\texttt{fermion\_scan\_standalone.py}) applies all filters and records
intermediate counts; permuting the filter order and comparing outputs
verifies independence.

\item \textbf{CPT symmetry.}\quad The CPT quotient (F7) identifies
each template with its charge conjugate.  The conjugation map
$3 \leftrightarrow \bar{3}$, $6 \leftrightarrow \bar{6}$,
$8 \to 8$, $1 \to 1$ is an involution.  The canonical representative
is the lexicographic minimum of the pair $\{T, \overline{T}\}$.  The
4 post-F6 survivors form 2 CPT pairs (listed in \S\ref{sec:waterfall}).

\item \textbf{Sensitivity to generation count.}\quad The scan assumes
$N_{\mathrm{gen}} = 3$ (derived in Paper~2, \S9).  With
$N_{\mathrm{gen}} = 1$ or $2$: the AF bounds loosen ($b_0$ terms
scale as $N_{\mathrm{gen}}$), admitting additional survivors.  With
$N_{\mathrm{gen}} \ge 4$: the AF bounds tighten, reducing or
eliminating survivors.  The $N_{\mathrm{gen}} = 3$ scan is therefore
not vacuous (it admits templates) and not trivial (it excludes most).

\item \textbf{Second survivor.}\quad The second CPT-inequivalent
survivor (DOF $= 48$, SM $+$ extra singlet) provides a
falsifiability test: if the SM contained 48 Weyl fermions (e.g., with
a right-handed neutrino), the framework would predict $C = 64$ and
$\Omega_\Lambda = 42/64 = 0.656$, excluded at $> 5\sigma$ by Planck
(\S\ref{sec:neutrinos}).
\end{enumerate}


% ══════════════════════════════════════════════════════════════
% BIBLIOGRAPHY (placeholder)
% ══════════════════════════════════════════════════════════════

\begin{thebibliography}{99}

\bibitem{BrookePaper1}
E.\,S.~Brooke, ``Paper~1: The Enforceability of Distinction,''
Admissibility Physics, 2025.
Submitted to \textit{Foundations of Physics}.

\bibitem{BrookeSupplement}
E.\,S.~Brooke, ``Technical Supplement: End-to-End Derivation from Finite
Enforcement Capacity to Quantum Structure,'' companion to Paper~1, 2025.

\bibitem{Humphreys1972}
J.\,E.~Humphreys, \textit{Introduction to Lie Algebras and Representation
Theory}, Springer, 1972.

\bibitem{FultonHarris1991}
W.~Fulton and J.~Harris, \textit{Representation Theory: A First Course},
Springer, 1991.

\bibitem{Witten1982}
E.~Witten, ``An SU(2) Anomaly,'' \textit{Phys.\ Lett.\ B} \textbf{117},
324--328 (1982).

\bibitem{PeskinSchroeder1995}
M.\,E.~Peskin and D.\,V.~Schroeder, \textit{An Introduction to Quantum
Field Theory}, Westview Press, 1995.

\bibitem{GrossWilczek1973}
D.\,J.~Gross and F.~Wilczek, ``Ultraviolet Behavior of Non-Abelian
Gauge Theories,'' \textit{Phys.\ Rev.\ Lett.} \textbf{30}, 1343 (1973).

\bibitem{Politzer1973}
H.\,D.~Politzer, ``Reliable Perturbative Results for Strong Interactions?''
\textit{Phys.\ Rev.\ Lett.} \textbf{30}, 1346 (1973).

\bibitem{Planck2020}
Planck Collaboration, ``Planck 2018 results.\ VI.\ Cosmological
parameters,'' \textit{Astron.\ Astrophys.} \textbf{641}, A6 (2020).

\bibitem{PDG2024}
Particle Data Group, ``Review of Particle Physics,''
\textit{Phys.\ Rev.\ D} \textbf{110}, 030001 (2024).

\bibitem{Hardy2001}
L.~Hardy, ``Quantum theory from five reasonable axioms,''
arXiv:quant-ph/0101012 (2001).

\bibitem{CDP2011}
G.~Chiribella, G.\,M.~D'Ariano, and P.~Perinotti,
``Informational derivation of quantum theory,''
\textit{Phys.\ Rev.\ A} \textbf{84}, 012311 (2011).

\bibitem{Hohn2017}
P.\,A.~H\"ohn, ``Quantum theory from questions,''
\textit{Phys.\ Rev.\ A} \textbf{95}, 012102 (2017).

\bibitem{Connes1996}
A.~Connes, ``Gravity coupled with matter and the foundation of
noncommutative geometry,''
\textit{Commun.\ Math.\ Phys.} \textbf{182}, 155--176 (1996).

\bibitem{Bousso1999}
R.~Bousso, ``A covariant entropy conjecture,''
\textit{JHEP} \textbf{07}, 004 (1999).

\bibitem{HKM1976}
S.\,W.~Hawking, A.\,R.~King, and P.\,J.~McCarthy,
``A new topology for curved space-time which incorporates the causal,
differential, and conformal structures,''
\textit{J.\ Math.\ Phys.} \textbf{17}, 174--181 (1976).

\bibitem{Malament1977}
D.\,B.~Malament, ``The class of continuous timelike curves determines
the topology of spacetime,''
\textit{J.\ Math.\ Phys.} \textbf{18}, 1399--1404 (1977).

\bibitem{Kolmogorov1933}
A.\,N.~Kolmogorov, \textit{Grundbegriffe der Wahrscheinlichkeitsrechnung},
Springer, 1933.

\bibitem{Lovelock1971}
D.~Lovelock, ``The Einstein tensor and its generalizations,''
\textit{J.\ Math.\ Phys.} \textbf{12}, 498--501 (1971).

\bibitem{BrookeFermionScan}
E.\,S.~Brooke, ``APF Fermion Template Scan---Standalone Verification
Script (v2),'' 2026.
Zenodo: \url{https://doi.org/10.5281/zenodo.19154197}.
GitHub: \url{https://github.com/Ethan-Brooke/APF-Paper-2-The-Structure-of-Admissible-Physics}.
Interactive walkthrough:
\url{https://colab.research.google.com/drive/1otO1Qwp_Lr3OsK6ykkeMRATanecPkKg6}.

\end{thebibliography}


\end{document}
